\xchapter{Preface}

THE
Works of which this Volume is composed, being of an historical
character, naturally require a Chronological Table of the principal
events recorded in them; but the difficulties of forming any
satisfactory statement, during the period to which they belong, are so
great, that any arrangement can be but hypothetical, and must be
accompanied with some notice of the difficulties themselves, and the
various expedients which have been adopted with the view of overcoming
them. Though such notice will be necessarily very imperfect, it shall
here be attempted.

\xsubsection{1. Interposition of Pope Julius in the affairs of the East}

It is certain, that both the
Eusebians and the Egyptian Bishops had recourse to Rome; that
Athanasius went thither; that a synodal judgment was passed there; and
that Legates went from S. Julius to Antioch; but the order and dates
of these events are variously determined. For the sake of perspicuity,
it will be necessary in the first place to take a view of the
transactions to which dates are to be assigned; though it is
impossible to do so, without prejudging some of the questions in
dispute.

It appears then, that shortly
after the return of S. Athanasius to Alexandria from his exile in
Gaul, the Eusebian party brought charges against him before the three
Emperors, (infr. pp. 18, 226,) and the Pope, (p. 37.) Their embassy or
legation to the latter consisted of Macarius, Martyrius, and Hesychius, (pp. 42, 47.) and they were met by a counter deputation
from S. Athanasius, (pp. 44, 226,) supported, (p. 48,) or preceded,
(p. 43,) by letters from many Catholic Bishops, (pp. 47, 70,) and by a
letter to the Pope, (p. 38,) which an Alexandrian Council of from
eighty, (p. 61,) to one hundred Bishops, (p. 14,) had written in his
favour, (pp. 14, 17, and 48.) The discussions which ensued at Rome
perhaps were held before a Council of Bishops then present, (p. 46,)
and ended in the defeat of the Eusebian legates, (p. 43,) one of whom
abruptly left the city in consequence, (p. 44.) Julius, however, did
not decide the matter at once, but at their suggestion, (pp. 39, 42,
226,) proposed a Council, (p. 11,) at which both Eusebians, (p. 54,)
and Athanasius should attend, (p. 40,) and the Alexandrians have the
choice of place, (p. 226.) Athanasius, who was otherwise disposed to
betake himself to Rome, in consequence of the outrages of Gregory whom
the Arian Council of the Dedication had sent to Alexandria in his
place, (p. 227,) promptly obeyed the call (p. 49); and on his arrival
at Rome, the Pope sent Elpidius and Philoxenus as legates to Antioch,
(p. 39,) with a letter to the Eusebians, (p. 46,) repeating the
invitation to a Council, (p. 41,) and fixing the day, (pp. 45, 227.)
There they were detained over the time, ibid. and at length came back
with a refusal on the part of the Orientals to attend (pp. 40, 46,
47); though the Eusebian legates had not only been the originators of
the measure, but had gone so far as to offer to submit the question to
the arbitration of the Pope, (p. 39.) Upon this Julius proceeded to
hold a Council of fifty Bishops, (pp. 14, 39, 230,) at which
Athanasius and others were pronounced innocent and admitted to
communion, ibid. and in the name of which, (pp. 39, 46,) the Pope,
eighteen months from the date of Athanasius's arrival, (p. 49,)
proceeded to address a letter of remonstrance to the Orientals, who
had written to him from Antioch.

This is a sketch of the
history, and now to proceed to its chronology. The only date which is
known for certain is that of the Eusebian Council of Antioch, held A.D. 341. This
we learn from Athanasius, de Syn. §. 25. “Ninety Bishops,” he
says, “met at the Dedication under the Consulate of Marcellinus and
Probinus, in the 14th of the Indiction;” L. F. vol. 8, p. 109. As,
in dating by the Indiction, the new year began in September, the
Council must have assembled during the spring or summer of 341; nay,
it would appear, in the first months of it, if Gregory, who was
appointed in it to the See of Alexandria, began his persecution at
Alexandria in that year. Gregory entered Alexandria during Lent, (infr.
p. 7.) that is, either in Lent 341 while the Council was still
sitting, or the Lent following. Upon Gregory's coming, Athanasius left
Alexandria for Rome, that is, after Easter; thus Athanasius's visit to
Rome commences in the spring of 341 or 342; unless indeed we suppose
with Mansi, that Gregory's invasion and Athanasius's flight were prior
to the Council of the Dedication, viz. in 340. He remained at Rome
three years, (p. 158.) and in the fourth year was called by Constans
to Milan. Now in the latter part of 345 the delegates of the Eusebians
also came to Milan, Eudoxius, Martyrius, and Macedonius, (vid. L. F.
vol. 8, p. 111.) with the Macrostich or Long Confession, which had
been drawn up at Antioch in the beginning of the year. They presented
themselves before a Council there, according to a letter of Liberius,
of the date of 354; which rejected them; and that, according to the
same letter, eight years before that date, which nearly agrees with
Athanasius's account of the publication of the Macrostich. It is
natural to connect this visit of the Eusebians to Milan with the
summons of Athanasius by Constans to that city, and to conclude that
the proceedings of the Council issued in the resolution which the
Emperor adopted at this time to treat with his brother for the meeting
of a General Council. If so, the date of Athanasius's journey to Rome
is 342. And it certainly seems much more probable that Gregory should
proceed to Alexandria the Lent after the Dedication, than that the
ecclesiastical and military acts and movements \footnote{1. Pagi after Schelstrate
contends, that the Confession of faith and the Canons preceded the
cause of Athanasius in the Council. Montfaucon and Tillemont, (with
the exception of the Canon, which was expressly levelled at Athanasius,
and which Montf. does not notice as a Canon,) place it first of all.
If there were at first orthodox Bishops at the Council, as is said, we
cannot suppose, that Athanasius was condemned till after their
departure. Schelstrate, who places matters of faith and discipline
first, in his task of vindicating the Catholicity of the Council, is
obliged to suppose its commencement in 340, in order to gain time for
Gregory's expedition by Lent 341.} which attended his expedition should be despatched between
January and Lent, which the date of 341 requires, i.e. did not
Athanasius's words p. 226. on the other hand shew that the Eusebians
were very much bent on the measure, and were likely to prosecute it
promptly. And Baronius and others date the Councils of the Macrostich
and of Milan at 344, not 345, which throws back the journey of
Athanasius to 341. And moreover if the Anonymus Maffeianus, relied on
by Mansi, be correct, the Council of Sardica was held at the end of
344, a date which may just allow time for a preliminary Council of
Milan (in 344.) between the Sardican Council and the end of three
years from May 341. In this uncertainty about the year of Athanasius's
journey to Rome, 341 may be more fitly taken than 342 or 340, as
having the suffrages of more critics in its favour. But in this
question does not consist the main difficulty of the chronology on the
point before us, which is internal to the documents which are to
follow, arising out of the relative not the absolute dates which they
contain.

It appears that S. Athanasius
was eighteen months at Rome before Pope Julius's letter, (p. 49;) that
is, the Council of Rome, in or upon which he wrote it, was ending or
just ended eighteen months after Athanasius's arrival, or in the month
of October or rather November, since he set out for Rome after Easter.
But the meeting of the Council was fixed for a day before the January
preceding that November; because the Pope's legates who were sent into
the East upon Athanasius's arrival at Rome are said, by being kept at
Antioch till January, to be kept over the time of meeting. Thus
we have an interval of eleven months between the meeting and the
termination. It follows then that the Council did not meet at the time
proposed, \emph{or} that it was continued for nearly a whole year, \emph{or}
that there were two Councils, one in December, the other in November.
Now as to the last supposition, it is most improbable that the same
Bishops of Italy should meet twice over at so short a period, and
Julius and Athanasius speak distinctly of but one synodal \emph{body},
(even supposing they are not clear about one meeting,) which both
pronounced the innocence of Athanasius and commissioned Julius to
write. Still less is it conceivable that the Council should be
prolonged for ten or eleven months. Nor can we easily conjecture, what
is at first sight plausible, a postponement of the day of meeting, for
Julius seems positively to say that they met at the very time for
which they had been convened. (p. 46.)

In this difficulty, which can
on no hypothesis perhaps be satisfactorily removed, some critics have
thrown the \emph{fault}, as it may be called, upon one place in the
history, others on another.

The form in which it has been
above exhibited is that which arises out of the arrangement of facts
and dates first suggested by Valesius, and adopted after him by
Schelstrate, Pagi, Montfaucon, Coustant, Du Pin, S. Basnage, and
others. It seems far more natural and less open to objections than any
other; and perhaps the readiest explanation of the difficulty, which
has been above described as attaching to it, is to consider the letter
of Pope Julius to be later than the Italian Council by eleven months,
and written in the ordinary Autumnal Synod (Baron. 342. 34.), to
which, on occasion of the delay of the Eusebians, the Italian Council
of December, might naturally delegate \footnote{2.
Tillemont will be found to make a similar suggestion, vol. 7. pp. 706,
7. He supplies parallel instances.}, as to a sort of Committee, the office of concluding
negociations with them and issuing the Council's sentence, whenever
the legates of the Pope should return. What makes this the more
probable is, that Julius speaks of Athanasius as being among the \emph{Romans}
eighteen months. “He continued here a year and six months, … his
presence overcame us all,” p. 49, words which properly belong to
Bishops residing in the neighbourhood, not to an Italian Council. It
is observable, moreover, that Julius says, “the sentiments I am
expressing are not those of myself alone, but of all the Bishops
throughout Italy, \emph{and in these parts},” [\emph{en toutois tois
meresi}], p. 46. (Baronius, however, adduces this passage in order
to shew that S. Julius's first letter issued from a Council.) And he
proceeds, “The Bishops \emph{now too}, [\emph{kai nun}], assembled
on the appointed day,” as if there had been a former appointment,
and that punctually kept; (though Valesius and Schelstrate understand
the words, “I \emph{again} write,” which follow, to refer to
Julius's former communication with the Eusebians before Athanasius's
coming, as we may understand it still.) And that a delay of some kind
was occasioned in the proceedings at Rome by the conduct of the
Eusebians, is plain, as various critics observe, from Julius's words,
p. 40, “I, when I had read your letter, after much consideration, \emph{kept
it to myself}, thinking that \emph{after all} some of you would
come … but when no one arrived, and it became \emph{necessary} that
the letter should be produced, \&c.” This passage too accounts
for the long interval between the departure of the legates from the
Eusebians in January, and the Pope's Letter to them of the November
following in answer.

Such is the disposition of the
dates which is the most satisfactory on the whole; but it must not be
concealed, that names of the greatest weight may be alleged in favour
of other chronological arrangements. Such is Baronius, who has been
followed by Labbe, Petavius, and others; such are Hermant, Papebroke,
and Tillemont, who adopt a third hypothesis. Such again is Mansi, who
follows an arrangement of his own, founded on a document which has
come to light since the time of his predecessors.

Baronius supposes two visits
of Athanasius to Rome, and two Italian Councils held there. He refers
to a statement of Socrates, as apparently the basis of the former of
these suppositions; though Socrates is so inextricably perplexed in
his account of the events and even of the names of persons which occur
in the history, that it is difficult to determine what he does and
what he does not say on this point. Baronius refers to Hist. ii. 11.
where no such statement occurs. He may be taken, however, to say,
(e.g. ii. 15.) that Athanasius after his acquittal at Rome returned to
Alexandria before the violent entrance of Gregory, upon which he
retired to Rome a second time. Accordingly, Baronius terminates the
eighteen months some time before Lent, 342, which he considers the
date of Gregory's entrance, or towards the close of 341, and places
their commencement, that is, the first journey of Athanasius in the
early part of 340, and the Council of Alexandria in 339. Further,
since the termination of the eighteen months must coincide with the
date of the Roman Council, which acquitted Athanasius, he supposes
that Council to have been held in 341, before the outrages of Gregory,
and before the return of the legates, whom he sends into the East in
340, previous to Athanasius' first journey, and brings back to Rome
not till 342, when Julius holds a second Council, in which he writes
his synodal letter.

Baronius urges in behalf of
his two Councils that Pope Julius notices in his Letter written from
the Council, the complaint of the Eusebians that Athanasius had been
admitted to communion, which was undeniably the act of the Council of
fifty Bishops. Valesius answers first by denying that Julius notices
any such complaint, next by arguing that the act of the Council of
fifty was not mere admission into communion, for Athanasius had never
been out of communion, and of this the Eusebians might be complaining,
but a formal recognition of his being, and deserving to be, in
communion with the Church. And hence Athanasius says, that they gave
him “the confirmation of their fellowship,” p. 39. [\emph{ekur\={o}san
t\={e}n koin\={o}nian}]. As to the question, which has been
raised, whether the Pope suspended communion with Athanasius, it is
treated of by Tillemont, vol. 8. p. 673.

Tillemont, though he agrees
with Baronius in supposing two journeys of S. Athanasius to Rome,
follows Papebroke in differing from him altogether in the dates at
which he places them. He argues that the Council at Rome must be dated
shortly after the Council of the Dedication at Antioch 341; after it,
because Julius complains that the Eusebians had anticipated him \footnote{3.
Schelstrate of course, whom Pagi follows, will not allow any
intentional anticipation on the part of the Council, which he
maintains to be in its beginnings Catholic, and to have assembled at
the end of 340 to dedicate the Aureum Dominicum.}, (p. 50.) and but shortly after, because they pleaded the
suddenness of the summons to Rome as a reason for not going, whereas
it had been sent them by the Pope's legates as far back as the
foregoing year. And he considers that the legates set out in the year
340, because in Athanasius's Encyclical Letter, written in the spring
of 341, mention is made (p. 11.) of an intention at Rome to hold a
Council for settling the existing troubles, an intention moreover the
news of which occasioned the Eusebians to assemble at Antioch in 341.
Accordingly he places the Council of Rome in June of that year; and
this, in spite of S. Julius's express statement that January, when the
legates were dismissed from Antioch, was about (because just beyond)
the time when the Council was held, meeting the difficulty by an
arbitrary alteration of the text, of June for January. And he supposes
the Council to continue by adjournment and representation till the
return of the legates, when S. Julius wrote his letter to the
Eusebians. Athanasius's eighteen months therefore terminated at this
date, i.e. in the autumn of 341; but, as agreeing with Valesius in
fixing Gregory's arrival at Alexandria in Lent of that year, Tillemont
is obliged to suppose that the eighteen months were not consecutive,
even if they were complete. He dates Athanasius's first coming as at
the end of 339 \footnote{4. The
words [\emph{monon akousas}] in Athanasius, infr. p. 227. §. 11 init.
are felt as a difficulty both by Tillemont and Montfaucon; by
Montfaucon, as if shewing that his flight was before Gregory's coming;
by Tillemont, as shewing that it was after Gregory's ordination.}; considers
that he returned to Alexandria in the course of 340 on the rumour
of the Eusebian movements at Antioch, and retired a second time to
Rome on the forcible entrance of Gregory during the Lent following.

Valesius argues against the
double journey of Athanasius from the strong negative fact that
Athanasius nowhere speaks of more than one, (vid. infr. pp. 39,
\&c. 158, 227, \&c.) He considers too that he could not have
returned to Alexandria without formal Letters from Constantius, which
there is no appearance of his obtaining.

Mansi differs from other
critics in this, that he rejects the testimony of Socrates, \&c.
upon which it rests that Gregory's appointment proceeded from the
Council of the Dedication, and considers his violences at Alexandria
to have taken place in Lent 340. He argues from the language of
Athanasius in his Encyclical Letter and elsewhere that Gregory
certainly was not elected by Bishops, and therefore not in a Council,
(vid. infr. pp. 5, 64, 229, \&c.) Yet surely, according to
Socrates, \&c. Athanasius was deposed by the Council “because he
had violated a rule which they themselves then passed,” viz. that he
had exercised his episcopal office without the formal leave of a
Council of Bishops; and it can hardly be supposed that, when the
Eusebians took the pains to be thus formal, they had already
despatched Gregory to take possession of the Alexandrian See. And Pope
Julius's letter too, p. 50 fin. implies that the Council passed some
act against Athanasius. Hence Schelstrate and Pagi maintain that he
was not deposed till after the question of faith and at least some
canons had been settled. Mansi, however, relies upon a document
discovered by Maffei in the Veronese Library, presently to be
mentioned, which anticipates the date of Athanasius's return after the
Council of Sardica by some years, placing it on Oct. 21, 346. and
assigning six years and six months for the length of his exile. In
consequence he fixes Athanasius's flight from Gregory and journey to
Rome at the beginning of 340, agreeing with Baronius and Papebroke
in supposing that it was preceded, as Sozomen reports Hist. ii.
9. by a time of concealment. He places the Council of Rome at the end
of the eighteen months after Athanasius's arrival, i.e. towards the
end of 341. And he argues that the Council of the Dedication was held
in the month of August, from the circumstance of St. Jerome's
assigning the Council in his Chronicon to the fifth year of the
Emperors, (as does Socrates Hist. ii. 8.) while the fourteenth of the
Indiction, which is also its date, ended with the beginning of
September. But the fifth year from Constantine's death began on May
22; and from the new Emperors' assumption of the title of Augustus,
not in August as Mansi states, (vid. Suppl. Cone. p. 175.) but on
Sept. 9. vid. Tillem. Emp. t. 4. p. 312. l'Art de verifier les Dates,
t. 1. p. 392.

The mention of the accession
of the sons of Constantine leads to the notice of one date in which
Schelstrate, Pagi, and Montfaucon, as well as Papebroke, and Tillemont,
side with Baronius against Valesius, who wishes to make 337 instead of
338 the year of S. Athanasius's return from Gaul. Valesius argues in
favour of 337, from the circumstances that Constantine the younger in
his letter to the Church of Alexandria, (infr. p. 121.) which is dated
June 17, designates himself as “Cæsar,” not by the title of
Augustus, which he assumed with his brothers the September after his
accession, i.e. Sept. 9, 337. Valesius adds, that while the brothers
were but Cæsars, Constantine would have the highest authority of the
three, as being the eldest; as if thus accounting for Constantine's
writing to the Alexandrians, not Constantius their sovereign.
Tillemont, after Schelstrate and Pagi, urges in reply the testimony of
Theodoret, who says that Athanasius was two years and four months at
Treves; and as he arrived there not before the end of 335, (Tillem.
Montf.) or in 336, (Baron. Schelstr) he did not leave till 338.
Moreover, Constantine's letter was written too soon after his father's
death, on the supposition of its belonging to 337, to allow even of
his hearing of that event, much less of his speaking, as he
does, of his father's wishes as regards Athanasius. It appears too
that the three brothers met in Pannonia in 338, where Athanasius tells
us, (infr. p. 159,) he had about this time an interview with
Constantius, viz. at Viminacium; it is natural then to suppose that
the letter of Constantine was the consequence of the meetings then and
there held. And while Athanasius, (infr. p. 225,) expressly says, that
his return was the joint act of the three brothers, it is known that
Constantius and Constans were at Viminacium in June 338, since one of
their laws bears this date and place; not to say that, according to
Epiphanius, Constantius's approbation of the return of Athanasius was
given when that Emperor was at Antioch, which he is known to have been
in October 338. (vid. Schelstrate, Pagi.) As to Valesius's difficulty
about Constantine's title, Pagi solves it by observing that
Constantine was writing to a Church under his brother's jurisdiction,
and in such case he would naturally drop the title Augustus, though he
was in possession of it. He refers to parallel instances. And as to
Constantine's writing at all, it is sufficient to answer that Treves
where Athanasius was staying was within his territory.

Valesius also maintains, that
the Encyclical Letter was written on occasion of the second attack on
the Alexandrian Church, by George in 356, not upon the first under
Gregory. He is misled by the faults in the text noticed infr. p. 1,
which Baronius had corrected from the necessity of the case, and which
Montfaucon has been able to set right from one of his Mss. To meet the
difficulty which the mention of Philagrius creates, of whose
connection with Gregory we are informed by Athanasius himself, infr.
p. 224, Pagi, who, as well as Schelstrate, follows Valesius in this
point, supposes that there were two Prefects of the name of Philagrius,
the second the son of the first. He supports this supposition by the
mention which occurs, (ibid.) of a Philagrius, Vicar of Cappadocia,
i.e. under the Prefect, and who cannot, he considers, be the man who
had served the higher office of Prefect of Egypt. In this way would be explained the praise bestowed upon a Philagrius by Nazianzen,
(vid. ibid. note b.) whom he supposes to be the second of the two.

\xsubsection{2. The council of Sardica}

If any period in the life of
S. Athanasius might at first sight be considered free from
chronological difficulties, it would be that which lies between his
second and his third exiles. Baronius, Montfaucon, and Tillemont,
whose dates we have found so discordant in the foregoing years, have
hardly a subject of difference in those which follow. There is a
general consent among them and the critics which come between them
concerning the date of the Council of Sardica, the restoration of S.
Athanasius, and the irruption of Syrianus and his flight. The great
difficulties attaching to the Councils of Sirmium in these years
scarcely fall into the narrative of his life. Thus stands the matter,
if we confine ourselves to the discussions and researches of the
seventeenth century. But in the course of the eighteenth a fresh
source of information was discovered, which, while it added perplexity
to the perplexed period which has already come under review, brought
into serious difficulty the hitherto unquestioned dates of the Council
of Sardica, and of S. Athanasius's return to Alexandria consequent
upon it.

Maffei published from the
Library of Verona a fragment of the Latin Version of Annals of the
life of S. Athanasius, written apparently in Greek at Alexandria, and
not very long after the times which it records. The high value which
he sets upon this document, is confirmed by the judgment of Mansi and
the Ballerini, the latter of whom call it an “aureum opusculum,”
Observ. in Noris. p. 834. and the former has made it the basis of a
new chronological arrangement \footnote{5. Vid.
also Vallars. in Hieron. Chron. p. 793.}.
That it contains very great historical misstatements is evident at
first sight; but it is a question whether these may not be attributed
to the ignorance of the translator, errors in transcription, e.g. in
numerals, and other causes; while on the other hand, were the
mistakes even so numerous and flagrant, an apparent internal
consistency as well as plausible external support may be urged in
behalf of those particular statements, on which are founded the
corrections of the chronology of the historical period now under
review.

In the very passage which is
of main importance in the inquiry, and with which the fragment opens,
we find a glaring error, at variance too with the account which
follows. “Post Gregorii mortem Athanasius reversus est ex urbe Româ
… et remansit quietus apud Alexandriam annis xvi. et mens. vi.”
whereas it is notorious, as the Annalist himself goes on to say, that
he was driven into banishment again in little more than nine years.

In the paragraph that follows,
the Author speaks of the Consuls of the year 349, as \emph{Hypatius}
and Catulinus, instead of Limenius; and of Eusebius of Nicomedia as
then alive, who died in 341 or 342; and of the murder of Hermogenes at
Constantinople, which took place at the same date. Mansi, however, has
a very ingenious explanation of the mistake in the Consul's name.

Afterwards he speaks of
Constans for Constantius, and Gregory for George.

The statement in which we are
immediately concerned runs thus: “Et factus est, post Gregorii
mortem Athanasius reversus est ex urbe Româ et partibus Italiæ et
ingressus est Alexandriam, Phaophi xxiv. Consulibus Constantio iv. et
Constante iii. hoc est post annos vi.” The Consuls named belong to
346, and the Egyptian date, according to Mansi, corresponds to October
21; whereas the received date of Athanasius's return is 349, and is
computed thus:—Sozomen Hist. iii. 12. places the Council of Sardica
in the Consulate of Rufinus and Eusebius, that is, A.D.
347. From the Council an embassy or legation was sent by Constans to
his brother, consisting of Euphrates and Vincentius. What happened to
them at Antioch we read infr. p. 235, and it took place “at the
season of the most holy Easter,” which must be 348, Easter-Day
being April 3; now Gregory died “about ten months after,” p. 236;
that is, in February 349, upon which Athanasius was restored to his
see, ibid. But on the other hand, reckoning backwards, if his
restoration took place, as the Annalist would have it, in 346, then
Euphrates and Vincentius were at Antioch at Easter 345, and the
Council took place in 344.

In another place the anonymous
Annalist speaks of the irruption of Syrianus, infr. p. 206. as
occurring, “Mechir xiii. die per noctem supervenientem xiv.” or
February 9, which answers to the received account infr. p. 294. and
adds, ”Hoc factum est post annos ix, et menses iii, ac dies xix,
quam Italiâ reversus est Episcopus;” a period, which, reckoning
according to Alexandrian months of thirty days, consistently answers,
as Maffei and Mansi observe, to the interval between Oct. 21, 346. and
Feb. 9, 356. One cannot suppose then the date assigned, whatever be
its value, to have been altered in transcription or translation. It is
the date intended by the Author. Now in St. Jerome's Chronicon, the
year assigned for Athanasius's return, is the tenth year of Constans,
that is, this very year 346, though the date A.D.
is there otherwise marked, viz. as 350 (349). Theodoret too reckons
the length of Gregory's usurpation at six years, which, however
treated, cannot be made to reach to 349. Moreover, if Euphrates was
convicted of Arianism in 346, which is the date assigned to the
Council of Cologne, he could not have been a legate from the Council
of Sardica to Constantius in Easter 348; but this difficulty, so
celebrated in controversy, vanishes, if for 348 we substitute 345, as
the date of the visit of Euphrates to Antioch. It may be added, that
in Surius's Edition of the Council of Sardica, the Consuls of 344 are
named in the title; which is also the case in an ancient Ms. of the
Collection of Mercator formerly contained in the Jesuit Library at
Paris, though other chronological specifications are added
inconsistent with this date.

What alterations in the
chronology of the period seem to be required by this and other
notices contained in the fragment under consideration, will be seen by
inspecting Mansi's table, a specimen of which shall presently be
given. Here the dates set down by the Annalist himself shall be set
before the reader.

Entrance of
S. Athanasius into Alexandria on his return from Italy.
Oct. 21, 346.

Legation of
five Bishops from S. Athanasius to Constans [Constantius] at
Milan.
May 19, 353.

Montanus the
Palatine enters Alexandria, four days after, with Letters

from the Emperor to S. Athanasius
prohibitory of his legation.
May 23, 353.

Diogenes the
Notary comes to Alexandria with a view of driving S. Athanasius

from the city.

he was there 4 months from the intercalation (after
July) to Dec. 22.
end of July,

355.

Syrianus
enters Alexandria.
Jan. 5, 356.

Breaks into
the Church at night.
Feb. 9, 356.

George is
driven from Alexandria.
Oct. 2, 358.

Death of S.
Athanasius.
May 3, 373.

It does not fall within the
scope of this Preface to enter into the Chronology of the Councils of
Milan, upon which so much has been written. On the critics who have
treated the subject and their respective judgments, vid. Pagi, ann.
344. n. 4.

\xsubsection{3. Councils of Sirmium}

Something was said on the
subject of the Councils of Sirmium, in the eighth Volume of the
Library of the Fathers, p. 160, in course of enumerating the Sirmian
and other Confessions. Mansi, however, was scarcely referred to; and
Zaccaria who has written after him not at all. A few words will be
sufficient to supply the omission.

Socrates and Sozomen assign
the condemnation of Photinus at Sirmium to a Council held there in
351. Baronius, Sirmond, and Gothofred, consider them mistaken, and fix
it in the year 357, towards or at the end of which, Constantius came
to that place, and remained there through the greater part or
whole of 358, and part of 359, (Gothofred in Philost. p. 200. Mansi,
Suppl. Conc. p. 182. ed. 1748.) Petavius, Tillemont, S. Basnage,
\&c. speak of three Councils or Conferences of Sirmium, placing
them respectively in 351, 357, and 359. Gothofred three, in 357, 358,
359. Mansi three, in 358, 359, 359. Zaccaria makes in all five, viz.
in 349, (in which indeed he follows Petavius,) 351, 357, (at which
Hosius lapsed,) 357 (following Valesius and Pagi,) and 359. The main
point of dispute is, whether there are \emph{two} dates for Sirmian
Councils, 351, and 357-9, or but \emph{one}, and that, at the latter
period, the former date, though assigned by Socrates, being in that
case impossible; and the main argument in favour of Baronius and Mansi,
who assert that there was but one, is the improbability, be it great
or be it little, that there should have been two Councils or
Conferences in that city, of an ecumenical not local character, within
a few years of each other. There does not seem much more to be said
than this, against Petavius and other advocates for 351 and 357.

This is evident from the mode
in which Mansi draws out his argument. He urges that Socrates and
Sozomen, the two writers who date the Council at 351, nevertheless
state, that “George, Bishop of Alexandria,” was present at it,
that is, George of Cappadocia, who was not consecrated till 356, and
was not driven from Alexandria till the end of August, (or Oct. 2,
according to the Anonymus,) 358. The Council then was held towards the
end of that year, a date at which we happen to know that Constantius
was making a long stay at Sirmium. Such seems the utmost of Mansi's
argument. Tillemont had already urged the mention of George to shew
that there was a Sirmian Council at a later date, but it does follow
from thence, as Tillemont well understands, that still Photinus was
not condemned at an earlier Council held in 351. Now the reasons for
the latter opinion, with the replies made to them, are as follows: 1.
Socrates dates in this place by naming the Consuls (of the foregoing
year,—there were no Consuls in 351,) and is never wrong, according
to Petavius, when he dates by the Consuls. Mansi, however,
denies this, and Zaccaria concedes it, vid, also infr. p. xxi. 2. The
Council of Sirmium, says Tillemont, was composed of Bishops of the
East, yet held in Illyricum, all which agrees with the date 351, when
the West was under the power of usurpers; Mansi accounts for the fact
by alleging that the West had already declared its judgment in two
Councils held against Photinus at Rome and Milan. 3. Basil of Ancyra,
who was the life of the Council against Photinus, opposed himself at
Ancyra to the Council of 357 or 358; which obliges us to distinguish
between the two Councils. Mansi explains by stating, what was the
fact, that there were two parties, Arians and Semi-Arians, at the
Council, and that when the latter, of which Basil was the leader, left
it, the former stayed and passed the Confession which Hosius
subscribed, and Basil, \&c. at Ancyra repudiated. 4. Germinius, who
succeeded Photinus in the see of Sirmium, sat as Bishop as early as
the Council of Milan, 355; it is answered, that at least he was Bishop
of Cyzicus before the deposition of Photinus. 5. Theodore, who
subscribed the formulary against Photinus, was dead in 355, that is,
if the Theodore who subscribed was the Bishop of Heraclea, and this
formulary the confession which Liberius signed. vid. Hilar. Fragm. vi.
7. 6. Cecropius of Nicomedia, says Zaccaria against Mansi, though not
against Baronius, was present at the Council, but he was killed in the
earthquake in that city, August 28, 358. 7. Pagi too observes, that
the disputation between Basil and Photinus was taken down, according
to Epiphanius, Hær. 71. p. 829. by “Callicrates, registrar of
Rufinus the Prefect;” now if Prætorian Prefect be meant, Rufinus
was Prefect of Illyricum 349-352. Exceptores or registrars were
attached to all judges, Gothofr. Cod. Theod. t. 2. p. 459. but they
are especially connected with Prætorian Prefects by Gothofred, ibid.
Pancirollus Not. Dign. p. 36. and Lami Erud. Apost. p. 262.

\xsubsection{4. The year of S. Athanasius's death}

Though there is nothing in the
following Treatises which leads specially to a discussion of the year
of S Athanasius's death; yet since it is one of the principal points
of controversy in a history which, as we have seen, abounds in
chronological difficulties, and is closely connected with passages
which occur below, it will not be out of place here to set down the
opinions of various critics on the subject. Many of them are collected
together in Fontanini's Dissertation appended to his Historia
Literaria Aquileiensis.

Socrates places his death in
the Consulate of Gratian ii. and Probus, that is, in 371; in which he
is followed by Petavius; Hermant in his Life of S. Athanasius; P. F.
Chifflet, (upon Ep. Paulin. 29.) Paulin. Illustr. part. 2. c. 11. p.
150; Papebroke in vit. Ath. p. 248; and Sollerius (who answers Pagi
and Montfaucon in a very disagreeable tone) de Patriarchis
Alexandrinis, Act. SS. in t. 5. Jun.

Baronius; Valesius (Theod.
Hist. iv. 22.); Renaudot, Hist. Patriarch. Alex. p. 95; and Fontanini \emph{supr}.
adopt the date of 372, from the duration of his Episcopate being 46
years, (on which there is a general agreement,) and its commencement
in 326. Sollerius too confesses, that of the two he should prefer 372
to 373, de Patr. Alex. n. 213. and it can hardly be doubted, that this
date would have, what may be called, the second votes of the advocates
both of 371 and of 373.

Cardinal Noris in his Censur.
in Not. Garner. (Opp. t. 3. p. 1178.) in correction of a former
statement in his Hist. Pelag. in which he agreed with Baronius; his
Editors the Ballerini in their Obss. p.834; Bucherius (in. Victor.
Can. Pasch.); Pagi; Quesuel (Leon. Opp. t. 2. p. 1545. ed. Baller.);
Du Pin, making S. Athanasius's Episcopate “more than 48 years;”
Oudinus (in supplem. Script. Eccles.); Tillemont; Montfaucon; Ceillier
(Hist. des Aut. Eccles.); S. Basnage (Annal), Le Quien (Or. Christ. t.
2. p. 400.); Scip. Maffei (Osserv. Lett. t. 3.); and Mansi in
the Dissertation quoted above, (though he speaks respectfully of
Sollerius's objections, in Pag. Ann. 372. 9.) argue in favour of 373.
This last opinion, which Montfaucon is considered to have established,
in his Vit. Ath. and a “Dissertatio de tempore mortis Alex. Ep.
Alex. ac de anno ob. Athan. M.” (which has not fallen in the way of
the present writer,) is founded principally upon S. Proterius's
Paschal Epistle.

Little seems to be adduced in
favour of 371, beyond the circumstance of Socrates mentioning the
Consuls of that year, a mode of dating which, according to Baronius,
may ordinarily be trusted, (in Ann. 69. n. 36.) that is, in the case
of public acts or contemporary events, as Montfaucon observes, Fontan.
Diss. p. 444. Petavius, however, says, Socrates nunquam temere, aut
falso notas Consulares adhibet, de Phot. Hær. c. 2. p. 379; on this
point, however, something has occurred above, p. xix. After alleging
the evidence of Socrates, Sollerius, who is the latest of the above
advocates of the year 371, does little more than attempt to adjust
that date with other existing chronological data, and to refute
objections.

The most obvious difficulty in
his hypothesis is, that Socrates himself, in the very passage in which
he mentions the Consuls of 371, states that S. Athanasius was Bishop
for 46 years, which, since he did not succeed Alexander till 326, will
bring the date of his death to 372 or 373. A controversy follows,
whether his consecration was at the end of 326, or at the beginning.
S. Alexander died, according to the Coptite History, as late as April
17 (326); but according to Athanasius himself, infr. p. 88. and
Theodoret, within five months after the reception of the Meletians,
(which followed upon the termination of the Nicene Council, i.e. upon
Aug. 25, 325,) and therefore in the beginning of 326, or the end of
325. Epiphanius too reports, that S. Alexander died the year of the
Nicene Council, Hær. 69. 11. (though he adds what invalidates his
testimony, or rather turns it the other way;) and his Festival is
fixed in the Roman Martyrology on Feb. 26. Next comes the
question of the interval between Alexander's death and Athanasius's
ordination, which Sollerius of course wishes to curtail as much as
possible. With this view he refers to the words of the Alexandrian
Council, infr. p. 22, which he interprets to imply, that the vacancy
in the see was immediately filled, and he maintains, after Papebroke,
that the Greek Feast-Day of S. Athanasius, Jan. 18, was really the day
of his consecration, i.e. in 326. However, though this be granted for
argument's sake, even then the 46 years of S. Athanasius's Episcopate
extend to January 372, i.e. beyond May 2, (his day of death,) 371. Nor
can we suppose, that Socrates merely uses round numbers, when he
speaks of 46 years, for S. Cyril expressly tells us, that Athanasius's
Episcopate was “46 \emph{whole} years;” and Theodoret, Sozomen,
the Arabian writers, (Renaudot Hist. Patr. Alex. p. 96.) and others
say the same thing. Yet Rufinus, who was in Egypt about the time of
Athanasius's death, certainly says only, that he died in his 46th
year.

And here at first sight is an
argument in favour of 372, rather than 373; Papebroke and Fontanini
observe, that S. Athanasius would have been Bishop 47 not 46 years on
supposition of the latter date. But this depends on the time of year
at which his Episcopate commenced. Sollerius maintains above, that it
dates from January 18; but Montfaucon (in his Monitum in correction of
his Vit. Athan.) and Tillemont place the death of S. Alexander on the
17th or 18th of April, following the Jacobite Chronicon of Abraham
Eckellensis, as above cited, and other Coptite, as well as Abyssinian
Calendars. To the five months spoken of above by Athanasius and
Theodoret, must in this case be added, as indeed is reasonable, the
time consumed in the return of S. Alexander from Nicæa to Alexandria,
and the proceedings in reconciliation of the Meletians, which will
make up the whole interval between August 25, and the April following.
Again, S. Athanasius's consecration does not seem to have followed
immediately upon the death of his predecessor, infr. p. 22.
which will carry down the beginning of his Episcopate far into the
year 326; and if we date it from the middle or the end, and much more
if, as the Ballerini propose, we fix it on Jan. 18, 327, then 46 years
and some months, or as it is natural that S. Cyril should express it,
46 whole years, will bring us to May 2, (the received day of his
death,) 373. The known duration then of S. Athanasius's Episcopate
does not decide between 372 and 373, being consistent with the latter
date as well as with the former. Other arguments, decisive against
371, but available for both 372 and 373, are deducible from the date
of the coming of Valens to Antioch, where, as Socrates tells us, he
was staying at the time of S. Athanasius's death; and of Melania's
visit to Alexandria, when Athanasius gave her Macarius's
sheep-skin,—a proof, says Montfaucon, that Athanasius was not dead
then, a proof, says Fontanini, that he was dying.

The direct evidence in favour
of 373 has been mentioned above. It consists in the Paschal Epistle of
S. Proterius, a contemporary of S. Leo, which is contained in Petavius's
Doctr. Temp. t. 2. who, however, p. 889. ed. 1627. as Sollerius and
Fontanini after him, thinks the text corrupt and untrustworthy, as it
evidently is in part. Sollerius also argues against it as irrelevant
in its context, and unmeaning. It is confirmed by S. Jerome's
Chronicon, which places Athanasius's death in the 10th year of Valens;
and by the Coptite History, which, by dating it on a Thursday, fixes
it in 373; and especially by Maffei's fragment, of which so much has
been said above. Collateral evidence is gained from the date of the
consecration of S. Basil 370, who, when he was Bishop, corresponded
with S. Athanasius; which, under the circumstances, could hardly have
been the case, had Athanasius died in 372. Sollerius, however,
suggests, that the Athanasius addressed by S. Basil was Athanasius of
Ancyra, at one time an Arianizer, though afterwards zealous for
orthodoxy, ii. 250.

It only remains to exhibit the
historical events which have come under review according to the
respective chronologies which different critics have adopted.

Dates
according to Valesius, Schelstrate, Pagi, Montfaucon, Sam. Basnage

A.D.

S. Athanasius returns from
Gaul 337.
337. \emph{V.}\\
338. \emph{S.P.M.B.}

leaving
Treves end of June, \emph{M.}

Three Eusebian Legates sent
to Rome.
339. \emph{V.S.P.M.B.}

Council of Alexandria.
340. \emph{S.P.M.B.}

Council of the Dedication.
341. \emph{V.S.P.M.B.}

in beginning of Year, \emph{V}.
end of 340, till January

341, \emph{S}. before Sept.
\emph{P}. to anticipate Roman,

\emph{Bar}. not to anticipate Roman, \emph{S.P.}

Entrance of Gregory into
Alexandria. Lent.
341. \emph{V.P.M.B.}

Athanasius writes his
Encyclical Letter.

in concealment, \emph{M.}[in 356 according to \emph{V.S.P}.]
341. \emph{M.}

S. Athanasius escapes to
Rome.

March or April, \emph{S.P.} after Easter, (April 19,) \emph{V.}May \emph{M}. after Council of
Dedication, \emph{P}.
341. \emph{V.S.P.M.B.}

Legates set out from Rome to
the Eusebians.

before Athanasius arrives there, and in

beginning of Year, \emph{V.}\\
after Athanasius's arrival, in March or April,

\emph{S. P.} May, \emph{M.}
341. \emph{V.S.P.B.}

Legates arrive at Antioch.

in April or May or June, \emph{S.} in June,
\emph{P.}
341. \emph{S.P.}

Legates set out from
Antioch.

January, \emph{S.B.M.}they return in March or April, \emph{S.}
342. \emph{V.S.M.B.}

Council of Rome, in which
Athanasius is acquitted. 

October, \emph{S.B.} or November, \emph{M}.
342. \emph{V.S.P.M.B.}

The Pope's Letter to the
Eusebians.
342. \emph{V.M.B.M.}

Baronius
and Petavius

Athanasius returns from
Gaul.
338. \emph{B.P.}

The three Eusebian Legates,
Macarius, \&c. sent

to Rome.
339.
\emph{B.P.}

Council of Alexandria.
339. \emph{B.}

The Legates sent from the
Pope to the Eusebians.
340. \emph{B.}

Athanasius comes to
Rome (first time) beginning of
340. \emph{B.P.}

Council of the Dedication at
Antioch,

to anticipate Roman Council, \emph{B.}
341. \emph{B.P.}

First Council of Rome, in
which Athanasius is acquitted.
341. \emph{B.P.}

Athanasius returns
immediately to Alexandria, 

end of year, or beginning of next, \emph{B.}
341. \emph{B.P.}

Eusebians send back the
Legates.

after
the Council of Rome, \emph{B.} before it, \emph{P.}
341. \emph{B.P.}

Entrance of Gregory into
Alexandria, Lent.
342. \emph{B.P.}

Athanasius retreats from
Alexandria into a place

of concealment.
342. \emph{B.P.}

He writes his Encyclical
Letter.
342. \emph{B.}

The Pope's Legates return to
Rome.
342. \emph{B.}

Second Council of Rome.
342. \emph{B.}

The Pope's Letter to the
Eusebians.
342. \emph{B.}

Athanasius comes to Rome
(second time).
342. \emph{B.P.}

Papebroke,
Tillemont

S. Athanasius returns from
Gaul.
335. \emph{P.T.}

The three Eusebian Legates
sent to Rome.
339. \emph{T.}

Council of Alexandria.
339. \emph{P.T.}

S. Athanasius goes to Rome.

and
his 18 months begin, \emph{T.} September, \emph{P.}
339. \emph{P.T.}

The Legates sent from the
Pope to the Eusebians, 

immediately after Sept. 339. \emph{P.}
340. \emph{T.}

S. Athanasius returns to
Alexandria,
end of

340. \emph{P.T.}

Council of the Dedication.

beginning
of Year, \emph{T.}before
September, \emph{T.}
341. \emph{P.T.}

Entrance of Gregory into
Alexandria, Lent.
341. \emph{P.T.}

S. Athanasius writes his
Encyclical Letter.
341. \emph{P.T.}

He leaves Alexandria and
retreats to Rome.

after
Easter, \emph{T}.
341. \emph{P.T.}

The Pope's Legates leave
Antioch.

in
June not January, \emph{P.T.}
341. \emph{P.T.}

Council of Rome,

opened before return of Legates, \emph{P.}sitting
till August or September, \emph{T.}
June

341. \emph{P.T.}

The Pope's Letter to the
Eusebians.

August or September, \emph{T.} 341. \emph{T.}

Mansi

Entrance of Gregory into
Alexandria.
Lent, 340.

S. Athanasius leaves
Alexandria for a place 

of concealment.
May, 340.

He goes to Rome.
June, 340.

Council of the Dedication.
August, 341.

Council of Rome.
End of 341.

\emph{Baron.}
\emph{Pag.}
\emph{Mont.}
\emph{Tillem.}
\emph{Mans.}

Macrostich is
drawn up by Arian Council

of Antioch.
344.
345.
344.
345.
end of

343.

It is
rejected by the Westerns in the Council

of Milan.
344.
346.
345.
345.
344.

when the Arian Legates leave the Assembly

in anger.
344.
346.
345.
345.
346.

Council of
Sardica.
347.
347.
347.
347.
end of

344.

Sardican
Legates at Antioch.
Easter,

348.
348.
348.
348.
345.

Death of the
usurper Gregory,
Jan. or
Feb.

349.
349.
349.
349.
346.

Council of
Cologne deposes Euphrates.
346.
346.

346.
346.

Council of
Milan against Photinus, at which Valens 

and
Ursacius appear.
350.
347.
347.
347.
346.

Council of
Jerusalem.
350.
349.
349.
349.
346.

S. Athanasius
returns to Alexandria.
350.
349.
349.
349.
Oct. 21.,

346.

First Sirmian
Council against Photinus.
357.
351.
351.
351.
358.

Montanus
comes to Alexandria.
351.
353.
353 or

354.

end of

May 353.

Diogenes
the Notary attempts to drive S. Athanasius

from
Alexandria.
354.
355.
355.
355.
end of

July 355.

Irruption of
Syrianus into the Church,
Feb. 9.,

356.
356.
356.
356.
356.

George is
driven from Alexandria.
357.

358.
Oct 2.,

358.

Second
Sirmian Council or Conference, in

which was passed
the “blasphemia,” vol. 8. p 161.
Beg. of

357.
End of

357.
357.
357.
359.

Council of
Ancyra just before Easter.
357.
358.

358.
359.

Third Sirmian
Council or Conference.
357.
358.
May 22, 

359.
359.
359.

Council of
Ariminum, July 21.
359.
359.
359.
359.
359.

Death of S.
Athanasius, May 2.
372.
373.
373.
373.
373.

Before concluding, it is
necessary to observe, that in the references in the notes or margin,
S. Athanasius's Works are designated by their Latin titles for the
sake of clearness; and “Hist. Arian.” is the same work as “ad
Mon.” There is some unavoidable irregularity in the mode of
reference to former Volumes of this series, e.g. “Libr. F.” with
the Volume specified, is equivalent to “Oxf. Tr.'' or “O.T.”
or to the name of the Treatise with “Tr.” added. Also the
reference is sometimes made according to pages, sometimes according to
sections \&c. Consistency has not been thought of much consequence
in a matter of this kind, where clearness and conciseness of reference
were rather to be consulted in each particular case.

Also it may be right to refer
the reader to a Letter addressed to Montfaucon on the words [\emph{thall\={o}n}]
or “boughs,” infr. p. 270. in the Collectio Nova (t. ii. in Cosm.
p. 18); and to a note of Quesnel's on S. Leo, (t. 3. p. xlvii. ed.
Baller.) who observes, that Siscia, infr. p. 60. is not a province,
but the city of that name in Pannonia.

And. it should be added to
page 13, that Tillemont dates the Apologia contra Arian. not earlier
than A.D.
356. arguing from the mention of the banishment of Liberius and Hosius.
Also in note g, p. 49, justice is not done to Baronius's view of
Athanasius's double journey to Rome, as the foregoing pages will shew.
And in p. 76, note m, Thomassin is quoted not to corroborate Febronius's
interpretation, but principle.

Also in p. 46, Valesius Obss.
Eccles. i. 2. p. 174. understands Eusebius himself by [\emph{hoi peri
Eusebion}] §. 26. Montfaucon observes, that Eusebius alone is
spoken of in §. 1. He adds, “res hic in dubio versatur.” Baronius
adduces the phrase as used in the Encyclical Letter in proof that it
was written while Eusehius was still alive, but Valesius denies the
argument on grammatical grounds, Obss. Eccl. i. 7 fin. Montfaucon,
however, observes, in his Monitum prefixed to that Letter, that in
matter of fact the phrase is never used by S. Athanasius of
Eusebius's party after E.'s death, but always [\emph{koin\={o}noi t\={o}n
peri E.}] or [\emph{kl\={e}ronomoi t\={e}s asebeias tou E.}]

Also with reference to the
subject of note n, p. 77. it should be observed, that the majority of
critics side with Du Cange against Gothofred on the meaning of the
word Canalis. “Those Bishops,” says Baronius, were “in Canalio,
qui sedes haberent in cursu publico, viâ scilicet quâ equi publici
per stationes singulas dispositi essent ad iter agendum.” An. 347.
55. “Qui præerant sacris urbium, quæ regiæ viæ insidebant,”
says Noris, professing his agreement with Baronius, Opp. t. 4. p. 623.
Pitiscus also, “qui sedes habent in cursu publico,” \emph{in voc}.
So also Kiesling, adding, “intelliguntur hoc nomine urbes, seu
potius civitates, in quibus Episcopi sedem habuerunt fixam.” de
Discipl. Cleric. p. 13. Beveridge reports Zonaras and Balsamon as
furnishing the same interpretation; “cities which are in the public
ways, or canal, through which travellers pass without trouble, as
water flows in an aqueduct.” Pandect. t. 1. p. 507.

For the Translation, the
Editors have to express their acknowledgments to the Rev. MILES
ATKINSON,
M.A. late Fellow of Lincoln College.

J. H. N.

\emph{Dec}. 4,
1843.

\xchapter{I. Encyclical Epistle of the Blessed Athanasius, Bishop of Alexandria}

[S. Athanasius wrote the following Epistle in the
year 341. In that year the Eusebians held the famous Council of the
Dedication at Antioch, vid. Athan. de Syn. §. 25. (Libr. F. vol. 8.
p. 109, \&c.) Here they appointed Gregory to the see of Alexandria
in the place of Athanasius, whom they had already condemned and
denounced at the Synod of Tyre, A.D.
335. Gregory was by
birth a Cappadocian, and, (if Nazianzen speaks of the same Gregory,
which some critics doubt,) studied at Alexandria, where S. Athanasius
had treated him with great kindness and familiarity, though Gregory
afterwards took part in propagating the calumny against him of having
murdered Arsenius. Gregory was on his appointment dispatched to
Alexandria with Philagrius Prefect of Egypt, and their proceedings on
their arrival are related in the following Encyclical Epistle, which
S. Athanasius forwarded immediately upon his retreat from the city to
all the Bishops of the Catholic Church. It is less correct in style,
as Tillemont observes, than other of his works, as if composed in
haste. In the Editions previous to the Benedictine, it was called an ``Epistle
to the Orthodox every where;'' but Montfaucon has been able to restore
the true title. He has been also able from his MSS. to make a far more
important correction, which has cleared up some very perplexing
difficulties in the history. All the Editions previous to the
Benedictine read ``George'' throughout for ``Gregory,'' and ``Gregory'' in
the place where ``Pistus'' occurs. Baronius, Tillemont, \&c. had
already made the alterations from the necessity of the case.]

———————

To his fellow-Ministers \footnote{Margin Notes

1. [\emph{sulleitourgois}].

Return to text} in every place, beloved Lords, Athanasius sends health in the
Lord.

§. 1.

1. O\textsc{ur} sufferings have been
dreadful beyond endurance, and it is impossible to describe them in
suitable terms; but in order that the dreadful nature of the
events which have taken place may be more readily apprehended, I have
thought it good to bring to your notice a history out of the
Scriptures. It happened that a certain Levite was injured in the
person of his wife [Judg. xix. 19.]; and, when he considered the
exceeding greatness of the pollution, (for the woman was a Hebrew, and
of the tribe of Judah,) being astounded at the outrage which had been
committed against him, he divided his wife's body, as the Holy
Scripture relates in the Book of Judges, and sent a part of it to
every tribe in Israel, in order that it might be understood that an
injury like this pertained not to himself only, but extended to all
alike; and that, if the people sympathised with him in his sufferings,
they might as avenge him; or if they neglected to do so, might bear
the disgrace of being considered henceforth as themselves guilty of
the wrong. The messenger's whom he sent related what had happened; and
they that heard and saw it, declared that such things had never been
done from the day that the children of Israel came up out of Egypt. So
every tribe of Israel was moved, and all came together against the
offenders, as though they had themselves had been the sufferers; and
at last the perpetrators of this iniquity were destroyed in war, and
became a curse \footnote{[\emph{anathema}].

Return to text} in the
mouths of all: for the assembled people considered not their kindred
blood, but regarded only the crime they had committed. You know the
history, brethren, and the particular account of the circumstances
given in Scripture. I will not therefore describe them more in detail,
since I write to persons acquainted with them, and as I am anxious to
represent to your piety our present circumstances, which are even
worse than those to which I have referred. For my object in reminding
you of this history is this, that you may compare those ancient
transactions with what has happened to us now, and perceiving how
little these last exceed the other in cruelty, may be filled with
greater indignation on account of them, than were the people of old
against those offenders.

2. For the treatment we have undergone, surpasses
the bitterness of any persecution; and the calamity of the Levite was
but small, when compared with the enormities which have now been
committed against the Church; or rather such deeds as these were
never before heard of in the whole world, or the like experienced by
any one. In that case it was but a single woman that was injured, and
one Levite who suffered wrong; now the whole Church is injured, the
priesthood insulted, and worst of all, piety \footnote{[\emph{eusebeia}], orthodoxy, vid. vol. viii. p. 1, note A.

Return to text} is persecuted by impiety. On that occasion the tribes were
astounded, each at the sight of part of the body of one woman; but now
the members of the whole Church are seen divided from one another, and
are sent abroad some to you, and some to others, bringing word of the
insults and injustice which they have suffered. Be ye therefore also
moved, I beseech you, considering that these wrongs are done unto you
no less than unto us; and let every one lend his aid, as feeling that
he is himself a sufferer, lest shortly the Ecclesiastical Canons, and
the faith of the Church be corrupted. For both are in danger, unless
God shall speedily by your hands amend what has been done amiss, and
the Church be avenged on her enemies. For our Canons \footnote{vid. Beveridg. Cod. Can. Illustr. i. 3. §. 2.
who comments on this passage at length. Allusion is also made to the
Canons in Apol. contr. Arian. §. 69.

Return to text} and our forms were not given to the Churches at the present
day, but were wisely and safely transmitted to us from our
forefathers. Neither had our faith its beginning at this time, but it
came down to us from the Lord through his disciples \footnote{vid. Athan. de Syn. §. 4. (Oxf. Tr. p. 78. and note O.) Orat. i. §.
8. (ibid. p. 191.) Tertull. Præscr. Hær. §. 29. (O. T. p. 462, and
note C.)

Return to text}. That therefore the ordinances which have been preserved in the
Churches from old time until now, may not be lost in our days, and the
trust which has been committed to us required at our hands; rouse
yourselves, brethren, as being stewards of the mysteries of God, and
seeing them now seized upon by aliens. Further particulars of our
condition you will learn from the bearers of our letters; but I was
anxious myself to write you a brief account thereof, that you may know
for certain, that such things have never before been committed against
the Church, from the day that our Saviour, when He was taken up, gave
command to his disciples, saying, Go ye, and make disciples of all
nations, baptizing them in the name of the Father, and of the Son, and
of the Holy Ghost [Mat. xxviii. 19.].

§. 2.

3. Now the outrages which have been committed
against us, and against the Church are these. While we were holding
our assemblies in peace, as usual, and while the people were rejoicing
in them, and advancing in godly conversation, and while our
fellow-ministers in Egypt, and the Thebais, and Libya, were in love
and peace both with one another and with us; on a sudden the Prefect
of Egypt puts forth a public letter, bearing the form of an edict, and
declaring that one Gregory from Cappadocia was coming to be my
successor, supported by his own body-guard. This announcement
confounded every one, for such a proceeding was entirely novel, and
now heard of for the first time. The people however assembled still
more constantly in the Churches \footnote{Assembling in the Churches seems to have been a sort of protest or
demonstration, sometimes peaceably, but sometimes in a less
exceptionable manner;—peaceably, during Justina's persecution at
Milan. Ambros. Ep. i. 20. August. Confess. ix. 15. but at Ephesus
after the third Ecumenical Council the Metropolitan shut up the
Churches, took possession of the Cathedral, and succeeded in repelling
the imperial troops. Churches were asylums, vid. Cod. Theodos. ix. 45.
§. 4. \&c. at the same time arms were prohibited.

Return to text}, for they very well knew that neither they themselves, nor any
Bishop or Presbyter, nor in short any one had ever complained against
me; and they saw that Arians only were on his side, and were aware
also that he was himself an Arian, and was sent by the Eusebians to
the Arian party. For you know, brethren, that the Eusebians have
always been the supporters and associates of the impious heresy of the
Arian fanatics \footnote{[\emph{areiomanit\={o}n}]. Vid. Ath. Oxf. Tr. viii. p. 91, note Q.

Return to text}, by
whose means they have ever carried on their designs against me, and
were the authors of my banishment into Gaul.

4. The people, therefore, were justly indignant
and exclaimed against the proceeding, calling the rest of the
magistrates and the whole city to witness, that this novel and
iniquitous attempt was now made against the Church, not on the ground
of any charge brought against me by Ecclesiastical persons, but
through the wanton assault of the Arian heretics. For even if there
had been any complaint generally prevailing against me, it was not an
Arian, or one professing Arian doctrines, that ought to have been
chosen to supersede me but according to the Ecclesiastical Canons, and
the direction of Paul, when the people were \emph{gathered together, and
the} \emph{spirit} of them that ordain, \emph{with the power of
our Lord Jesus Christ}, all things ought to have been enquired into
and transacted canonically, in the presence of those among the laity
and clergy who demanded the change; and not that a person brought from
a distance by Arians, as if making a traffic of the title of Bishop,
should with the support and strong arm of heathen magistrates, thrust
himself upon those who neither demanded nor desired his presence, nor
indeed knew any thing of what had been done. Such proceedings tend to
the dissolution of all Ecclesiastical rules, and compel the heathen to
blaspheme, and to suspect that our appointments are not made according
to a divine rule, but as a matter of traffic and patronage \footnote{O. T. viii. p. 190, note C.

Return to text}.

§. 3.

5. Thus was this notable appointment of Gregory
brought about by the Arians, and such was the beginning of it. And
what outrages he committed on his entry into Alexandria, and of what
great evils that event was the cause, you may learn both from our
letters, and by enquiry of those who travel among you. While the
people were offended at such an unusual proceeding, and in consequence
assembled in the Churches, in order to prevent the impiety of the
Arians from mingling itself with the faith of the Church, Philagrius
who has long been a persecutor of the Church and her virgins, and is
now Prefect \footnote{The Prefect of Egypt was called Augustalis as having been first
appointed by Augustus, after his victories over Antony. He was of the
Equestrian, not, as other Prefects, of the Senatorian order. He was
the imperial officer, as answering to Proprætors in the Imperial
Provinces. vid. Hofman. in voc.

Return to text} of Egypt, an
apostate already, and a fellow-countryman of Gregory, a man too of no
respectable character, and moreover supported by the Eusebians, and
therefore full of zeal against the Church; this person, by means of
promises which he afterwards fulfilled, succeeded in gaining over the
heathen multitude, with the Jews and disorderly persons, and having
excited their passions, sent them in a body with swords and clubs into
the Churches to attack the people.

6. What followed upon this it is by no means easy
to describe: indeed it is not possible to set before you a just
representation of the circumstances, nor even could one recount a
small part of them without tears and lamentations. Have such deeds as
these ever been made the subjects of tragedy among the ancients?
or has the like ever happened before in time of persecution or of war?
The Church and the holy Baptistery were set on fire, and straightway
groans, shrieks, and lamentations, were heard through the city; while
the citizens in their indignation at these enormities, cried shame
upon the governor, and protested against the violence used to them.
For the holy and undefiled virgins \footnote{The sister of S. Antony was one of the earliest known inmates at a
nunnery, vit. Ant. §. 2. 3. They were called by the Catholic Church
by the title, ``Spouse of Christ.'' Apol. ad Const. §. 33.

Return to text} were stripped naked, and suffered treatment which is not to be
named, and if they resisted, they were in danger of their lives. Monks
were trampled under foot and perished; some were hurled headlong;
others were destroyed with swords and clubs; others were wounded and
beaten. And oh! what deeds of impiety and iniquity were committed upon
the Holy Table! They offered birds and pine cones \footnote{The [\emph{thuos}] or suffitus of Grecian sacrifices generally
consisted of portions of odoriferous trees. vid. Potter. Antiqu. ii.
4. Some translate the word here used, ([\emph{strobilous}],) ``shell-fish.''

Return to text} in sacrifice, singing the praises of their idols, and
blaspheming even in the very Churches our Lord and Saviour Jesus
Christ, the Son of the living God. They burned the books of Holy
Scripture which they found in the Church; and the Jews, the murderers
of our Lord, and the godless heathen entering irreverently (O strange
boldness!) the holy Baptistery, stripped themselves naked, and acted
such a disgraceful part, both by word and deed, as one is ashamed even
to relate. Certain impious men also, following the examples set them
in the bitterest persecutions, seized upon the virgins, and widows,
and having tied their hands together, dragged them along, and
endeavoured to make them blaspheme and deny the Lord; and when they
refused to do so, they beat them violently and trampled them under
foot.

§. 4.

7. In addition to all this, after such a notable
and illustrious entry into the city, the Arian Gregory, taking
pleasure in these calamities, and as if desirous to secure to the
heathens and Jews, and those who had wrought these evils upon us, a
prize and price of their iniquitous success, gave up the Church to be
plundered by them. Upon this licence of iniquity and disorder, their
deeds were worse than in time of war, and more cruel than those of
robbers. Some of them plundered whatever fell in their way; others
divided among themselves the sums which individuals had laid up there
\footnote{Churches, as heathen temples before them, were used for deposits. At
the sack of Rome, Alaric spared the Churches and their possessions;
nay, he himself transported the costly vessels of St. Peter into his
Church.

Return to text}; the wine, of which there
was a large quantity, they either drank or emptied out or carried
away; they plundered the store of oil, and every one took as his spoil
the doors and chancel rails; the candlesticks they forthwith laid
aside in the wall \footnote{[\emph{en toi toichioi}. (?) \emph{toicharchoi}] for deacons. Apost.
Const. ii. 57. Clement. p. 615. from idea of navis or nave.

Return to text}, and
lighted the candles of the Church before their idols; in a word,
rapine and death pervaded the Church.

8. And the impious \footnote{[\emph{dussebeis}].

Return to text} Arians, so far from feeling shame that such things should be
done, added yet further outrages and cruelty. Presbyters and laymen
had their flesh torn, virgins were stripped of their veils \footnote{[\emph{apomaphorizomenai}].

Return to text}, and led away to the tribunal of the governor, and then cast
into prison; others had their goods confiscated, and were scourged;
the bread of the ministers and virgins was intercepted. And these
things were done even during the holy season of Lent \footnote{Lent and Passion Week was the season during which Justina's
persecution of St. Ambrose took place, and the proceedings against St.
Chrysostom at Constantinople. On the Paschal Vigils, vid. Tertull. ad
Uxor. ii. 4. p. 426, note N. Oxf. Tr.

Return to text}, about the time of Easter; a time when the brethren were
keeping fast, while this notable person Gregory exhibited the
disposition of a Caiaphas, and, together with that Pilate the
Governor, furiously raged against the pious worshippers of Christ.
Going into one of the Churches on the Preparation \footnote{[\emph{paraskeu\={e}}], i.e. Good Friday. The word was used for
Friday generally as early as S. Clem. Alex. Strom. vii. p. 877. ed.
Pott. vid. Constit. Apostol. v. 13. Pseudo-Ign. ad Philipp.13.

Return to text

}, in company with the Governor and the heathen multitude, when
he saw that the people regarded with abhorrence his forcible entry
among them, he caused that most cruel person, the Governor, publicly
to scourge in one hour, four and thirty virgins and married women, and
men of rank, and to cast them into prison. Among whom there was one
virgin, who, being fond of reading, had the Psalter in her hands, at
the time when he caused her to be publicly scourged: the book was
seized by the officers, and the virgin herself shut up in prison.

§. 5.

9. When all this was done, they did not stop even
here; but consulted how they might act the same part in the other Church, where I principally abode during those days; and they were
eager to extend their fury \footnote{[\emph{manian}].

Return to text}
to this Church also, in order that they might hunt out and dispatch
me. And this would have been my fate, had not the grace of Christ
assisted me, if it were only that I might escape to relate these few
particulars concerning their conduct. For seeing that they were
exceedingly mad against me, and being anxious that the Church should
not be injured, nor the virgins that were in it suffer, nor additional
murders be committed, nor the people again outraged, I withdrew myself
from among them, remembering the words of our Saviour, \emph{If they
persecute you in this city, flee ye into another} [Mat. x. 23.]. I
judged from the mischief they had done to one Church, that there was
no outrage they would forbear to perpetrate against the other,
especially since they had not reverenced even the Lord's day \footnote{Easter Day.

Return to text} on this holy Festival, but on that day when our Lord delivered
all men from the bonds of death, they had shut up in prison the people
of his Church; and Gregory and his associates, as if fighting against
our Saviour, and depending upon the support of the Governor, had
turned into mourning this day of liberty to the servants of Christ.
The heathens were rejoiced to do this, for they abhor that day; and
Gregory perhaps did but fulfil the commands of the Eusebians, when he
forced the Christians to mourn under the infliction of bonds.

10. With these acts of violence has the Governor
seized upon the Churches, and has given them up to Gregory and the
Arian fanatics. Thus, those persons who were excommunicated by us for
their impiety, now glory in the plunder of our Churches; while the
people of God, and the Clergy of the Catholic Church are compelled
either to have communion with the impiety of the Arian heretics, or
else to forbear entering into them. Moreover, by means of the
Governor, Gregory has exercised no small violence towards the captains
of ships and others who pass over sea, torturing and scourging some,
putting others in bonds, and casting them into prison, in order to
oblige them not to resist his iniquities, and to convey letters \footnote{i.e. letters of communion.

Return to text} from him. And not satisfied with all this, that he may glut
himself with my blood, he has caused his savage associate the
Governor, to prefer an indictment against me, as in the name of
the people, before the most religious Emperor Constantius, which
contains such odious charges, that if they were true, I ought not only
to be banished, but should deserve to suffer a thousand deaths. The
person who drew it up is an apostate from Christianity, and a
shameless worshipper of idols, and they who subscribed it are
heathens, and keepers of idol temples, and others of them Arians. In
short, not to make my letter tedious to you, a persecution rages here,
and such a persecution as was never before raised against the Church.
For in former instances a man at least might pray while he fled from
his persecutors, and be baptized while he lay in concealment. But now
their extreme cruelty has imitated the godless conduct of the
Babylonians. For as they falsely accused Daniel, so does the notable
Gregory now accuse before the Governor those who pray in their houses,
and watches every opportunity to insult their ministers, so that
through his violent conduct, the souls of many are endangered from
missing baptism, and many who are in sickness and sorrow have no one
to visit them, a calamity which they bitterly lament, accounting it
worse than their sickness. For while the ministers of the Chinch are
under persecution, the people who condemn the impiety of the Arian
heretics choose rather thus to be sick and to run the risk, than that
a hand of the Arians should come upon their heads.

§. 6.

11. Gregory then is an Arian, and has been sent
to the Arian party; for none demanded him, but they only; and
accordingly as a hireling and a stranger, he makes use of the Governor
to inflict these dreadful and cruel deeds upon the people of the
Catholic Churches, as not being his own. For since Pistus, whom the
Eusebians formerly appointed over the Arians, was justly anathematized
and excommunicated for his impiety by you the Bishops of the Catholic
Church, as you all know, on our writing to you concerning him, they
have now, therefore, in like manner sent this Gregory to them; and
lest they should a second time be put to shame, by our again writing
against them, they have employed foreign force against me, in order
that, having obtained possession of the Churches, they may seem to
have escaped all suspicion of being Arians. But in this too they have
been mistaken, for none of the people of the Church are with
them, except the heretics only, and those who have been excommunicated
for their crimes, and such as have been compelled by the Governor to
dissemble.

12. This then is the plot of the Eusebians, which
they have long been devising and bringing to bear; and now have
succeeded in accomplishing through the false charges which they have
made against me before the Emperor. Notwithstanding, they are not yet
content to be quiet, but even now seek to kill me; and they make
themselves so formidable to my friends, that they are all driven into
banishment, and expect death at their hands. But you must not for this
stand in awe of their iniquity, but on the contrary avenge: and shew
your indignation at this their unprecedented conduct against me. For
if when one member suffers all the members suffer with it, and,
according to the blessed Apostle, we ought to weep with them that
weep, let every one, now that so great a Church as this is suffering,
avenge its wrongs, as though he were himself the sufferer. For we have
a common Saviour, who is blasphemed by them, and Canons belonging to
us all, which they are transgressing. If while any of you had been
sitting in your Church, and while the people were assembled with you,
without any blame, some one had suddenly come under plea of an edict
to be your successor, and had acted the same part towards you, would
you not have been indignant? would you not have demanded to be
righted? If so, then it is right that you should be indignant now,
lest if these things be passed over unnoticed, the same mischief shall
by degrees extend itself to every Church, and so our schools of
religion be turned into a market-house and an exchange.

§. 7.

13. You are acquainted with the history of the
Arian fanatics, beloved, for you have often, both individually and in
a body, condemned their impiety; and you know also that the Eusebians,
as I said before, are engaged in the same heresy; for the sake of
which they have long been carrying on a conspiracy against me. And I
have represented to you, what has now been done, both for them and by
them, with greater cruelty than is usual even in time of war, in order
that after the example set before you in the history which I related at the beginning, you may entertain a zealous hatred of their
wickedness, and reject those who have committed such enormities
against the Church. If the brethren at Rome last year, before these
things had happened, and on account of their former misdeeds, wrote
letters to call a Council, that these evils might be set right,
(fearing which, the Eusebians took care previously to throw the Church
into confusion, and desired to destroy me, in order that they might
thenceforth be able to act as they pleased without fear, and might
have no one to call them to account;) how much more ought you now to
be indignant at these outrages, and to condemn them, seeing they have
added this to their former misconduct.

14. I beseech you, overlook not such proceedings,
nor suffer the famous Church of the Alexandrians to be trodden down by
heretics. In consequence of these things the people and their
ministers are separated from one another, as one might expect,
silenced by the violence of the Prefect, yet abhorring the impiety of
the Arian fanatics. If therefore Gregory shall write unto you, or any
other in his behalf, receive not his letters, brethren, but tear them
in pieces and put the bearers of them to shame, as the ministers of
impiety and wickedness. And even if he presume to write to you after a
friendly fashion, nevertheless receive them not. Those who bring his
letters convey them only from fear of the Governor, and on account of
his frequent acts of violence. And since it is probable that the
Eusebians will write to you concerning him, I was anxious to admonish
you beforehand, so that you may herein imitate God, who is no
respecter of persons, and may drive out from before you those that
come from them because for the sake of the Arian fanatics they caused
persecutions, rape of virgins, murders, plunder of the Church's
property, burnings and blasphemies in the Churches, to be committed by
the heathens and Jews at such a season. The impious and mad Gregory
cannot deny that he is an Arian, being proved to be so by the person
who writes his letters. This is his secretary Ammon, who was cast out
of the Church long ago by my predecessor the blessed Alexander for his
many crimes and for his impiety.

15. For all these reasons, therefore, vouchsafe
to send me a reply, and condemn these impious men; so that even
now the ministers and people of this place, seeing your orthodoxy and
hatred of wickedness, may rejoice in your concord in the Christian
faith, and that those who have been guilty of these lawless deeds
against the Church may be reformed by your letters, and brought at
last, though late, to repentance. Salute the brotherhood that is among
you. All the brethren that are with me salute you. Fare ye well, and
remember me, and the Lord preserve you continually, most truly beloved
Lords.

\xchapter{II. An Apology of our Holy Father Athanasius, Archbishop of Alexandria, against the Arians}

[The following Apology, or Defence of his
conduct, was written by S. Athanasius between A.D. 349-352, after his
return from his second exile upon the Council of Sardica. It is
scarcely more than a collection of exculpatory documents, which might
serve as a record of his innocence. These documents extend from A.D.
300, to A.D. 350, of which those between 340 and 350, are placed
first. ``This Apology,'' says Montfaucon, ``is the most authentic source
of the history of the Church in the first half of the fourth century.
Athanasius is far superior to any other historians of the period, both
from his bearing for the most part a personal testimony to the facts
he relates, and from his great accuracy and use of actual documents.
On the other hand, Ruffinus, Socrates, Sozomen, Theodoret, must not be
used without extreme caution, unless they adduce documents, which is
seldom the case.'' He proceeds to give instances; for this reason it
will not be worth while in this work, nor was it in the foregoing, to
compare Athanasius's statements with those of other historians, or to
use the latter except in connecting the line of the narrative. The
charges which he notices are as follow:—that he had been
clandestinely consecrated; that he had imposed a duty on Egyptian
linen; that he had assisted Philumenus with money, when in rebellion
against the Emperor; that he had sanctioned the overthrow of a
Communion Table and breaking of one of the Communion Vessels; that he
had killed a Meletian Bishop named Arsenius; that he had been the
cause of many executions or murders after his return from Gaul; that
he had sold for his own benefit the corn bestowed by Constantine on
the widows of the Church, and that he had stopped the supplies of corn
intended for Constantinople.]

———————

\xsection{Introduction}

1. I supposed that, after so many proofs of my
innocence had been given, my enemies would have shrunk from further
enquiry, and would now have condemned themselves for their false
accusations of others. But as they are not yet abashed, though they
have been so clearly convicted, but, as insensible to shame, persist
in their slanderous reports against me, professing to think that the
whole matter ought to be tried over again, (not that they may have
judgment passed on them, for that they avoid, but in order to harass
me, and to disturb the minds of the simple;) I therefore thought it
necessary to make my defence unto you, that you may listen to their
murmurings no longer, but may denounce their wickedness and base
calumnies. And it is only to you, who are men of sincere minds, that I
offer a defence; as for the contentious, I appeal confidently to the
decisive proofs which I have against them. For my cause needs not a
second judgment; which has already been given, and not once or twice
only, but many times. First of all, it was tried in my own country in
an assembly of nearly one hundred of its Bishops \footnote{The Council of Sardica says eighty; which is a
usual number in Egyptian Councils. (vid. Tillemont, vol. 8. p. 74.)
There were about ninety Bishops in Egypt, the Thebais, and Libya. The
present Council was held in 339, or 340. Its Synodal Epistle is
contained below, §. 3. and is particularly addressed to Pope Julius,
§. 20.

Return to text} a second time at Rome, when, in consequence of Letters from
Eusebius, both they and we were summoned, and more than fifty Bishops
met \footnote{This was held in 341. Julius's Letter is found below, §. 21.

Return to text}; and a third time in
the great Council assembled at Sardica \footnote{In A.D.
347, though Marsi, contrary to other writers, maintains its date to be
344. vid. §. 44. infr.

Return to text} by order of the most religious Emperors Constantius and
Constans, when my enemies were degraded as false accusers, and the
sentence that was passed in my favour received the suffrages of more
than three hundred Bishops, out of the provinces of Egypt, Libya, and
Pentapolis, Palestine, Arabia, Isauria, Cyprus, Pamphylia, Lycia,
Galatia, Dacia, Mysia, Thrace, Dardania, Macedonia, Epirus, Thessaly,
Achaia, Crete, Dalmatia, Siscia, Pannonia, Noricum, Italy, Picenum,
Tuscany, Campania, Calabria, Apulia, Bruttia, Sicily, the whole of
Africa, Sardinia, Spain, Gaul, and Britain.

2. Added to these was the testimony \footnote{vid. infr. §. 58. This was A.D. 349.

Return to text} of Ursacius and Valens, who had formerly calumniated me, but
afterwards changed their minds, and not only gave their assent to the
sentence that was passed in my favour, but also confessed that
they themselves and the rest of my enemies were false accusers; for
men who make such a change and such a recantation of course reflect
upon the Eusebians, for with them they had contrived the plot against
me. Now after a matter has been examined and decided on such clear
evidence by so many eminent Bishops, every one will confess that
further discussion is unnecessary; else, if an investigation be
instituted at this time, it may be again discussed and again
investigated, and there will be no end of this trifling.

§. 2.

3. Now the decision of so many Bishops was
sufficient to confound those who would still fain pretend some charge
against me. But when my enemies also bear testimony in my favour and
against themselves, declaring that the proceedings against me were a
conspiracy, who is there that would not be ashamed to doubt any
longer? The law requires that in the mouth of two or three witnesses
judgments shall be settled, and we have here this great multitude of
witnesses in my favour, with the addition of the proofs afforded by my
enemies; so much so that those who still continue opposed to me no
longer attach any importance to their own arbitrary \footnote{[\emph{h\={o}s \={e}thel\={e}san}]. vid. infr. §. 14.de Decr.
§. 3. de Syn. §. 13. ad Ep. Ag §. 5.

Return to text} judgment, but now have recourse to violence, and in the place
of fair reasoning seek to injure \footnote{This implies that Valens and Ursacius were subjected to some kind of
persecution, which is natural. They relapsed in 351, when Constantius
on the death of Constans came into possession of his brother's
dominions; and professed to have been forced to their former
recantation by the latter Emperor.

Return to text} those by whom they were exposed. For this is the chief cause of
vexation to them, that the measures they carried on in secret,
contrived by themselves in a corner, have been brought to light and
disclosed by Valens and Ursacius; for they are well aware that their
recantation not only clears those whom they have injured, but condemns
them.

4. Indeed this led to their degradation in the
Council of Sardica, as mentioned before; and with good reason; for, as
the Pharisees of old, when they undertook the defence of Paul, gave
clear judgment against the conspiracy which they and the Jews had
formed against him; and as the blessed David was proved to be
persecuted unjustly when the persecutor confessed, \emph{I have
sinned, my son David} [1 Sam. xxvi. 21.]; so it was with these men;
being overcome by the truth they became suppliants, and addressed a
letter to that effect to Julius Bishop of Rome. They wrote also to me
desiring to be on terms of peace with me, though they have spread such
reports concerning me; and probably even now they are covered with
shame, on seeing that those whom they sought to destroy by the grace
of the Lord are still alive. Consistently also with this conduct they
anathematized Arius and his heresy; for knowing that the Eusebians had
conspired against me in behalf of their own misbelief, and of nothing
else, as soon as they had determined to confess their calumnies
against me, they immediately renounced also that antichristian heresy
for the sake of which they had falsely asserted them.

\xsection{Chapter 1. Encyclical Letter of the Council of Egypt}

1. T\textsc{he} following are the
letters written in my favour by the Bishops in the several Councils;
and first the letter of the Egyptian Bishops.

The holy Council assembled at Alexandria, out of Egypt, the Thebais,
Libya, and Pentapolis, to the Bishops of the Catholic Church every
where, brethren beloved and greatly longed for, sendeth health in the
Lord.

§. 3.

Dearly beloved brethren, we might have put forth
a defence of our brother Athanasius \footnote{Margin Notes

1. [\emph{sulleitourgou}].

Return to text}, as respects the conspiracy of the Eusebians against him, and
complained of his sufferings at their hands, and have exposed all
their false charges, either at the beginning of their conspiracy or
upon his arrival at Alexandria. But circumstances did not permit it
then, as you also know; and lately, after the return of the Bishop
Athanasius, we thought that they would be confounded and covered with
shame at their manifest injustice: in consequence we prevailed with
ourselves to remain silent. Since, however, after all his severe
sufferings, after his retirement into Gaul, after his sojourn in a
foreign and far distant country in the place of his own, after his
narrow escape from death through their calumnies, but for the clemency
of the Emperor,—distress which would have satisfied even the most
cruel enemy,—still they are insensible to shame, and are again
raging against the Church and Athanasius; and from indignation at his
deliverance venture on still more atrocious schemes against him, and
are ready with any accusation, fearless of the words in holy
Scripture, \emph{A false witness shall not be unpunished} [Prov. xix.
5.]; and, \emph{The mouth that} \emph{belieth slayeth the soul} [Wisd.
i. 11.]; we therefore are unable longer to hold our peace, being
amazed at their wickedness and at the insatiable love of contention
displayed in their treacherous proceedings.

2. For see, they cease not to disturb the ear of
royalty with fresh reports against us; they cease not to write letters
of deadly import, for the destruction of the Bishop who is the enemy
of their impiety. For again have they written to the Emperors against
him; again are they conspiring against him, charging him with a
butchery which has never taken place; again they wish to shed his
blood, accusing him of a murder that never was committed, (for at that
former time would they have murdered him by their calumnies, had we
not found favour with the Emperor;) again they are urgent, to say the
least, that he should be sent into banishment, while they pretend to
lament the miseries of the exiles, as though they had been exiled by
him. They lament before us things that have never been done, and, not
satisfied with what has been done to him, desire to add thereto other
and more cruel treatment.

3. So mild are they and merciful, and of so just
a disposition; or rather (for the truth shall be spoken) so wicked are
they and malicious; obtaining respect through fear and by threats,
rather than by their piety and justice, as becomes Bishops. They have
dared in their letters to the Emperors to pour forth language such as
no contentious person would employ even among those that are without;
they have charged him with a number of murders and butcheries, and
that not before a Governor, or any other superior officer, but before
the three Augusti; nor shrink they from any journey however long,
provided only all the greater courts may be filled with their
accusations. For indeed, dearly beloved, their business consists in
accusations, and that of the most solemn character, forasmuch as the
tribunals to which they make their appeal are the most solemn of any
upon earth. And what other end do they propose by these
investigations, except to move the Emperor to capital punishment?

§. 4.

4. Their own conduct therefore, and not that of
Athanasius, is the fittest subject for lamentation and mourning, and
one would more properly lament them, for such actions ought to be
bewailed, since it is written, \emph{Weep ye not for the dead, neither
bemoan him: but weep sore for him that goeth away, for he shall return
no more} [Jer. xxii. 10.]. For their whole letter speaks of nothing
but his death; and their endeavour is to kill, whenever they may be
permitted, or if not, to drive into exile. And this they were
permitted to do by the most religious father of the Emperors, who
gratified their fury by the banishment of Athanasius, though not by
his death. Now that this is not the conduct even of ordinary
Christians, (nay, even of heathens,) much less of Bishops, who profess
to teach others righteousness, we suppose that your Christian
consciences must at once perceive. How can they forbid others to
accuse their brethren, who themselves become their accusers, and that
to the Emperors? How can they teach compassion for the misfortunes of
others, who cannot rest satisfied even with our banishment? For there
was confessedly a general sentence of banishment against us Bishops,
and we all looked upon ourselves as banished men: and now again we
consider ourselves as restored with Athanasius to our native country,
and in the place of our former lamentations and mourning over him, as
having the greatest encouragement and grace,—which may the Lord
continue to us, nor suffer the Eusebians to destroy!

5. Even if their charges against him were true,
here is a certain charge against them, that against the precept of
Christianity, and after his banishment and trials, they have assaulted
him again, and accuse him of murder, and butchery, and other crimes,
which they sound in the royal ears against the Bishops. But how
exceeding manifold is their wickedness, and what manner of men think
you them, when every word they speak is false, every charge they bring
a calumny, and there is no truth whatever either in their speeches or
their writings! However, let us now enter upon these matters, and meet
their last charges. This will prove, that in their former
representations in the Council and at the trial their conduct was
dishonourable, or rather their words untrue, besides exposing them for
what they have now advanced.

§. 5.

6. We are indeed ashamed to make any defence
against such charges. But since our reckless accusers lay hold of any
charge, and allege that murders and butcheries were committed
after the return of Athanasius, we beseech you to bear with our answer
though it be somewhat long; for circumstances constrain us. No murder
was committed either by Athanasius or on his account, since our
accusers, as we said before, compel us to enter upon this strange
apology. Slaughter and imprisonment are foreign to our Church. No one
did Athanasius commit into the hands of the executioner; and the
prison, so far as he was concerned, was never disturbed. Our
sanctuaries are now, as they have always been, pure, and honoured only
with the Blood of Christ and his pious worship. Neither Presbyter nor
Deacon was destroyed by Athanasius; he perpetrated no murder, he
caused the banishment of no one. Would that they had never caused the
like to him, nor given him actual experience of it! No one here was
banished on his account; no one at all except Athanasius himself the
Bishop of Alexandria, whom they banished, and whom, now that he is
restored, they again seek to entangle in the same or even a more cruel
plot than before, setting their tongues to speak all manner of false
and deadly words against him.

7. For, behold, they now attribute to him the
acts of the magistrates; and although they plainly confess in their
letter that the Prefect of Egypt passed sentence upon certain persons,
they now are not ashamed to impute this sentence to Athanasius; and
that, though he had not at the time entered Alexandria, but was yet on
his return from his place of exile. Indeed he was then in Syria; since
we must needs adduce in his defence his length of way from home, that
a man may not be responsible for the actions of a Governor or Prefect
of Egypt. But supposing Athanasius had been in Alexandria, what were
the proceedings of the Prefect to Athanasius? However, he was not even
in the country; and what the Prefect of Egypt did was not done on
ecclesiastical grounds, but for reasons which you will learn from the
records, which, after we understood what they had written, we made
diligent enquiry for, and have transmitted to you. Since then they now
raise a cry against certain things which were never done either by him
or for him, as through they had certainly taken place, and testify
against such evils as though they were assured of their existence; let
them inform us from what Council they obtained their knowledge of
them, from what proofs, and in the course of what investigation? But
if they have no such evidence to bring forward, and nothing but their
own mere assertion, we leave it to you to consider as regards their
former charges also, how the things took place, and why they so speak
of them. In truth, it is nothing but calumny, and a plot of our
enemies, and anger full of atrocious projects, and an impiety in
behalf of the Arian fanatics \footnote{or Ario-maniacs, \emph{passim}.

Return to text},
which is frantic against true godliness, and desires to root out the
orthodox, so that henceforth the advocates of impiety may preach
without fear whatever doctrines they please. The history of the matter
is as follows: —

§. 6.

8. When Arius, from whom the heresy of the Arian
fanatics has its name, was cast out of the Church for his impiety by
Bishop Alexander, of blessed memory, the Eusebians, who are the
disciples and partners of his impiety, considering themselves also to
have been ejected, wrote frequently to the Bishop Alexander,
beseeching him not to keep the heretic Arius out of the Church. But
when Alexander in his piety towards Christ refused to admit that
impious man, they directed their resentment against Athanasius, who
was then a Deacon, because in their busy enquiries they had heard that
he was much in the familiarity of Alexander the Bishop, and much
honoured by him. And their hatred of him was greatly increased after
they had experience of his piety \footnote{i.e. orthodoxy, \emph{passim}.

Return to text} towards Christ, in the Council assembled at Nicæa, wherein he
spoke boldly against the impiety of the Arian fanatics. But when God
raised him to the Episcopate, their long-cherished malice burst forth
into a flame, and fearing his orthodoxy and resistance of their
impiety, they (and especially Eusebius, who was smitten with a
consciousness of his own evil doings,) engaged in all manner of
treacherous designs against him. They prejudiced the Emperor against
him; they frequently threatened him with Councils; and at last
assembled at Tyre; and to this day they cease not to write against
him, and are so implacable that they even find fault with his
appointment to the Episcopate \footnote{The Eusebians alleged that, fifty-four Bishops
of the two parties of S. Alexander and Meletius being assembled for
the election, and having sworn to elect by the common voice, six or
seven of these broke their oaths in favour of S. Athanasius, whom no
one had thought of, and consecrated him in secret to the great
surprise and scandal of both ecclesiastical and lay persons. vid. Socr.
ii. 17. Philostorgius (A.D. 423.) adds particulars, explanatory or
corrective of this statement, of which the Bishops in the text do not
seem to have heard; sic, that Athanasius with his party one night
seized on the Church of St. Dionysius, and compelled two Bishops whom
he found there to consecrate him against their will; that he was in
consequence anathematized by all the other Bishops, but that,
fortifying himself in his position, he sent in his election to the
Emperor, and by this means obtained its confirmation. Hist. ii. 16. It
appears, in matter of fact, that S. Athan. was absent at the time of
his election; as Socrates says, in order to avoid it, or as Epiphanius,
on business at the Court; these reasons are compatible.

Return to text},
taking every means of shewing their enmity and hatred towards
him, and spreading false reports for the sole purpose of thereby
destroying his character.

9. However, the very misrepresentations which
they now are making, do but convict their former statements of being
falsehoods, and a mere conspiracy against him. For they say, that ``after
the death of the Bishop Alexander, a certain few having mentioned the
name of Athanasius, six or seven Bishops elected him clandestinely in
a secret place:'' and this is what they wrote to the Emperors, having
no scruple about asserting the greatest falsehoods. Now that the whole
multitude and all the people of the Catholic Church assembled together
as with one mind and body, and cried, shouted, that Athanasius should
be Bishop of their Church, made this the subject of their public
prayers to Christ, and conjured us to grant it for many days and
nights, neither departing themselves from the Church, nor suffering us
to do so; of all this we are witnesses, and so is the whole city, and
the province too. Not a word did they speak against him, as these
persons represented, but gave him the most excellent titles they could
devise, calling him the good, the pious, Christian, an ascetic \footnote{It is contested whether S. Athan. was ever one of S. Antony's monks,
the reading of a passage in the commencement of his Vit. Ant., which
would decide the question, varying in different MSS. The word ``ascetic''
is used of those who lived a life, as afterwards followed in
Monasteries, in the Ante-Nicene times.

Return to text}, a genuine Bishop. And that he was elected by a majority of our
body in the sight and with the acclamations of all the people, we who
elected him also testify, who are surely more credible witnesses than
those who were not present, and now spread these false accounts.

10. But yet Eusebius \footnote{Of Nicomedia.

Return to text} finds fault with the appointment of Athanasius,—he who
perhaps never received any appointment to his office at all; or
if he did, has himself rendered it invalid \footnote{The Canons of Nicæa and Sardica were absolute against translation,
but, as Bingham observes, Antiqu. vi. 4. §. 6. only as a general
rule. The so-called Apostolical Canons except ``a reasonable cause'' and
the sanction of a Council; one of the Councils of Carthage prohibit
them when subserving ambitious views, and except for the advantage of
the Church. Vid. list of translations in Socr. Hist. vii. 36.
Cassiodor. Hist. xii. 8. Niceph. Hist. xiv. 39. Cotelier adds others
ad Can. Apost. 14.

Return to text}. For he had first the See of Berytus, but leaving that he came
to Nicomedia. He left the one contrary to the law, and contrary to the
law invaded the other; he deserted his own See for he failed in
affection, and took possession of another's though he failed in a
plea; he lost his love for the first in his lust for another, nor
retained that love for the second which his lust had occasioned. For,
behold, withdrawing himself from the second, again he takes possession
of another's, casting an evil eye all around him upon the cities of
other men, and thinking that godliness \footnote{[\emph{eusebeian}], orthodoxy, (vid. 1 Tim. vi. 5.)

Return to text} consists in wealth and in the greatness of cities, and making
light of the heritage of God to which he had been appointed; not
knowing that \emph{where} even \emph{two or three are gathered in the
name} of the Lord, \emph{there} is the Lord \emph{in the midst of them};
not considering the words of the Apostle, \emph{I will not boast in
another man's labours}; not perceiving the charge which he has
given, \emph{Art thou bound unto a wife? seek not to be loosed} [Mat.
xviii. 20. 2 Cor. x. 15. 1 Cor. vii. 27.]. For if this expression
applies to a wife, how much more does it apply to a Church, and to the
same Episcopate; to which whosoever is bound ought not to seek
another, lest he prove an adulterer according to holy Scripture.

§. 7.

11. But though conscious of these his own
misdoings, he has boldly undertaken to arraign the appointment of
Athanasius, to which honourable testimony has been borne by all; and
he ventures to reproach him with his deposition, though he has been
deposed himself, and has a standing proof of his deposition in the
appointment of another. How could either he or Theognius \footnote{Or Theognis; he was, as well as Eusebius, a pupil of Lucian's, and was
deposed together with him after the Nicene Council for communicating
with Arians. Constantine banished them to Gaul; they were recalled in
the course of two or three years. He was dead by the date of the
Council of Sardica.

Return to text} degrade another, after they had been degraded themselves, which
is sufficiently proved by the appointment of others in their room? For
you know very well that there were appointed instead of them
Amphion to Nicomedia and Chrestus to Nicæa, in consequence of their
own impiety and connection with the Arian fanatics, who were rejected
by the Ecumenic Council. But while they desire to set aside that true
Council, they endeavour to give that name to their own unlawful
combination \footnote{Eusebian Council of Tyre, A.D. 335.

Return to text}; while they
are unwilling that the decrees of the Council should be enforced, they
desire to enforce their own decisions; and they use the name of a
Council, while they refuse to submit themselves to one so great as
this. Thus they care not for Councils, but only pretend to do so in
order that they may root out the orthodox, and annul the decrees of
the true and great Council against the Arians, in support of whom,
both now and heretofore, they have ventured to assert these falsehoods
against the Bishop Athanasius. For their former statements resembled
those they have now made, viz. that disorderly meetings were held at
his entrance \footnote{On his return from Gaul, A.D. 338.

Return to text}, with
lamentation and mourning, the people indignantly refusing to receive
him. Now such was not the case, but, on the other hand, joy and
cheerfulness prevailed, and the people ran together, hastening to
obtain the desired sight of him. The Churches were full of rejoicings,
and thanksgivings were offered up to the Lord every where; and all the
Ministers and Clergy beheld him with such feelings, that their souls
were possessed with delight, and they esteemed that the happiest day
of their lives. Why need we mention the inexpressible joy that
prevailed among us Bishops, for we have already said that we counted
ourselves to have been partakers in his sufferings?

§. 8.

12. Now this being confessedly the truth of the
matter, although it is very differently represented by them, what
weight can be attached to that Council or trial of which they make
their boast? Since they presume thus to controvert the circumstances
of a case which they did not witness, which they have not examined,
and for which they did not meet, and to write as though they were
assured of the truth of their statements, how can they claim credit
respecting those matters for the consideration of which they say that
they did meet together? Will it not rather be believed that they have
acted both in the one case and in the other out of enmity to us?
For what kind of a Council of Bishops was then held? Was it an
assembly which aimed at the truth? Was not almost every one among them
our enemy? Did not the attack of the Eusebians upon us proceed from
their zeal for the Arian fanaticism? Did they not urge on the others
of their party? Have we not always written against them as professing
the doctrines of Arius? Was not Eusebius of Cæsarea in Palestine
accused by our confessors of sacrificing to idols \footnote{At the Council of Tyre, Potamo an Egyptian Bishop and Confessor asked
Eusebius what had happened to \emph{him} in prison during the
persecution, Epiph. Hær. 68, 7. as if hinting at his cowardice. It
appears that Eusebius was prisoner at Cæsarea with S. Pamphilus; yet
he never mentions the fact himself, which is unlike him, if it was
producible.

Return to text}? Was not George proved to have been degraded by the blessed
Alexander \footnote{George, Bishop of Laodicea, had been degraded when a Priest by S.
Alexander, for his profligate habits as well as his Arianism. Athan.
speaks of him elsewhere as reprobated even by his party. de Fug. 26.

Return to text}? Were not they
charged with various offences, some with this, some with that?

13. How then could such men entertain the purpose
of holding a meeting against us? How can they have the boldness to
call that a Council, at which a single Count presided, which an
executioner attended, and where a chief jailor instead of the Deacons
of the Church introduced us into Court; and where the Count only
spoke, and all present held their peace, or rather obeyed his
directions? The removal of those Bishops who seemed to deserve it, was
prevented at his desire; and when he gave the order we were dragged
about by soldiers;—or rather the Eusebians gave the order, and he
was subservient to their will. In short, dearly beloved, what kind of
Council was that, the object of which was banishment and murder at the
pleasure of the Emperor? And of what nature were their charges?—for
here is matter of still greater astonishment. There was one Arsenius
whom they declared to have been murdered; and they also complained
that a chalice belonging to the sacred mysteries had been broken.

14. Now Arsenius is alive, and prays to be
admitted to our communion. He waits for no other testimony to prove
that he is still living, but himself confesses it, writing in his own
person to our brother Athanasius, whom they positively asserted
to be his murderer. The impious wretches were not ashamed to accuse
him of having murdered a man who was at a great distance from him,
being separated by an immense tract both of land and water, and whose
abode at that time no one knew. Nay, they even had the boldness to
remove him out of sight, and place him in concealment, though he had
suffered no injury; and, if it had been possible, they would have
transported him to another world, nay, or have taken him from life in
earnest, so that either by a true or false statement of his murder
they might in as good earnest destroy Athanasius. But thanks to divine
Providence for this also, which permitted them not to succeed in their
injustice, but presented Arsenius alive to the eyes of all men, who
has clearly proved their conspiracy and calumnies. He does not
withdraw from us as murderers, nor hate us as having injured him, (for
indeed he has suffered no evil at all;) but he desires to hold
communion with us; he wishes to be admitted among us, and has written
to this effect.

§. 9.

15. Nevertheless they laid their plot against
Athanasius, accusing him of having murdered a person who was still
alive; and those same men are the authors of his banishment \footnote{by Constantine into Gaul, A.D. 335.

Return to text}. For it was not the father of the Emperors, but their
calumnies, that sent him into exile. Consider whether this is not the
truth. When nothing was discovered to the prejudice of our brother
Athanasius, but still the Count threatened him with violence, and was
very zealous against him, the Bishop \footnote{The circumstances of this appeal, which are related by Athan. below,
§. 86. are thus summed up by Gibbon; ``Before the final sentence could
he pronounced at Tyre, the intrepid primate threw himself into a bark
which was ready to hoist sail for the imperial city. The request of a
formal audience might have been opposed or eluded; but Athanasius
concealed his arrival, watched the moment of Constantine's return from
an adjacent villa, and boldly encountered his angry sovereign as he
passed on horseback through the principal street of Constantinople. So
strange an apparition excited his surprise and indignation; and the
guards were ordered to remove the importunate suitor; but his
resentment was subdued by involuntary respect; and the haughty spirit
of the Emperor was awed by the courage and eloquence of a Bishop, who
implored his justice and awakened his conscience.'' Hist. xxi. Athan.
was a small man in person.

Return to text}, in order to avoid this violence, went up \footnote{i.e. to Constantinople.

Return to text} to the most religious Emperor, where he protested against the
Count and their conspiracy against him, and requested either that a
lawful Council of Bishops might be assembled, or that the Emperor
would himself receive his defence concerning the charges they brought
against him. Upon this the Emperor wrote in anger, summoning them
before him, and declaring that he would hear the cause himself, and
for that purpose he also ordered a Council to be held. Whereupon the
Eusebians went up and charged Athanasius, not with the same offences
which they had published against him at Tyre, but with an intention of
detaining the vessels laden with corn, as though Athanasius had been
the man to pretend that he could stop the exports of corn from
Alexandria to Constantinople.

16. Certain of our friends were present at the
palace with Athanasius, and heard the threats of the Emperor upon
receiving this report. And when Athanasius exclaimed against the
calumny, and positively declared that it was not true; (for how, he
argued, should he a poor man, and in a private station, be able to do
such a thing?) Eusebius did not hesitate publicly to repeat the
charge, and swore that Athanasius was a rich man, and powerful, and
able to do any thing; from which it might be supposed that he had used
this language. Such was the accusation these venerable Bishops
proffered against him. But the grace of God proved superior to their
wickedness, for it moved the pious Emperor to mercy, who instead of
death passed upon him the sentence of banishment. Thus their
calumnies, and nothing else, were the cause of this. For the Emperor,
in the letter which he previously wrote, complained of their
conspiracy, censured their machinations, and condemned the Meletians
as unrighteous and deserving of execration; in short, expressed
himself in the severest terms concerning them. For he was greatly
moved when he heard the story of the dead alive; he was moved at
hearing of this murder of one who lived after it without loss of life.
We have sent you the letter.

§. 10.

17. But these marvellous Eusebians, to make a
show of refuting the truth of the case, and the statements contained
in this letter, put forward the name of a Council, and ground its
proceedings upon the authority of the Emperor. Hence the attendance of
a Count at their meeting, and the soldiers as guards of the Bishops,
and royal letters compelling the attendance of any persons whom
they required. But observe here the strange character of their
machinations, and the inconsistency of their bold measures, so that by
some means or other they may take Athanasius away from us. For if as
Bishops they claimed for themselves alone the judgment of the case,
what need was there for the attendance of a Count and soldiers? or how
was it that they assembled under the sanction of royal letters? Or if
they required the Emperor's countenance and wished to derive their
authority from him, why did they then entrench upon his judgment? and
when he declared in the letter which he wrote, that the Meletians were
profligate calumniators, and that Athanasius was most innocent, and
enlarged upon the pretended murder of the living, how was it that they
determined that the Meletians had spoken the truth, and that
Athanasius was guilty of the offence; and were not ashamed to make the
living dead, living both after the Emperor's judgment, and at the time
when they met together, and who even until this day is amongst us? So
much concerning the case of Arsenius.

§. 11.

18. And as for the chalice belonging to the
mysteries, what was it, or where was it broken by Macarius? for this
is the report which they spread up and down. But for Athanasius, even
his accusers would not have ventured to blame him, had they not been
suborned by them. However, they attribute the origin of the offence to
him; although it ought not to be imputed even to Macarius who is clear
of it. And they are not ashamed to parade the sacred mysteries before
Catechumens, and worse than that, even before heathens \footnote{This period, when Christianity was acknowledged by the state but not
embraced by the population, is just the time when we hear most of this
Reserve as a principle. While Christians were but a sect, persecution
enforced a discipline, and when they were commensurate with the
nation, faith made it unnecessary. We are now returned to the state of
the fourth century.

Return to text}: whereas, they ought to attend to what is written, \emph{It is
good to keep close the secret of a king}[Tob. xii. 7.]; and as the
Lord has charged us, \emph{Give not that which is holy unto the dogs,
neither cast ye your pearls before swine} [Matt. vii. 6.]. We ought
not then to parade the holy mysteries before the uninitiated, lest the
heathen in their ignorance deride them, and the Catechumens being
over-curious be offended. However, what was the chalice, and where and
before whom was it broken? It is the Meletians who make the
accusation, who are not worthy of the least credit, for they have been
schismatics and enemies of the Church, not of a recent date, but from
the times of the blessed Peter, Bishop and Martyr \footnote{Meletius, Bishop of Lycopolis in the Thebaid, being deposed for
lapsing in the Dioclesian Persecution, separated from the Catholic
Church and commenced a succession of his own in Egypt. In the same
persecution S. Peter suffered.

Return to text}. They formed a conspiracy against Peter himself; they
calumniated his successor Achillas; they accused Alexander even before
the Emperor; and being thus well versed in these arts, they have now
transferred their enmity to Athanasius, acting altogether in
accordance with their former wickedness. For as they slandered those
that have been before him, so now they have slandered him. But their
calumnies and false accusations have never prevailed against him until
now, that they have got the Eusebians for their assistants and
patrons, on account of the impiety \footnote{i.e. heresy, \emph{passim}.

Return to text} which these have adopted from the Arian fanatics, which has
led them to conspire against many Bishops, and among the rest
Athanasius.

19. Now the place where the say the chalice was
broken, was not a Church; there was no Presbyter in occupation of the
place; and the day on which they say that Macarius did the deed, was
not the Lord's day. Since then there was no Church there; since there
was no one to perform the priest's office; and since the day did not
require the use of it \footnote{This seems to imply that the Holy Communion was only celebrated on
Sundays in the Egyptian Churches.

Return to text};
what was this sacred chalice, and when, or where was it broken? There
are many cups, it is plain, both in private houses, and in the public
market; and if a person breaks one of them, he is not guilty of
impiety. But the chalice which belongs to the mysteries, and which if
it be broken intentionally, makes the perpetrator of the deed an
impious person, is found only among those who are lawfully appointed
to preside over the Church. This is the only description that can be
given of this kind of chalice; there is none other; of this you drink
prior to the people; this you have received according to the canon of
the Church \footnote{vid. Can. Ap. 65.

Return to text}; this
belongs only to those who preside over the Catholic Church, for
to you only it appertains to have the first taste \footnote{[\emph{propinein}].

Return to text} of the Blood of Christ, and to none besides. But as he who
breaks a sacred cup is an impious person, much more impious is he who
treats the Blood of Christ with contumely: and he does so who performs
this mystical rite contrary to the rule of the Church;—(we say this,
not as if a chalice even of the schismatics was broken by Macarius,
for there was no chalice there at all; how should there be? where
there was neither Lord's house nor any one belonging to the Church,
nay, it was not the time of the celebration of the mysteries;)—now
such a person is the notorious Ischyras, who was never appointed to
his office by the Church, and when Alexander admitted the Presbyters
that had been ordained by Meletius, he was not even numbered amongst
them; and therefore did not receive ordination even from that quarter.

§. 12.

20. By what means then did Ischyras become a
Presbyter \footnote{Vid. Bp. Taylor, Episcop. Assert. §. 32. Potter on Church Gov. ch. v.

Return to text}? who was it
that ordained him? was it Colluthus? for this is the only supposition
that remains. But it is well known, and no one has any doubt about the
matter, that Colluthus died a Presbyter, and that every ordination of
his was invalid, and that all that were ordained by him during the
schism were reduced to the condition of laymen, and in that rank
appear in the congregation. How then can it be believed that a private
person, occupying a private house, had in his possession a sacred
chalice? But the truth is, they gave the name of Presbyter at the time
to a private person, and gratified him with this title to support him
in his iniquitous conduct towards us; and now as the reward of his
accusations they procure for him the erection of a Church. So that
this man had then no Church; but as the reward of his malice and
subserviency to them in accusing us, he receives now what he had not
before; nay, perhaps they have even remunerated his services with the
Episcopate, for so he goes about reporting, and accordingly behaves
towards us with great insolence. Thus are such rewards as these now
bestowed by Bishops upon accusers and calumniators; though indeed it
is reasonable, in the case of an accomplice, that as they have
made him a partner in their proceedings, so they should also make him
their associate in their own Episcopate. But this is not all; give ear
yet further to their proceedings at that time.

§. 13.

21. Being unable to prevail against the truth,
though they had thus set themselves in array against it, and Ischyras
having proved nothing at Tyre, except that he was a calumniator, and
the calumny ruining their plot, they defer proceedings until they
obtain fresh evidence, and propose to send to the Mareotis certain of
their party to enquire diligently into the matter. Accordingly they
dispatched secretly, with the assistance of the civil power, persons
to whom we openly objected on many accounts, as being of the party of
Arius, and therefore our enemies; namely, Diognius, Maris, Theodorus,
Macedonius, and two others, young both in years and mind \footnote{Vid. also Athan. ad Ep. Æg. 7. Euseb. Vit. c. iv. 43. Hilar. ad
Const. i. 5. Fragm. ii. 12.

Return to text}, Ursacius and Valens from Pannonia; who, after they had
undertaken this long journey for the purpose of obtaining justice
against their enemy, set out again from Tyre for Alexandria. They did
not shrink from becoming witnesses themselves, although they were the
judges, but openly adopted every means of furthering their design, and
undertook any labour or journey whatsoever in order to bring to a
successful issue the conspiracy which was in progress. They left the
Bishop Athanasius detained in a foreign country while they themselves
entered their enemy's city, as if to have their revel both against his
Church and against his people. And what was more outrageous still,
they took with them the accuser Ischyras, but would not permit
Macarius, the accused person, to accompany them, but left him in
custody at Tyre. For ``Macarius the Presbyter of Alexandria'' was made
answerable for the charge far and near.

§. 14.

22. They therefore entered Alexandria alone with
the accuser, their partner in lodging, board, and wine-cup; and taking
with them Philagrius the Prefect of Egypt they proceeded to the
Mareotis, and there carried on the investigation by themselves, all
their own way, with the forementioned person. Although the Presbyters
frequently begged that they might be present, they would not
permit them. The Presbyters both of the city and of the whole country
desired to attend, that they might detect who and whence the persons
were who were suborned by Ischyras. But they forbade the Ministers to
be present, while they carried on the examination concerning the
Church, the chalice, the table, and the holy things, before the
heathen; nay, worse than that, they summoned heathen witnesses during
the enquiry concerning the sacred chalice; and those persons who they
affirmed were taken out of the way by Athanasius by means of the
summons of the Receiver-general, and they knew not where in the world
they were, these same individuals they brought forward before
themselves and the Prefect only, and avowedly used their testimony,
whom they affirmed without shame to have been secreted by the Bishop
Athanasius.

23. But here too their only object is to effect
his death, and so they again pretend that persons are dead who are
still alive, following the same method they adopted in the case of
Arsenius. For the men are living, and are to be seen in their own
country; but to you who are at a great distance from the spot they
give a tragical representation of the matter as though they had
disappeared, in order that, as the evidence is so far removed from
you, they may falsely accuse our brother-minister, as though he used
violence and the civil power; whereas they themselves have in all
respects acted by means of that power and the countenance of others.
For their proceedings in the Mareotis were parallel to those at Tyre;
and as there a Count attended with military assistance, and would
permit nothing either to be said or done contrary to their pleasure,
so here also the Prefect of Egypt was present with a band of men,
frightening all the members of the Church, and permitting no one to
give true testimony. And what was the strangest thing of all, the
persons who came, whether as judges or witnesses, or, what was more
likely, in order to serve their own purposes and those of Eusebius,
lived in the same place with the accuser, even in his house, and there
seemed to carry on the investigation as they pleased.

§. 15.

24. We suppose you are not ignorant what outrages
they committed at Alexandria; for they are reported every where. They attacked the holy virgins and brethren with naked swords; they
beat with scourges their persons, esteemed honourable in the sight of
God, so that their feet were lamed by the stripes, whose souls are
whole and sound in purity and all good works \footnote{Hist. Arian. 12.

Return to text}. The trades \footnote{[\emph{ergasiai}] (?)

Return to text}
were excited against them; and the heathen multitude was set to strip
them naked, to beat them, wantonly to insult them, and to threaten
them with their altars and sacrifices. And one coarse fellow, as
though license had now been given them by the Prefect in order to
gratify the Bishops, took hold of a virgin by the hand, and dragged
her towards an altar that happened to be near, imitating the practice
of compelling to offer sacrifice in time of persecution. When this was
done, the virgins took to flight, and a shout of laughter was raised
by the heathen against the Church; the Bishops being in the place, and
occupying the very house where this was going on; and from which, in
order to obtain favour with them, the virgins were assaulted with
naked swords, and were exposed to all kinds of danger, and insult, and
wanton violence. And this treatment they received during a season of
fasting \footnote{supr. p. 7.

Return to text}, and at the
hands of persons who themselves were feasting with the Bishops in that
house.

§. 16.

25. Foreseeing these things, and reflecting that
the entrance of enemies into a place is no ordinary calamity, we
protested against this commission. And Alexander \footnote{This Alexander had been one of the Nicene Fathers, in 325, and had the
office of publishing their decrees in Macedonia, Greece, \&c. He
was at the Council of Jerusalem ten years after, at which the Church
of the Holy Sepulchre was consecrated, and afterwards Arius admitted
to communion. His influence with the Court party seems to have been
great, judging from Count Dionysius's tone in speaking of him, infr.
§. 81.

Return to text}, Bishop of Thessalonica, considering the same, wrote to the
people residing there, discovering the conspiracy, and testifying of
the plot. They indeed reckon him to be one of themselves, and account
him a partner in their designs; but they only prove thereby the
violence they have exercised towards him. For even the profligate
Ischyras himself was only induced by fear and violence to proceed in
the matter, and was obliged by force to undertake the accusation. As a
proof of this, he wrote himself to our brother Athanasius \footnote{infr. §. 64.

Return to text}, confessing that nothing of the kind that was alleged had
taken place there, but that he was suborned to make a false
statement. This declaration he made, though he was never admitted by
Athanasius as a Presbyter, nor received that title from him as a boon,
nor was entrusted by way of recompense with the erection of a Church,
nor expected the bribe of a Bishopric; all of which he obtained from
them in return for undertaking the accusation. Moreover, his whole
family held communion with us \footnote{vid. infr. §. 63 fin. §. 85 fin.

Return to text}, which they would not have done had they been injured in the
slightest degree.

§. 17.

26. Now to prove that these things are facts and
not mere assertions, we have the testimony \footnote{infr. §. 74.

Return to text} of all the Presbyters of the Mareotis \footnote{The district, called Mareotis from a neighbouring lake, lay in the
territory and diocese of Alexandria, to the west. It consisted of
various large villages, with handsome Churches, and resident Priests,
and of hamlets which had none; of the latter was ``the Peace of
Secontaruri,'' (infr. §. 85.) where Ischyras lived.

Return to text}, who always accompany the Bishop in his visitations, and who
also wrote at the time against Ischyras. But neither those of them who
came to Tyre were allowed to declare the truth \footnote{infr. §. 79.

Return to text}, nor could those who remained in the Mareotis obtain
permission to refute the calumnies of Ischyras \footnote{§. 72 fin.

Return to text}. Copies also of the letters of Alexander, and of the
Presbyters, and of Ischyras, will prove the same thing. We have sent
also the letter of the father of the Emperors, in which he expresses
his indignation that the murder of Arsenius was charged upon any one
while the man was still alive; as also his astonishment at the
variable and inconsistent character of their accusations with respect
to the chalice; since at one time they accused the Presbyter Macarius,
at another the Bishop Athanasius, of having broken it with his hands.
He declares also on the one hand that the Meletians are calumniators,
and on the other that Athanasius is perfectly innocent.

27. And are not the Meletians calumuniators, and
above all John \footnote{Arcaph. infr. 65 fin. head of the Melitians.

Return to text}, who
after coming into the Church, and communicating with us, after
condemning himself, and no longer taking any part in the proceedings
respecting the chalice, when he saw the Eusebians zealously supporting
the Arian fanatics, though they had not the daring to cooperate with
them openly, but were attempting to employ others as their masks,
undertook a character, as an actor in the heathen theatres \footnote{vid. infr. §. 37. 46. vol. 8. p. 127, note G.

Return to text}? The subject of the drama was the contest of the Arians;
the real design of the piece being their success, but John and his
partizans being appended and playing the parts, in order that under
colour of these, the supporters of the Arians, in the garb of judges,
might drive away the enemies of their impiety, firmly establish their
impious doctrines, and bring the Arians into the Church. And those who
wish to drive out true godliness \footnote{[\emph{eusebeia}], \&c. vid. supr. p. 3. ref. 1.

Return to text} strive all they can to prevail by ungodliness ; they
who have chosen the part of that impiety \footnote{[\emph{eusebeia}], \&c. vid. supr. p. 3. ref. 1.

Return to text}
which wars against Christ, endeavour to destroy the enemies thereof,
as though they were impious \footnote{[\emph{eusebeia}], \&c. vid. supr. p. 3. ref. 1.

Return to text} persons;
and they impute to us the breaking of the chalice, for the purpose of
making it appear that Athanasius, equally with themselves, is guilty
of impiety \footnote{[\emph{eusebeia}], \&c. vid. supr. p. 3. ref. 1.

Return to text} towards Christ.

28. For what means this mention of the sacred
chalice by them? Whence comes this religious \footnote{[\emph{eusebeia}], \&c. vid. supr. p. 3. ref. 1.

Return to text} regard for the chalice among those who support impiety \footnote{[\emph{eusebeia}], \&c. vid. supr. p. 3. ref. 1.

Return to text} towards Christ? Whence comes it that Christ's chalice is known
to them who know not Christ? How can they who profess to honour that
chalice, dishonour the God of the chalice? or how can they who lament
over the chalice, seek to murder the Bishop who celebrates the
mysteries therewith? for they would have murdered him, had it been in
their power. And how can they who lament the loss of the throne that
was Episcopally covered \footnote{cathedræ velatæ, Aust. ap. Bingh. viii. 6. §. 10.

Return to text},
seek to destroy the Bishop that sat upon it, to the end that both the
throne may be without its Bishop, and that the people may be deprived
of godly \footnote{[\emph{eusebeia}], \&c. vid. supr. p. 3. ref. 1.

Return to text} doctrine? It was not then
the chalice, nor the murder, nor any of those portentous deeds they
talk about, that induced them to act thus; but the forementioned
heresy of the Arians, for the sake of which they conspired against
Athanasius and other Bishops, and still continue to wage war against
the Church.

29. Who are they that have really been the cause
of murders and banishments? Are not these? Who are they that, availing
themselves of external support, conspire against the Bishops? Are not
the Eusebians they, and not Athanasius, as in their letters they
pretend? Both he and others have suffered at their hands. Even at the
time of which we speak, four Presbyters \footnote{vid. their names infr. §. 40.

Return to text} of Alexandria, though they had not even proceeded to Tyre,
were banished by their means. Who then are they whose conduct calls
for tears and lamentations? Does not theirs, who after they have
been guilty of one course of persecution, do not scruple to add to it
a second, but have recourse to all manner of falsehood, in order that
they may destroy a Bishop who will not give way to their impious
heresy? hence arises the enmity of the Eusebians; hence their
proceedings at Tyre; hence their pretended trials; hence also now the
letters which they have written even without any trial, expressing the
utmost confidence in their statements; hence their calumnies before
the father of the Emperors, and before the most religious Emperors
themselves.

§. 18.

30. For it is necessary that you should know what
is now reported to the prejudice of our brother Athanasius, in order
that you may thereby be led to condemn their wickedness, and may
perceive that they desire nothing else but to murder him. A quantity
of corn was given by the father of the Emperors for the support of
certain widows, some to be of Libya, and some out of Egypt. They have
all received it up to this time, Athanasius getting nothing therefrom,
but the trouble of assisting them. But now, although the recipients
themselves make no complaint, but acknowledge that they have received
it, Athanasius has been accused of selling all the corn, and
appropriating the profits to his own use: and the Emperor wrote to
this effect about it, charging him with the offence in consequence of
the calumnies which had been raised against him. Now who are they that
have raised these calumnies? Is it not those who after they have been
guilty of one course of persecution, scruple not to set on foot
another? Who are the authors of those letters which are said to have
come from the Emperor? Are not the Arians who are so zealous against
Athanasius, and scruple not to speak and write any thing? No one would
pass over persons who have acted as they have done, in order to
entertain suspicion of others. Nay, the proof of their calumny appears
to be most evident, for they are anxious under cover of it, to take
away the corn from the Church, and to give it to the Arians. And this
circumstance more than any other, brings the matter home to the
authors of this design and their principals, who scrupled neither to
set on foot a charge of murder against Athanasius, and as a base means
of prejudicing the Emperor against him, nor yet to take away from
the Clergy \footnote{[\emph{t\={o}n kl\={e}r\={o}n}], f. [\emph{ch\={e}r\={o}n}],
Montf. §. 19.

Return to text} of the
Church the subsistence of the poor, in order that in fact they might
make gain for the heretics.

§. 19.

31. We have sent also the testimony of our
brother ministers in Libya, Pentapolis, and Egypt, from which likewise
you may learn the false accusations which have been brought against
Athanasius. And these things they do, in order that, the professors of
true godliness being henceforth induced by fear to remain quiet, the
heresy of the impious Arians may be introduced in place of the truth.
But thanks be to your piety, dearly beloved, that you have frequently
anathematized the Arians in your letters, and have never given them
admittance into the Church. The exposure of the Eusebians is also
easy, and ready at hand. For behold, after their former letters
concerning the Arians, of which also we have sent you copies, they now
openly stir up the Arian fanatics against the Church, though the whole
Catholic Church has anathematized them; they have appointed a Bishop \footnote{Pistus.

Return to text} over them; they distract the Churches with threats and alarms,
that they may gain assistants in their impiety in every part.
Moreover, they send Deacons to the Arians, who openly join their
assemblies; they write letters to them, and receive answers from them,
thus making schisms in the Church, and holding communion with them;
and they send to every part, commending their heresy, and repudiating
the Church, as you will perceive from the letters they have addressed
to the Bishop of Rome \footnote{vid. infr. §. 21.

Return to text},
and perhaps to yourselves also. You perceive therefore, dearly
beloved, that these things are not undeserving of vengeance: they are
indeed dreadful and alien from the doctrine of Christ.

32. Wherefore we have assembled together, and
have written to you, to request of your Christian wisdom to receive
this our declaration and sympathize with our brother Athanasius, and
to shew your indignation against the Eusebians who have essayed such
things, in order that such malice and wickedness may no longer prevail
against the Church. We call upon you to be the avengers of such
injustice, reminding you of the injunction of the Apostle, \emph{Put away
from among yourselves that wicked person} [1 Cor. v. 13.]. Wicked
indeed is their conduct, and unworthy of your communion.
Wherefore give no further heed to them, though they should again write
to you against the Bishop Athanasius; (for all that proceeds from them
is false;) not even though they subscribe their letter with names \footnote{The Eusebians availed themselves of the subscriptions of the Meletians,
as at Philippopolis, Hilar. Fragm. 3.

Return to text

} of Egyptian Bishops. For it is evident that it will not be we
who write, but the Meletians \footnote{infr. §. 73.

Return to text}, who have ever been schismatics, and who even unto this day
make disturbances and raise factions in the Churches. For they ordain
improper persons, and all but heathens; and they are guilty of such
actions as we are ashamed to set down in writing, but which you may
learn from those whom we have sent unto you, and who will deliver to
you our letter.

§. 20.

33. Thus wrote the Bishops of Egypt to all
Bishops, and to Julius Bishop of Rome.

\xsection{Chapter 2. Letter of Pope Julius to the Eusebians at Antioch}

1. T\textsc{he} Eusebians also wrote
to Julius, and thinking to frighten me, requested him to call a
Council, and to be himself the judge, if he so pleased \footnote{Margin Notes

1. A.D.
340. vid. Hist. Arian. §. 11.

Return to text}. When therefore I went up to Rome, Julius wrote to the
Eusebians, as was suitable, and sent moreover two of his own
Presbyters \footnote{Vito and Vincentius, Presbyters, had
represented Silvester at Nicæa. Liberius sent Vincentius, Bishop, and
Marcellus, Bishop, to Constantius; and again Lucifer, Bishop, and
Eusebius, Bishop. St. Basil suggests that Damasus should send Legates
into the East, Ep. 69. The Council of Sardica, Can. 5. recognised the
Pope's power of sending Legates into foreign Provinces to hear certain
appeals; ``ut de \emph{Latere suo} Presbyterum mittat.'' vid. Thomassin.
de Eccl. Disc. Part 1. ii. 117.

Return to text}, Elpidius and
Philoxenus \footnote{May, A.D. 341.

Return to text}. But they,
when they heard of me, were thrown so into confusion, as not expecting
my going up thither; and they declined the proposed Council, alleging
unsatisfactory reasons for so doing, but in truth they were afraid
lest the things should be proved against them which Valens and
Ursacius afterwards confessed \footnote{infr. §. 58.

Return to text}. However, more than fifty Bishops assembled, in the place where
the Presbyter Vito held his congregation \footnote{[\emph{sun\={e}gen}].

Return to text}; and they acknowledged my defence, and gave me the confirmation
\footnote{vid. infr. p. 60, ref. 2.

Return to text} both of their
fellowship and their loving hospitality. On the other hand, they
expressed great indignation against the Eusebians, and requested that
Julius would write to the following effect to those of their number
who had written to him. Which accordingly he did, and sent it by the
hand of Count Gabianus.

2. The Letter of Julius \footnote{A.D.
342, but 341. Tillem. \&c.

Return to text}

Julius to his dearly beloved brethren \footnote{By Danius, which had been considered the same name as Dianæus, Bishop
of Cæsarea in Cappadocia, Montfaucon in loc. understands the
notorious Arian, Bishop of Nicæa, called variously Diognius, (supr.
§. 13.) Theognius, (infr. §. 28.) Theognis, (Philostorg. Hist. ii.
7.) Theogonius, (Theod. Hist. i. 19.) and assigns some ingenious and
probable reasons for his supposition. vid. supr. p. 23, note d.
Flacillus, Arian Bishop of Antioch, as Athan. names him, is called
Placillus,(in St. Jerome's Chronicon, p.785.) Placitus, (Soz. iii. 5,)
Flacitus, (Theod. Hist. i. 21.) Theodorus was Arian Bishop of Heraclea,
whose Comments on the Psalms are supposed to be those which bear his
name in Corderius's Catena.

Return to text}, Danius, Flacillus, Narcissus, Eusebius, Maris, Macedonius,
Theodorus, and their friends, who have written to me from
Antioch, sends health in the Lord.

§. 21.

I have read your letter \footnote{Some of the topics contained in the Eusebian Letter are specified in
Julius's answer. It acknowledged, besides, the high dignity of the See
of Rome, as being ``The School ([\emph{phrontist\={e}rion}]) of the
Apostles and the Metropolis of orthodoxy from the beginning,'' but
added that ``doctors came to it from the east; and that they ought not
themselves to hold the second place, for they were superior in virtue,
though not in their Church.'' And they said that they would hold
communion with Julius if he would agree to their depositions and
substitutions in the Eastern Sees. Soz. iii. 8.

Return to text} which was brought to me by my Presbyters Elpidius and
Philoxenus, and I am surprised to find that, whereas I wrote to you in
charity and with conscious sincerity, you have replied to me in an
unbecoming and quarrelsome temper; for the pride and arrogance of the
writers is plainly exhibited in that letter. Yet such feelings are
inconsistent with the Christian faith; for what was written in a
charitable spirit ought likewise to be answered in a spirit of charity
and not of contention. And was it not a token of charity to send
Presbyters to sympathize with them that are in suffering, and to
desire those who had written to me to come hither, that the questions
at issue might obtain a speedy settlement, and all things be duly
ordered, so that our brethren might no longer be exposed to suffering,
and that you might escape further imputation? But something seems to
show that your temper is such, as to force us to conclude that the
terms in which you appear to pay honour \footnote{[\emph{timan}].

Return to text} to us, are with some dissimulation modified in their meaning.
The Presbyters also whom we sent to you, and who ought to have
returned rejoicing, did on the contrary return sorrowful on account of
the proceedings they had witnessed among you. And I, when I had read
your letter, after much consideration, kept it to myself, thinking
that after all some of you would come, and there would be no need to
bring it forward, lest if it should be openly exhibited, it should
grieve many of our brethren here. But when no one arrived, and it
became necessary that the letter should be produced, I declare to you,
they were all astonished, and were hardly able to believe that such a
letter had been written by you at all; for it is expressed in terms of
strife rather than of charity.

3. Now if the author of it wrote with an ambition
of exhibiting his power of language, such a practice surely is more
suitable for other subjects: in ecclesiastical matters, it is not
a display of eloquence that is needed, but the observance of Apostolic
Canons, and an earnest care not to offend one of the little ones of
the Church. For it were better for a man, according to that
ecclesiastical sentence, that a millstone were hanged about his neck,
and that he were drowned in the sea, than that he should offend even
one of the little ones [vid. Matt. xviii. 6.]. But if such a letter
was written, because certain persons through a narrow feeling \footnote{[\emph{mikropsuchian}], jealousy. vid. Chrys. in Eph. O. T. pp. 119,
326, 331.

Return to text} took offence among themselves, (for I will not impute it to
all); it were better not to entertain any such feeling of offence at
all, at least not to let the sun go down upon their vexation; and
certainly not to give it room to exhibit itself in writing.

§. 22.

4. Yet what has been done that is a just cause of
offence? or in what respect was my letter to you such? Was it, that I
invited \footnote{[\emph{proetrepsametha}].

Return to text} you to be
present at a Council? You ought rather to have received the proposal
with joy. Those who have confidence in their proceedings, or as they
choose to term them, in their decisions, are not wont to be angry, if
such decision is enquired into by others; they rather shew all
boldness, seeing that if they have given a just decision, it can never
prove to be the reverse. The Bishops who assembled in the great
Council of Nicæa agreed, not without the will of God, that the
decisions of one Council should be examined in another \footnote{As this determination does not find a place among the now received
Canons of the Council, the passage in the text becomes of great moment
in the argument in favour of the twenty Canons extant in Greek being
but a portion of those passed at Nicæa. vid. Alber. Dissert. in Hist.
Eccles. vii. Abraham Ecchellensis has argued on the same side, (apud
Colet. Concil. t. ii. p. 399. Ed. Ven. 1728.) also Baronius, though
not so strongly, Ann. 325. nn. 157, \&c. and Montfaucon in loc.
Natalis Alexander, Sæc. 4. Dissert. 28. argues against the larger
number, and Tillemont, Mem. t. 6. p. 674.

Return to text}, to the end that the judges, having before their eyes that
other trial which was to follow, might be led to investigate matters
with the utmost caution, and that the parties concerned in their
sentence might have assurance that the judgment they received was
just, and not dictated by the enmity of their former judges. Now if
you are unwilling that such a practice should be adopted in your own
case, though it is of ancient standing, and has been noticed and
recommended by the great Council, your refusal is not becoming; for it
is unreasonable that a custom which has once obtained in the
Church, and been established by Councils, should be set aside by a few
individuals.

5. For a further reason they cannot justly take
offence in this point. When the persons whom you the Eusebians
dispatched with your letters, I mean Macarius the Presbyter, and
Martyrius and Hesychius the Deacons, arrived here, and found that they
were unable to withstand the arguments of the Presbyters who came from
Athanasius, but were confuted and exposed on all sides, they then
requested me to call a Council together \footnote{A.D.
340.

Return to text}, and to write to Alexandria to the Bishop Athanasius, and also
to the Eusebians, in order that a just judgment might be given in the
presence of all parties. And they undertook in that case to move all
the charges which had been brought against Athanasius. For Martyrius
and Hesychius had been publicly detected by us, and the Presbyters of
the Bishop Athanasius had withstood them with great confidence:
indeed, if one must tell the truth, the part of Martyrius had been
utterly overthrown; and thus it was that led them to desire that a
Council might be held. Now supposing that they had not desired a
Council, but that I had been the person to propose it, in
discouragement of \footnote{[\emph{skulai}].

Return to text}
those who had written to me, and for the sake of our brethren who
complain that they have suffered injustice; even in that case the
proposal would have been reasonable and just, for it is agreeable to
ecclesiastical practice, and well pleasing to God. But when those
persons, whom you the Eusebians considered to be trustworthy, when
even they wished me to call the brethren together, it was inconsistent
in the parties invited to take offence, when they ought rather to have
shewn all readiness to be present. These considerations shew that the
display of anger in the offended persons is unreasonable, and their
refusal to meet the Council is unbecoming, and has a suspicious
appearance. Does any one find fault, if he sees that done by another,
which he would allow if done by himself? If, as you write, the decrees
of any Council have an irreversible force, and he who has given
judgment on a matter is dishonoured, if his sentence is examined by
another; consider, dearly beloved, who are they that dishonour
Councils? who are setting aside the decisions of former judges?

6. Not to inquire at present into every
individual case, lest I should appear to press too heavily on
certain parties, the last instance that has occurred, and which every
one who hears it must shudder at, will be sufficient in proof of the
others which I omit. §. 23. The Arians who were excommunicated for
their impiety by Alexander, the late Bishop of Alexandria, of blessed
memory, were not only proscribed by the brethren in the several
cities, but were also anathematized by the whole body assembled
together in the great Council of Nicæa. For theirs was no ordinary
offence, neither had they sinned against man, but against our Lord
Jesus Christ Himself, the Son of the living God. And yet these persons
who were proscribed by the whole world, and branded in every Church,
are said now to have been admitted to communion again; which I think
you ought to hear with indignation. Who then are the parties who
dishonour Councils? Are not they who have set at nought the votes of
the Three hundred \footnote{The number of the Fathers at the Nicene Council is generally
considered to have been 318, the number of Abraham's servants, Gen.
xiv. 14. Anastasius (Hodeg. 3. fin.) referring to the first three
Ecumenical Councils, speaks of the faith of the 318, the 150, and the
200.

Return to text}, and
have preferred impiety to godliness?

7. The heresy of the Arian fanatics \footnote{[\emph{areiomanit\={o}n}], supr. p. 4. ref. 1.

Return to text} was condemned and proscribed by the whole body of Bishops
every where; but the Bishops Athanasius and Marcellus have many
supporters who speak and write in their behalf. We have received
testimony in favour of Marcellus, that he resisted the advocates of
the Arian doctrines in the Council of Nicæa; and in favour of
Athanasius, that at Tyre nothing was brought home to him, and that in
the Mareotis, where the Reports against him are said to have been
drawn up, he was not present. Now you know, dearly beloved, that \emph{ex
parte} proceedings are of no weight, but bear a suspicious
appearance. Nevertheless, these things being so, we, in order to be
accurate, and neither showing any prepossession in favour of
yourselves, nor of those who wrote in behalf of the other party,
invited those who had written to me to come hither; that, since there
were many who wrote in their behalf, all things might be enquired into
in a Council, and neither the guiltless might be condemned, nor the
guilty be accounted innocent. Who then are not the parties who
dishonour Councils, but they who at once and recklessly have received
the Arians whom all had condemned, and contrary to the decision
of the judges. The greater part of those judges have now departed, and
are with Christ; but some of them are still in this life of trial, and
are indignant at learning that certain persons have set aside their
judgment.

§. 24.

8. We have also been informed of the following
circumstance by those who were at Alexandria. A certain Carpones, who
had been excommunicated by Alexander for Arianism, was sent hither \footnote{A.D.
341.

Return to text} by one Gregory with certain others, also excommunicated for
the same heresy. However, I had learnt the matter also from the
Presbyter Macarius, and the Deacons Martyrius and Hesychius \footnote{A.D.
339.

Return to text}. For before the Presbyters of Athanasius arrived, they urged
me to send letters to one Pistus at Alexandria, though at the same
time the Bishop Athanasius was there. And when the Presbyters of the
Bishop Athanasius came, they informed me that this Pistus was an
Arian, and that he had been excommunicated \footnote{vid. infr. Depos. Ar.

Return to text} by the Bishop Alexander and the Council of Nicæa, and then
ordained by one Secundus, whom also the great Council excommunicated
as an Arian. This statement the party of Martyrius did not gainsay,
nor did they deny that Pistus had received his ordination from
Secundus. Now consider, after this who are most justly liable to
blame? I, who could not be prevailed upon to write to the Arian Pistus;
or those, who advised me to do dishonour to the great Council, and to
address the impious as if they were godly persons? Moreover, when the
Presbyter Macarius, who had been sent hither by Eusebius with
Martyrius and the rest, heard of the opposition which had been made by
the Presbyters of Athanasius, while we were expecting his appearance
with Martyrius and Hesychius, he decamped in the night, in spite of a
bodily ailment; which leads us to conjecture that his departure arose
from shame on account of the exposure which had been made concerning
Pistus. For it is impossible that the ordination of the Arian Secundus
should be considered valid \footnote{[\emph{ischusai}].

Return to text} in the Catholic Church. This would indeed be dishonour to the
Council, and to the Bishops who composed it, if the decrees they
framed, as in the presence of God, with such extreme earnestness and
care, should be set aside as nugatory.

§. 25.

9. If, as you write \footnote{vid. also Hilar. Fragm. iii. 26.

Return to text}, the decrees of all Councils ought to be of force,
according to the precedent in the case of Novatus \footnote{i.e. Novatian.

Return to text} \footnote{The instance of Novatian makes against the Eusebians, because for some
time after Novatian was condemned in the West, his cause was not
abandoned in the East. Tillemont, Mem. t. 7. p. 277.

Return to text} and Paul of
Samosata, certainly the sentence of the Three hundred ought not to be
reversed, certainly a Catholic Council ought not to be set at nought
by a few individuals. For the Arians are heretics as they, and the
like sentence has been passed both against the one and the other. And,
after such bold proceedings as these, who are they that have lighted
up the flame of discord? for in your letter you blame us for having
done this. Have we, who have sympathized with the sufferings of the
brethren, and have acted in all respects according to the Canon; or
they who contentiously and contrary to the Canon have set aside the
sentence of the Three hundred, and dishonoured the Council in every
way? For not only have the Arians been received into communion, but
Bishops also have adopted the practice of removing from one place to
another \footnote{vid. supr. p. 23.

Return to text}. Now if you
really believe that all Bishops have the same and equal authority \footnote{Cyprian. de Unit. Eccl. 4. O. T.

Return to text}, and you do not, as you assert, account of them according to
the magnitude of their cities; he that is entrusted with a small city
ought to abide in the place committed to him, and not from disdain of
his trust to remove to one that has never been put under him;
despising that which God has given him, and making much of the vain
applause of men. You ought then, dearly beloved, to have come and not
declined, that the matter may be brought to a conclusion; for this is
what reason demands.

10. But perhaps you were prevented by the time
fixed upon for the Council, for you complain in your letter that the
interval before the day we appointed \footnote{[\emph{prothesmia}], vid. Cyr. Cat. O. T. pp. 3, 246.

Return to text} was too short. But this, dearly beloved, is a mere excuse. Had
certain of you set out to come, and the day arrived before them, the
interval allowed would then have been proved to be too short. But when
persons do not wish to come, and detain even my Presbyters up to the
month of January \footnote{A.D.
342. Tillem. reads June.

Return to text}, it
is the mere excuse of those who have no confidence in their cause;
otherwise, as I said before, they would have come, not regarding the
length of the journey, not considering the shortness of the time,
but trusting to the justice and reasonableness of their cause. But
perhaps they did not come on account of the aspect of the times \footnote{the Persian war. Hist. Arian. §. 11.

Return to text}, for again you declare in your letter, that we ought to have
considered the present circumstances of the East, and not to have
desired you to come. Now if as you say you did not come because the
times were such, you ought to have considered such times beforehand,
and not to have become the authors of schism, and of mourning and
lamentation in the Churches. But as the matter stands, men, who have
been the cause of these things, shew that it is not the times that are
to blame, but the determination of those who will not meet a Council.

§. 26.

11. But I wonder also how you could ever have
written that part of your letter, in which you say, that I alone
wrote, and not to all of you, but to the Eusebians only. In this
complaint one may discover more of readiness to find fault than of
regard for truth. I received the letters against Athanasius from none
other than those connected with Martyrius and Hesychius, and I
necessarily wrote to them who had written against him. Either then the
Eusebians ought not alone to have written, apart from you all, or else
you, to whom I did not write, ought not to be offended that I wrote to
them who had written to me. If it was right that I should address my
letter to you all, you also ought to have written with them; but now,
considering what was reasonable, I wrote to them who had addressed
themselves to me, and had given me information. But if you were
displeased because I alone wrote to them, it is but consistent that
you should also be angry, because they wrote to me alone. But for this
also, dearly beloved, there was a fair and reasonable cause.
Nevertheless it is necessary that I should acquaint you that, although
I only wrote, yet the sentiments I expressed were not those of myself
alone, but of all the Bishops throughout Italy and in these parts. I
indeed was unwilling to cause them all to write, lest the others
should be overpowered by their number. The Bishops however assembled
on the appointed day, and agreed in these opinions, which I again
write to signify to you; so that, dearly beloved, although I alone
address you, yet you may be assured that these are the sentiments of
all. Thus much for the excuses, not reasonable, but unjust and
suspicious, which some of you have alleged for your conduct.

§. 27.

12. Now although what has already been said were
sufficient to shew that we have not admitted to our communion our
brothers Athanasius and Marcellus either too readily, or unjustly, yet
it is but fair briefly to set the matter before you. Eusebius's
friends wrote formerly against the friends of Athanasius, as you also
have written now; but a great number of Bishops out of Egypt and other
provinces wrote in his favour. Now in the first place, your letters
against him are inconsistent with one another, and the latter have no
sort of agreement with the former, but in many instances the former
are answered by the latter, and the latter are impeached by the
former. Now where there is this contradiction in letters, no credit
whatever is due to the statements they contain. In the next place, if
you require us to believe what you have written, it is but consistent
that we should not refuse credit to those who have written in his
favour \footnote{vid. supr. §. 3.-§. 19.

Return to text}; especially,
considering that you write from a distance, while they are on the
spot, are acquainted with the man, and the events which are occurring
there, and testify in writing to his manner of life, and positively
affirm that he has been the victim of a conspiracy throughout.

13. Again, a certain Bishop Arsenius was said at
one time to have been destroyed by Athanasius, but we have learned
that he is alive, nay, that he is on terms of friendship with him. He
has positively asserted that the Reports drawn up in the Mareotis were
\emph{ex parte} ones; for that neither the Presbyter Macarius, the
accused party, was present, nor yet his Bishop, Athanasius himself.
This we have learnt, not only from his own mouth, but also from the
Reports which Martyrius and Hesychius brought to us \footnote{vid. infr. §. 83 fin.

Return to text}; for we found on reading them, that the accuser Ischyras was
present there, but neither Macarius, nor the Bishop Athanasius; and
that the Presbyters of Athanasius desired to attend, but were not
permitted. Now, dearly beloved, if the trial was to be conducted
honestly, not only the accuser, but the accused also ought to have
been present. As the accused party Macarius attended at Tyre, as well
as the accuser Ischyras, when nothing was proved against him, so not
only ought the accuser to have gone to the Mareotis, but also the
accused, so that he might be present when he was convicted, or if he
was acquitted, might have opportunity to expose the calumny. But now,
as this was not the case, but the accuser only went out thither, with
those to whom Athanasius objected, the proceedings wear a suspicious
appearance.

§. 28.

14. And he complained also that the persons who
went to the Mareotis went against his wish, for that Theognius, Maris,
Theodorus, Ursacius, Valens, and Macedonius, who were the persons they
sent out, were of suspected character. This he shewed not by his own
assertion merely, but from a letter of Alexander who was Bishop of
Thessalonica; for he produced a letter written by him to Dionysius \footnote{infr. §. 80.

Return to text}, the Count who presided in the Council, in which he shews most
clearly that there was a conspiracy on foot against Athanasius. He has
also brought forward a genuine document, all in the handwriting of the
accuser Ischyras himself \footnote{§. 64.

Return to text},
in which he calls God Almighty to witness that no chalice was broken,
nor table overthrown, but that he had been suborned by certain persons
to invent these accusations. Moreover, when the Presbyters of the
Mareotis arrived \footnote{§. 74.

Return to text},
they positively affirmed that Ischyras was not a Presbyter of the
Catholic Church, and that Macarius had not committed any such offence
as the other had laid to his charge. The Presbyters and Deacons also
who came to us testified in the fullest manner in favour of the Bishop
Athanasius, strenuously asserting that none of those things which were
alleged against him were true, but that he was the victim of a
conspiracy.

15. And all time Bishops of Egypt and Libya wrote
and protested \footnote{supr. §. 6. p. 22.

Return to text} that
his ordination was lawful and strictly ecclesiastical, and that all
that you had advanced against him was false, for that no murder had
been committed, nor any persons despatched on his account, nor any
chalice broken, but that all was false. Nay, the Bishop Athanasius
also shewed from the \emph{ex parte} Reports drawn up in the Mareotis,
that a Catechumen was examined and said \footnote{infr. §. 83.

Return to text}, that he was within with Ischyras, at the time when they say
Macarius the Presbyter of Athanasius burst into the place; and that
others who were examined said,—one, that Ischyras was in a small
cell \footnote{[\emph{en kelli\={o}i}].

Return to text},—and another,
that he lay behind the door, being sick at that very time, when
they say Macarius came thither. Now from these representations of his,
we are naturally led to ask the question, How was it possible that a
man who was lying behind the door sick could get up, conduct the
service, and offer the Oblations? and how could it be that Oblations
were offered in the presence of Catechumens \footnote{Bingh. Ant. x. 5. §. 8.

Return to text}? for if there were Catechumens present, it was not yet the
time for presenting; the Oblations. These representations, as I said,
were made by the Bishop Athanasius, and he shewed from the Reports,
what was also positively affirmed by those who were with him, that
Ischyras has never been a Presbyter at all in the Catholic Church, nor
has ever appeared as a Presbyter in the assemblies of the Church; for
not even when Alexander admitted those of the Meletian schism, by the
indulgence of the great Council, was he named by Meletius among his
Presbyters, as they deposed \footnote{infr. §. 71.

Return to text}; which is the strongest argument possible that he was not even
a Presbyter of Meletius; for otherwise, he would certainly have been
numbered with the rest. Besides, it was shewn also by Athanasius from
the Reports, that Ischyras had spoken falsely in other instances: for
he set up a charge respecting the burning of certain books, when, as
they pretend, Macarius burst in upon them, but was convicted of
falsehood by the witnesses he himself brought to prove it.

§. 29.

16. Now when these things were thus represented
to us, and so many witnesses appeared in his favour, and so much was
advanced by him in his own justification, what did it become us to do?
what did the Canon \footnote{pp. 3, 45, 55.

Return to text}
of the Church require of us, but that we should not condemn him, but
rather receive him and treat him as a Bishop, as we have done?
Moreover, besides all this he continued here a year and six months \footnote{Valesius, Montfaucon, and Coustant, consider these eighteen months to
run from about May 341, upon Gregory's usurpation, to October or
November 342, when the Council of Rome terminated, as Schelstrate also
thinks. Baronius and Tillemont follow Socrates in supposing two
journeys of Athan. to Rome, and that the eighteen months began in 339
or 340, and had a break in them, during which he returned to
Alexandria.

Return to text}, expecting the arrival of yourselves and of whoever chose to
come. His presence overcame us all, for he would not have been here,
had he not felt confident in his cause; and he came not of his own
accord, but on a summons \footnote{[\emph{kl\={e}theis}].

Return to text}
by letter from us, in the manner in which we wrote to you. But
still you complain after all of our transgressing the Canons. Now
consider; who are they that have so acted? we who received this man
after such ample proof of his innocence, or they who being at Antioch
at the distance of six and thirty posts \footnote{or rather, halts, [\emph{monai}]. They are enumerated in the Itinerary
of Antonius, and are set down on Montfaucon's plate. The route passes
over the Delta to Pelusium, and then coasts all the way to Antioch.
These [\emph{monai}] were day's journeys, Coustant in Hilar. Psalm
118, Lit. 5. 2. or half a day's journey, Herman. Ibid; and were at
unequal intervals, Ambros. in Psalm 118, Serm. 5. §. 5. Gibbon says
that by the government conveyances, ``it was easy to travel an 100
miles in a day along the Roman roads.'' ch. ii. [\emph{Mon\={e}}] or
mansio properly means the building, where soldiers or other public
officers rested at night, (hence its application to monastic houses.)
Such buildings included granaries, stabling, \&c. vid. Cod. Theod.
ed. Gothofr. 1665. t. 1. p. 47. t. 2. p. 507. Ducange Gloss. t. 4. p.
426. Col. 2.

Return to text}; appointed a stranger to be Bishop, and sent him to Alexandria
with a military force; a thing which was not done even when Athanasius
was banished into Gaul, though it would have been done then, had he
been really proved guilty of the offence. But when he returned, of
course he found his Church unoccupied and waiting for him.

§. 30.

17. But now I am ignorant under what colour these
proceedings have been conducted. In the first place, if the truth must
be spoken, it was not right, when we had written to summon a Council,
that any persons should anticipate its decisions \footnote{The Eusebians kept the Pope's legates, and hastened their own Council
of the Dedication by way of anticipating him in their decision.

Return to text}: and in the next place, it was not fitting that such novel
proceedings should be adopted against the Church. For what Canon of
the Church \footnote{p. 41, note D, p. 55.

Return to text}, or what
Apostolical tradition warrants this, that when the Church was at
peace, and so many Bishops were in unanimity with Athanasius the
Bishop of Alexandria, Gregory should be sent thither, a stranger to
the city, not having been baptized there, nor known to the general
body, and desired neither by Presbyters, nor Bishops, nor Laity—that
he should be ordained at Antioch, and sent to Alexandria, accompanied
not by Presbyters, nor by Deacons of the city, nor by Bishops of
Egypt, but by soldiers? for they who came hither complained that this
was the case.

18. Even supposing that Athanasius was in the
position of a criminal after the Council, this appointment ought not
to have been made thus illegally and contrary to the Canon of the
Church, but the Bishops of the province ought to have ordained
one in that very Church, of that very Priesthood, of that very Clergy
\footnote{vid. Bingh. Ant. ii. 11. §. 7.

Return to text}; and the Canons \footnote{pp. 3, 50.

Return to text} received from the Apostles ought not thus to be set side. Had
this offence been committed against any one of you, would you not have
exclaimed against it, and demanded justice as for the transgression of
the Canons? Dearly beloved, we speak honestly, as in the presence of
God, and declare, that this proceeding was neither pious, nor lawful,
nor ecclesiastical. Moreover, the account which is given of the
conduct of Gregory on his entry into the city, plainly shews the
character of his appointment. In such peaceful times, as those who
came from Alexandria declared them to have been, and as the Bishops
also represented in their letters, the Church was set on fire \footnote{supr. p. 6.

Return to text}; Virgins were stripped; Monks were trodden under foot;
Presbyters and many of the people were scourged and suffered violence;
Bishops were cast into prison; multitudes were dragged about from
place to place; the holy Mysteries \footnote{Athan. only suggests this, supr. p. 6. S. Hilary says the same of the
conduct of the Arians at Toulouse; ``Clerks were beaten with clubs;
Deacons bruised with lead; nay, even \emph{on Christ Himself} (the
Saints understand my meaning) hands were laid.'' Contr. Coustant. 11.

Return to text}, about which they accused the Presbyter Macarius, were seized
upon by heathens and cast upon the ground; and all to constrain
certain persons to admit the appointment of Gregory. Such conduct
plainly shews who they are that transgress the Canons. Had the
appointment been lawful, he would not have had recourse to illegal
proceedings to compel the obedience of those who in a legal way
resisted him. And notwithstanding all this, you write that perfect
peace prevailed in Alexandria and Egypt. Surely not, unless the works
of peace are entirely changed, and you call such doings as these
peace.

§. 31.

19. I have also thought it necessary to point out
to you this circumstance, viz. that Athanasius positively asserted
that Macarius was kept at Tyre under a guard of soldiers, while only
his accuser accompanied those who went to the Mareotis \footnote{p. 31.

Return to text}; and that the Presbyters who desired to attend the inquiry
were not permitted, while the said inquiry respecting the chalice and
the Table was carried on before the Prefect and his band, and in the
presence of Heathens and Jews. This at first seemed incredible,
but it was proved to have been so from the Reports; which caused great
astonishment to us, as I suppose, dearly beloved, it does to you also.
Presbyters, who are the ministers of the Mysteries, are not permitted
to attend, but an enquiry concerning Christ's Blood and Christ's Body
is carried on before an external \footnote{[\emph{ex\={o}tikou}].

Return to text} judge, in the presence of Catechumens, nay, worse than that,
before heathens and Jews, who have so bad a name in regard to
Christianity. Even supposing that an offence had been committed, it
should have been investigated legally in the Church and by the Clergy,
not by heathens who abhor the Word and know not the Truth. I am
persuaded that both you and all men must perceive the nature and
magnitude of this sin. Thus much concerning Athanasius.

§. 32.

20. With respect to Marcellus \footnote{Julius here acquits Marcellus; but it would seem that he did not
eventually preserve himself from heretical notions, even if he
deserved a favourable judgment at this time. Athan. sides with him, de
Fug. 3. Hist. Arian. 6. but Epiphanius records, that on his asking
Athanasius what he (Athan.) thought of Marcellus, a smile came on his
face and he implied that there was some unsoundness in Marcellus's
views which perhaps he did not like to expose. Hær. 72. n. 4. And S.
Hilary says that Athan. separated him from his communion, as agreeing
with Photinus his disciple, Fragm. ii. 23. Sulpicius says the same. He
is considered heretical by S. Epiphanius, \emph{loc. cit}. S. Basil,
Epp. 69, 125, 263, 265. S. Chrysostom in Hebr. Hom. ii. 2. Theodoret,
Hær. ii. 10. vid. Petav. de Trin. i. 13. who condemns him, and Bull
far more strongly. Def. F. N. ii. 1. §. 9. Montfaucon defends him,
(in a special Dissertation, Collect. Nov. tom. 2.) and Tillemont, Mem.
tom. 7. p. 513. and Natalis Alex. Sæc. iv. Dissert. 30.

Return to text}, forasmuch as you have charged him also of impiety towards
Christ, I am anxious to inform you, that when he was here, he
positively declared that what you had written concerning him was not
true; but being nevertheless requested by us to give an account of his
faith, he answered in his own person with the utmost boldness, so that
we were obliged to acknowledge that he maintains nothing except the
truth. He made a confession \footnote{vid. Epiph. Hær. 72. 2, 3. and p. 73. infr.

Return to text} of the same godly doctrines concerning our Lord and Saviour
Jesus Christ as the Catholic Church confesses; and he affirmed that he
had held these opinions for a very long time, and had not recently
adopted them: as indeed our Presbyters \footnote{Vicentius and Vito.

Return to text}, who were at a former date present at the Council of Nicæa,
testified to his orthodoxy; for he maintained then, as he has done
now, his opposition to Arianism, (on which point it is right to
admonish you, lest any of you admit such heresy, instead of
abominating it as alien from \emph{sound doctrine} [1 Tim. i. 10.].)
Seeing then that he professed orthodox opinions, and had
testimony to his orthodoxy, what, I ask again in his case, ought we to
have done, except to receive him as a Bishop, as we did, and not
reject him from our communion?

21. These things I have written, not so much for
the purpose of defending their cause, as in order to convince you,
that we acted justly and canonically \footnote{pp. 5. 29.

Return to text} in receiving these persons, and that you are contentious
without a cause. But it is your duty to use your anxious endeavours
and to labour by every means to correct the irregularities which have
been committed contrary to the Canon, and to secure the peace of the
Churches; so that the peace of our Lord which has been given to us may
remain, and the Churches may not be divided, nor you incur the charge
of being authors of schism. For I confess, your past conduct is an
occasion of schism rather than of peace.

§. 33.

22. For not only the Bishops Athanasius and
Marcellus came hither and complained of the injustice that had been
done them, but many other Bishops also \footnote{The names of few are known; perhaps Marcellus, Asclepas, Paul of
Constantinople, Lucius of Adrianople. vid. Montf. in loc. Tillem. Mem.
tom. 7. p. 272.

Return to text}, from Thrace, from Cœle-Syria, from Phœnicia and Palestine,
and Presbyters not a few, and others from Alexandria and from other
parts, were present at the Council here, and in addition to their
other statements, lamented before all the assembled Bishops the
violence and injustice which the Churches had suffered, and affirmed
that similar outrages to those which had been committed in Alexandria
had occurred in their own Churches, and in others also. Again, there
lately came Presbyters with letters from Egypt and Alexandria, who
complained that many Bishops and Presbyters who wished to come to the
Council were prevented; for they said that, since the departure of
Athanasius \footnote{These outrages took place immediately on the dismission of Elpidius
and Philoxenus, the Pope's legates, from Antioch. Athan. Hist. Ar. 12.

Return to text} even up to
this time, Bishops who are confessors \footnote{e.g. Saparammon and Potamo, both Confessors, who were of the number of
the Nicene Fathers, and had defended Athan. at Tyre, were, the former
banished, the latter beaten to death. vid. infr. Hist. Ar. 12.

Return to text} have been beaten with stripes, that others have been cast into
prison, and that but lately aged men, who have been an exceedingly
long period in the Episcopate, have been given up to be employed
in the public works, and nearly all the Clergy of the Catholic Church
with the people are the objects of plots and persecutions. Moreover
they said that certain Bishops and other brethren had been banished
for no other reason than to compel them against their will to
communicate with Gregory and his Arian associates. We have heard also
from others, what is confirmed by the testimony of the Bishop
Marcellus, that a number of outrages, similar to those which were
committed at Alexandria, have occurred also at Ancyra in Galatia \footnote{The Pseudo-Sardican Council, i.e. the Eusebians at Philippopolis,
retort this accusation on the party of Marcellus; Hilar. Fragm. iii.
9. but the character of the outrages fixes them on the Arians. vid.
infr. p. 71, note H.

Return to text}. And in addition to all this, those who came to the Council
reported against some of you (for I will not mention names) certain
charges of so dreadful a nature that I have declined setting them down
in writing: perhaps you also have heard them from others. It was for
this cause especially that I wrote to desire \footnote{[\emph{protrepomenos}].

Return to text} you to come, that you might be present to hear them, and that
all irregularities might be corrected and differences healed. And
those who were called for these purposes ought not to have refused,
but to have come the more readily, lest by failing to do so they
should be suspected of what was alleged against them, and be thought
unable to prove what they had written.

§. 34.

23. Now according to these representations, since
the Churches are thus afflicted and treacherously assaulted, as our
informants positively affirmed, who are they that have lighted up the
flame of discord \footnote{vid. supr. p. 45.

Return to text}?
We, who grieve for such a state of things and sympathize with the
sufferings of the brethren, or those who have brought these things
about? While then such extreme confusion existed in every Church,
which was the cause why those who visited us came hither, I wonder how
you could write that unanimity prevailed in the Churches. These things
tend not to the edification of the Church, but to her destruction; and
those who rejoice in them are not sons of peace, but of confusion: but
our \emph{God is not} a God \emph{of confusion, but of peace} [1 Cor.
xiv. 33.]. Wherefore, as the God and Father of our Lord Jesus Christ
knows, it was from a regard for your good name, and with prayers
that the Churches might not fall into confusion, but might continue as
they were regulated \footnote{[\emph{ekanonisth\={e}}], vid. p. 51, infr. §. 69.

Return to text}
by the Apostles, that I thought it necessary to write thus unto you,
to the end that you might at length discountenance those who through
the effects of their mutual enmity have brought the Churches to this
condition. For I have heard, that it is only a certain few \footnote{ad Ep. Æg. 5. de Syn. 5.

Return to text} who are the authors of all these things.

24. Now, as having bowels of mercy, take ye care
to correct, as I said before, those irregularities which have been
committed contrary to the Canon, so that if any mischief has already
befallen, it may be healed through your zeal. And write not that I
have preferred the communion of Marcellus and Athanasius to yours, for
such like complaints are no indications of peace, but of
contentiousness and hatred of the brethren. For this cause I have
written the foregoing, that you may understand that we acted not
unjustly in admitting them to our communion, and so may cease this
strife. If you had come hither, and they had been condemned, and had
appeared unable to produce reasonable evidence in support of their
cause, you would have done well in writing thus. But seeing that, as I
said before, we acted agreeably to the Canon, and not unjustly, in
holding communion with them, I beseech you for the sake of Christ,
suffer not the members of Christ to be torn asunder, neither trust to
prejudices, but seek rather the peace of the Lord. It is neither holy
nor just, in order to gratify the narrow-spirit \footnote{[\emph{mikropsuchian}] supr. p. 41.

Return to text} of a few persons, to reject those who have never been
condemned, and thereby to grieve the Spirit. But if you think that you
are able to prove any thing against them, and to confute them face to
face, let those of you who please come hither: for they also promised
that they would be ready to establish completely the truth of those
things which they have reported to us.

§. 35.

25. Give us notice therefore of this, dearly
beloved, that we may write both to them, and to the Bishops who will
have again to assemble, so that the guilty may be condemned in the
presence of all, and confusion no longer prevail in the Churches. What
has already taken place is enough: it is enough surely that Bishops
have been sentenced to banishment in the presence of Bishops; of
which it behoves me not to speak at length, lest I appear to press too
heavily on those who were present on those occasions. But if one must
speak the truth, matters ought not to have proceeded so far; their
private feelings \footnote{[\emph{mikropsuchias}] p. 55.

Return to text}
ought not to have been suffered to reach their present pitch. Let us
grant the ``removal,'' as you write, of Athanasius and Marcellus, from
their own places, yet what must one say of the case of the other
Bishops and Presbyters who, as I said before, came hither from various
parts, and who complained that they also had been forced away, and had
suffered the like injuries? O dearly beloved, the decisions of the
Church are no longer according to the Gospel, but tend only to
banishment and death \footnote{Hist. Arian. §. 67.

Return to text}.
Supposing, as you assert, that some offence rested upon those persons,
the case ought to have been conducted against them, not after this
manner, but according to the Canon of the Church \footnote{p. 53.

Return to text}. Word should have been written of it to us all \footnote{Coustant in loc. fairly insists on the word ``all,'' as shewing that S.
Julius does not here claim the prerogative of judging \emph{by himself}
all Bishops whatever, and that what follows relates merely to the
Church of Alexandria.

Return to text}, that so a just sentence might proceed from all. For the
sufferers were Bishops, and Churches of no ordinary note, but those
which the Apostles themselves had governed in their own persons \footnote{St. Peter (Greg. M. Epist. vii. Ind. 15. 40 ) or St. Mark (Leo, Ep.
9.) at Alexandria. St. Paul at Ancyra in Galatia, (Tertull. contr.
Marcion. iv. 5.) vid. Coustant, in loc.

Return to text}.

26. And why was nothing said to us concerning the
Church of the Alexandrians in particular? Are you ignorant that the
custom has been for word to be written first to us, and then for a
just sentence to be past from this place \footnote{Socrates says somewhat differently, ``Julius wrote back ... that they
acted against the Canons, because they had not called him to a
Council, the Ecclesiastical Canon commanding that the Churches ought
not to make Canons beside the will of the Bishop of Rome.'' Hist. ii.
17. Sozomen in like manner, ``for it was a sacerdotal law, to declare
invalid whatever was transacted beside the will of the Bishop of the
Romans.'' Hist. iii. 10. vid. Pope Damasus ap. Theod. Hist. v. 10.
Leon. Epist. 14. \&c. In the passage in the text the prerogative of
the Roman see is limited, as Coustant observes, to the instance of
Alexandria; and we actually find in the third century a complaint
lodged against its Bishop Dionysius with the Pope.

Return to text}? If then any such suspicion rested upon the Bishop there,
notice thereof ought to have been sent to the Church of this place;
whereas, after neglecting to inform us, and proceeding on their own
authority as they pleased, now they desire to obtain our
concurrence in their decisions, though we never condemned him. Not so
have the Constitutions \footnote{[\emph{diataxeis}]. St. Paul says [\emph{hout\={o}s en tais ekkl\={e}siais
diatassomai}]. 1 Cor. vii. 17. [\emph{ta de loipa diataxomai}].
Ibid. xi. 34. vid. Pearson, Vind. Ignat. p. 298. Hence Coustant in
loc. Athan. would suppose Julius to refer to 1 Cor. v. 4. which Athan.
actually quotes, Ep. Encycl. §. 2. supr. pp. 4. 5. Pearson \emph{loc.
cit}. considers the [\emph{diataxeis}] of the Apostles, as a
collection of regulations and usages, which more or less represented,
or claimed to represent, what may be called St. Paul's rule, or St.
Peter's rule, \&c. Cotelier considers the [\emph{diataxeis}] as the
same as the [\emph{didaxai}] the ``doctrine'' or ``teaching'' of the
Apostles. Præfat. in Const. Apost. So does Beveridge, Cod. Can.
Illustr. ii. 9. §. 5.

Return to text} of
Paul, not so have the traditions of the Fathers directed; this is
another form of procedure, a novel practice. I beseech you, readily
bear with me: what I write is for the common good. For what we have
received from the blessed Apostle Peter \footnote{[Petri] in Sede suâ vivit potestas et excellit auctoritas. Leon. Serm.
iii. 3. vid. contra Barrow on the Supremacy, p. 116. ed. 1836. ``not
one Bishop, but all Bishops together through the whole Church, do
succeed St. Peter, or any other Apostle.''

Return to text

}, that I signify to you; and I should not have written this, as
deeming that these things were manifest unto all men, had not these
proceedings so disturbed us. Bishops are forced away from their sees
and driven into banishment, while others from different quarters are
appointed in their place; others are treacherously assailed, so that
the people have to grieve for those who are forcibly taken from them,
while, as to those who are sent in their room, they are obliged to
give over seeking the man whom they desire, and to receive those they
do not.

27. I ask \footnote{[\emph{axi\={o}}].

Return to text} of you, that such things may no longer be, but that you will
denounce in writing those persons who attempt them; so that the
Churches may no longer be afflicted thus, nor any Bishop or Presbyter
be treated with insult, nor any one be compelled to act contrary to
his judgment, as they have represented to us, lest we become a
laughing-stock among the heathen, and above all, lest we excite the
wrath of God against us. For every one of us shall give account in the
Day of judgment of the things which he has done in this life. May we
all be possessed with the mind of God! so that the Churches may
recover their own Bishops, and rejoice evermore in Jesus Christ our
Lord; through Whom to the Father be glory, for ever and ever. Amen.

I pray for your health in the Lord, brethren
dearly beloved and greatly longed for.

§. 36

28. Thus wrote the Council of Rome by Julius
Bishop of Rome.

\xsection{Chapter 3. Letters of the Council of Sardica to the Churches of Egypt and of Alexandria, and to all Churches}

1. B\textsc{ut} when,
notwithstanding, the Eusebians proceeded without shame, disturbing the
Churches, and plotting the ruin of many, the most religious Emperors
Constantius and Constans being informed of this, commanded \footnote{Margin Notes

1. [\emph{ekeleusen}].

Return to text}
the Bishops from both the West and East to meet together in the city
of Sardica. In the mean time Eusebius \footnote{of
Nicodemia.

Return to text} died: but a great
number assembled from all parts, and we challenged the associates of
Eusebius to submit to a trial. But they, having before their eyes the
things that they had done, and perceiving that their accusers had come
up to the Council, were afraid to do this; but, while all beside met
with honest intentions, they again brought with them the Counts \footnote{Hist. Ar. 15.

Return to text} Musonianus \footnote{Musonian was originally of Antioch, and his
name Strategius; he had been promoted and honoured with a new name by
Constantine, for whom he had collected information about the Manichees.
Amm. Marc. xv. 13. §. 1. In 354, he was Prætorian Prefect of the
East. (vid. Libr. of F. O. T. vol. viii. p. 73, note a.) Libanius
praises him.

Return to text} and Hesychius the Castrensian
\footnote{The Castrensians were the officers of the
palace; castra, as [\emph{stratopedon}] infr. §. 86. being at this
time used for the Imperial Court. vid. Gothofred in Cod. Theod. vi.
30. p. 218. Ducange in voc.

Return to text}, that,
as their custom was, they might effect their own aims by their
authority. But when the Council met without the Counts, and no
soldiers were permitted to be present, they were confounded, and
conscience-stricken, because they could no longer obtain what judgment
they wished, but such only as reason and truth \footnote{[\emph{ho t\={e}s al\={e}theias logos}],
vid. p. 72.

Return to text} required. We,
however, frequently repeated our challenge, and the Council of Bishops
called upon them to come forward, saying, ``You have come for the
purpose of undergoing a trial; why then do you now withdraw
yourselves? Either you ought not to have come, or having come, not to
conceal yourselves. Such conduct will prove your greatest
condemnation. Behold, Athanasius and his friends are here, whom you
accused while absent; if therefore you think that you have any thing
against them, you may convict them face to face. But if you pretend to
be unwilling to do so, while in truth you are unable, you plainly shew
yourselves to be calumniators, and the Council will give sentence
against you accordingly.'' When they heard this they were
self-condemned, (for they were conscious of their machinations and
fabrications against us,) and were ashamed to appear, thereby proving
themselves to have been guilty of many base calumnies.

2. The holy Council therefore denounced their
indecent and suspicious flight \footnote{to
Philippopolis.

Return to text}, and admitted us to make our
defence and when we had related their conduct towards us, and proved
the truth of our statements by witnesses and other evidence, they were
filled with astonishment, and all acknowledged that our opponents had
good reason to be afraid to meet the Council, lest their guilt should
be proved before their faces. They said also, that probably they had
come from the East, supposing that Athanasius and his friends would
not appear, but that, when they saw them confident in their cause, and
challenging a trial, they fled. They accordingly received us as
injured persons who had been falsely accused, and confirmed \footnote{[\emph{ekur\={o}san k.t.l}.] p. 38, ref. 5.

Return to text}
yet more towards us their fellowship and loving hospitality \footnote{[\emph{agap\={e}n}].

Return to text}.
But they deposed Eusebius's associates in wickedness, who had become
even more shameless than himself, viz. Theodorus \footnote{p. 39, note B.

Return to text} of Heraclea,
Narcissus of Neronias, Acacius \footnote{vol. 8, p. 7. p.

Return to text} of Cæsarea, Stephanus \footnote{Hist. Arian. §. 20.

Return to text} of Antioch, Ursacius and Valens of Pannonia, Menophantus of
Ephesus, and George \footnote{p. 25. F.

Return to text} of Laodicæa; and they wrote to the
Bishops in all parts of the world, and to the diocese \footnote{[\emph{paroikiai}].

Return to text} of
each of the injured persons, in the following terms.

3. Letter of the council of Sardica to the Church
of Alexandria

The holy Council, by the grace of God assembled
at Sardica, from \footnote{vid. supr. p. 14. where Isauria, Thessaly,
Sicily, Britain \&c. added. Also Theod. Hist. ii. 6. vid. p. 78. r.
1.

Return to text} Rome, Spain, Gaul, Italy, Campania,
Calabria, Apulia, Africa, Sardinia, Pannonia, Mysia, Dacia, Noricum,
Siscia, Dardania, the other Dacia, Macedonia, Thessaly, Achaia,
Epirus, Thrace, Rhodope, Palestine, Arabia, Crete, and Egypt, to
their dearly beloved brethren, the Presbyters and Deacons, and to all
the Holy Church of God abiding at Alexandria, sends health in the
Lord.

§. 37.

We were not ignorant, but the fact was well known
to us, even before we received the letters of your piety, that the
supporters of the abominated heresy of the Arians were practising many
dangerous machinations, rather to the destruction of their own souls,
than to the injury of the Church. For this has ever been the object of
their unprincipled craft; this is the deadly design in which they have
been continually engaged; viz. how they may best expel from their
places and persecute all who are to be found any where of orthodox
sentiments, and maintaining the doctrine of the Catholic Church, which
was delivered to them from the Fathers. Against some they have laid
false accusations; others they have driven into banishment; others
they have destroyed by the punishments inflicted on them. Thus also
they endeavoured by violence and tyranny to surprise the innocence of
our brother and fellow Bishop Athanasius, and therefore conducted
their enquiry into his case without any scrupulous care, without any
faith, without any sort of justice. Accordingly having no confidence
in the part they had played on that occasion, nor yet in the reports
they had circulated against him, but perceiving that they were unable
to produce any certain evidence respecting them, when they came to the
city of Sardica, they were unwilling to meet the Council of all the
holy Bishops. From this it became evident that the decision of our
brother and fellow-Bishop Julius was a just one \footnote{vid. infr. p. 80, note P.

Return to text}; for after
cautious deliberation and care he had determined, that we ought not to
hesitate at all about holding communion with our brother Athanasius.
For he had the credible testimony of eighty Bishops, and was also able
to advance this fair argument in his support, that by the mere means
of our dearly beloved brethren his own Presbyters, and by
correspondence, he had defeated the designs of the Eusebians, who
relied more upon violence than upon a judicial enquiry.

4. Wherefore all the Bishops from all parts
determined upon holding communion with Athanasius on the ground that
he was innocent. And let your charity also observe, that when he
came to the holy Council assembled at Sardica, the Bishops of the East
were informed of the circumstance, as we said before, both by letter,
and by injunctions conveyed by word of mouth, and were summoned \footnote{[\emph{ekl\={e}th\={e}san}], vid. p.
49. r. 3.

Return to text} by us to be present. But, being condemned by their own conscience,
they had recourse to unbecoming excuses, and set themselves to avoid
the enquiry. They demanded that an innocent man should be rejected
from our communion, just as if he had been guilty, not considering how
unbecoming, or rather how impossible, such a proceeding was. And as
for the Reports which were framed in the Mareotis by certain most
wicked and most profligate youths \footnote{supr. p. 31, note M.

Return to text}, to whose hands one would
not commit the very lowest office of the ministry, it is certain that
they were \emph{ex parte} statements. For neither was our brother the
Bishop Athanasius present on the occasion, nor the Presbyter Macarius
who was accused by them. And besides, their enquiry, or rather their
falsification of facts, was attended by the most disgraceful
circumstances. Sometimes heathens, sometimes Catechumens, were
examined, not that they might declare what they knew, but that they
might assert those falsehoods which they had been taught by others.
And when you Presbyters, who were anxious in the absence of your
Bishop, desired to be present at the enquiry, in order that you might
shew the truth, and disprove falsehood, no regard was paid to you;
they would not permit you to be present, but drove you away with
insult.

5. Now although their calumnies have been most
plainly exposed before all men by these circumstances; yet we found
also, on reading the Reports, that that most iniquitous person,
Ischyras, who has obtained from them the empty title of Bishop as his
reward for the false accusation, had convicted himself of calumny. He
declares in the Reports that at the very time when, according to his
positive assertions, Macarius entered his cell, he lay there sick;
whereas the Eusebians have had the boldness to write that Ischyras was
standing up and offering the oblations, when Macarius came in \footnote{pp. 30, 48.

Return to text}.

§. 38.

6. The base and slanderous charge which they next
alleged against him, has become well-known to all men. They raised a
great outcry, affirming that Athanasius had committed murder, and had
destroyed one Arsenius a Meletian Bishop, whose loss they
pretended to deplore with feigned lamentations and untrue tears, and
demanded that the body of a living man, as if a dead one, should be
given up to them. But their fraud was easily detected: one and all
knew that the person was alive, and was numbered among the living \footnote{pp. 26, 47.

Return to text}.

7. And when these men, who are ready upon any
opportunity, perceived their falsehoods detected, (for Arsenius shewed
himself alive, and so proved that he had not been destroyed, and was
not dead,) yet they would not rest, but proceeded to other calumnies \footnote{vid. supr. 36. infr. §. 87.

Return to text}, and to slander Athanasius by a fresh expedient. Well; our
brother, dearly beloved, was not confounded, but again in the present
case also with great boldness challenged them to the proof, and we too
prayed and exhorted them to come to the trial, and if they were able,
to establish their charge against him. O great arrogance! O dreadful
pride! or rather, if one must say the truth, O evil and guilt-stricken
conscience! for this is the view which all men take of it.

8. Wherefore, dearly beloved brethren, we
admonish and exhort you, above all things to maintain the right faith
of the Catholic Church. You have undergone many severe and grievous
trials; many are the insults and injuries which the Catholic Church
has suffered, but \emph{he that endureth to the end the same shall be
saved} [Matt. x. 22.]. Wherefore even though they shall still
recklessly assail you, let your tribulation be unto you for joy. For
such afflictions have a share in martyrdom, and such confessions and
tortures as yours will not be without their reward, but ye shall
receive the prize from God. Therefore strive above all things in
support of the sound faith, and of the innocence of your Bishop and
our brother Athanasius. We also have not held our peace, nor been
negligent of what concerns your comfort, but have deliberated and done
whatsoever the claims of charity demand. We sympathize with our
suffering brethren, and their afflictions we consider as our own.

§. 39.

9. Accordingly we have written to beseech our
most religious and godly Emperors, that their Graces would give orders
for the release of those who are still suffering from affliction and
oppression, and would command that none of the magistrates, whose
duty it is to attend only to civil causes, give judgment upon Clergy \footnote{vid. Bingham Antiqu. v. 2. §. 5. \&c.
Gieseler Eccl. Hist. vol. 1. p. 242. Bassi. Biblioth. Jur. t. 1. p.
276. Bellarm. de Cleric. 28.

Return to text}, nor henceforward in any way, on pretence of providing for
the Churches, attempt any thing against the brethren; but that every
one may live, as he prays and desires to do, free from persecution,
from violence and fraud, and in quietness and peace may follow the
Catholic and Apostolic Faith. As for Gregory, who has the reputation
of being illegally ordained by the heretics, and has been sent by them
to your city, we wish your unanimity to understand, that he has been
degraded by a judgment of the whole sacred Council, although indeed he
has never at any time been considered to be Bishop at all. Wherefore
receive gladly your Bishop Athanasius, for to this end we have
dismissed him in peace. And we exhort all those who either through
fear, or through the intrigues of certain persons, have held communion
with Gregory, that now being admonished, exhorted, and persuaded by
us, they withdraw from that his accursed communion, and straightway
unite themselves to the Catholic Church.

§. 40.

10. Forasmuch as we have learnt that Aphthonius,
Athanasius the son of Capito, Paul, and Plutio, our fellow Presbyters
\footnote{supr. p. 35.

Return to text}, have also suffered from the machinations of the Eusebians,
so that some of them have had trial of exile, and others have lied on
peril of their lives, we have in consequence thought it necessary to
make this known unto you, that you may understand that we have
received and acquitted them also, being aware that whatever has been
done by the Eusebians against the Orthodox has tended to the glory and
commendation of those who have been attacked by them. It were fitting
that your Bishop and our brother Athanasius should make this known to
you respecting them, to his own respecting his own; but as for more
abundant testimony he wished the holy Council also to write to you, we
deferred not to do so, but hastened to signify this unto you, that you
may receive them as we have done, for they also are deserving of
praise, because through their piety towards Christ they have been
thought worthy to endure violence at the hands of the heretics.

11. What decrees have been past by the holy
Council against those who are at the head of the Arian heresy, and
have offended against you, and the rest of the Churches, you will
learn from the subjoined documents \footnote{vid. Encycl. Letter, infr. p. 69.

Return to text}. We have sent them to
you, that you may understand from them that the Catholic Church will
not overlook those who offend against her.

12. Letter of the Council of Sardica to the
Bishops of Egypt and Libya

The holy Council, by the grace of God assembled
at Sardica, to the Bishops of Egypt and Libya, their fellow ministers
and dearly beloved brethren, sends health in the Lord.

§. 41.

We were not ignorant \footnote{It will be observed that this Letter is nearly
a transcript of the foregoing. It was first printed in the Benedictine
Edition.

Return to text}, but the fact was
well known to us, even before we received the letters of your piety,
that the supporters of the abominated heresy of the Arians were
practising many dangerous machinations, rather to the destruction of
their own souls, than to the injury of the Church. For this has ever
been the object of their craft and villainy: this is the deadly design
in which they have been continually engaged, viz. how they may best
expel from their places and persecute all who are to be found any
where of orthodox sentiments, and maintaining the doctrine of the
Catholic Church, which was delivered to them from the Fathers. Against
some they have laid false accusations; others they have driven into
banishment; others they have destroyed by the punishments inflicted on
them. Thus also they endeavoured by violence and tyranny to surprise
the innocence of our brother and fellow Bishop Athanasius, and
therefore conducted their enquiry into his case without any scrupulous
care, without any faith, without any sort of justice. Accordingly,
having no confidence in the part they had played on that occasion, nor
yet in the reports they had circulated against him, but perceiving
that they were unable to produce any certain evidence respecting them,
when they came to the city of Sardica, they were unwilling to meet the
Council of all the holy Bishops. From this it became evident that the
decision of our brother and fellow Bishop Julius was a just one;
for after cautious deliberation and care he had decided, that we ought
not to hesitate at all about holding communion with our brother
Athanasius. For he had the credible testimony of eighty Bishops, and
was also able to advance this fair argument in his support, that by
the mere means of our dearly beloved brethren his own Presbyters, and
by correspondence, he had defeated the designs of the Eusebians, who
relied more upon violence, than upon a judicial enquiry.

13. Wherefore all the Bishops from all parts
determined upon holding communion with Athanasius on the ground that
he was innocent. And let your charity also observe, that when he came
to the holy Council assembled at Sardica, the Bishops of the East were
informed of the circumstance, as we said before, both by letter, and
by injunctions conveyed by word of mouth, and were invited by us to be
present. But, being condemned by their own conscience, they had
recourse to unbecoming excuses, and began to avoid the enquiry. They
demanded that an innocent man should be rejected from our communion,
just as if he had been guilty, not considering how unbecoming, or
rather how impossible, such a proceeding was. And as for the reports
which were framed in the Mareotis by certain most wicked and abandoned
youths, to whose hands one would not commit the very lowest office of
the ministry, it is certain that they were \emph{ex parte} statements.
For neither was our brother the Bishop Athanasius present on the
occasion, nor the Presbyter Macarius, who was accused by them. And
besides, their enquiry, or rather their falsification of facts, was
attended by the most disgraceful circumstances. Sometimes Heathens,
sometimes Catechumens, were examined, not that they might declare what
they knew, but that they might assert those falsehoods which they had
been taught by others. And when you Presbyters, who were anxious in
the absence of your Bishop, desired to be present at the enquiry, in
order that you might shew the truth, and disprove falsehood, no regard
was paid to you; they would not permit you to be present, but drove
you away with insult.

14. Now although their calumnies have been most
plainly exposed before all men by these circumstances; yet we found also, on reading the Reports, that that most iniquitous person
Ischyras, who has obtained from them the empty title of Bishop as his
reward for the false accusation, had convicted himself of calumny. He
declares in the Reports, that at the very time when, according to his
positive assertions, Macarius entered his cell, he lay there sick;
whereas the Eusebians have had the boldness to write that Ischyras was
standing up and offering the oblations, when Macarius came in.

§. 42.

15. The base and slanderous charge which they
next alleged against him has become well known unto all men. They
raised a great outcry, affirming that Athanasius had committed murder,
and destroyed one Arsenius a Meletian Bishop, whose loss they
pretended to deplore with feigned lamentations, and untrue tears, and
demanded that the body of a living man, as if a dead one, should be
given up to them. But their fraud was easily detected; one and all
knew that the person was alive, and was numbered among the living.

16. And when these men, who are ready upon any
opportunity, perceived their falsehood detected, (for Arsenius shewed
himself alive, and so proved that he had not been destroyed, and was
not dead,) yet they would not rest, but proceeded to add other to
their former calumnies, and to slander Athanasius by a fresh
expedient. Well: our brother, dearly beloved, was not confounded, but
again in the present case also with great boldness challenged them to
the proof, and we too prayed and exhorted them to come to the trial,
and if they were able, to establish their charge against him. O great
arrogance! O dreadful pride! or rather, if one must say the truth, O
evil and guilt-stricken conscience! for this is the view which all men
take of it.

17. Wherefore, dearly beloved brethren, we
admonish and exhort you, above all things, to maintain the right faith
of the Catholic Church. You have undergone many severe and grievous
trials; many are the insults and injuries which the Catholic Church
has suffered, but \emph{he that endureth to the end, the same shall be
saved} [Mat. x. 22.]. Wherefore, even though they shall still
recklessly assail you, let your tribulation be unto you for joy. For
such afflictions have a share in martyrdom, and such confessions and
tortures as yours will not be without their reward, but ye shall
receive the prize from God. Therefore, strive above all things in
support of the sound Faith, and of the innocence of your Bishop and
our brother Athanasius. We also have not held our peace, nor been
negligent of what concerns your comfort, but have deliberated and done
whatever the claims of charity demand. We sympathize with our
suffering brethren, and their afflictions we consider as our own, and
have mingled our tears with yours. And you, brethren, are not the only
persons who have suffered: many others also of our brethren in
ministry have come hither, bitterly lamenting these things.

§. 43.

18. Accordingly, we have written to beseech our
most religious and godly Emperors, that their Graces would give orders
for the release of those who are still suffering from affliction and
oppression, and would command that none of the magistrates, whose duty
it is to attend only to civil causes, give judgment upon Clergy, nor
henceforward in any way, on pretence of providing for the Churches,
attempt any thing against the brethren, but that every one may live,
as he prays and desires to do, free from persecution, from violence
and fraud, and in quietness and peace may follow the Catholic and
Apostolic Faith. As for Gregory who has the reputation of being
illegally ordained by the heretics, and who has been sent by them to
your city, we wish your unanimity to understand, that he has been
degraded by the judgment of the whole sacred Council, although indeed
he has never at any time been considered to be a Bishop at all.
Wherefore receive gladly your Bishop Athanasius; for to this end we
have dismissed him in peace. And we exhort all those, who either
through fear, or through the intrigues of certain persons, have held
Communion with Gregory, that being now admonished, exhorted, and
persuaded by us, they withdraw from his accursed communion, and
straightway unite themselves to the Catholic Church.

19. What decrees have been passed by the holy
Council against Theodorus, Narcissus, Stephanus, Acacius, Menophantus,
Ursacius, Valens, and George \footnote{p. 60.

Return to text}, who are the heads of the Arian
heresy, and have offended against you and the rest of the Churches,
you will learn from the subjoined documents. We have sent them to you,
that your piety may assent to our decisions, and that you may
understand from them, that the Catholic Church will not overlook those
who offend against her.

20. Encyclical Letter of the Council of Sardica

The holy Council \footnote{vid. Theod. Hist. ii. 6. Hil. Fragm. ii.

Return to text}, by the grace of God,
assembled at Sardica, to their dearly beloved brethren, the Bishops
and fellow-Ministers of the Catholic Church every where, sends health
in the Lord.

§. 44.

The Arian fanatics have dared repeatedly to
attack the servants of God, who maintain the right faith; they
attempted to substitute a spurious doctrine, and to drive out the
orthodox; and at last they made so violent an assault against the
Faith, that it became known to the piety of our most religious
Emperors. Accordingly, the grace of God assisting them, our most
religious Emperors have themselves assembled us together out of
different provinces and cities, and have permitted this holy Council
to be held in the city of Sardica; to the end that all dissension may
be done away, and all false doctrine being driven from us, Christian
godliness may alone be maintained by all men. The Bishops of the East
also attended, being exhorted to do so by the most religious Emperors,
chiefly on account of the reports they have so often circulated
concerning our dearly beloved brethren and fellow-ministers Athanasius
Bishop of Alexandria, and Marcellus Bishop of Ancyro-Galatia. Their
calumnies have probably already reached you, and perhaps they have
attempted to disturb your ears, that you may be induced to believe
their charges against those innocent men, and that they may obliterate
from your minds any suspicious respecting their own wicked heresy. But
they have not been permitted to effect this to any great extent; for
the Lord is the Defender of His Churches, who endured death for their
sakes and for us all, and provided access to heaven for us all through
Himself. When therefore the Eusebians wrote long ago to Julius our
brother and Bishop of the Church of the Romans, against our
fore-mentioned brethren, that is to say, Athanasius, Marcellus, and
Asclepas \footnote{Asclepas, or Asclepius of Gaza, Epiph. Hær.
69. 4. was one of the Nicene Fathers, and according to Theod. Hist. i.
27. was at the Council of Tyre, which Athan. also attended, but only
by compulsion. According to the Eusebians at Philippopolis, they had
deposed him about 330, if the Council of Sardica was held 347. They
state, however, at the same time, that he had been condemned by
Athanasius and Marcellus. vid. Hilar. Fragm. iii. 13. Sozomen, Hist.
iii. 8. says that they deposed him on the charge of having overturned
an altar; and, after Athan. infr. §. 47. that he was acquitted at
Sardica on the ground that Eusebius of Cæsarea and others had
reinstated him in his see, (before 339.) There is mention of a Church
built by him in Gaza, ap. Bolland. Febr. 26. Vit. S. Porphyr. n. 20.
p. 648.

Return to text}, the Bishops from the other parts wrote also,
testifying to the innocence of our fellow-minister Athanasius, and
declaring that the representations of the Eusebians were nothing else
but mere falsehood and calumny.

21. And indeed their calumnies were clearly
proved by the fact that, when they were called \footnote{[\emph{kl\={e}thentas}].

Return to text} to a Council
by our dearly beloved fellow-minister Julius, they would not come, and
also by what was written to them by Julius himself. For had they had
confidence in the measures and the acts in which they were engaged
against our brethren, they would have come. And besides, they gave a
still more evident proof of their conspiracy by their conduct in this
great and holy Council. For when they arrived at the city of Sardica,
and saw our brethren Athanasius, Marcellus, Asclepas, and the rest,
they were afraid to come to a trial, and though they were repeatedly
invited to attend, they would not obey the summons. Although all we
Bishops met together, and above all that man of an happy old age,
Hosius, one who on account of his age, his confession, and the many
labours he has undergone, is worthy of all reverence; and although we
waited and besought them to come to the trial, that in the presence of
our fellow-ministers they might establish the truth of those charges
which they had circulated and written against them in their absence;
yet they would not come, when they were thus called, as we said
before, thus giving proof of their calumnies, and almost proclaiming
to the world by this their refusal, the plot and conspiracy in which
they have been engaged. They who are confident of the truth of their
assertions are able to make them good against their opponents face to
face. But as they would not meet us, we think that no one can now
doubt, however they may again have recourse to their bad practices,
that they possess no proof against our brethren, but calumniate them
in their absence, while they avoid their presence.

§. 45.

22. They fled, dearly beloved brethren, not only
on account of the calumnies they had uttered, but because they saw
that those had come who had various charges to advance against them.
For chains and iron were brought forward which they had used; persons
appeared who had returned from banishment; there came also our
brethren, kinsmen of those who were still detained in exile, and
friends of such as had perished through their means. And what was the
most weighty ground of accusation, Bishops were present, one \footnote{Perhaps Lucius of Hadrianople, says
Montfaucon, referring to Apol. de Fug. §. 3. vid. also Hist. Arian.
19.

Return to text}
of whom brought forward the iron and the chains which they had caused
him to wear, and others testified to the deaths which had been brought
about by their calumnies. For they had proceeded to such a pitch of
madness, as even to attempt to destroy Bishops; and would have
destroyed them, had they not escaped their hands. Our fellow-minister,
Theodulus of blessed memory \footnote{Theodulus, Bishop of Trajanopolis in Thrace,
who is here spoken of as deceased, seems to have suffered this
persecution from the Eusebians upon their retreat from Sardica, vid.
Athan. Hist. Arian. §. 19. We must suppose then with Montfaucon, that
the Council, from whom this letter proceeds, sat some considerable
time after that retreat, and that the proceedings spoken of took place
in the interval. Socrates, however, makes Theodulus survive Constans,
who died 350. Hist. ii. 26.

Return to text}, died during his flight from
their false accusations, orders having been given in consequence of
these to put him to death. Others also exhibited sword-wounds; and
others complained that they had been exposed to the pains of hunger
through their means. Nor were they ordinary persons who testified to
these things, but whole Churches, in whose behalf legates appeared \footnote{The usual proceeding of the Arians was to
retort upon the Catholics the charges which they brought against them,
supr. p. 54, note P. Accordingly, in their Encyclical from
Philippopolis, they say that ``a vast multitude had congregated at
Sardica, of wicked and abandoned persons, from Constantinople and
Alexandria; who lay under charges of murder, blood, slaughter,
robbery, plunder, spoiling, and all nameless sacrileges and crimes;
who had broken altars, burnt Churches, ransacked private houses,
\&c. \&c. Hil. Fragm. iii. 19.

Return to text}, and told us of soldiers sword in hand, of multitudes armed
with clubs, of the threats of judges, of the use of forged letters.
For there were read certain forged letters of Theognius against our
fellow ministers Athanasius, Marcellus, and Asclepas, written with the
design of exasperating the Emperors against them; and those who had
then been Deacons of Theognius proved the fact. In addition to these
things, we heard of virgins stripped naked, Churches burnt,
ministers in custody, and all for no other end, but only for the sake
of the accursed heresy of the Arian fanatics, whose communion whoso
refused was forced to suffer these things.

23. When they perceived then how matters lay,
they were in a strait what course to choose. They were ashamed to
confess all that they had done, but were unable to conceal it any
longer. They therefore came to the city of Sardica, that by their
appearance there they might seem to remove suspicion from themselves
of the guilt of such things. But when they saw those whom they had
calumniated, and those who had suffered at their hands; when they had
before their eyes their accusers and the proofs of their guilt, they
were unwilling to come forward, though invited by our fellow-ministers
Athanasius, Marcellus, and Asclepas, who with great freedom complained
of their conduct, and urged and challenged them to the trial,
promising not only to refute their calumnies, but also to bring proof
of the offences which they had committed against their Churches. But
they were seized with such terrors of conscience, that they fled; and
in doing so they exposed their own calumnies, and confessed by running
away the crimes of which they had been guilty.

§. 46.

24. But although their malice and their calumnies
have been plainly manifested on this as well as on former occasions,
yet that they may not devise means of practising a further mischief in
consequence of their flight, we have considered it advisable to
examine the part they have played according to the principles of truth
\footnote{supr. p. 59. ref. 2. Orat. 1. O.T. p. 227 \emph{init}.

Return to text}; this has been our purpose, and we have found them
calumniators by their acts, and authors of nothing else than a plot
against our brethren in ministry. For Arsenius, who they said had been
murdered by Athanasius, is still alive, and is numbered among the
living; from which we may infer that the reports they have circulated
on other subjects are fabrications also. And whereas they spread
abroad a rumour concerning a chalice, which they said had been broken
by Macarius the Presbyter of Athanasius, those who came from
Alexandria, the Mareotis, and the other parts, testified that nothing
of the kind had taken place. And the Egyptian Bishops \footnote{p. 30.

Return to text} who
wrote to Julius our brother in ministry, positively affirmed that
there did not exist among them even any suspicion whatever of
such a thing.

25. Moreover, the Reports, which they say they
have to produce against him, are, as is notorious, \emph{ex parte}
statements; and even in the formation of these very Reports, Heathens
and Catechumens were examined; one of whom, a Catechumen, said \footnote{pp. 48, 49.

Return to text} in his examination that he was present in the room, when Macarias
broke in upon them; and another declared, that Ischyras of whom they
speak so much, lay sick in his cell at the time; from which it appears
that the Mysteries were never celebrated at all, because Catechumens
were present, and also that Ischyras was not there, but was lying sick
on his bed. Besides, this wicked wretch Ischyras, who has falsely
asserted, as he was convicted of doing, that Athanasius had burnt some
of the sacred books, has himself confessed that he was sick, and was
lying in his bed when Macarius came; from which it is plain that he is
a slanderer. Nevertheless, as a reward for these his calumnies, they
have given to this very Ischyras the title of Bishop, although he has
never been even a Presbyter. For two Presbyters, who were once
associated with Meletius, but were afterwards received by Alexander,
Bishop of Alexandria, of blessed memory, and are now with Athanasius,
appeared before the Council, and testified that he was not even a
Presbyter of Meletius, and that Meletius never had either Church or
Minister in the Mareotis. And yet this man, who has never been even a
Presbyter, they have now brought forward as a Bishop, that by this
name they may have a means of overpowering those who are within
hearing his calumnies.

§. 47.

26. The book of our brother Marcellus was also
read, by which the fraud of the Eusebians were plainly discovered. For
what Marcellus had advanced by way of enquiry \footnote{vid. de Decr. O.T. vol. 8. p. 44, e.

Return to text}, they falsely
represented as his professed opinion; but when the subsequent parts of
the book were read, and the parts preceding these queries, his faith
was found to be correct. He had never pretended, as they positively
affirmed \footnote{de Syn. O.T. p. 110, note r.

Return to text}, that the word of God had His beginning from holy
Mary, nor that His kingdom had an end; on the contrary he had written
that His kingdom was both without beginning and without end.

Our brother Asclepas also produced Reports which
had been drawn up at Antioch in the presence of his accusers and
Eusebius of Cæsarea, and proved that he was innocent by the sentence
of the Bishops who judged his cause \footnote{p. 70. E.

Return to text}. They had good reason
therefore, dearly beloved brethren, for disobeying our frequent
summons, and for deserting the Council. They were driven to this by
their own consciences; but their flight only confirmed the proof of
their calumnies, and caused those things to be believed against them,
which their accusers, who were present, were asserting and arguing.
But besides all these timings, they had not only received those who
were formerly degraded and ejected on account of the Arian heresy, but
had even promoted them to a higher station, advancing Deacons to the
Presbytery, and of Presbyters making Bishops, for no other end, but
that they might disseminate and spread abroad impiety, and corrupt the
orthodox faith.

§. 48.

27. Their present leaders are, after Eusebius,
Theodorus of Heraclea, Narcissus of Neronias in Cilicia, Stephanus of
Antioch, George of Laodicea, Acacius of Cæsarea in Palestine,
Menophantus of Ephesus in Asia, Ursacius of Singidonum in Mysia, and
Valens of Mursia in Pannonia \footnote{Vid. supr. p. 31, note M. p. 60. ref. 4.
\&c. vol. 8. p. 74, note D. About Stephanus, vid. infr. Hist.
Arian. §. 20.

Return to text}. These men would not permit
those who came with them from the East to meet the holy Council, nor
even to approach the Church of God, but as they were coming to Sardica,
they held Councils in various places by themselves, and made an
engagement under threats, that when they came to Sardica, they would
not at all appear at the trial, nor attend the assembling of the holy
Council, but simply coming, and making known their arrival as a matter
of form, would speedily take to flight. This we have been able to
ascertain from our brethren in ministry, Macarius of Palestine and
Asterius of Arabia \footnote{These two Bishops were soon after the Council
banished by Eusebian influence into upper Libya, where they suffered
extreme ill usage. vid. infr. Hist. Arian. §. 18.

Return to text}, who after coming in their company,
separated themselves from their unbelief. These came to the holy
Council, and complained of the violence they had suffered, and said
that no orthodox act proceeded from them; adding that there were many
among them who adhered to the true doctrine, but were prevented
by those men from coming hither, by means of the threats and promises
which they held out to those who wished to separate from them. On this
account it was that they were so anxious that all should abide in one
dwelling, and would not suffer them to be by themselves even for the
shortest space of time.

§. 49.

28. Since then it became us not to hold our
peace, nor to pass over unnoticed their calumnies, imprisonments,
murders, scourgings, conspiracies by means of forged letters,
outrages, stripping of the virgins, banishments, destruction of the
Churches, burnings, translations from small cities to larger dioceses,
and above all, the rising of the accursed Arian heresy by their means
against the orthodox faith; we have therefore pronounced our dearly
beloved brethren and fellow-ministers Athanasius, Marcellus, and
Asclepas, and those who minister to the Lord with them, to be innocent
and clear of offence, and have written to the diocese of each, that
the people of each Church may know the innocence of their own Bishop,
and may esteem him as their Bishop and expect his coming.

29. And as for those who like wolves \footnote{vid. Acts xx. 29.

Return to text}
have invaded their Churches, Gregory at Alexandria, Basil at Ancyra,
and Quintianus at Gaza, let them neither give them the title of
Bishop, nor hold any communion at all with them, nor receive letters \footnote{p. 8. ref. 3.

Return to text} from them, nor write to them. And for Theodorus, Narcissus,
Acacius, Stephanus, Ursacius, Valens, Menophantus, and George,
although the last from fear did not come from the East, yet because he
was degraded by the blessed Alexander, and because both he and the
others were connected with the Arian fanaticism, as well as on account
of the charges which lie against them, the holy Council has
unanimously deposed them from the Episcopate, and we have decided that
they not only are not Bishops, but that they are unworthy of holding
communion with the faithful.

30. For they who separate the Son and alienate
the Word from the Father, ought themselves to be separated from the
Catholic Church and to be alien from the Christian name. Let them
therefore be anathema to you, because they have adulterated the word
of truth. It is an Apostolic injunction, \emph{If any man preach any
other Gospel unto you than that ye have} \emph{received, let him
be accursed} [Gal. i. 9.]. Charge your people that no one hold
communion with them, for there is no \emph{communion of light with
darkness}; put away from you all these, for there is no \emph{concord
of Christ with Belial} [2 Cor. vi. 14. 15.]. And take heed, dearly
beloved, that ye neither write to them, nor receive letters from them;
but desire rather, brethren and fellow-ministers, as being present in
spirit with our Council, to assent to our judgments by your
subscriptions \footnote{In like manner the Council of Chalcedon was
confirmed by as many as 470 subscriptions, according to Ephrem, (Phot.
Bibl. p. 801.) by 1600 according to Eulogius, (ibid. p. 877.) i.e. of
Bishops, Archimandrites, \&c.

Return to text}, to the end that concord may be preserved by
all our fellow-ministers every where. May Divine Providence protect
and keep you, dearly beloved brethren, in sanctification and joy.

I, Hosius, Bishop, have subscribed this, and all
the rest likewise.

31. This is the letter which the Council of
Sardica sent to those who were unable to attend, and they on the other
hand gave their judgment in accordance; and the following are the
names both of those Bishops who subscribed in the Council, and of the
others also.

§. 50.

Hosius of Spain \footnote{Hosius is called by Athan. the father and the
president of the Council. Hist. Arian. 15. 16. Roman controversialists
here explain why Hosius does not sign himself as the Pope's legate, De
Marc. Concord. v. 4. Alber. Dissert. ix. and Protestants why his
legates rank before all the other Bishops, even before Protogenes,
Bishop of the place. Basnage, Ann. 347. 5. Febronius considers that
Hosius signed here and at Nicæa, as a sort of representative of the
civil, and the Legates of the ecclesiastical supremacy. de Stat. Eccl.
vi. 4. And so Thomassin, ``Imperator velut exterior Episcopus: præfuit
autem summus Pontifex, ut Episcopus interior.'' Dissert. in Conc. x.
14. The Pope never attended in person the Eastern Councils. St. Leo
excuses himself on the plea of its being against usage. Epp. 37. and
93.

Return to text}, Julius of Rome by his
Presbyters Archidamus and Philoxenus, Protogenes of Sardica,
Gaudentius, Macedonius, Severus \footnote{of Ravenna.

Return to text}, Prælextatus \footnote{of Barcelona.

Return to text},
Ursicius \footnote{of
Brescia.

Return to text}, Lucillus \footnote{of Verona.

Return to text}, Eugenius, Vitalius, Calepodius,
Florentius \footnote{of
Merida.

Return to text}, Bassus, Vincentius
\footnote{of
Capua.

Return to text}, Stercorius,
Palladius, Domitianus, Chalbis, Gerontius, Protasius \footnote{of Milan.

Return to text},
Eulogus, Porphyrius \footnote{of Philippi.

Return to text}, Dioscorus, Zozimus, Januarius, Zozimus,
Alexander, Eutychius, Socrates, Diodorus, Martyrius, Eutherius,
Eucarpus, Athenodorus, Irenæus, Julianus, Alypius, Jonas, Aetius \footnote{of
Thesalonica.

Return to text}, Restitutus, Marcellinus, Aprianus, Vitalius, Valens,
Hermogenes, Castus, Domitianus, Fortunatius \footnote{of
Aquilea.

Return to text}, Marcus,
Annianus, Heliodorus, Musæus, Asterius, Paregonius, Plutarchus,
Hymenæus, Athanasius, Lucius, Amantius, Arius, Asclepius, Dionysius,
Maximus \footnote{of
Lucca.

Return to text}, Tryphon, Alexander, Antigonus, Ælianus, Petrus,
Symphorus, Musonius, Eutychus, Philologius, Spudasius, Zozimus,
Patricius, Adolius, Sapricius.

From Gaul the following; Maximianus \footnote{of
Treves.

Return to text},
Verissimus \footnote{of Lyons.

Return to text}, Victurus, Valentinus
\footnote{of Arles.

Return to text}, Desiderius,
Eulogius, Sarbatius, Dyscolius, Severinus \footnote{of
Sens.

Return to text}, Satyrus, Martinus,
Paulus, Optatianus, Nicasius, Victor \footnote{of Worms.

Return to text}, Sempronius, Valerinus,
Pacatus, Jesses, Ariston, Simplicius, Metianus, Amantus \footnote{of
Strasbourgh.

Return to text},
Amillianus, Justinianus, Victorinus \footnote{of Paris.

Return to text}, Saturnilus, Abundantius,
Donatianus, Maximus.

From Africa; Nessus, Gratus \footnote{of Carthage.

Return to text}, Megasius,
Coldæus, Rogatianus, Consortius, Rufinus, Manninus, Cessilianus,
Herennianus, Marianus, Valerius, Dynamius, Myzonius, Justus,
Celestinus, Cyprianus, Victor, Honoratus, Marinus, Pantagathus, Felix,
Bandius, Liber, Capito, Minervalis, Cosmus, Victor, Hesperio, Felix,
Severianus, Optantius, Hesperus, Fidentius, Salustius, Paschasius.

From Egypt; Liburnius, Amantius, Felix,
Ischyrammon, Romulus, Tiberinus, Consortius, Heraclides, Fortunatius,
Dioscorus, Fortunatianus, Bastamon, Datyllus, Andreas, Serenus, Arius,
Theodorus, Evagoras, Helias, Timotheus, Orion, Andronicus, Paphnutius,
Hermias, Arabion, Psenosiris, Apollonius, Muis, Sarapampon \footnote{p. 53, note O. and §. 78.

Return to text},
Philo, Philippus, Apollonius, Paphnutius, Paulus, Dioscorus, Nilammon,
Serenus, Aquila, Aotas, Harpocration, Isac, Theodorus, Apollos,
Ammonianus, Nilus, Heraclius, Arion, Athas, Arsenius, Agathammon,
Theon, Apollonius, Helias, Paninuthius, Andragathius, Nemesion,
Sarapion, Ammonius, Ammonius, Xenon, Gerontius, Quintus, Leonides,
Sempronianus, Philo, Heraclides, Hieracys, Rufus, Pasophius,
Macedonius, Apollodorus, Flavianus, Psaes, Syrus, Apphus, Sarapion,
Esaias, Paphnutius, Timotheus, Elurion, Gaius, Musæus, Pistus,
Heraclammon, Hero, Helias, Anagamphus, Apollonius, Gaius, Philotas,
Paulus, Tithoes, Eudæmon, Julius.

Those in the cross roads \footnote{[\emph{hoi en t\={o}i kanali\={o}i t\={e}s
Italias}]. ``Canalis est, non via regia aut militaris, verum via
transversa, quæ in regiam seu basilicam influit, quasi aquæ canalis
in alveum.'' Gothofred. in Cod. Theod. vi. de Curiosis, p. 196. who
illustrates the word at length. Du Cange on the contrary, \emph{in voc}.
explains it of ``the high road.'' Tillemont professes himself unable to
give a satisfactory sense to it. vol. viii. p. 685.

Return to text} of Italy are, Probatius,
Viator, Facundinus, Joseph, Numedius, Sperantius, Severus,
Heraclianus, Faustinus, Antoninus, Heraclius, Vitalius, Felix,
Crispinus, Paulianus.

From Cyprus; Auxibius, Photius, Gerasius,
Aphrodisius, Irenicus, Nunechius, Athanasius, Macedonius, Triphyllius,
Spyridon, Norbanus, Sosicrates.

From Palestine; Maximus, Aetius, Arius,
Theodosius, Germanus, Silvanus, Paulus, Claudius, Patricius, Elpidius,
Germanus, Eusebius, Zenobius, Paulus, Petrus.

These are the names of those who subscribed to
the acts of the Council; but there are very many beside, out of Asia,
Phrygia, and Isauria \footnote{p. 60. r. 9.

Return to text}, who wrote in my behalf before this
Council was held, and whose names, nearly sixty-three in number, may
be found in their own letters. They amount altogether to three hundred
and forty-four \footnote{There is great uncertainty what was the actual
number of Bishops present at the Council. Athan. Hist. Arian. §. 15.
says 170, while Theodoret names 250. Hist. ii. 6. If the Western
Bishops, whose signatures are given by Athan. in the text to the
number of 163, were all present, it might have been conjectured that
he was speaking of the Western only; but he expressly includes the
Eastern. In that case, subtracting the 73 or 80 Eusebians, so small a
majority of orthodox remains, that it is incredible, considering the
notorious dexterity and unscrupulousness of the Eusebians in Synodal
meetings, that they should have been obliged to secede. Athan. says,
supr. §. 1. that the Letter of the Council was signed in all by more
than 300. It will be observed, that Athan.'s numbers in the text do
not accurately agree with each other. The subscriptions enumerated are
284, to which 63 being added, make a total of 347, not 344.

Return to text

}.

\xsection{Chapter 4. Imperial and Ecclesiastical Acts in consequence of the decision of the Council of Sardica}

§. 51.

1. W\textsc{hen} the most religious
Emperor Constantius heard of these things, he sent for me, having
written privately to his brother Constans of blessed memory, and to me
three several times in the following terms.

2. Constantius Victor Augustus to Athanasius

Our benignant clemency will not suffer you to be
any longer tempest-tossed by the wild waves of the sea; for our
unwearied piety has not lost sight of you, while you have been bereft
of your native home, deprived of your goods, and have been wandering
in savage wildernesses. And although I have for a long time deferred
expressing by letter the purpose of my mind concerning you,
principally because I expected that you would appear before us of your
own accord, and would seek a relief of your sufferings; yet forasmuch
as fear, it may be, has prevented you from fulfilling your intentions,
we have therefore addressed to your fortitude letters full of our
bounty, to the end that you may use all speed and without fear present
yourself in our presence, thereby to obtain the enjoyment of your
wishes, and that, having experience of our grace, you may be restored
again to your friends. For this purpose I have besought my Lord and
brother Constans Victor Augustus in your behalf; that he would give
you permission to come, in order that you may be restored to your
country with the consent of us both, receiving this as a pledge of our
favour.

3. The Second Letter

Although we made it very plain to you in a former
letter that you may without hesitation come to our Court, because we greatly wished to send you home, yet, we have further sent this
present letter to your fortitude, to exhort you without any distrust
or apprehension, to place yourself in the public conveyances \footnote{Margin Notes

1. Gothof. in Cod. Theod. viii. 5. p. 507.

Return to text}, and to hasten to us, that you may enjoy the fulfilment of your
wishes.

4. The Third Letter

Our pleasure was, while we abode at Edessa, and
your Presbyters were there, that, on one of them being sent to you,
you should make haste to come to our Court, in order that you might
see our face, and straightway proceed to Alexandria. But as a long
period has elapsed since you received letters from us, and you have
not yet come, we are therefore desirous to remind you again, that you
may endeavour to present yourself before us with all speed, and so may
be restored to your country, and obtain the accomplishment of your
prayers. And for your fuller information we have sent Achitas the
Deacon, from whom you will be able to learn our earnest desires
concerning you, and that you may now secure the objects of your
prayers.

5. Such was the tenour of the Emperor's letter;
on receiving which I went up to Rome to bid farewell to the Church and
the Bishop: for I was at Aquileia when it was written. The Church was
filled with all joy, and the Bishop Julius rejoiced with me in my
return and wrote to the Church \footnote{``They acquainted Julius the Bishop of Rome
with their case and he, according to the prerogative ([\emph{pronomia}])
of the Church in Rome, fortified them with letters in which he spoke
his mind, and sent them back to the East, restoring each to his own
place, and remarking on those who had violently deposed them. They
then set out from Rome, and on the strength ([\emph{tharrhountei}]) of
the letters of Bishop Julius, take possession of their Churches.'' Socr.
ii. 15. It must be observed, that in the foregoing sentence Socrates
has spoken of ``\emph{imperial} Rome,'' Sozomen says, ``Whereas the care
of all ([\emph{k\={e}demonias}]) pertained to him on account of the
dignity of his see, he restored each to his own Church, iii. 8. ``I
answer,'' says Barrow, ``the Pope did not restore them \emph{judicially},
but \emph{declaratively}, that is, declaring his approbation of their
right and innocence, did admit them to communion ... Besides, the Pope's
proceeding was taxed, and protested against, as irregular; … and,
lastly, the restitution of Athanasius and the other Bishops had no
complete effect, till it was confirmed by the synod of Sardica, backed
by the imperial authority.'' Suprem. p. 360. ed. 1836.

Return to text};
and as I passed along, the Bishops of every place sent me on my way in
peace. The letter of Julius was as follows.

§. 52.

6. Julius to the Presbyters, Deacons, awl people
abiding at Alexandria.

I congratulate you, beloved brethren, that you
now behold the fruit of your faith before your eyes; for any one may
see that such indeed is the case with respect to my brother and
fellow-Bishop Athanasius, whom for the innocency of his life, and by
reason of your prayers, God hath restored to you again. Wherefore it
is easy to perceive, that you have continually offered up to God pure
prayers and full of love. Being mindful of the heavenly promises, and
of the conversation that leads to them, which you have learnt from the
teaching of this my brother, you knew certainly and were persuaded by
the right faith that is in you, that he, whom you always had as
present in your most pious minds, would not be separated from you for
ever. Wherefore there is no need that I should use many words in
writing to you; for your faith has already anticipated whatever I
could say to you, and has by the grace of God procured the
accomplishment of the common prayers of you all. Therefore, I repeat
again, I congratulate you, because you have preserved your souls
unconquered in the faith; and I also congratulate no less my brother
Athanasius, in that, though he has endured many afflictions, he has at
no time been forgetful of your love and earnest desires towards him.
For although for a season he seemed to be withdrawn from you in body,
yet has he continued to live as always present with you in spirit \footnote{Athan. here omits a paragraph in his own praise. vid. Socr. ii. 23.

Return to text}.

§. 53.

7. Wherefore he returns to you now more
illustrious than when he went away from you. Fire tries and purifies
the precious metals, gold and silver: but how can one describe the
worth of such a man, who, having passed victorious through the perils
of so many tribulations, is now restored to you, being pronounced
innocent not by my voice only, but by the voice of the whole Council \footnote{p. 56, note S. p. 80, note P.

Return to text}? Receive therefore, dearly beloved brethren, with all godly
honour and rejoicing, your Bishop Athanasius, together with those who
have been partners with him in so many labours. And rejoice that you
have now obtained the fulfilment of your prayers, after that in your
salutary writings, you have given meat and drink to your Pastor, who,
so to speak, longed and thirsted after your godliness. For while
he sojourned in a foreign land, you were his consolation; and you
refreshed him during his persecutions by your most faithful minds and
spirits. And it delights me now to conceive and figure to my mind the
joy of every one of you at his return, and the pious greetings of the
multitude, and the glorious festivity of those that run to meet him.
What a day will that be to you, when my brother comes back again, and
your former sufferings terminate, and his much-prized and desired
return inspires you all with an exhilaration of perfect joy! The like
joy it is mine to feel in a very great degree, since it has been
granted me by God, to be able to make the acquaintance of so eminent a
man.

8. It is fitting therefore that I should conclude
my letter with a prayer \footnote{[\emph{euch\={e}n}].

Return to text}.
May Almighty God, and His Son our Lord and Saviour Jesus Christ,
afford you continual grace, giving you a reward for the admirable
faith which you displayed in your noble confession in behalf of your
Bishop, that He may impart unto you and unto them that are with you,
both here and hereafter, those better things, which \emph{eye hath not
seen, nor ear heard, neither hath entered into the heart of man, the
things which God hath prepared for them that love Him} [1 Cor. ii.
9.]; through our Lord Jesus Christ, through whom to Almighty God be
glory for ever and ever. Amen. I pray dearly beloved brethren, for
your health and strength in the Lord.

§. 54.

9. The Emperor, when I came to him with these
letters, received me kindly, and sent me forward to my country and
Church, addressing the following to the Bishops, Presbyters and
People.

10. Victor Constantius, Maximus, Augustus, to the
Bishops and Presbyters of the Catholic Church

The most reverend Athanasius has not been
deserted by the grace of God, but although for a brief season he was
subjected to trials to which human nature is liable, he has obtained
from the superintending Providence such an answer to his prayers as
was meet, and is restored by the will of the Most High, and by our
sentence, at once to his country and to the Church, over which by
divine permission he presided. Wherefore, in accordance with
this, it is fitting that it should be provided by our clemency, that
all the decrees which have heretofore been passed against those who
held communion with him, be now consigned to oblivion, and that all
suspicions respecting them be henceforward set at rest, and that an
immunity, such as the Clergy who are associated with him formerly
enjoyed, be duly confirmed to them. Moreover to our other acts of
favour towards him we have thought good to add the following, that all
persons of the sacred catalogue \footnote{vid. Bingh. Antiqu. i. 5. §. 10.

Return to text} should understand, that an assurance of safety is given to all
who adhere to him, whether Bishops, or other Clergy. And union with
him will be a sufficient guarantee, in the case of any person, of an
upright intention. For whoever, acting according to a better judgment
and part, shall choose to hold communion with him, we order, in
imitation of that Providence which has already gone before, that all
such should have the advantage of the grace which by the will of the
Most High is now offered to them from us. May God preserve you.

11. The Second Letter

Victor Constantius, Maximus, Augustus, to the
people of the Catholic Church at Alexandria.

§. 55.

Desiring as we do your welfare in all respects,
and knowing that you have for a long time been deprived of episcopal
superintendence, we have thought good to send back to you your Bishop
Athanasius, a man known to all men for the uprightness that is in him,
and for his personal deportment. Receive him, as you are wont to
receive every one, in a suitable manner, and, putting him forth as
your succour in your prayers to God, endeavour to preserve continually
that unanimity and peace according to the order of the Church, which
is at the same time becoming in you, and most advantageous for us. For
it is not becoming that any dissension or faction should be raised
among you, so subversive of the prosperity of our times. We desire
that this offence may be altogether removed from you, and we exhort
you to continue stedfastly in your accustomed prayers, and to make
him, as we said before, your advocate and helper towards God. So that,
when this your determination, dearly beloved, has influenced the
prayers of all men, even the heathen who are still addicted to the
false worship of idols may eagerly desire to come to the knowledge of
our sacred worship.

12. Again therefore we exhort you to continue in
these timings, and gladly to receive your Bishop, who is sent back to
you by the decree of the Most High, and by our desire, and determine
to greet him cordially with all your soul and with all your mind. For
this is what is both becoming in you, and agreeable to our clemency.
In order that all occasion of excitement and sedition may be taken
away from those who are maliciously disposed, we have by letter
commanded the magistrates who are among you to subject to the
vengeance of the law all whom they find to be factious. Wherefore
taking into consideration both these things, our desire in accordance
with the will of the Most High, and our regard for you and for concord
among you, and the punishment that awaits the disorderly, observe such
things as are proper and suitable to the order of our sacred religion,
and receiving the fore-mentioned Bishop with all reverence and honour,
take care to offer up with him your prayers to God, the Father of all,
in behalf of yourselves, and for the well-being of your whole lives.

§. 56.

13. Having written these letters, he also
commanded that the decrees, which he had formerly sent out against me
in consequence of the calumnies of the Eusebians, should be abolished,
and removed from out the Orders of the Duke and the Prefect of Egypt;
and Eusebius the Decurion \footnote{member of the Curia or Council.

Return to text}
was sent to withdraw them from the Order-books. His letter on this
occasion was as follows.

11. Victor, Constantius, Augustus, to Nestorius \footnote{Prefect of Egypt, vid. p. 5, note D.

Return to text}

(And in the same terms, to the Governors of
Augustamnica, the Thebais, and Libya.)

Whatever Orders are found to have been passed
heretofore, tending to the injury and dishonour of those who hold
communion with the Bishop Athanasius, we wish them to be now erased.
For we desire that whatever immunities his Clergy possessed before,
they should again possess the same. And we wish this our Order to
be observed, that when the Bishop Athanasius is restored to his
Church, those who hold communion with him may enjoy the immunities
which they have always enjoyed, and which the rest of the Clergy
enjoy; so that they may have the satisfaction of being on an equal
footing with others.

§. 57.

15. Being thus set forward on my journey, as I
passed through Syria, I met with the Bishops of Palestine, who when
they had called a Council \footnote{Hist. Arian. 25.

Return to text}
at Jerusalem, received me courteously, and themselves also sent me on
my way in peace, and addressed the following letter to the Church and
the Bishops.

16. The Holy Council, assembled at Jerusalem, to
the brethren in ministry in Egypt and Libya, and to the Presbyters,
Deacons, and People at Alexandria, dearly beloved brethren, and
greatly longed for, sends health in the Lord.

We cannot give worthy thanks to the God of all,
dearly beloved, for the wonderful things which He has done at all
times, and especially at this time with respect to your Church, in
restoring to you your pastor and lord \footnote{[\emph{kurion}], infr. p. 86.

Return to text}, and our fellow-minister Athanasius. For who ever hoped that
his eyes would see what you are now actually enjoying? Of a truth,
your prayers have been heard by the God of all, who cares for His
Church, and has looked upon your tears and groans, and has therefore
heard your petitions. For ye were as sheep scattered and fainting, not
having a shepherd. Wherefore the true Shepherd, who careth for His own
sheep, has visited you from heaven, and has restored to you him whom
you desire. Behold, we also, being ready to do all things for the
peace of the Church, and being prompted by the same affection as
yourselves, have saluted him before you; and communicating with you
through him, we send you these greetings, and our offering of
thanksgiving, that you may know that you are united in one bond of
love with him and with us. You are bound to pray also for the piety of
our most religious Emperors, who, when they knew your earnest longings
after him, and his innocency, determined to restore him to you with
all honour. Wherefore receive him with uplifted hands, and take good
heed that you offer up due thanksgivings on his behalf to God who has
bestowed these blessings upon you; so that you may continually rejoice
with God and glorify our Lord, in Christ Jesus our Lord, through
whom to the Father be glory for ever. Amen.

17. I have set down here the names of those who
subscribed this letter, although I have mentioned them before \footnote{p. 78.

Return to text}. They are these; Maximus, Aetius, Arius, Theodorus \footnote{Theodosius, supr.

Return to text}, Germanus, Silvanus, Paulus, Patricius, Elpidius, Germanus,
Eusebius, Zenobius, Paulus, Macrinus \footnote{not supr.

Return to text}, Petrus, Claudius.

§. 58.

18. When Ursacius and Valens witnessed these
proceedings, they forthwith condemned themselves for what they had
done, and going up to Rome, confessed their crime, declared themselves
penitent, and sought forgiveness \footnote{vid. p. 15, note F.

Return to text}, addressing the following letters to Julius Bishop of ancient
Rome, and to myself. Copies of them were sent to me from Paulinus,
Bishop of Tibur \footnote{[\emph{Tiber\={o}n}], Paul infr. p. 239. Paulinus, supr. p. 78?

Return to text}.

19. A Translation from the Latin of a Letter \footnote{Hist Arian. 25. 26.

Return to text} to Julius, concerning the recantation of Ursacius and Valens \footnote{``I have always entertained some doubts,'' says Gibbon, ``concerning the
retractation of Ursacius and Valens. Their Epistles to Julius Bishop
of Rome, and to Athanasius himself, are of so different a cast from
each other, that they cannot both be genuine. The one speaks the
language of criminals, who confess their guilt and infamy; the other
of enemies, who solicit on equal terms an honourable reconciliation.''
ch. xxi. note 118. Surely this is just the difference of tone in which
an apology is made to a superior, and to an equal ([\emph{adelph\={o}i}]).
except by very generous, or by deeply repentant, persons. Athan.'s
account of it, infr. p. 239, r. 2. is quite in accordance. It will be
observed too that they appear to have made their peace with Rome with
the view of being defended by the Pope against Athanasius.

Return to text

}

Ursacius and Valens to the most blessed Lord \footnote{[\emph{kuri\={o}i}], infr. p. 87.

Return to text}, Pope Julius.

Whereas it is well known that we have heretofore
in letters laid many grievous charges against the Bishop Athanasius,
and whereas, when we were corrected by the letters of your Goodness \footnote{[\emph{chr\={e}stot\={e}ti}].

Return to text}, we were unable to render an account of our conduct, by reason
of the circumstance which we notified unto you; we do now confess
before your Goodness, and in the presence of all the Presbyters our
brethren, that all the reports which have heretofore come to your
hearing respecting the case of the aforesaid Athanasius, are
falsehoods and fabrications, and are utterly inconsistent with his
character. Wherefore we earnestly desire communion with the aforesaid
Athanasius, especially since your Piety, with your characteristic
generosity, has vouchsafed to pardon our error. But we also
declare, that if at any time the Eastern Bishops, or Athanasius
himself, with an evil intent, should wish to bring us to judgment for
this offence, we will not attend contrary to your judgment and desire.
And as for the heretic Arius and his supporters, who say that once the
Son was not, and that the Son is made of that which was not, and who
deny that Christ is God \footnote{not in Latin.

Return to text}
and the Son of God before the worlds, we anathematize them both now
and for evermore, as also we set forth in our former declaration at
Milan \footnote{A.D.
346, 7, or 8.

Return to text}. We have
written this with our own hands, and we profess again, that we have
renounced for ever, as we said before, the Arian heresy and its
authors.

I Ursacius subscribed this my confession in
person; and likewise I Valens.

20. Ursacius and Valens, Bishops, to their Lord \footnote{[\emph{kuri\={o}i}], vid. infr. p. 95.

Return to text} and Brother, the Bishop Athanasius

Having an opportunity of sending by our brother
and fellow Presbyter Musæus, who is coming to your Charity, we salute
you affectionately, dearly beloved brother, through him, from Aquileia,
and pray you, being as we trust in health, to read our letter. You
will also give us confidence, if you will return to us an answer in
writing. For know that we are at peace with you, and in communion with
the Church, of which the salutation prefixed to this letter is a
proof. May Divine Providence preserve you, my Lord \footnote{[\emph{kurie}].

Return to text}, our dearly beloved brother!

21. Such were their letters, and such the
sentence and the judgment of the Bishops in my behalf. But in order to
prove that they did not act thus to ingratiate themselves, or under
compulsion \footnote{p. 15, note F.

Return to text}, in any
quarter, I desire, with your permission, to recount the whole matter
from the beginning, so that you may perceive that the Bishops wrote as
they did with upright and just intentions, and that Ursacius and
Valens, though they were slow to do so, at last confessed the truth.

\xsection{Chapter 5. Documents connected with the Charges of the Meletians against St. Athanasius}

§. 59. \footnote{Margin Notes

1. Second part of Apology.

Return to text}

1. P\textsc{eter} was Bishop among
us before the persecution, and during the course of it he suffered
martyrdom. When Meletius, who held the title of Bishop in Egypt, was
convicted of many crimes, and among the rest of offering sacrifice to
idols, Peter deposed him in a general Council of the Bishops.
Whereupon Meletius did not appeal to another Council, or attempt to
justify himself before those who should come after, but made a schism,
so that they who espoused his cause are even yet called Meletians
instead of Christians \footnote{vol. viii. p. 180, note f.

Return to text}.
He began immediately to revile the Bishops, and made false
accusations, first against Peter himself, and after him against
Achillas, and after Achillas against Alexander \footnote{ad Ep. Æg. §. 22. supr. p. 29.

Return to text}. And he thus practised craftily, following the example of
Absalom, to the end that, as he was disgraced by his deposition, he
might by his calumnies mislead the minds of the simple. While Meletius
was thus employed, the Arian heresy arose, and in the Council of Nicæa,
when that heresy was anathematized, and the Arians were
excommunicated, the Meletians on whatever grounds \footnote{Meletius had the name of Bishop secured to
him, but was interdicted from all Episcopal functions. Those who had
been ordained by him were received to communion and allowed to
continue in ministerial duties, on condition that they gave precedence
in their own Church or Diocese to those whom Alexander had ordained,
and performed no ecclesiastical act without leave of the Catholic
Bishop; but when the Catholic Bishop in each place died, they were to
be considered capable of succeeding. Athan. speaks more openly against
this arrangement. infr. §. 71. vid. vol. viii. p. 181, note g.

Return to text} (for it is not necessary now to mention the reasons of this
proceeding) were received into the Church. Five months however had not
elapsed when the blessed Alexander died, and the Meletians, who ought
to have remained quiet, and to have been grateful that they were
received on any terms, like dogs unable to forget their vomit [vid. 2
Pet. ii. 22.], began again to trouble the Churches.

2. Upon learning this, Eusebius, who had the lead
in the Arian heresy, sends and bribes the Meletians with large
promises, becomes their secret friend, and arranges with them for
their assistance on any occasion when he might wish for it. At first
he sent to me, urging me to admit the Arians to communion \footnote{ad Ep. Æg. 19.

Return to text}, and threatening me in his verbal communications, which he
requested me in his letters. And when I refused, declaring that it was
not right that those who had invented heresy contrary to the truth,
and had been anathematized by the Ecumenical \footnote{supr. §. 7. and vol. 8. p. 49, note o.

Return to text} Council, should be admitted to communion, he caused the Emperor
also, Constantine, of blessed memory, to write to me, threatening me,
in case I should not receive the Arians, with those afflictions, which
I have before undergone, and which I am still suffering. The following
is a part of his letter. Syncletius and Gaudentius, officers of the
palace \footnote{[\emph{palatinoi}], vid. Apol. Ad Const. §. 19.

Return to text}, were the
bearers of it.

3. Part of a Letter from the Emperor Constantine

Having therefore knowledge of my will, grant free
admission to all who wish to enter into the Church. For if I learn
that you have hindered or excluded any who claim to be admitted into
communion with the Church, I will immediately send some one who shall
depose you by my command, and shall remove you from your place.

§. 60.

4. When upon this I wrote and endeavoured to
convince the Emperor, that that anti-Christian \footnote{[\emph{christomach\={o}i}], vol. 8. p. 6, note n.

Return to text} heresy had no communion with the Catholic Church, Eusebius
forthwith, availing himself of the occasion which he had agreed upon
with the Meletians, writes and persuades them to invent some pretext,
so that, as they had practised against Peter and Achillas and
Alexander, they might also lay a plot for me, and might spread abroad
reports to my prejudice. Accordingly, after seeking for a long time,
and finding nothing, they at last agree together, with the advice of
the Eusebians, and fabricate their first accusation by means of Ision,
Eudæmon, and Callinicus \footnote{infr. §. 71 fin.

Return to text},
respecting the linen vestments \footnote{[\emph{sticharia}], ecclesiastical, vid. Du Cange.

Return to text}, to the effect that I had imposed a law upon the Egyptians, and
had required its observance of them first. But when certain Presbyters
of mine were found to be present, and the Emperor took cognizance of
the matter, they were condemned, (the Presbyters were Apis and
Macarion,) and the Emperor wrote, condemning Ision, and ordering me to
appear before him. His letters were as follows \footnote{they are lost.

Return to text}. * *
*

5. Eusebius, having intelligence of this,
persuades them to wait; and when I arrive, they next accuse Macarius
of breaking the chalice, and bring against me the most heinous
accusation possible, viz. that, being an enemy of the Emperor, I had
sent a purse of gold to one Philamenus. The Emperor therefore heard us
on this charge also in Psammathia \footnote{suburb of Nicomedia, infr. §. 65.

Return to text}, when they, as usual, were condemned, and driven from the
presence; and, as I returned, he wrote the following letter to the
people.

6. Constantine Maximus, Augustus, to the people
of the Catholic Church at Alexandria

§. 61.

Dearly beloved brethren, I greet you well,
calling upon God, who is the chief witness of my good-will towards
you, and on the Only-begotten, the Author of our Law, who is Sovereign
over the lives of all men, and who hates dissensions. But what shall I
say to you? That I am in good health? Nay, but I should be able to
enjoy better health and strength, if you were possessed with mutual
love one towards another, and had rid yourselves of your enmities,
through which, in consequence of the storms excited by contentious
men, we have left the haven of brotherly love. Alas! what perverseness
is this! What evil consequences are produced every day by the tumult
of envy which has been stirred up among you! hence it is that an evil
character attaches to the people of God. Whither has the faith of
righteousness departed? For we are so involved in the mists of
darkness, not only through manifold errors, but through the faults of
ungrateful men, that we bear with those who favour folly, and though
we are aware of them, take no heed of those who beat down goodness and
truth. What strange inconsistency is this! We do not convict our
enemies, but we follow the example of robbery which they set us,
whereby the most pernicious errors, finding no one to oppose them,
easily, if I may so speak, make a way for themselves. Is there no
understanding among us, for the credit of our common nature,
since we are thus neglectful of the injunctions of the Law?

7. But some one will say, that that mutual love
which nature prompts is exercised among us. But, I ask, how is it that
we who have the law of God for our guide, in addition to the light of
nature, thus tolerate the disturbances and disorders raised by our
enemies, who set every thing in a flame, as it were, with firebrands?
How is it, that having eyes, we see not, neither understand, though we
are surrounded by the intelligence of the law? What a stupor has
seized upon our senses, that we are thus neglectful of ourselves,
although God admonishes us of these things! Is it not an intolerable
calamity? and ought we not to esteem such men as our enemies, and not
the household and people of God? For they are infuriated against us,
desperate as they are: they lay grievous crimes to our charge, and
persecute us as enemies.

§. 62.

8. And I would have you yourselves to consider
with what exceeding madness they do this. The foolish men carry their
maliciousness at their tongues’ end. They carry about with them a
sort of sullen anger, so that, by way of retaliation, they smite one
another, and give us a share in the punishment which they inflict upon
themselves. The good teacher is accounted an enemy, while he who
clothes himself with the vice of envy, contrary to all justice makes
his gain of the gentle temper of the people; he ravages, and consumes,
he decks himself out, and recommends himself with false praises; he
subverts the truth, and corrupts the faith, until he finds out a hole
and hiding place for his conscience. Thus their very perverseness
makes them wretched, while they impudently prefer themselves to places
of honour, however unworthy they may be. Ah! what a mischief is this!
they say, “Such an one is too old; such an one is a mere boy; the
office belongs to me; it is due to me, since it is taken away from
him. I will gain over all men to my side, and then I will endeavour
with my power to ruin him.” Plain indeed is this proclamation of
their madness to all the world; the sight of companies, and
gatherings, and rowers under command \footnote{[\emph{archieresian}].

Return to text} in their offensive cabals. Alas! what preposterous conduct is
ours, if I may say it! Do they make an exhibition of their folly in
the Church of God? And are they not yet ashamed of themselves? Do
they not yet blame themselves? Are they not smitten in their
consciences, so that they now at length shew that they entertain a
proper sense of their deceit and contentiousness? Theirs is the mere
force of envy, supported by those baneful influences which naturally
belong to it. But those wretches have no power against your Bishop.
Believe me, brethren, their endeavours will have no other effect than
this, after they have worn down our days, to leave to themselves no
place of repentance in this life.

9. Wherefore I beseech you, lend help to
yourselves; receive kindly our love, and with all your strength drive
away those who desire to obliterate from among us the grace of
unanimity; and looking unto God, love one another. I received
graciously your Bishop Athanasius, and addressed him in such a manner,
as being persuaded that he was a man of God. It is for you to
understand these things, not for me to judge of them. I thought it
becoming that the most Reverend Athanasius himself should convey my
salutation to you, knowing his kind care of you, which, in a manner
worthy of that peaceable faith which I myself profess, is continually
engaged in the good work of declaring saving knowledge, and will be
furnished with a word of exhortation for you. May God preserve you,
dearly beloved brethren.

Such was the letter of Constantine.

§. 63.

10. After these occurrences the Meletians
remained quiet for some time, but afterwards shewed their hostility
again, and contrived the following plot, with the aim of pleasing
those who had hired their services. The Mareotis is a region of
Alexandria, in which Meletius was not able to make a schism. Now while
the Churches still existed within their appointed limits, and all the
Presbyters had congregations in them, and while the people were living
in peace, a certain person named Ischyras \footnote{supr. p. 30. 48. 62.

Return to text}, who was not a Clergyman, but depraved in his habits,
endeavoured to lead astray the people of his own village, declaring
himself to be a Clergyman. Upon learning this, the Presbyter of the
place, informed me of it when I was going through my visitation of the
Churches, and I sent Macarius the Presbyter with him to summon
Ischyras. The found him sick and lying in his cell, and charged
his father to admonish his son not to continue any such practices as
had been reported against him. But when he recovered from his
sickness, being prevented by his friends and his father from pursuing
the same course, he fled over to the Meletians; and they communicate
with the Eusebians, and at last that calumny is invented by them, that
Macarius had broken a chalice, and that a certain Bishop named
Arsenius had been murdered by me. Arsenius they placed in concealment,
in order that he might seem taken off, when he did not make his
appearance; and they carried about a hand pretending that he had been
cut to pieces. As for Ischyras, whom they did not even know, they
began to spread a report that he was a Presbyter, in order that what
he said about the chalice might mislead the people. Ischyras, however,
being censured by his friends, came to me weeping, and said that no
such thing as they had reported had been done by Macarius, and that
himself had been suborned by the Meletians to invent this calumny. And
he wrote the following letter.

§. 64.

11. To the Blessed Pope \footnote{vid. vol. 8. p. 96, note g.

Return to text} Athanasius, Ischyras sends health in the Lord.

As when I came to you, my Lord \footnote{[\emph{kurie}], supr. p. 86.

Return to text} Bishop, desiring to be received into the Church, you reproved
me for what I formerly said, as though I had proceeded to such lengths
of my own free choice, I therefore submit to you this my apology in
writing, in order that you may understand, that violence was used
towards me, and blows inflicted on me by Isaac and Heraclides, and
Isaac of Letopolis, and those of their party. And I declare, and take
God as my witness in this matter, that of none of the things which
they have stated, do I know you to be guilty. For no breaking of a
chalice or overturning of the holy Table ever took place, but they
compelled me by their violent usage to assert all this. And this
defence I make and submit to you in writing, desiring and claiming for
myself to be admitted among the members of your congregation. I pray
that you may have health in the Lord.

12. I submit this my handwriting to you the
Bishop Athanasius in the presence of the Presbyters, Ammonias of
Dicella, Heraclius of Phascus, Boccon of Chenebris, Achillas of
Myrsine, Didymus of Taphosiris, and Justus from Bomotheus; and of the
Deacons, Paul, Peter, and Olympius, of Alexandria, and Ammonius,
Pistus, Demetrius, and Gaius, of the Mareotis.

§. 65.

13. Notwithstanding this statement of Ischyras,
they again spread abroad the same charges against me every where, and
also reported them to the Emperor Constantine. He had heard before of
the affair of the chalice in Psammathia \footnote{vid. §. 60.

Return to text}, when I was there, and had detected the falsehood of my
enemies. But now he wrote to Antioch to Dalmatius \footnote{Dalmatius is as the name of father and son, the brother and nephew of
Constantine, Socrates, Hist. i. 27. gives the title of Censor to the
son; but the Alexandrian Chronicon (according to Tillemont, Empereurs,
vol. 4. p. 657.) gives it to the father. Valesius, and apparently
Tillemont, think Socrates mistaken. The younger Dalmation was created
Cæsar by Constantine a few years before his death; and, as well as
his brother Hannibalian, and a number of other relatives, was put to
death by Constantius, or his ministers and the soldiery, on the death
of his father. vid. Athan. Hist. Mon. 69.

Return to text} the Censor, requiring him to institute a judicial enquiry
respecting the murder. Accordingly the Censor sent me notice to
prepare for my defence against the charge. Upon receiving his letters,
although at first I paid no regard to the thing, because I knew that
nothing of what they said was true, yet seeing that the Emperor was
moved, I wrote to my brethren in Egypt, and sent a deacon, desiring to
learn something of Arsenius, for I had not seen the man for five or
six years. Well, not to relate the matter at length, Arsenius was
found in concealment, in the first instance in Egypt, and at last my
friends discovered him still in concealment at Tyre. And what was most
remarkable, even when it was discovered he would not confess that he
was Arsenius, until he was convicted in court before Paul, who was
then Bishop of Tyre, and at last out of very shame he could not deny
it.

14. This he did in order to fulfil his contract
with the Eusebians, lest, if he were discovered, the game they were
playing should at length be broken up; which in fact came to pass. For
when I wrote the Emperor word, that Arsenius was discovered, and
reminded him of what he had heard in Psammathia concerning Macarius
the Presbyter, he stopped the proceedings of the Censor’s court, and
wrote condemning the proceedings against me as calumnious, and
commanded the Eusebians to return, who were coming into the East
to appear against me. Now in order to shew that they accused me of
having murdered Arsenius, (not to bring forward the letters of many
persons on the subject,) it shall be sufficient only to produce one
from Alexander the Bishop of Thessalonica, from which the tenor of the
rest may be inferred. He then being acquainted with the reports which
Archaph, who is also called John, circulated against me on the subject
of the murder, and having heard that Arsenius was alive, wrote as
follows.

15. Letter of Alexander

To his dearly beloved son and brother
like-minded, the Lord \footnote{[\emph{kuri\={o}i}], supr. p. 93.

Return to text}
Athanasius, Alexander the Bishop sends health in the Lord.

§. 66.

I congratulate the most excellent Serapion, that
he is striving so earnestly to adorn himself with holy habits, and is
thus advancing to higher praise the memory of his father. For, as the
Holy Scripture somewhere says, \emph{though his father die, yet he is as
though he were not dead}[Ecclus. xxx. 4.]: for he has left behind
him a memorial of his life. What my feelings are towards the
ever-memorable Sozon, you yourself, my lord \footnote{[\emph{despota}]. Theod. Hist. i. 5. init.

Return to text}, are not ignorant, for you know the sacredness of his memory,
as well as the excellent disposition of the young man. I have received
only one letter from your reverence, which I had by the hands of this
youth. I mention this to you, my lord, that you may know that I have
received it. Our dearly beloved brother and deacon Macarius, afforded
me great pleasure by writing to me from Constantinople, that the false
accuser Archaph had met with disgrace, for having given out before all
men that a live man had been murdered. That he will receive from the
righteous Judge, together with all the tribe of his associates, that
punishment which his crimes deserve, the infallible Scriptures assure
us. May the Lord of all preserve you for very many years, my most
excellent lord \footnote{[\emph{kurie}].

Return to text}.

§. 67.

16. And they who lived with Arsenius bear
witness, that he was kept in concealment for this purpose, that they
might pretend his death; for in searching after him we found the
following person, and he in consequence wrote the following letter to John, who supported this false accusation against me.

17. To his dearly beloved brother John, Pinnes,
Presbyter of the Monastery of Ptemencyrcis, in the district of
Anteopolis, sends greeting.

I wish you to know, that Athanasius sent his
deacon into the Thebais, to search every where for Arsenius; and
Pecysius the Presbyter, and Sylvanus the brother of Helias, and
Tapenacerameus, and Paul monk of Hypsele, whom he first fell in with,
confessed that Arsenius was with us. Upon learning this we caused him
to be put on board a vessel, and to sail to the lower countries with
Helias the monk. Afterwards the deacon returned again suddenly with
certain others, and entered our monastery, in search of the same
Arsenius, and him they found not, because, as I said before, we had
sent him away to the lower countries; but they conveyed me together
with Helias the monk, who took him out of the way, to Alexandria, and
brought us before the Duke \footnote{According to the system of government introduced by Dioclesian and
Constantine, there were thirty-five military commanders of the troops,
under the Magistri militum, and all of these bore the name of duces or
dukes; the comites, or counts, were ten out of the number, who were
distinguished as companions of the Emperor. vid. Gibbon, ch. 17. Three
of these dukes were stationed in Egypt.

Return to text};
when I was unable to deny, but confessed that he was alive, and had
not been murdered: the monk also who took him out of the way confessed
the same. Wherefore I acquaint you with these things, Father, lest you
should determine to accuse Athanasius; for I said that he was alive,
and had been concealed with us, and all this is become known in Egypt,
and it cannot any longer be kept secret.

I, Paphnutius, monk of the same monastery \footnote{[\emph{mon\={e}s}].

Return to text}, who wrote this letter, heartily salute you. I trust that you
are in health.

18. The following also is the letter which the
Emperor wrote when he learnt that Arsenius was found to be alive.

19. Victor, Constantine, Maximus, Augustus, to
the Pope \footnote{vid. supr. p. 93.

Return to text} Athanasius.

§. 68.

Having read the letters of your wisdom, I felt
the inclination to write in return to your gravity, and to exhort you
that you would endeavour to restore the people of God to tranquillity, and to merciful feelings. For in my own mind I hold
these things to be of the greatest importance, that we should
cultivate truth, and ever keep righteousness in our thoughts, and have
pleasure especially in those who walk in the right way of life. But as
concerning those who are deserving of all execration, I mean the most
perverse and ungodly Meletians, who have at last stultified themselves
by their folly, and are now raising unreasonable commotions by envy,
uproar, and tumult, thus making manifest their own ungodly
dispositions, I will say thus much. You see that those who they
pretended had been slain with the sword, are still amongst us, and in
the enjoyment of life. Now what could be a stronger presumption
against them, and one so manifestly and clearly tending to their
condemnation, as that those whom they declared to have been murdered,
are yet in the enjoyment of life, and accordingly will be able to
speak for themselves?

20. But this further accusation was advanced by
these same Meletians. They positively affirmed that you, rushing in
with lawless violence, had seized upon and broken a chalice, which was
deposited in the most Holy Place; than which there certainly could not
be a more serious charge, nor a more grievous offence, had such a
crime actually been perpetrated. But what manner of accusation is
this? What is the meaning of this change and variation and difference
in the circumstances of it, insomuch that they now transfer this same
accusation to another person \footnote{pp. 48, 49.

Return to text}, a fact which makes it clearer, so to speak, than the light
itself that they designed to lay a plot for your wisdom? After this
who can be willing to follow them, men that have fabricated such
charges to the injury of another, seeing too that they are hurrying
themselves on to ruin, and are conscious that they are accusing you of
false and feigned crimes? Who then, as I said, will follow after them,
and thus go headlong in the way of destruction; in that way in which
it seems they alone suppose that they have hope of safety and of help?
But if they were willing to walk according to a pure conscience, and
to be directed by the best wisdom, and to go in the way of a sound
mind, they would easily perceive that no help can come to them from
Divine Providence, while they are given up to such doings, and
tempt their own destruction. I should not call this a harsh judgment
of them, but the simple truth.

21. And finally, I will add, that I wish this
letter to be read frequently by your wisdom in public, that it may
thereby come to the knowledge of all men, and especially reach the
ears of those who thus act, and thus raise disturbances; for the
judgment which is expressed by me according to the dictates of equity
is confirmed also by real facts. Wherefore, seeing that in such
conduct there is so great criminality, let them understand that I so
judge of them; and that I have come to this determination, that if
they excite any further commotion of this kind, I will myself in
person take cognizance of the matter, and that not according to the
ecclesiastical, but according to the civil laws, and so I will find
them out, because they seem to be offenders not only against human
kind, but against the divine doctrine itself. May God ever preserve
you, dearly beloved brother!

§. 69.

22. But that the wickedness of the calumniators
might be more fully displayed, behold Arsenius also wrote to me after
he was discovered in his place of concealment; and as the letter which
Ischyras had written confessed the falsehood of their accusation, so
that of Arsenius proved their maliciousness still more completely.

23. To the blessed Pope Athanasius, Arsenius,
Bishop of those who were heretofore under Meletius in the city of the
Hypselites, together with the Presbyters and Deacons, wishes much
health in the Lord.

Being earnestly desirous of peace and union with
the Catholic Church, over which by the grace of God you are appointed
to preside, and wishing to submit ourselves to the Canon of the
Church, according to the ancient rule \footnote{vid. supr. p. 3, note A; the (so called) Apostolical Canon apparently
referred to here, is Can. 27. according to Beveridge.

Return to text}, we write unto you, dearly beloved Pope, and declare in the
name of the Lord, that we will not for the future hold communion with
those who continue in schism, and are not at peace with the Catholic
Church, its Bishops, Presbyters, and Deacons. Neither will we take
part with them if they wish to establish any thing in a Council;
neither will we send letters of peace \footnote{p. 8, r. 3.

Return to text} unto them nor receive such from them; neither yet without the
consent of you our Metropolitan, will we publish any decree concerning
Bishops, or on any other general Ecclesiastical question; but we will
yield obedience to all the Canons that have heretofore been ordained,
after the example of the Bishops \footnote{i.e. Meletian Bishops who had conformed; or, since they are not in the
list, §. 71. Catholic Bishops with whom the conforming party were
familiar; or Meletians after the return of Meletius. vid. Tillemont,
Mem. vol. 8. p. 658.

Return to text} Ammonian, Tyrannus, Plusian, and the rest. Wherefore we beseech
your goodness to write to us speedily in answer, and likewise to our
fellow-ministers concerning us, informing them that we will henceforth
abide by the fore-mentioned resolution and will be at peace with the
Catholic Church, and at unity with our fellow-ministers in every part.
And we are persuaded that your prayers, being acceptable unto God,
will so prevail with Him, that this peace shall be firm and
indissoluble unto the end, according to the will of God the Lord of
all, through Jesus Christ our Lord.

24. The sacred Ministry that is under you, we and
those that are with us salute. Very shortly, if God permit, we will
come unto your goodness. I, Arsenius, pray that you may be strong in
the Lord for many years, most blessed Pope.

§. 70.

But a stronger and clearer proof of the calumny
is the recantation of John, of which the most godly Emperor
Constantine of blessed memory is a witness, for knowing how John had
accused himself, and having received letters from him expressing his
repentance, he wrote to him as follows.

25. Constantine Maximus Augustus to John

The letters which I have received from your
prudence were extremely pleasing to me, because I learned from them
what I very much longed to hear, that you had laid aside every narrow
feeling \footnote{[\emph{mikropsuchian}], vid. p. 41, ref. 1.

Return to text}, had joined
the communion of the Church as became you, and were now in perfect
concord with the most reverend Bishop Athanasius. Be assured therefore
that so far I entirely approve of your conduct; because, dismissing
all occasions of quarrel, you have done that which is pleasing
to God, and have embraced the unity of His Church. In order therefore
that you obtain the accomplishment of your wishes, I have thought it
right to grant you permission to enter the public conveyance \footnote{On the “cursus publicus,” vid. Gothofred. in Cod. Theod. viii.
tit. 5. It was provided for the journeys of the Emperor, for persons
whom he summoned, for magistrates, ambassadors, and for such private
persons as the Emperor indulged in the use of it, which was gratis.
The use was granted by Constantine to the Bishops who were summoned to
Nicæa, as far as it went, in addition though aliter Valesius in loc.
to other means of travelling. Euseb. v. Const. iii. 6. The cursus
publicus brought the Bishops to the Council of Tyre. ibid. iv. 43. In
the conference between Liberius and Constantius, Theod. Hist. ii. 13.
it is objected that the cursus publicus is not sufficient to convey
Bishops to the Council which Liberius proposes; he answers that the
Churches are rich enough to convey their Bishops as far as the sea.
Thus S. Hilary was compelled, (datâ evectionis copiâ, Sulp. Sev.
Hist. ii. 57.) to attend at Seleucia, as Athan. at Tyre. Julian
complains of the abuse of the cursus publicus, perhaps with an
allusion to these Councils of Constantius. vid. Cod. Theod. viii. tit.
5. l. 12. where Gothofred quotes Liban. Epitaph. in Julian. (vol. i.
p. 569. ed. Reiske.) Vid. the well-known passage of Ammianus, who
speaks of the Councils being the ruin of the res vehicularia. Hist.
xxi. 16. The Eusebians at Philippopolis say the same thing. Hilar.
Fragm. iii. 25. The Emperor provided board and perhaps lodging for the
Bishops at Ariminum; which the Bishops of Aquitaine, Gaul, and
Britain, declined, except three British from poverty. Sulp. Hist. ii.
56. Hunneric in Africa, after assembling 466 Bishops at Carthage,
dismissed them without modes of conveyance, provision, or baggage.
Victor. Utic. Hist. iii. init. In the Emperor’s letter previous to
the assembling of the sixth Ecumenical Council, A.D. 678, (Harduin. Conc. t. 3. p. 1048 fin.)
he says he has given orders for the conveyance and maintenance of its
members. Pope John VIII. reminds Ursus, Duke of Venice, (A.D. 876.) of the same duty of providing for
the members of a Council, “secundum pios principes, qui in talibus
munificè semper erant intenti.” Colet. Concil. (Ven. 1730.) t. xi.
p. 14.

Return to text}, and to come to the court \footnote{[\emph{stratopedon}]. vid. Chrys. on the Statues, p. 118, note d.
Gothofr. in Cod. Theod. vi. 32. l. 1. Castra sunt ubi Princeps est.
ibid. 35. l. 15. also Kiesling. de Discipl. Cler. i. 5. p. 16.
Beveridge in Can. Apost. 83. interprets [\emph{strateia}] of any civil
engagement as opposed to clerical.

Return to text

} of my clemency. Let it then be your care to make no delay; but
as this letter gives you authority to use the public conveyance, come
to me immediately, that you may have your desires fulfilled, and by
appearing in my presence may enjoy that pleasure which it is fit for
you to receive. May God preserve you continually, dearly beloved
brother.

\xsection{Chapter 6. Documents connected with the Council of Tyre}

§. 71.

1. THUS
ended the conspiracy. The Meletians were repulsed and covered with
shame; but notwithstanding this the Eusebians still did not remain
quiet, for it was not for the Meletians but for the Arians that they
cared, and they were afraid lest, if the proceedings of the former
should be stopped, they should no longer find persons to play the
parts \footnote{Margin Notes

1. p. 34, r. 6.

Return to text}, by whose
assistance they might bring in that heresy. They therefore again
stirred up the Meletians, and persuaded the Emperor to give orders
that a Council should be held afresh at Tyre, and Count Dionysius was
despatched thither, and a military guard was given to the Eusebians.
Macarius also was sent as a prisoner to Tyre under a guard of
soldiers; and the Emperor wrote to me, and laid a peremptory command
upon me, so that, however unwilling, I was obliged to go. The whole
conspiracy may be understood from the letters which the Bishops of
Egypt wrote; but it will be necessary to relate how it was contrived
by them in the outset, that so may be perceived the malice and
wickedness that was exercised against me.

2. There are in Egypt, Libya, and Pentapolis,
nearly one hundred Bishops; none of whom laid any thing to my charge;
none of the Presbyters found any fault with me; none of the people
spoke aught against me; but it was the Meletians who were ejected by
Peter, and the Arians, that divided the plot between them, while the
one party claimed to themselves the right of accusing me, the other of
sitting in judgment on the case. I objected to the Eusebians as being
my enemies on account of the heresy; next, I shewed in the following
manner that the person who was called my accuser was not a Presbyter
at all. When Meletius was admitted into communion, (would that he had
never been so admitted \footnote{p. 88, note R.

Return to text}!)
the blessed \footnote{[\emph{makarit\={e}s}], infr. pp. 161, 162.

Return to text} Alexander
who knew his craftiness required of him a catalogue of the Bishops
whom he said he had in Egypt, and of the Presbyters and Deacons
that were in Alexandria itself, and if he had any in the country
adjoining. This the Pope Alexander did, lest Meletius, assuming full
liberty of action in the Church, should sell ordination to many, and
thus continually, by a fraudulent procedure, put in whatever ministers
he pleased. Accordingly he made out the following catalogue of those
in Egypt.

3. A catalogue presented by Meletius to the
Bishop Alexander.

I, Meletius of Lycus, Lucius of Antinopolis,
Phasileus of Hermopolis, Achilles of Cusæ, Ammonius of Diospolis.

In Ptolemais, Pachymes of Tentyræ.

In Maximianopolis, Theodorus of Coptus.

In Thebais, Cales of Hermethes, Colluthus of
Upper Cynus, Pelagius of Oxyrynchus, Peter of Heracleopolis, Theon of
Nilopolis, Isaac of Letopolis, Heraclides of Niciopolis, Isaac of
Cleopatris, Melas of Arsenoitis.

In Heliopolis, Amos of Leontopolis, Ision of
Athribis.

In Pharbethus, Harpocration of Bubastus, Moses of
Phacusæ, Callinicus of Pelusium, Eudæmon of Tanis, Ephraim of Thmuis.

In Sais, Hermæon of Cynus and Busiris,
Soterichus of Sebennytus, Pininuthes of Phthenegys, Chronius of
Metelis, Agathammon of the district of Alexandria.

In Memphis, John who was ordered by the Emperor
to be with the Archbishop \footnote{Meletius?

Return to text}.
These are those of Egypt.

And the Clergy that he had in Alexandnia were
Apollonius Presbyter,
Irenæus Presbyter, Dioscorus Presbyter, Tyrannus Presbyter. And
Deacons; Timotheus Deacon, Antinous Deacon, Hephæstion Deacon. And
Macarius Presbyter of Parembole \footnote{a village on the Mareotic lake. vid. Socr. iv. 23. Athan. Opp. ed.
Pat. t. 3. pp. 86-89.

Return to text}.

§. 72.

4. These Meletius presented in person to the
Bishop Alexander, but he made no mention of the person called Ischyras,
nor ever professed at all that he had any Clergy in the Mareotis.
Notwithstanding our enemies did not desist from their attempts, but
still he that was no Presbyter was feigned to be one, for there was
the Count ready to use compulsion towards us, and soldiers hurried us
about \footnote{p. 25. vol. 8. p. 4, note H.

Return to text}. But even
then the grace of God prevailed: for they could not convict Macarius
in the matter of the chalice; and Arsenius, whom they reported to have
been murdered by me, stood before them alive and shewed the falseness
of their accusation. When therefore they were unable to convict
Macarius, the Eusebians, who became enraged that they had lost the
prey of which they had been in pursuit, persuaded the Count Dionysius
who is one of them to send to the Mareotis, in order to see whether
they could not find out something there against the Presbyter, or
rather that they might at a distance patch up their plot as they
pleased in my absence: for this was their aim. However, when I
represented that the journey to the Mareotis was a superfluous
undertaking, (for that they ought not to pretend that statements were
defective which they had been employed upon so long, and ought not now
to defer the matter; for they had said whatever they thought they
could say, and now being at a loss what to do, they were making
pretences,) or if they must needs go to the Mareotis, that at least
the suspected parties should not be sent,—the Count was convinced by
my reasoning, with respect to the suspected persons but they did any
thing rather than what I proposed, for the very persons whom I
objected against on account of the Arian heresy, these were they who
specially went, viz. Diognius, Maris, Theodorus, Macedonius, Ursacius,
and Valens. Again, letters were written to the Prefect of Egypt, and a
military guard was provided; and, what was remarkable and altogether
most suspicious, they caused Macarius the accused party to remain
behind under a guard of soldiers, while they took with them the
accuser \footnote{supr. p. 31.

Return to text}.

5. Now who after this does not see through this
conspiracy? Who does not clearly perceive the wickedness of these
Eusebians? For if a judicial enquiry must needs take place in the
Mareotis, the accused ought also to have been sent thither. But if
they did not go for the purpose of such an enquiry, why did they take
the accuser? It was enough that he had not been able to prove the
fact. But this they did in order that they might carry on their
designs against the absent Presbyter, whom they could not convict when
present, and might concoct a plan as they pleased. For when the
Presbyters of Alexandria and of the whole district found fault with
them because they were there by themselves, and required that they too
might be present at their proceedings, (for they said that they knew
both the circumstances of the case, and the history of the person
named Ischyras,) they would not allow them; and although they had with
them Philagrius the Prefect of Egypt, who was an apostate, and heathen
soldiers, during an enquiry which it was not becoming even for
Catechumens to witness, they would not admit the Clergy, lest there as
well as at Tyre there might be those who would expose them.

§. 73.

6. But in spite of these precautions they were
not able to escape detection: for the Presbyters of the City and of
the Mareotis, perceiving their evil designs, addressed to them the
following protest.

7. To Theognius, Maris, Macedonius, Theodorus,
Ursacius, and Valens, the Bishops who have come from Tyre, these from
the Presbyters and Deacons of the Catholic Church of Alexandria under
the most reverend Bishop Athanasius.

It was incumbent upon you when you came hither
and brought with you the accuser, to bring also the Presbyter Macarius;
for trials are appointed by holy Scripture to be so constituted, that
the accuser and accused may stand up together. But since neither you
brought Macarius, nor our most reverend Bishop Athanasius came with
you, we claimed for ourselves the right of being present at the
investigation, that we might see that the enquiry was conducted
impartially, and might ourselves be convinced of the truth. But when
you refused to allow this, and wished, in company only with the
Prefect of Egypt and the accuser, to do whatever you pleased, we
confess that we entertained all evil suspicion of the affair, and
perceived that your coming was only the act of a cabal and a
conspiracy. Wherefore we address to you this letter, to be a testimony
before a genuine Council, that it may be known to all men, that you
have carried on an \emph{ex parte} proceeding and for your private
ends, and have desired nothing else but to form a conspiracy against
us. A copy of this, lest it should be kept secret by you, we have
transmitted also to Palladius the Controller \footnote{Curiosus; the Curiosi (in curis agendis) were
properly the overseers of the public roads, Ducange in voc. but they
became in consequence a sort of imperial spy, and were called the
Emperor's eyes. Gothofr. in Cod.Theod. t. 2. p. 194. ed. 1665.
Constantius confined them to the school of the Agentes in rebus, (infr.
Apol. ad Const. §. 10.) under the Master of the Offices. Gothofr.
ibid. p. 192.

Return to text} of Augustus. For what you have already done causes us to
suspect you, and to reckon on the like conduct from you hereafter.

I Dionysius Presbyter have delivered this letter,
Alexander Presbyter, Nilaras Presbyter, Longus Presbyter, Aphthonius
Presbyter, Athanasius Presbyter, Amyntius Presbyter, Pistus Presbyter,
Plution Presbyter, Dioscorus Presbyter, Apollonius Presbyter, Serapion
Presbyter, Ammonius Presbyter, Gaius Presbyter, Rhinus Presbyter, Œthales
Presbyter.

Deacons; Marcellinus Deacon, Appianus Deacon,
Theon Deacon, Timotheus Deacon, a second Timotheus Deacon.

§. 74.

8. This is the letter, and these the names of the
Clergy of the city; and the following was written by the Clergy of the
Mareotis, who know the character of the accuser, and who were with me
in my visitation.

9. To the holy Council of blessed Bishops of the
Catholic Church, all the Presbyters and Deacons of the Mareotis send
health in the Lord.

Knowing that which is written, \emph{Speak that
thine eyes have seen}, and, \emph{A false witness shall not be
unpunished} [Prov. xxv. 7. Sept. xix. 5.]; we testify what we have
seen, especially since the conspiracy which has been formed against
our Bishop Athanasius has made our testimony necessary. We wonder how
Ischyras ever came to be reckoned among the number of the Ministers of
the Church, which is the first point we think it necessary to mention.
Ischyras never was a Minister of the Church; but when formerly he
represented himself to be a Presbyter of Colluthus, he found no one to
believe him, except only his own relations. For he never had a Church,
nor was ever considered a Clergyman, by those who lived but a short
distance from his village, except only, as we said before, by his own
relations. But, notwithstanding he assumed this designation, he was
deposed in the presence of our Father Hosius at the Council which
assembled at Alexandria, and was reduced to the condition of a
layman, and so he continued subsequently, being deprived of his
pretended claim to the priesthood. Of his character we think it
unnecessary to speak, as all men have it in their power to become
acquainted therewith. But since he has falsely accused our Bishop
Athanasius of breaking a chalice and overturning a table, we are
necessarily obliged to address you on this point.

10. We have said already that he never had a
Church in the Mareotis; and we declare before God as our witness, that
no chalice was broken, nor table overturned by our Bishop, nor by any
one of those who accompanied him; but all that is alleged respecting
this affair is mere calumny. And this we say, not as having been
absent from the Bishop, for we are all with him when he makes his
visitation of the Mareotis, and he never goes about alone, but is
accompanied by all the Presbyters and Deacons, and by a considerable
number of the people. Wherefore we make these assertions, as having
been present with him during the whole of the visitation which he made
amongst us, and testify that neither was a chalice ever broken, nor
table overturned, but the whole story is false, as the accuser himself
also witnesses under his own hand \footnote{supr. p. 93.

Return to text}. For when, after he had withdrawn with the Meletians, and had
reported these things against our Bishop Athanasius, he wished to be
admitted to communion, he was not received, although he wrote and
confessed under his own hand that none of these things were true, but
that he had been suborned by certain persons to say so.

§. 75.

11. Wherefore also Theognius, Theodorus, Maris,
Macedonius, Ursacius, and Valens, came into the Mareotis, and when
they found that none of these things were true, but it was likely to
be discovered that they had framed a false accusation against our
Bishop Athanasius, the party of Theognius being themselves his
enemies, caused the relations of Ischyrus and certain Arian fanatics \footnote{supr. p. 4, r. 1.

Return to text} to say whatever they wished. For none of the people spoke
against the Bishop but these persons, through a dread of Philagrius
the Prefect of Egypt, and by threats and with the support of the Arian
fanatics \footnote{supr. p. 4, r. 1.

Return to text}, accomplished whatever they
desired. For when we came to disprove the calumny, they would
not permit us, but cast us out, while they admitted whom they pleased
to a participation in their schemes, and concerted matters with them,
influencing them by fear of the Prefect Philagrius. Through his means
they prevented us from being present, that we might discover whether
those who were suborned by them were members of the Church or Arian
fanatics. And you also, dearly beloved Fathers, know, as you teach us,
that the testimony of enemies avails nothing. That what we say is the
truth the handwriting \footnote{[\emph{cheir}], infr. Apol. ad Const. §. 11.

Return to text}
of Ischyras testifies, as do also the facts themselves, because when
we were conscious that no such thing as was pretended had taken place,
they took with them Philagrius, that through fear of the sword and by
threats they might frame whatever plots they wished. These things we
testify as in the presence of God; we make these assertions as knowing
that there will be a judgment held by God; desiring indeed all of us
to come to you, but being content with these letters which we send to
you, that they may be instead of the presence of those who cannot
come.

I, Ingenius Presbyter, pray that you may be
strong in the Lord, dearly beloved Fathers. Theon P. Ammonas P.
Heraclius P. Boccon P. Tryphon P. Peter P. Hierax P. Serapion P.
Marcus P. Ptollarion P. Gaius P. Dioscorus P. Demetrius P. Thyrsus P.

Deacons; Pistus D. Apollos D. Serras D. Pistus D.
Polynicus D. Ammonius D. Maurus D. Hephæstus D. Apollos D. Metopas D.
Apollos D. Serapas P. Meliphthmongus D. Lucius D. Gregoras D.

§. 76.

12. The same to the Controller, and to Philagrius,
at that time Prefect of Egypt.

To Flavius Philagrius, and to Flavius Palladius,
Ducenary \footnote{On the different kinds of Ducenaries, vid. Gothofr. in Cod. Theod. xi.
7. leg. 1. Here, as in Euseb. Hist. vii. 30. the word stands for a
Procurator, whose annual pay amounted to 200 sestertia. vid. Salmas.
in Hist. Aug. t. 1. p. 533. In like manner a Centenary is one who
receives 100.

Return to text}, Officer of the
Palace \footnote{vid. p. 89, r. 3.

Return to text}, and
Controller, and to Flavius Antoninus, Commissary of Provisions \footnote{[\emph{biarch\={o}i}].

Return to text}, and Centenary of my Lords, the most illustrious Prefects of
the sacred Prætorium, these from the Presbyters and Deacons of the
Mareotis, a district of the Catholic Church which is under the
most Reverend Bishop Athanasius, we address this testimony by those
whose names are under-written:—

Whereas Theognius, Maris, Macedonius, Theodorus,
Ursacius, and Valens, as if sent by all the Bishops who assembled at
Tyre, came into our Diocese alleging that they had received orders to
investigate certain ecclesiastical affairs, among which they spoke of
the breaking of a chalice belonging to the Lord, of which information
was given them by Ischyras, whom they brought with them, and who says
that he is a Presbyter, although he is not,—for he was ordained by
the Presbyter Colluthus who pretended to the Episcopate, and was
afterwards ordered by a whole Council, by Hosius and the Bishops that
were with him, to take the place of a Presbyter, as he was before; and
accordingly all that were ordained by Culluthus, resumed the same rank
which they held before, and so Ischyras himself proved to be a
layman,—and the Church, which he says he has, never was a Church at
all, but a small dwelling house belonging to an orphan boy of the name
of Ision;—for this reason we have offered this testimony, adjuring
you by Almighty God, and by our Lords Constantine Augustus, and the
most illustrious Cæsars his sons, to bring these things to the
knowledge of their piety. For neither is he a Presbyter of the
Catholic Church, nor does he possess a Church, nor has a chalice ever
been broken, but the whole story is false and an invention.

Dated in the Consulship \footnote{A.D.
335.

Return to text} of Julius Constantius the most illustrious Patrician \footnote{The title Patrician was revived by Constantine as a personal
distinction. It was for life, and gave precedence over all the great
officers of state except the Consul. It was usually bestowed on
favourites, or on ministers as a reward of services. Gibbon, Hist. ch.
17. This Julius Constantius, who was the father of Julian, was the
first who bore the title, with L. Optatus, who had been consul the
foregoing year. Illustrissimus was the highest of the three ranks of
honour. ibid.

Return to text}, brother of the most religious Emperor Constantine Augustus,
and of Rufinus Albinus, most illustrious men, on the tenth day of the
month Thoth \footnote{August.

Return to text}.

These were the letters of the Presbyters.

§. 77.

13. The following also are the letters and
protests of the Bishops who came with us to Tyre, when they discovered
the conspiracy and plot.

14. To the Bishops assembled at Tyre, most
honoured Lords, those of the Catholic Church who have come from Egypt
with Athanasius send health in the Lord.

We suppose that the conspiracy which has been
formed against us by Eusebius, Theognius, Marus, Narcissus, Theodorus,
and Patrophilus, is no longer uncertain. From the very beginning we
all demurred, through our fellow-minister Athanasius, to the holding
of the inquiry in their presence, knowing that the presence of even
one enemy only, much more of many, is able to disturb and injure the
hearing of a cause. And you also yourselves know the enmity which they
entertain, not only towards us, but towards all the orthodox, how that
for the sake of the fanaticism of Arius, and his impious doctrine,
they direct their assaults, they form conspiracies against all. And
when, being confident in the truth, we desired to shew the falsehood,
which the Meletians had employed against the Church, the Eusebians
endeavoured by some means or other to interrupt our representations,
and strove eagerly to set aside our testimony, threatening those who
gave an honest judgment, and insulting others, for the sole purpose of
carrying out the design they had against us. Your divinely inspired \footnote{[\emph{entheos}].

Return to text} piety, most honoured Lords, was probably ignorant of their
conspiracy, but we suppose that it has now been made manifest.

15. For indeed they have themselves plainly
disclosed it; for they desired to send to the Mareotis those of their
party who are suspected by us, so that, while we were absent and
remained here, they might disturb the people and accomplish what they
wished. They knew that the Arian fanatics, and Colluthians \footnote{Colluthus formed a schism on the doctrine that God was not the cause
of any sort of evil, e.g. did not inflict pain and suffering. Though a
Priest, he took on himself to ordain, even to the Priesthood. vid.
supr. p. 30, note L. St. Alexander even seems to imply that he did so
for money. Theod. Hist. i. 3.

Return to text} and Meletians, were enemies of the Catholic Church, and
therefore they were anxious to send them, that in the presence of our
enemies they might devise against us whatever schemes they pleased.
And those of the Meletians who are here, even four days previously,
(as they knew that this inquiry was about to take place,) despatched
at evening certain of their party, as a post, for the purpose of
collecting Meletians out of Egypt into the Mareotis, because there
were none at all there, and Colluthians and Arian fanatics, from
other parts, and to prepare them to speak against us. For you also
know that Ischyras himself confessed before you, that he had not more
than seven persons in his congregation. When therefore we heard that,
after they had made what preparations they pleased against us, and had
sent these suspected persons, they were going about to each of you,
and requiring your subscriptions, in order that it might appear as if
this had been done with the consent of you all; for this reason we
hastened to write to you, and to present this our testimony; declaring
that we are the objects of a conspiracy under which we are suffering
by and through them, and demanding that having the fear of God in your
minds, and condemning their conduct in sending whom they pleased
without our consent, you would refuse your subscriptions, lest they
pretend that those things are done by you, which they are contriving
only among themselves.

16. Surely it becomes those who are in Christ,
not to regard men, but to prefer the truth before all things. And be
not afraid of their threatenings, which they employ against all, nor
of their plots, but rather fear God. If it was at all necessary that
persons should be sent to the Mareotis, we also ought to have been
there with them, in order that we might convict the enemies of the
Church, and point out those who were aliens, and that the
investigation of the matter might be impartial. For you know that the
Eusebians contrived that a letter should be presented, as coming front
the Colluthians, the Meletians, and Arians, and directed against us:
but it is evident that these enemies of the Catholic Church speak
nothing that is true concerning us, but say every thing against us.
And the law of God forbids an enemy to be either a witness or a judge.
Wherefore as you will have to give an account in the day of judgment,
receive this testimony, and recognising the conspiracy which has been
framed against us, beware, if you are requested by them, of doing any
thing against us, and of taking part in the designs of the Eusebians.
For you know, as we said before, that they are our enemies, and are
aware why Eusebius of Cæsarea became such last year. We pray that you
may be in health, greatly beloved Lords \footnote{[\emph{kurioi}].

Return to text}.

17. To the most illustrious Count Flavius
Dionysius, from the Bishops of the Catholic Church in Egypt who have
come to Tyre \footnote{nearly verbatim as the foregoing.

Return to text}.

§. 78.

We suppose that the conspiracy which has been
formed against us by Eusebius, Theognius, Maris, Narcissus, Theodorus,
and Patrophilus, is no longer uncertain. From the very beginning we
all demurred, through our fellow-minister Athanasius, to the holding
of the inquiry in their presence, knowing that the presence of even
one enemy only, much more of many, is able to disturb and injure the
hearing of a cause. For their enmity is manifest which they entertain,
not only towards us, but also towards all the orthodox, because they
direct their assaults, they form conspiracies against all. And when,
being confident in the truth, we desired to shew the falsehood which
the Meletians had employed against the Church, the Eusebians
endeavoured by some means or other to interrupt our representations,
and strove eagerly to set aside our testimony, threatening those who
gave a honest judgment and insulting others, for the sole purpose of
carrying out the design they had against us. Your goodness was
probably ignorant of the conspiracy which they have formed against us,
but we suppose that it has now been made manifest.

18. For indeed they have themselves plainly
disclosed it; for they desired to send to the Mareotis those of their
party who are suspected by us, so that, while we are absent, and
remained here, they might disturb the people and accomplish what they
wished. They knew that Arian fanatics, Colluthians, and Meletians were
enemies of the Church, and therefore they were anxious to send them,
that in the presence of our enemies, they might devise against us
whatever schemes they pleased. And those of the Meletians who are
here, even four days before, (as they knew that this inquiry was about
to take place,) despatched at evening two individuals of their own
party, as a post, for the purpose of collecting Meletians out of Egypt
into the Mareotis, because there were none at all there, and
Colluthians, and Arian fanatics, from other parts, and to prepare them
to speak against us. And your goodness knows that he himself confessed before you, that he had not more than seven persons in his
congregation. When therefore we heard that, after they had made what
preparations they pleased against us, and had sent these suspected
persons, they were going about to each of the Bishops and requiring
their subscriptions, in order that it might appear that this was done
with the consent of them all; for this reason we hastened to refer the
matter to your honour, and to present this our testimony, declaring
that we are the objects of a conspiracy, under which we are suffering
by and through them, and demanding of you that having in your mind the
fear of God, and the pious commands of our most religious Emperor, you
would no longer tolerate these persons, but condemn their conduct in
sending whom they pleased without our consent.

I Adamantius Bishop have subscribed this letter,
Ischyras, Ammon, Peter, Ammonianus, Tyrannus, Taurinus, Sarapammon, Œlurion,
Harpocration, Moses, Optatus, Anubion, Saprion, Apollonius, Ischyrion,
Arbæthion, Potamon, Paphnutius, Heraclides, Theodorus, Agathammon,
Gaius, Pistus, Athas, Nicon, Pelagius, Theon, Paninuthius, Nonnus,
Ariston, Theodorus, Irenæus, Blastammon, Philippus, Apollos,
Dioscorus, Timotheus of Diospolis, Macarius, Heraclammon, Cronius,
Muis, James, Ariston, Artemidorus, Phinees, Psais, Heraclides.

[§. 79.?]

19. Another from the same.

The Bishops of the Catholic Church who have come
from Egypt to Tyre, to the most illustrious Count Flavius Dionysius.

Perceiving that many conspiracies and plots are
being formed against us through the machinations of Eusebius,
Narcissus, Flacillus, Theognius, Maris, Theodorus, and Patrophilus,
(against whom we wished at first to enter an objection, but were not
permitted,) we are constrained to have recourse to the present appeal.
We observe also that great zeal is exerted in behalf of the Meletians,
and that a plot is laid against the Catholic Church in Egypt in our
persons. Wherefore we address this letter to you, beseeching you to
bear in mind the Almighty Power of God, who defends the kingdom of our
most religious and godly Emperor Constantine, and to reserve the
bearing of the affairs which concern us for the most religious Emperor
himself. For it is but reasonable, since you were commissioned by his
Majesty, that you should reserve the matter for him upon our appealing
to his piety. We can no longer endure to be the objects of the
treacherous designs of the fore-mentioned Eusebians, and therefore we
demand that the case be reserved for the most religious and godly
Emperor, before whom we shall be able to set forth our own and the
Church's just claims. And we are convinced that when his piety shall
have heard our cause, he will not condemn us. Wherefore we again
adjure you by Almighty God, and by our most religious Emperor, who,
together with the children of his piety, has thus ever been victorious
\footnote{pp. 79 and 96, p. 119, r. 2.

Return to text} and prosperous these
many years, that you proceed no further, nor suffer yourself to move
at all in the Council in relation to our affairs, but reserve the
hearing of them for his piety. We have likewise made the same
representations to my Lords \footnote{[\emph{kuriois mou}].

Return to text} the orthodox Bishops.

§. 80.

20. Alexander \footnote{p. 33, note N.

Return to text}, Bishop of Thessalonica, on receiving these letters, wrote to
the Count Dionysius as follows.

21. The Bishop Alexander to my Lord \footnote{[\emph{despot\={e}i mou}].

Return to text} Dionysius.

I see that a conspiracy has evidently been formed
against Athanasius; for they have determined, I know not on what
grounds, to send all those to whom he has objected, without giving any
information to us, although it was agreed that we should consider
together who ought to be sent. Take care therefore that nothing be
done rashly, (for they have come to me in great alarm, saying that the
wild beasts have already roused themselves, and are going to rush upon
them; for they had heard it reported, that John had sent certain,)
lest they be beforehand with us, and concoct what schemes they please.
For you know that the Colluthians \footnote{p. 109, note D.

Return to text} who are enemies of the Church, and the Arians, and Meletians,
are all of them leagued together, and are able to work much evil.
Consider therefore what is best to be done, lest some mischief befal,
and we be subject to censure, as not having judged the matter fairly.
Great suspicions are also entertained of these persons, lest, as being
devoted to the Meletians, they should go through those Churches
whose Bishops are here \footnote{at Tyre.

Return to text},
and raise an alarm amongst them, and so disorder the whole of Egypt.
For they see that this is already taking place to a great extent.

22. In consequence of this, the Count Dionysius
wrote to the Eusebians as follows.

§. 81.

23. This is what I have already mentioned to my
lords associated with Flacillus \footnote{perhaps president of Council, vid. p. 39, note B.

Return to text}, that Athanasius has come forward and complained that those
very persons have been sent whom he objected to; and crying out that
he has been wronged and deceived. Alexander the lord of my soul has
also written to me on the subject; and that you may perceive that what
his Excellence has said is reasonable, I have subjoined his letter to
be read by you. Remember also what I wrote to you before: I impressed
upon your Excellences, my lords, that the persons who were sent ought
to be commissioned by the general vote and decision of all. Take care
therefore lest our proceedings fall under censure, and we give just
grounds of blame to those who are disposed to find fault with us. For
as the accuser's side ought not to suffer any oppression, so neither
ought the defendant's. And I think that there is no slight ground of
blame against us, when my lord Alexander appears to disapprove of what
we have done.

§. 82.

24. While matters were proceeding thus we
withdrew from them, as from \emph{an assembly of treacherous men} [Jer.
ix. 2.], for whatsoever they pleased they did, whereas there is no man
in the world but knows that \emph{ex parte} proceedings cannot stand
good. This the divine law determines; for when the blessed Apostle was
suffering under a similar conspiracy and was brought to trial, he
demanded, saying, The Jews from Asia ought to have been here before
thee, and object, if they had ought against me [Acts xxiv. 18. 19.].
On which occasion Festus also, when the Jews wished to lay such a plot
against him, as these men have now laid against me, said, \emph{It is not
the manner of the Romans to deliver any man to die, before that he
which is accused have the accuser face to face}, \emph{and have
licence to answer for himself concerning the crime laid against him}
[Acts xxv. 16.]. But the Eusebians have both had the boldness to
pervert the law, and have acted more unjustly even than those unjust
persons. For they did not proceed privately at the first, but when in
consequence of our being present they found themselves weak, then they
straightway went out, like the Jews, and took counsel together alone,
how they might destroy us and bring in their heresy, as they demanded
Barabbas. For this purpose it was, as they have themselves confessed,
that they did all these things.

§. 83.

25. Although these circumstances were amply
sufficient for our vindication, yet in order that the wickedness of
these men and the freeness of the truth might be more fully exhibited,
I have not felt averse to repeat them again, in order to shew that
they have acted in a manner inconsistently with themselves, and as men
scheming in the dark have fallen foul upon one another, and while they
desired to destroy us have like insane persons wounded themselves. For
in their investigation of the subject of the Mysteries, they
questioned Jews, they examined Catechumens \footnote{vid. p. 73.

Return to text}; ``Where were you,'' they said, ``when Macarius came and
overturned the Table?'' They answered, ``We were present within doors;''
whereas there could be no oblation if Catechumens were present. Again,
although they had written word every where, that Macarius came and
overthrew every thing, while the Presbyter was standing and
celebrating the Mysteries, yet when they questioned whomsoever they
pleased, and asked them, ``Where was Ischyras when Macarius rushed in?''
those persons answered that he was lying sick in a cell. Now he that
was lying could not be standing, nor could one that lay sick in his
cell offer the oblation. Besides whereas Ischyras said that certain
books had been burnt by Macarius, the witnesses who were suborned to
give evidence, declared that nothing of the kind had been done, but
that Ischyras spoke falsely. And what is most remarkable, although
they had again written word every where, that those who were able to
give evidence had been concealed by us, yet these persons made their
appearance, and they questioned them, and were not ashamed to find it
proved on all sides that they were slanderers, and had acted in
this matter clandestinely, and according to their pleasure. For they
prompted the witnesses by signs, while the Prefect threatened them,
and the soldiers pricked them with their swords; but the Lord revealed
the truth, and shewed them to be slanderers. Therefore also they
concealed the Records of their proceedings, which they retained
themselves, and charged those who wrote them to keep out of sight, and
to communicate to no one whomsoever. But in this too also they were
disappointed; for the person who wrote them was Rufus, who is now
public executioner in the Augustalian \footnote{vid. p. 5, note D.

Return to text} prefecture, and is able to testify to the truth of this; and
the Eusebians sent them to Rome by the hands of their own friends, and
Julius the Bishop transmitted them to me. And now they are mad with
rage, because we have obtained and read what they wished to conceal.

§. 84.

26. As such was the character of their
machinations, so they very soon shewed plainly the reasons of their
conduct. For when they went away, they took the Arians with them to
Jerusalem, and there admitted them to communion, having sent out a
letter concerning them, part \footnote{vid. de Syn. §. 21. (vol. 8, p. 103.)

Return to text} of which, and the beginning, is as follows.

27. The holy Council by the grace of God
assembled at Jerusalem, to the Church of God which is in Alexandria,
and to the Bishops, Presbyters, and Deacons, in all Egypt, the Thebais,
Libya, Pentapolis, and throughout the world, sends health in the Lord.

Having come together out of different Provinces
to a great meeting which we have held for the consecration of the
Martyry of the Saviour, which has been appointed to the service of God
the King of all and of His Christ, by the zeal of our most religious
Emperor Constantine, the grace of God hath afforded us more abundant
rejoicing of heart; which our most religious Emperor himself hath
occasioned us by his letters, wherein he hath stirred us up to do that
which is right, putting away all envy from the Church of God, and
driving far from us all malice, by which the members of God have been
heretofore torn asunder, and that we should with single and peaceable
minds receive the Arians, whom envy, that enemy of all goodness, has
caused for a season to be excluded from the Church. Our most
religious Emperor has also in his letter testified to the correctness
of their faith, which he has ascertained from themselves, himself
receiving the profession of it from them by word of mouth, and has now
made manifest to us by subjoining a written declaration of their
orthodox belief.

§. 85.

28. Every one that hears of these things must see
through their treachery. For they made no concealment of what they
were doing; unless perhaps they confessed the truth without wishing
it. For if I was the hindrance to the admittance of the Arians into
the Church, and if they were received while I was suffering from their
plots, what other conclusion can be arrived at, than that these things
were done on their account, and that all their proceedings against me,
and the story which they fabricated about the breaking of the chalice
and the murder of Arsenius, were for the sole purpose of introducing
impiety into the Church, and of preventing their being condemned as
heretics? For this was what the Emperor threatened long ago in his
letters to me. And they were not ashamed to write in the manner they
did, and to affirm that those persons whom the whole Ecumenical
Council anathematized held orthodox sentiments. And as they undertook
to say and do any thing without scruple, so they were not afraid to
meet together in a corner, in order to overthrow, as far as was in
their power, the authority of so great a Council.

29. Moreover, the price which they paid for false
testimony yet more fully manifests their wickedness and impious
intentions. The Mareotis, as I have already said, is a district of
Alexandria, in which there has never been either a Bishop or a
Chorepiscopus \footnote{That Chorepiscopi were real Bishops, vid. Bevereg. in Conc. Ancyr.
Can. 13. Routh in Conc. Neocæs. Can. 13. referring to Rhabanus Maurus.
Thomassin on the other hand denies that they were Bishops, Discipl.
Eccl. i. 2. c. 1.

Return to text}; but the
Churches of the whole district are subject to the Bishop of
Alexandria, and each Presbyter has under his charge one of the largest
villages, which are about ten or more in number \footnote{Ten under each Presbyter. Vales. ad Socr. Hist. i. 27. Ten altogether,
Montfacon in loc. with more probability; and so Tillemont, vol. 8. p.
20.

Return to text}. Now the village in which Ischyras lives, is a very small one,
and possesses so few inhabitants, that there has never been a Church
built there, but only in the adjoining village. Nevertheless,
they determined, contrary to ancient usage \footnote{It was against the Canon of Sardica, and doubtless against ancient
usage, to ordain a Bishop for so small a village, vid. Bingham, Antiqu.
ii. 12. who, howver, maintains by instances, that at least small towns
might be sees. Also it was against usage that a layman, as Ischyras,
should be made a Bishop. ibid. 10. §. 4, \&c. St. Hilary, however,
makes him a Deacon. Fragm. ii. 16.

Return to text}, to nominate a Bishop for this place, and not only so, but even
to appoint one, who was not so much as a Presbyter. Knowing as they
did the unusual nature of such a proceeding, yet being constrained by
the promises they had given in return for his false impeachment of me,
they submitted even to this, lest that abandoned person, if he were
ill-treated by them, should disclose the truth, and thereby show the
wickedness of the Eusebians. Notwithstanding this, he has no Church,
nor a people to obey him, but is scouted by them all, like a dog \footnote{Dogs without owners, and almost in a wild state, abound, as is well
known, in Eastern cities; vid. Psalm lix. 6, 14, 15. 2 Kings ix. 35,
36. and for the view taken in Scripture of dogs, vid. Bochart, Hieroz.
ii. 56.

Return to text}, although they have even caused the Emperor to write to the
Receiver-General \footnote{Catholicus, p. 32. Apol. ad Const. §. 10.

Return to text},
(for every thing is in their power,) commanding that a Church should
be built for him, that being possessed of that, his statement may
appear credible about the chalice and the table. They caused him
immediately to be nominated a Bishop, because if he were without a
Church, and not even a Presbyter, he would appear to be a false
accuser, and a fabricator of the whole matter. Nevertheless he
possesses but an empty title, as he has no people \footnote{pp. 108, 110.

Return to text}, and even his own relations are not obedient to him, and the
letter also has failed to accomplish its purpose, remaining only as a
convincing proof of the utter wickedness of himself and the Eusebians.
It runs as follows.

30. The Letter of the Receiver-General.

Flavius Hemnerius sends health to the
Tax-collector \footnote{Exactor.

Return to text} of the
Mareotis.

Ischyras the Presbyter having petitioned the
piety of our Lords, Cæsars Augusti, that a Church might be built in
the district of the Peace of Secontarurus \footnote{p. 34, note O.

Return to text}, their divinity has commanded that this should be done as soon
as possible. Take care therefore, as soon as you receive the copy of
the sacred Edict, which with all due veneration is placed above, and
the Reports which have been formed before my sanctity, that you quickly make an abstract of them, and transfer them to the Order book,
so that the sacred command may be put in execution.

§. 86.

31. While they were thus plotting and scheming, I
went up \footnote{p. 26, r. 2.

Return to text} and
represented to the Emperor the unjust conduct of the Eusebians, for he
it was who had commanded the Council to be held, and his Count
presided at it. When he heard my report, he was greatly moved, and
wrote to them as follows.

32. Victor \footnote{Euseb. v. Const. ii. 48.

Return to text}, Constantine, Maximus, Augustus, to the Bishops assembled at
Tyre.

I know not what the decisions are which you have
arrived at in your Council amidst noise and tumult; but somehow the
truth seems to have been perverted in consequence of certain
confusions and disorders, in that you, through your mutual
contentiousness, which you are resolved should prevail, have failed to
perceive what is pleasing to God. However, it will rest with Divine
Providence to disperse the mischiefs which manifestly are found to
arise from this contentious spirit, and to shew plainly to us, whether
you, while assembled in that place, have had any regard for the truth,
and whether you have made your decisions uninfluenced by either favour
or enmity. Wherefore I wish you all to assemble with all speed before
my piety, in order that you may render in person a true account of
your proceedings.

33. The reason why I have thought good to write
thus to you, and why I summon you before me by letter, you will learn
from what I am going to say. As I was entering on a late occasion our
all-happy home of Constantinople, which bears our name, (I chanced at
the time to be on horseback,) on a sudden the Bishop Athanasius, with
certain others whom he had with him, approached me in the middle of
the road, so unexpectedly, as to occasion me much amazement. God, who
knoweth all things, is my witness, that I should have been unable at
first sight even to recognise him, had not some of my attendants, on
my naturally enquiring of them, informed me both who it was, and under
what injustice he was suffering. I did not however enter into any
conversation with him at that time, nor grant him an interview; but
when he requested to be heard I refused, and all but gave orders
for his removal: when with increasing boldness he claimed only this
favour, that you should be summoned to appear, that he might have an
opportunity of complaining before me in your presence, of the
ill-treatment which he has met with. As this appeared to me to be a
reasonable request, and suitable to the times, I willingly ordered
this letter to be written to you, in order that all of you, who
constituted the Council which was held at Tyre, might hasten without
delay to the Court \footnote{[\emph{stratopedon}], p. 100, note Z.

Return to text}
of my clemency, so as to prove by facts that you had passed an
impartial and uncorrupt judgment. This, I say, you must do before me,
whom not even you will deny to be a true servant of God \footnote{[\emph{theraponta}].

Return to text}.

34. For indeed through my devotion \footnote{[\emph{latreias}].

Return to text} to God, peace is preserved every where, and the Name of God is
truly worshipped even by the barbarians, who have hitherto been
ignorant of the truth. And it is manifest, that he who is ignorant of
the truth, does not know God. Nevertheless, as I said before, even the
barbarians have now come to the knowledge of God, by means of me, His
true servant \footnote{``Once in an entertainment, at which he (Constantine) received Bishops,
he made the remark that he too was a Bishop using pretty much these
words in my hearing, You are Bishops of matters within the Church, I
am appointed by God to be Bishop of matters external to it.'' Euseb.
Vit. Const. iv. 24. vid. supr. p. 76, note M.

Return to text}, and have
learned to fear Him whom they perceive from actual facts to be my
shield and protector every where. And from this chiefly they have come
to know God, whom they fear through the dread which they have of me.
But we, who profess to set forth (for I will not say to guard) the
holy mysteries of His Goodness, we, I say, engage in nothing but what
tends to dissension and hatred, and, in short, whatever contributes to
the destruction of mankind. But hasten, as I said before, and all of
you with all speed come to us, being persuaded that I shall endeavour
with all my might to amend what is amiss, so that those things
specially may be preserved and firmly established in the law of God,
to which no blame nor dishonour may attach; while the enemies of the
law, who under pretence of His holy Name bring in manifold and divers
blasphemies, shall be scattered abroad, and entirely crushed, and
utterly destroyed.

§. 87.

35. When the Eusebians read this letter, being
conscious of what they had done, they prevented the rest of the
Bishops from going up, and only themselves went, viz. Eusebius,
Theognius, Patrophilus, the other Eusebius, Ursacius, and Valens. And
they no longer said any thing about the chalice and Arsenius, (for
they had not the boldness to do so,) but inventing another accusation
which concerned the Emperor himself, they declared before him, that
Athanasius had threatened that he would cause the corn to be withheld
which was sent from Alexandria to his own home \footnote{Constantinople.

Return to text}. The Bishops Adamantius, Anubion, Agathammon, Arbethion, and
Peter, were present and heard this. It was proved also by the anger of
the Emperor; for although he had written the preceding letter, and had
condemned their injustice, as soon as he heard such a charge as this,
he was immediately incensed, and instead of granting me a hearing, he
sent me away into Gaul. And this again shows their wickedness further:
for when the younger Constantine, of blessed memory, sent me back
home, remembering what his father had written, he also wrote as
follows.

36. Constantine Cæsar, to the people of the
Catholic Church of the city of Alexandria

I suppose that it has not escaped the knowledge
of your pious minds, that Athanasius, the interpreter of the adorable
Law, was sent away into Gaul for a time, with the intent that, as the
savageness of his bloodthirsty and inveterate enemies persecuted him
to the hazard of his sacred life \footnote{[\emph{kephal\={e}s}].

Return to text}, he might thus escape suffering some irremediable calamity,
through the perverse dealing of those evil men. In order therefore to
escape this, he was snatched out of the jaws of his assailants, and
was ordered to pass some time under my government, and so was supplied
abundantly with all necessaries in this city, where he lived, although
indeed his celebrated virtue, relying entirely on divine assistance,
set at nought the sufferings of adverse fortune. Now seeing that it
was the fixed intention of our Lord \footnote{[\emph{despot\={e}s}].

Return to text} Constantine Augustus, my Father, to restore the said Bishop to
his own place, and to your most beloved piety, but he was taken away
by that fate which is common to all men, and went to his rest before
he could accomplish his wish; I have thought proper to fulfil
that intention of the Emperor of sacred memory which I have inherited
from him. When he comes to present himself before you, you will learn
with what reverence he has been treated. Indeed it is not wonderful,
whatever I have done on his behalf; for the thoughts of your longing
desire for him, and the appearance of so great a man, moved my mind,
and urged me thereto. May Divine Providence continually preserve you,
dearly beloved brethren.

Dated from Treves the 15th before the Calends of
July \footnote{June 17. A.D. 338.

Return to text}.

§. 88.

37. This being the reason why I was sent away
into Gaul, who, I ask again, does not plainly perceive the intention
of the Emperor, and the murderous spirit of the Eusebians, and that
the Emperor did this in order to prevent their forming some more
desperate scheme? for he listened to them with a sincere purpose \footnote{[\emph{ep\={e}kouse gar hapl\={o}s}]. Montfaucon in Onomast. (Athan.
t. 2. ad calc.) points out some passages in his author, where [\emph{epakeuein}]
like [\emph{hypakouein}] means ``to answer.'' vid. Apol. ad Const. §.
16. init. Orat. iii. 27 fin.

Return to text}. Such were the practices of the Eusebians, and such their
machinations against me. Who that has witnessed them will deny that
nothing has been done in my favour out of partiality, but that that
great number of Bishops both individually and collectively wrote as
they did in my behalf and condemned the falsehood of my enemies
justly, and in accordance with the truth? Who that has observed such
proceedings as these will deny that Valens and Ursacius had good
reason to condemn themselves, and to write as they did, to accuse
themselves on their repentance, choosing rather to suffer shame for a
short time, than to undergo the punishment of false accusers for ever
and ever \footnote{Here ends the second part of the Apology, as is evident by turning
back to §. 58. (supra, p. 87 fin.) to which this paragraph is an
allusion. The express object of the second part was to prove, what has
now been proved by documents, that Valens and Ursacius did but succumb
to plain facts which they could not resist. It is observable too from
this passage that the Apology was written before their relapse, i.e.
before A.D. 351, or 352. The remaining two
sections are written after 357, as they mention the fall of Liberius
and Hosius, and speak of Constantius in different language from any
which has been found above. vid. Libr. F. vol. 8. p. 90, note p.

Return to text

}?

§. 89.

38. Wherefore also my blessed brothers in
ministry, acting justly and according to the laws of the Church, while
certain affirmed that my case was doubtful, and endeavoured to
compel them to annul the sentence which was passed in my favour, have
now endured all manner of sufferings, and have chosen rather to be
banished than to see the judgment of so many Bishops reversed. Now if
those genuine Bishops had withstood by words only those who plotted
against me, and wished to undo all that had been done in my behalf; or
if they had been ordinary men, and not the Bishops of illustrious
cities, and the heads of great Churches, there would have been room to
suspect that in this instance they too had acted contentiously and in
order to gratify me. But when they not only endeavoured to convince by
argument, but also endured banishment, and one of them is Liberius
Bishop of Rome, (for although he did not endure to the end the
sufferings of banishment, yet he remained in his exile for two years,
being aware of conspiracy formed against me,) and since there is also
the great Hosius, together with the Bishops of Italy, and of Gaul, and
others from Spain, and from Egypt, and Libya, and all those from
Pentapolis, (for although for a little while, through fear of the
threats of Constantius, he seemed not to resist them, yet the great
violence and tyrannical power exercised by Constantius, and the many
insults and stripes inflicted on him, prove that it was not because he
gave up my cause, but through the weakness of old age, being unable to
bear the stripes, that he yielded to them for a season,) therefore I
say, it is altogether right that all, as being fully convinced, should
hate and abominate the injustice and the violence which they have used
towards me; especially as it is well known that I have suffered these
things on account of nothing else but the Arian impiety.

§. 90.

39. Now if any one wishes to become acquainted
with my case, and the falsehood of the Eusebians, let him read what
has been written in my behalf, and let him hear the witnesses, not
one, or two, or three, but that great number of Bishops; and again let
him attend to the witnesses of these proceedings, Liberius and Hosius,
and their associates, who when they saw the attempts made against me,
chose rather to endure all manner of sufferings than to give up the
truth, and the judgment which had been pronounced in my favour. And this they did with an honourable and righteous intention, for what
they suffered proves to what straits the other Bishops were reduced.
And they are memorials and records against the Arian heresy, and the
wickedness of the false accusers, and afford a pattern and model for
those who come after, to contend for the truth unto death, and to
abominate the Arian heresy which fights against Christ \footnote{[\emph{christomachon}], vol. 8. p. 6, note N.

Return to text}, and is a fore-runner of Antichrist \footnote{ibid. p. 79, note Q.

Return to text}; and not to believe those who attempt to speak against me. For
the defence put forth, and the sentence given, by so many Bishops of
high character, are a trustworthy and sufficient testimony in my
behalf.

\xchapter{III. Encyclical Epistle of our holy Father Athanasius, Archbishop of Alexandria, to the Bishops of Egypt and Libya, against the Arians}

———————

[The Circular Epistle which follows was addressed
by S. Athanasius to the Bishops of his Patriarchate in the beginning
of 356, immediately after his flight from Egypt on the outrages
committed against the Church by Syrianus. Some indeed have referred it
to the year 361, with some plausibility, on the ground of a passage in
§. 22, where he speaks of the Arians being ``declared heretics 36
years ago and cast out of the Church by decree of the whole Ecumenical
Council;'' i.e. 325. However, if a stop is placed after ``ago,'' the
former clause may be made to refer to S. Alexander's condemnation of
them, as Montfaucon observes. On the other hand it is plainly proved
from §. 7, that it was written just as the Arians were sending George
of Cappadocia to Alexandria, i.e. before Easter 356, and after Feb. 9,
the date of Athanasius's leaving Alexandria. The stress too which is
laid upon maintaining the Nicene Creed, and the notice of the Arian
appeal to Scripture, and the respectful language he uses of
Constantius, all agree with the date 356, if corroboration is
necessary. There is very little in this Epistle which is not contained
in his other Treatises, and a considerable portion is of a doctrinal
character. It was written on occasion of an attempt made by the Arians
to seduce the Bishops addressed into subscribing one of the specious
Creeds of which so much is read in the history of the times; but
nothing can be gathered of the circumstances from collateral sources.
The Treatise was formerly put at the head of the Orations against the
Arians, and numbered as the first of them.]

———————

\xsection{Chapter 1.}

§. 1.

1. A\textsc{ll} things whatsoever
our Lord and Saviour Jesus Christ, as Luke hath written, \emph{both did
and taught} [Acts i. 1.], He did for our salvation, for which He
appeared in the world; for He came, as John saith, not to
condemn the world, but that the world through Him might be saved [John
iii. 17.]. And among the rest we have especially to admire this
instance of his goodness, that He was not silent concerning those who
should fight against us, but plainly told us beforehand, that, when
those things should come to pass, we might straightway be found with
minds established by his teaching. For He said, There shall arise
false prophets and false Christs, and shall shew great signs and
wonders; insomuch that, if it were possible, the very elect shall be
deceived. Behold, I have told you before [Mat. xxiv. 24. 25.].
Manifold indeed and beyond human conception are the instructions and
gifts of grace which He has laid up in us; as the pattern of heavenly
conversation, power against devils, the adoption of sons, and that
exceeding great and singular grace, the knowledge of the Father and of
the Word Himself, and the gift of the Holy Ghost. But the mind of man
is prone to evil exceedingly; moreover, our adversary the devil,
envying us the possession of such great blessings, goeth about seeking
to snatch away the seed of the word which is sown within us. Wherefore
as if by his prophetic warnings He would seal up His instructions in
our hearts as His own peculiar treasure, the Lord said, \emph{Take heed
that no man deceive you: for many shall come in my name, saying, I am
he; and the time draweth near; and they shall deceive many: go ye not
therefore after them} [Luke xxi. 8.].

2. This is a great gift which the Word has
bestowed upon us, that we should not be deceived by appearances, but
that, howsoever these things are concealed, we should distinguish them
by the grace of the Spirit. For whereas the inventor of wickedness and
great spirit of evil, the devil, is utterly hateful, and as soon as he
shews himself is rejected \footnote{Margin Notes

1. [\emph{balletai}], vid. vol. 8. p.53, note f.

Return to text}
of all men,—as a serpent, as a dragon, as a lion seeking whom he may
seize upon and devour,—therefore he conceals and covers what he
really is, and craftily personates that Name which all men desire, so
that deceiving by a false appearance, he may bind fast in his chains
those whom he has led astray. And as if one that desired to kidnap the
children of others during the absence of their parents, should
personate their appearance, and so putting a cheat on the affections
of the offspring, should carry them far away and destroy them in
like manner this evil and wily spirit the devil, having no confidence
in himself, and knowing the love which men bear to the truth, puts on
the resemblance thereof, and so spreads his poison among those that
follow after him.

§. 2.

3. Thus he deceived Eve, not speaking his own,
but artfully adopting the words of God, and perverting their meaning.
Thus he suggested evil to the wife of Job, persuading her to feign
affection for her husband, while he taught her to blaspheme God. Thus
does the crafty spirit mock men by false appearances, deluding and
drawing each into his own pit of wickedness. When of old he deceived
the first man Adam, thinking that through him he should have all men
subject unto him, he exulted with great boldness and said, \emph{My hand
hath found as a nest the riches of the people; and as one gathereth
eggs that are left, have I gathered all the earth; and there was none
that moved the wing, or opened the mouth, or peeped} [Is. x. 14.].
But when the Lord came upon earth, and the enemy made trial of his
human economy, being unable to deceive the flesh which He had taken
upon Him, from that time forth He, who promised Himself the occupation
of the whole world, is for His sake mocked even by children: that
proud one is mocked as a sparrow \footnote{vid. Job xl. 24. Sept.

Return to text}. For now the infant child lays his hand upon the hole of the
asp, and laughs at him that deceived Eve; and all that rightly believe
in the Lord tread under foot him that said, \emph{I will ascend above the
heights of the clouds: I will be like the Most High} [Is. xiv.
14.].

4. Thus he suffers and is dishonoured; and
although he still ventures with shameless confidence to disguise
himself, yet now, wretched spirit, he is detected the rather by them
that bear the Sign on their foreheads; yea, more, he is rejected of
them, and is humbled, and put to shame. For even if, now that he is a
creeping serpent, he shall transform himself into an angel of light,
yet his deception will not profit him; for we have been taught that \emph{though
an angel from heaven preach unto us any other gospel than that we have
received, he shall be accursed} [vid. Gal. i. 8, 9.]. §. 3. And
although, again, he conceal his natural falsehood, and pretend to
speak truth with his lips; yet are we \emph{not ignorant of his devices}[2
Cor. ii. 11.], but are able to answer him in the words spoken by the
Spirit against him; But \emph{unto the ungodly, why dost thou
preach My law?} [Ps. l. 16.] and, \emph{Praise is not seemly in the
mouth of a sinner} [Ecclus. xv. 9.]. For even though he speak the
truth, the deceiver is not worthy of credit.

5. And whereas Scripture has shewn this, when
relating his wicked artifices against Eve in Paradise, so the Lord
also reproved him,—first in the mount, when he laid open \emph{the
folds of his breast-plate} [Job xli. 4. Sept.] \footnote{vid. vol. 8, p. 178.

Return to text}, and shewed who the crafty spirit was, and proved that it was
not one of the saints \footnote{or sacred writers, [\emph{hagi\={o}n}].

Return to text},
but Satan that was tempting Him. For He said, \emph{Get thee behind me,
Satan; for it is written, Thou shalt worship the Lord thy God, and Him
only shalt thou serve} [Mat. iv. 10.]. And again, when He put a
curb in the mouths of the devils that cried after Him from the tombs.
For although what they said was true, and they lied not then, saying, \emph{Thou
art the Son of God} [Mat. viii. 29.], and \emph{the holy One of God}
[Mark i. 24.]; yet He would not that the truth should proceed from an
unclean mouth, and especially from such as them, lest under pretence
thereof they should mingle with it their own malicious devices, and
sow \footnote{vol. 8, p. 5, note k.

Return to text} them while men
slept. Therefore He suffered them not to speak these words, neither
would He have us to suffer such, but hath charged us by His own mouth,
saying, \emph{Beware of false prophets, which come to you in sheeps'
clothing, but inwardly they are ravening wolves} [Mat. vii. 15.];
and by the mouth of His holy Apostles, \emph{Believe not every spirit}
[1 John iv. 1.].

6. Such is the method of our adversary's
operations; and of the like nature are all these inventions of
heresies, each of which has for the father of its own device the
devil, who changed and became a murderer and a liar from the
beginning. But being ashamed to profess his hateful name, they usurp
the glorious Name of our Saviour \emph{which is above every name}
[Phil. ii. 9.], and deck themselves out in the language of Scripture,
speaking indeed the words, but stealing away the true meaning thereof;
and so disguising by some artifice their false inventions, they also
become the murderers of those whom they have led astray. For to what
benefit do Marcion and Manichæus receive the Gospel while they reject
the Law \footnote{vol. 8, p. 189, r. 1.

Return to text}? For the New
Testament arose out of the Old, and bears witness to the Old; if then
they reject this, how can they receive that which proceeds from it?
Thus Paul was an Apostle of the Gospel, \emph{which God promised}
\emph{afore by His prophets in the holy Scriptures} [Rom. i.
2.]; and our Lord Himself said, \emph{Search the Scriptures, for they are
they which testify of Me} [John v. 39.]. How then shall they
confess the Lord, unless they first search the Scriptures which are
written concerning Him? And the disciples say that they have found
Him, \emph{of whom Moses and the Prophets did write} [John i. 45.].

[§. 4.?]

7. And to what end do the Sadducees retain the
Law, if they receive not the Prophets \footnote{vid. Prideaux, Conn. Ii. 5. (vol. 3. p. 474. ed. 1725.)

Return to text}? For God who gave the Law, Himself promised in the Law that He
would raise up Prophets also, so that the same is Lord both of the Law
and of the Prophets, and he that denies the one must of necessity deny
the other also. And again, how can the Jews receive the Old Testament,
unless they acknowledge the Lord whose coming was expected according
to it? For had they believed the writings of Moses, they would have
believed the words of the Lord; for He said, \emph{He wrote of Me}
[John v. 46.]. Moreover, what are the Scriptures to Paul \footnote{vol. 8. p. 16, note i.

Return to text} of Samosata, who denies the Word of God and His incarnate
Presence \footnote{ibid. p. 252, note g.

Return to text}, which is
signified and declared both in the Old and New Testament? And of what
use are the Scriptures to the Arians also, and why do they bring them
forward, men who say that the Word of God is a creature, and like the
Gentiles, \emph{serve the creature more than} God \emph{the Creator}
[Rom. i. 25.]? Thus each of these heresies, in respect of the peculiar
impiety of its invention \footnote{[\emph{epinoias}].

Return to text},
has nothing in common with the Scriptures. And their advocates are
aware of this, that the Scriptures are very much, or rather
altogether, opposed to the doctrines of every one of them; but for the
sake of deceiving the more simple sort, (such as are those of whom it
is written in the Proverbs, \emph{The simple believeth every word} [Prov.
xiv. 15.],) they pretend like their \emph{father the devil} [John
viii. 5.] \footnote{Orat. ii. 73, 74. vol. 8. p.9, note 3.

Return to text} to study
and to quote the language of Scripture, in order that they may appear
by their words to have a right belief, and so may persuade their
wretched followers to believe contrary to the Scriptures \footnote{ibid. p. 189.

Return to text}.

8. Assuredly in every one of these heresies the
devil has thus disguised himself, and has suggested to them words full
of craftiness. The Lord spake concerning them, that \emph{there shall
arise false Christs and false prophets, so that they shall deceive
many} [Mat. xxiv. 24.]. Accordingly the Devil has come, speaking by each and saying, ``I am Christ, and the truth is with me;''
and he has made them, one and all, to be liars like himself. And
strange it is, that while all heresies are at variance with one
another concerning the mischievous inventions which each has framed,
they are united together only by the common purpose of lying \footnote{vol. 8, p. 187, note b. vid. Orat. ii. §. 18.

Return to text}. For they have one and the same father that has sown in them
all the seeds of falsehood. Wherefore the faithful Christian and true
disciple of the Gospel, having grace to discern spiritual things, and
having built the house of his faith upon a rock, stands continually
firm and secure from their deceits. But the simple person, as I said
before, that is not thoroughly grounded in knowledge, such an one,
considering only the words that are spoken and not perceiving their
meaning \footnote{p. 134, r. 4.

Return to text}, is
immediately drawn away by their wiles. Wherefore it is good and
needful for us to pray that we may receive the gift of discerning
spirits, so that every one may know, according to the precept of John,
whom he ought to reject and whom to receive as friends and of the same
faith. Now one might write at great length concerning these things, if
one desired to go into details respecting them; for the impiety and
perverseness of heresies will appear to be manifold and various, and
the craft of the deceivers to be very terrible. But since holy
Scripture is of all things most sufficient \footnote{vol. 8, p. 81, r. 4.

Return to text} for us, therefore recommending to those who desire to know
more of these matters, to read the Divine word, I now hasten to set
before you that which most claims attention, and for the sake of which
principally I have written these things.

§. 5.

9. I have heard during my sojourn in these parts
\footnote{\emph{f}. Palestine, Tillem. vol. 8, p. 186.

Return to text}, (and they were true
and orthodox brethren that informed me,) that certain professors of
Arian opinions have met together, and have drawn up a confession of
faith to their own liking, and that they intend to send word to you,
that you must either subscribe to what pleases them, or rather to what
the Devil has inspired them with, or in case of refusal must suffer
banishment. They have indeed already begun to molest the Bishops of
these parts; and thereby have plainly manifested their disposition.
For inasmuch as they have framed this document only for the purpose of
inflicting banishment or other punishments, what does such
conduct prove them to be, but enemies of the Christians, and friends
of the Devil and his angels? and especially, since they spread abroad
what they like contrary to the mind of that gracious Prince, our most
religious Emperor Constantius \footnote{vol. 8, p. 90, note p.

Return to text}. And this they do with great craftiness, and, as appears to
me, chiefly with two ends in view; first, that by obtaining your
subscriptions, they may seem to remove the evil repute that rests upon
the name of Arius, and may escape notice themselves as if not
professing his opinions; and again, that by putting forth these
statements they may cast a shade over the Council of Nicæa \footnote{ibid. p. 84, note b.

Return to text}, and the confession of faith which was then put forth against
the Arian heresy.

10. But this proceeding does but prove the more
plainly their own maliciousness and heterodoxy. For had they believed
aright, they would have been satisfied with the confession put forth
at Nicæa by the whole Ecumenic \footnote{ibid. p. 49, note o, de Syn. passim.

Return to text} Council; and had they considered themselves calumniated and
falsely called Arians, they ought not to have been so eager to
innovate upon what was written against Arius, lest what was directed
against him might seem to be aimed at them also. This however is not
the course they pursue, but they conduct the struggle in their own
behalf; just as if they were Arians. Observe how entirely they
disregard the truth, and how every thing they say and do is for the
sake of the Arian heresy. For in that they dare to question those
sound definitions of the faith, and take upon themselves to produce
others contrary to them, what else do they but accuse the Fathers, and
stand up in defence of that heresy which they opposed and protested
against? And what they now write proceeds not from any regard for the
truth, as I said before, rather they do it as in mockery and by an
artifice, for the purpose of deceiving men; that by sending about
their letters they may engage the ears of the people to listen to
these notions, and so put off the time when they will be brought to
trial; and that by concealing their impiety \footnote{p. 35, r. 1.

Return to text} from observation, they may have room to extend their heresy,
which \emph{like a gangrene} [2 Tim. ii. 17.] eats its way every
where.

§. 6.

11. Accordingly they disturb and disorder every
thing, and yet are never satisfied with their own proceedings. For
every year, as if they were going to draw up a contract, they
meet together and pretend to write about the faith, whereby they
expose themselves the more to ridicule and disgrace, because their
expositions are rejected, not by others, but by themselves \footnote{vol. 8. p. 2, note c.

Return to text}. For had they had any confidence in their statements, they
would not have desired to draw up others; nor again, rejecting these
last, would they now have set down the one in question, which no doubt
true to their custom, they will again alter, after a very short
interval, and as soon as they shall find a pretence for their
customary plotting against certain persons. For when they have a
design against any, then it is that they make a great shew of writing
about the faith; that, as Pilate washed his hands, so they by a like
proceeding may destroy those who rightly believe in Christ, hoping
that, as making definitions about the faith, they will appear, as I
have repeatedly said, to be free from the charge of false doctrine.

12. But they will not be able to hide themselves,
nor to escape; for they continually become their own accusers \footnote{ibid. p. 6, note o.

Return to text} even while they defend themselves. Justly so, since instead of
answering those who bring proof against them, they do but persuade
themselves to believe whatever they wish. And when is an acquittal
obtained, upon the criminal becoming his own judge? hence it is that
they are always writing, and always altering their own previous
statements, and thus they shew an uncertain faith \footnote{ibid. p. 76, note k. p. 81, note t.

Return to text}, or rather a manifest unbelief and perverseness. And this, it
appears to me, must needs be the case with them; for since, having
fallen away from the truth, and desiring to overthrow that sound
confession of faith which was drawn up at Nicæa, they have, in the
language of Scripture, \emph{loved to wander, and have not refrained
their feet} [Jer. xiv. 10.]; therefore, like Jerusalem of old, they
labour and toil in these their changes, sometimes writing one thing,
and sometimes another, but only for the sake of gaining time, and that
they may continue enemies of Christ \footnote{[\emph{christomachoi}], p. 134. §. 8. init. p. 135, r. 3.

Return to text}, and deceivers of mankind.

§. 7.

13. Who then, that has any real regard for truth,
will be willing to suffer these men any longer? who will not justly
reject their expositions? who will not denounce their audacity, that
being but few \footnote{supr. p. 55, r. 2. vol. 8. p. 80, note s. (i.e. r.)

Return to text} in
number, they would have their decisions to prevail over every
thing, and as desiring the supremacy of their own meetings, held in
corners and suspicious in their circumstances, would forcibly cancel
the decrees of an uncorrupt, pure, and Ecumenic Council? Men who have
been promoted by the Eusebians for advocating this Antichristian
heresy, venture to define articles of faith, and while they ought to
be brought to judgment as criminals, like Caiaphas, they take upon
themselves to judge. They compose a Thalia \footnote{vol. 8. p. 94, \&c.

Return to text}, and would have it received as a standard of faith, while they
are not yet themselves determined what they believe.

14. Who does not know that Secundus \footnote{ibid. pp. 88, 89. supr. p. 44.

Return to text} of Pentapolis, who was several times degraded long ago, was
received by them for the sake of the Arian fanaticism; and that George
\footnote{supr. p. 25, note F.

Return to text} now of Laodicea, and
Leontius the Eunuch, and before him Stephanus, and Theodorus of
Heraclea \footnote{supr. p. 60.

Return to text}, were
promoted by them? Ursacius and Valens also, who from the first were
instructed by Arius as young men \footnote{supr. p. 31, note M.

Return to text}, though they had been formerly degraded from the Priesthood,
afterwards got the title of Bishops on account of their impiety; as
did also Acacius, Patrophilus \footnote{omitted and rightly (?) supr. p. 68.

Return to text}, and Narcissus, who have been most forward in all manner of
impiety. These were degraded in the great Council of Sardica;
Eustathius also now of Sebastea, Demophilus and Germinius \footnote{vol. 8. pp. 85, 86.

Return to text}, Eudoxius and Basil, who are supporters of that impiety, were
advanced in the same manner. Of Cecropius \footnote{of Nicomedia.

Return to text}, and him they call Auxentius, and of Epictetus \footnote{vid. Hist. Mon. §. 74 fin.

Return to text} the stage-player, it were superfluous for me to speak, since
it is manifest to all men, in what manner, on what pretexts, and by
what enemies of ours these were promoted, that they might play their
false charges against the orthodox Bishops who were the objects of
their designs. For although they resided at the distance of eighty
posts \footnote{supr. p. 50, note H.

Return to text}, and were
unknown to the people, yet on the ground of their impiety they were
able to procure for themselves the title of Bishop. For the same
reason also they have now \footnote{p. 125.

Return to text}
hired one George of Cappadocia, whom they wish to impose upon you. But
no respect is due to him any more than to the rest; for there is a
report in these parts that he is not even a Christian, but is devoted
to the worship of idols; and he has a hangman's temper \footnote{vol. 8. p. 134, note f.

Return to text}. And this person, such as he is described to be, they have
taken into their ranks, that they may be able to injure, to
plunder, and to slay; for in these things he is a great proficient,
but is ignorant of the very principles of the Christian faith.

§. 8.

15. Such are the machinations of these men
against the truth: but their designs are manifest to all the world,
though they attempt in ten thousand ways, like eels, to elude the
grasp, and to escape detection as enemies of Christ. Wherefore I
beseech you, let no one among you be deceived, no one seduced by them;
rather, considering that a sort of judaical impiety is invading the
Christian faith, be ye all zealous for the Lord; hold fast, every one,
the faith we have received from the Fathers, which they who assembled
at Nicæa recorded in writing \footnote{vol. 8, p. 49, note p.

Return to text}, and endure not those who endeavour to innovate thereon. And
however they may quote phrases out of the Scripture, endure not their
compositions; however they may speak the language of the orthodox, yet
attend not to what they say; for they speak not with an upright mind,
but putting on such language like sheeps' clothing, in their hearts
they think with Arius, after the manner of the devil \footnote{supr. p. 129.

Return to text} who is the author of all heresies. For he too made use of the
words of Scripture, but was put to silence by our Saviour. For if he
had indeed meant them as he used them, he would not have fallen from
heaven; but now having fallen through his pride \footnote{Cypr. Treat. Tr. p. 24, note a.

Return to text}, he artfully dissembles in his speech, and oftentimes
maliciously endeavours to lead men astray by the subtleties and
sophistries of the Gentiles.

16. Had these expositions of theirs proceeded
from the orthodox \footnote{vol. 8, p. 17, note m.

Return to text},
from such as the great Confessor Hosius, and Maximinus \footnote{supr. p. 77, r. 2.

Return to text} of Gaul, or his successor, or from such as Philogonius and
Eustathius \footnote{at Nicæa as most of the others.

Return to text}, Bishops
of the East \footnote{i.e. of Antioch.

Return to text}, or
Julius and Liberius of Rome, or Cyriacus of Mysia \footnote{of Paphos? Leont. In Nest. I. p. 550. [ed. Can.]

Return to text}, or Pistus and Aristæus of Greece, or Silvester and
Protogenes of Dacia, or Leontius and Eupsychius of Cappadocia, or
Cecilian of Africa, or Eustorgius of Italy, or Capito of Sicily, or
Macarius of Jerusalem, or Alexander of Constantinople, or Pederos of
Heraclea, or those great Bishops Meletius, Basil, and Longianus, and
the rest from Armenia and Pontus, or Lupus and Amphion from Cilicia,
or James and the rest from Mesopotamia, or our own blessed Alexander,
with others of the same sentiments as these;—there would then
have been nothing to suspect in their statements, for the minds of
apostolical men are sincere and incapable of fraud. §. 9. But when
they proceed from those who are lured to advocate the cause of heresy,
and since, according to the divine proverb, \emph{The words of the wicked
are to lie in wait}, and \emph{The mouth of the wicked poureth out
evil things}, and \emph{The counsels of the wicked are deceit} [Prov.
xii. 6. xv. 28. xii. 5.]: it becomes us to watch and be sober,
brethren, as the Lord has said, lest any deception arise from subtlety
of speech and craftiness; lest any one come and pretend to say, 'I
preach Christ,' and after a little while he be found to be Antichrist.
These indeed are Antichrists, whosoever come to you in the cause of
the Arian fanaticism.

17. For what defect is there among you, that any
one need to come to you from without? Or, of what do the Churches of
Egypt and Libya and Alexandria stand so much in need, that these men
should make a purchase \footnote{Ap. ad Const. §. 28. Hist. Arian. §. 73, supr.

Return to text}
of the Episcopate as of wood and goods, and intrude into Churches
which do not belong to them? Who is not aware, who does not perceive
clearly, that they do all this in order to support their impiety?
Wherefore although they should make themselves mute, or although they
should bind on their garments larger borders than the Pharisees, and
pour themselves forth in long speeches, and practise the tones of
their voice \footnote{vid. Basil. Ep. 223. 3.

Return to text}, they
ought not to be believed; for it is not the mode of speaking, but the
intentions of the heart and a godly conversation that recommend the
faithful Christian. And thus the Sadducees and Herodians, although
they had the law in their mouths, were put to rebuke by our Saviour,
who said unto them, \emph{Ye do err, not knowing the Scriptures, nor the
power of God} [Mat. xxii. 29.]: and all men witnessed the exposure
of those who pretended to quote the words of the Law, as being in
their minds heretics and enemies of God \footnote{[\emph{theomachoi}].

Return to text}. Others indeed they deceived by these professions, but when
our Lord became man they were not able to deceive Him; for \emph{the Word
was made flesh} [John i. 14.], who \emph{knoweth the thoughts of men
that they are vain} [Ps. xciv. 11.]. Thus He exposed time evasions
of the Jews, saying, \emph{If God were your Father, ye would love Me, for
I proceeded forth from the Father, and am come to you} [John viii.
42. xvi. 28.] \footnote{[\emph{h\={e}k\={o}}], vid. Hipp. contr. Noet. 16. and vol. 8,
p. 98, r. 1.

Return to text}. In
like manner these men seem now to act; for they disguise their real
sentiments, and make use of the language of Scripture in their
writings, which they hold forth as a bait for the ignorant, that they
may inveigle them into their own wickedness.

§. 10.

18. Consider, whether this be not so. If, when
there is no reason for their doing so, they write confessions of
faith, it is a superfluous, and perhaps also a mischievous proceeding,
because, when no question is proposed for consideration, they give
occasion for controversy of words, and unsettle the simple hearts of
the brethren, disseminating among them such notions as have never
entered into their minds. And if they profess to clear themselves in
regard to the Arian heresy, they ought first to remove the seeds of
those evils which have sprung up, and to proscribe those who produced
them, and then in the room of former statements to set forth others
which are sound; or else let them openly vindicate the opinions of
Arius, that they may no longer covertly but openly shew themselves
enemies of Christ \footnote{p. 132, r. 4.

Return to text},
and that all men may fly from them as from the sight of a serpent. But
now they keep back those opinions, and for a pretence write on other
matters; just as if a surgeon, when summoned to attend a person
wounded and suffering, should upon coming in to him say not a word
concerning his wounds, but proceed to discourse about his sound limbs.
Such an one would be chargeable with utter stupidity, for saying
nothing on the matter for which he came, but discoursing on those
other points in which he was not needed. Yet just in the same manner
these men omit those matters which concern their heresy, and take upon
themselves to write on other subjects; whereas, if they had any regard
for the Faith, or any love for Christ, they ought first to remove out
of the way those blasphemous expressions uttered against Him, and then
in the room of them to speak and to write sound words. But this they
neither do themselves, nor permit those that desire to do so, whether
it be from ignorance, or through craft and artifice.

§. 11.

19. If they do this from ignorance they must be
charged with rashness, because they affirm positively concerning
things that they know not; but if they dissemble knowingly, their
condemnation is the greater, because while they overlook nothing in
consulting for their own interests, in writing about faith in
our Lord they make a mockery, and do any thing rather than speak the
truth; they keep back those particulars respecting which their heresy
is accused, and merely bring forward passages out of the Scriptures.
Now this is a manifest robbery of the truth, and a practice full of
all iniquity; and so I am sure your piety will readily perceive it to
be from the following illustrations. No person being accused of
adultery defends himself as innocent of theft; nor would any one in
prosecuting a charge of murder suffer the accused parties to defend
themselves by saying, 'We have not committed perjury, but have
preserved the deposit which was entrusted to us.' This would be mere
child's play, instead of a refutation of the charge and a
demonstration of the truth. For what has murder to do with a deposit,
or adultery with theft? The crimes are indeed related to each other as
proceeding from the same evil heart; yet in respect to the refutation
of an alleged offence, they have no connection with each other.

20. Accordingly as it is written in the Book of
Jesus the son of Nave, when Achan was charged with theft, he did not
excuse himself with the plea of his zeal in the wars; but being
convicted of the offence was stoned by all the people. And when Saul
was charged with negligence and a breach of the law, he did not
benefit his cause by alleging his conduct on other matters. For a
defence in one sort will not operate to obtain an acquittal in another
sort; but if all things should be done according to law and justice, a
man must defend himself in those particulars wherein he is accused,
and must either disprove the past, or else confess it with the promise
that he will do so no more. But if he is guilty of the crime, and will
not confess, but in order to conceal the truth speaks on other points
instead of the one in question, he shews plainly that he has acted
amiss, nay, and is conscious of his delinquency. But what need of many
words, seeing that these persons are themselves the accusers of the
Arian heresy? For since they have not the boldness to speak out, but
conceal their blasphemous expressions, it is plain that they know that
this heresy is separate and alien from the truth. But since they
conceal this and are afraid to speak, it is necessary for me to strip
off the veil from their impiety, and to expose the heresy to
public view, knowing as I do the statements which the Arians formerly
made, and how they were cast out of the Church, and degraded from the
Clergy. But here first I ask for pardon \footnote{vol. 8, p. 216, note c.

Return to text} of the foul words to which I am about to give utterance, since
I use them, not because I thus think, but in order to convict the
heretics.

\xsection{Chapter 2.}

§. 12.

1. N\textsc{ow} the Bishop Alexander
of blessed memory cast Arius out of the Church for holding and
maintaining the following sentiments \footnote{Margin Notes

1. vol. 8. pp. 10, 94, 185.

Return to text}: ``God
was not always a Father: The Son was not always: But whereas all
things were made out of nothing, the Son of God also was made out of
nothing: And since all things are creatures, He also is a creature and
a production \footnote{[\emph{poi\={e}ma}].

Return to text}: And since all things
once were not, but were afterwards made, there was a time when the
Word of God Himself was not; and He was not before He was begotten \footnote{[\emph{genn\={e}th\={e}nai}], vol. 8, p. 272.

Return to text}, but He had a beginning \footnote{[\emph{arch\={e}n}].

Return to text} of
existence: For He was then begotten when God determined to produce \footnote{[\emph{d\={e}miourg\={e}sai}]

Return to text} Him: For He also is one among the rest of His works. And since He
is by nature changeable \footnote{[\emph{treptos}], vid. vol. 8. p. 230, note a.

Return to text}, and only
continues good because He chooses by his own free will, He is capable
of being changed, as are all other things, whenever he wishes. And
therefore God, as foreknowing that He would be good, gave Him by
anticipation that glory which He would have obtained afterwards by His
virtue; and He is now become good by His works which God foreknew.''
Accordingly they say, that Christ is not truly God, but that He is
called God on account of His participation in God's nature, as are all
other creatures. And they add, that He is not that Word which is by
nature in the Father, and is proper to His Substance, nor is He His
proper wisdom by which He made this world; but that there is another
Word \footnote{ibid. p. 101.

Return to text} which is properly \footnote{[\emph{idios}].

Return to text} in the Father, and another Wisdom which is properly in the Father,
by which Wisdom also He made this Word; and that the Lord Himself is
called the Word by a fiction \footnote{[\emph{kat' epinoian}].

Return to text} in regard
of things endued with reason \footnote{Orat. ii. 38.

Return to text}, and is
called the Wisdom fictitiously in regard of things endued with wisdom.
Nay, they say that as all things are in substance separate and
alien from the Father, so He also is in all respects separate and
alien from the substance of the Father, and properly belongs to things
made and created, and is one of them; for He is a creature, and a
production, and a work.

2. Again, they say that God did not create us for
His sake, but Him for our sakes. For they say, ``God was alone, and the
Word was not with Him, but afterwards when He would create us \footnote{vol. 8. p. 12.

Return to text}, then He made Him; and from the time He was made, He called Him
the Word, and the Son, and the Wisdom, in order that He might create
us by Him. And as all things subsisted by the will of God, and did not
exist before; so He also was made by the will of God, and did not
exist before. For the Word is not the proper and natural Offspring of
the Father, but was Himself made by grace: for God who existed before
made by His will the Son who did not exist, by which will also He made
all things, and produced, and created, and willed them to be.'' \footnote{[\emph{genesthai}].

Return to text} Moreover they say also, that Christ is not the natural and true
power of God ; but as \emph{the locust} and \emph{the cankerworm}
[Joel ii. 25.] are called a power \footnote{ibid. p. 100.

Return to text},
so also He is called the power of the Father. Furthermore he said,
that the Father cannot be described by the Son, and that the Son can
neither see nor know the Father perfectly and exactly \footnote{ibid. p. 96.

Return to text}. For having a beginning of existence, He cannot know Him that is
without beginning; but what He knows and sees, He knows and sees in a
measure proportionate to His capacity \footnote{ibid. p. 95.

Return to text}, as we also know and see in proportion to our powers. And he added
also, that the Son not only does not know His own Father exactly, but
that He does not even know His own nature \footnote{[\emph{ousian}].

Return to text}.

§. 13.

3. For maintaining these and the like opinions
Arius was declared a heretic; for myself, while I have merely been
writing them down, I have been cleansing myself \footnote{p. 138, r. 1.

Return to text} by thinking of the contrary doctrines, and by possessing my mind
with the idea of the true faith. For the Bishops who all assembled
from all parts at the Council of Nicæa, stopped their ears when they
heard these statements, and all with one voice condemned this heresy
on account of them, and anathematized it, declaring it to be alien and
estranged from the faith of the Church. It was no necessity which led
the judges to this decision, but they all by free choice
vindicated the truth \footnote{``Know,'' says St. Athan. to Jovian, ``that these
things have been preached from the beginning, and this Creed the
Fathers who assembled at Nicæa confessed, and to these have been
awarded the suffrages of all the Churches every where in their
respective places … And thou knowest that, should there be some few
who are in opposition to this faith, they cannot create any prejudice
against it, the whole world maintaining the Apostolical Creed.'' Athan.
Ep. ad Jov. §. 2. ``Whether it be persecutions or afflictions or
threats from our sovereign, or cruelties from persons in office, …
we endured it on behalf of the Apostolical faith, \&c.'' Theod. Hist.
v. 9. vid. Keble on Primitive Trad. p. 122. 10. ``Let each boldly set
down his faith in writing, having the fear of God before his eyes.''
Conc. Chalced. Sess. 1. Hard. t. 2. p. 273. ``Give diligence without
fear, favour, or dislike, to set out the faith in its purity.'' ibid.
p. 285.

Return to text}: and they did so
justly and rightly. For infidelity is coming in through these men, or
rather a Judaism beside the Scriptures, which has close upon it
Gentile superstition, so that he who holds these opinions can no
longer be called a Christian, for they are all contrary to the
Scriptures.

4. John, for instance, saith, \emph{In the beginning
was the Word} [John i. 1.]; but these men say, ``He was not, before
He was begotten.'' And again he has written, \emph{And we are in Him that
is true, even in His Son Jesus Christ; this is the true God, and
eternal life} [1 John v. 20.]; but these men, as if in
contradiction to this, allege that Christ is not the true God, but
that He is only called God, as are other creatures, in regard of his
participation in the divine nature. And the Apostle blames the
Gentiles, because they worship creatures, saying, \emph{They served the
creature more than God the Creator} [Rom. i. 25.] \footnote{supr. p. 129, vol. 8. p. 191, note d.

Return to text}. But if these men say that the Lord is a creature, and worship Him
as a creature, how do they differ from the Gentiles? If they hold this
opinion, is not this passage also against them; and does not the
blessed Paul write as blaming them? The Lord also says, \emph{I and the
Father are One}: and \emph{He that hath seen Me, hath seen the Father}
[John x. 30; xiv. 9.] \footnote{ibid. p. 229, note g.

Return to text}; and the
Apostle who was sent by Him to preach, writes, \emph{Who being the
Brightness of His glory, and the express Image of His Person} [Heb.
i. 3.]. But these men dare to separate them, and to say that He is
alien from the substance and eternity of the Father; and impiously to
represent Him as changeable, not perceiving, that by speaking thus,
they make Him to be, not one with the Father, but one with created
things. Who does not see, that the brightness cannot be separated from
the light \footnote{ibid. p. 48.

Return to text}, but that it is by nature
proper to it, and co-existent with it, and is not produced after
it? Again, when the Father says, \emph{This is My beloved Son} [Mat.
xvii. 5.], and when the Scriptures say that \emph{He is the Word} of
the Father, by whom \emph{the heavens were established} [Ps. xxxiii.
6.], and in short, \emph{All things were made by Him} [John i. 3.];
these inventors of new doctrines and fables represent that there is
another Word, and another Wisdom of the Father, and that He is only
called the Word and the Wisdom by a fiction in regard of things endued
with reason, while they perceive not the absurdity of this \footnote{p. 139.

Return to text}.

§. 14.

5. But if He be styled the Word and the Wisdom by
a fiction on our account, what He really is they cannot tell \footnote{vol. 8. p. 11, note u.

Return to text}. For if the Scriptures affirm that the Lord is both these, and yet
these men will not allow Him to be so, it is plain that in their
impious opposition to the Scriptures they would deny His existence
altogether. The faithful are able to conclude this truth both from the
voice of the Father Himself; and from the Angels that worshipped Him,
and from the Saints that have written concerning Him; but these men,
as they have not a pure mind, and cannot bear to hear the words of
holy men who teach of God, may be able to learn something even from
the devils who resemble them, for they spoke of Him, not as if there
were many beside, but, as knowing Him alone, they said, \emph{Thou art
the Holy One of God} [Mark i. 24.], and \emph{the Son of God} [Mat.
viii. 29.]. He also who suggested to them this heresy \footnote{supr. p. 129, r. 5.

Return to text}, while tempting Him in the mount, said not, 'If thou also be a Son
of God,' as though there were others beside Him, but, \emph{If Thou be
the Son of God} [Luke iv. 3.], as being the only one. But as the
Gentiles, having renounced the notion of one God, have sunk into
polytheism, so these wonderful men, not believing that the Word of the
Father is one, have come to adopt the idea of many words, and they
deny Him that is really God and the true Word, and have dared to
conceive of Him as a creature, not perceiving how full of impiety is
such an opinion. For if He be a creature, how is He at the same time
the Creator of creatures? or how the Son and the Wisdom and the Word?
For the Word is not created, but begotten; and a creature is not a
Son, but a production. And if all creatures were made by Him, and He
is also a creature, then by whom was He made? Productions must of
necessity proceed from some one; as in fact they proceeded from
the Word; because He was not Himself a production, but the Word of the
Father. And again, if the Wisdom in the Father be beside the Lord,
then there is a Wisdom in a Wisdom: and if the Word of God be the
Wisdom of God, then there is a Word in a Word: and if the Word of God
be the Son of God, then there is a Son produced in the Son.

§. 15.

6. How is it that the Lord has said, \emph{I am in
the Father, and the Father in Me} [John xiv. 10.], if there be
another in the Father, by whom the Lord Himself also was made? And how
is it that John, passing over that other, relates of this One, saying,
\emph{All things were made by Him; and without Him was not any thing made}
[John i. 3.] \footnote{vol. 8. p. 208, note a.

Return to text}? If all things that
were made by the will of God were made by Him, how can He be Himself
one of the things that were made? And when the Apostle says, \emph{For
whom are all things, and by whom are all things} [Heb. ii. 10.],
how can these men say, that we were not made for Him, but He for us?
If it be so, He ought to have said, ``For whom the Word was made;'' but
He saith not so, but, \emph{For whom are all things, and by whom are all
things}, thus proving these men to be heretical and false.

7. But further, as they have had the boldness to
say that there is another Word in God, and since they cannot bring any
clear proof of this from the Scriptures, let them but shew one work of
His, or one work of the Father that was made without this Word; so
that they may seem to have some ground at least for this their
imagination \footnote{[\emph{epinoias}].

Return to text}. The works of the true
Word are manifest to all, and according to the evidence they afford is
He known by them. For as, when we see the creation, we conceive of God
as the Creator of it; so when we see that nothing is without order
therein, but that all things move and continue with order and design,
we have an idea of a Word of God who is over all and governs all. This
too the holy Scriptures testify, declaring that He is the Word of God,
and that \emph{all things were made by Him, and without Him was not any
thing made} [John i. 3.]. But of that other Word, of whom they
speak, there is neither word nor work that they have to shew. Nay,
even the Father Himself when He says, \emph{This is My beloved Son}
[Mat. xvii. 5.], signifies that besides Him there is none other.

§. 16.

8. It appears then that so far as these doctrines
are concerned, these wonderful men have now joined themselves to the
Manichees. For these also confess the existence of a good God, so far
as the mere name goes, but they are unable to point out any of His
works either visible or invisible. But inasmuch as they deny Him who
is truly and indeed God, the Maker of heaven and earth, and of all
things invisible, they are mere inventors of fables. And this appears
to me to be the case with these evil-minded men. They see the works of
the true Word who alone is in the Father, and yet they deny Him, and
make to themselves another Word \footnote{vid. passage in Orat. ii. 39 fin.

Return to text},
whose existence they are unable to prove either by His works or by the
testimony of others. Unless it be that they have adopted a fabulous
notion of God, that He is a compound being like man, speaking and then
changing His words, and as a man exercising understanding and wisdom;
not perceiving to what absurdities \footnote{[\emph{alogian}].

Return to text}
they are reduced by such an opinion. For if God has a succession of
words \footnote{vol. 8. p. 26, note g.

Return to text}, they certainly must consider
Him as a man. And if those words proceed from Him and then vanish
away, they are guilty of a greater impiety, because they resolve into
nothing what proceeds from the self-existent God. If they conceive
that God doth at all beget, it were surely better and more religious
to say that He is the Father of One Word, who is the fulness of His
Godhead, in whom are hidden the treasures of all knowledge, and that
He is co-existent with His Father, and that all things were made by
Him; rather than to suppose God to be the Father of many words which
are no where to be found, or to represent Him who is simple in His
nature as compounded of many \footnote{ibid. p. 37, note y.

Return to text}, and as being subject to human
passions \footnote{[\emph{anthr\={o}popath\={e}}], ibid. p. 16, r. 1.

Return to text} and variable.

9. Next, whereas the Apostle says, \emph{Christ the
power of God and the wisdom of God} [1 Cor. i. 24.], these men
reckon Him but as one among many powers; nay, worse than this, they
compare Him, transgressors as they are, with the cankerworm and other
irrational creatures which are sent by Him for the punishment of men.
Next, whereas the Lord says, \emph{No one knoweth the Father, save the
Son} [Mat. xi. 27.]; and again, \emph{Not that any man hath seen the
Father, save He which is of the Father}[John vi. 46.]; are not
these indeed enemies of God \footnote{supr. p. 135, r. 3.

Return to text} which
say that the Father is neither seen nor known of the Son
perfectly? If the Lord says, \emph{As the Father knoweth Me, even so know
I the Father} [John x. 15.], and if the Father knoweth not the Son
partially, are they not mad to pretend that the Son knoweth the Father
only partially, and not fully? Next, if the Son has a beginning of
existence, and all things likewise have a beginning, let them say,
which is prior to the other. But indeed they have nothing to say,
neither can they with all their craft prove such a beginning of the
Word. For He is the true and proper Offspring of the Father, and \emph{in
the beginning was the Word, and the Word was with God, and the Word
was God} [John i. 1.]. With regard to their assertion, that the Son
knows not His own nature \footnote{[\emph{ousian}].

Return to text}, it is
superfluous to reply to it, except only so far as to condemn their
madness; for how does not the Son know Himself; when He imparts to all
men the knowledge of His Father and of Himself and blames those who
know Them not?

§. 17.

10. But it is written \footnote{Orat. ii. 18-72.

Return to text}, say they, \emph{The Lord created Me in the beginning of His ways for
His works} [Prov. viii. 22.]. O untaught and insensate that ye are!
He is called also in the Scriptures, \emph{servant}, and \emph{son of a
handmaid}, and \emph{lamb}, and \emph{sheep} [Ps. cxvi. 16.
\&c.], and it is said that He suffered toil, and thirst, and was
beaten, and endured pain. But there is plainly a reasonable ground and
cause \footnote{vol. 8. p. 22.

Return to text}, why such representations as
these are given of Him in the Scriptures; and it is because He became
man and the Son of man, and took upon Him the form of a servant, which
is the human flesh: for \emph{the Word}, says John, \emph{was made flesh}
[John i. 14.]. And since He became man, no one ought to be offended at
such expressions; for it is proper to man to be created, and born, and
formed, to suffer toil and pain, to die and to rise again from the
dead. And as, being the Word and Wisdom of the Father, He has all the
attributes of the Father, His eternity, and His unchangeableness, and
is like Him in all respects and in all things \footnote{vol. 8. p. 115, note e.

Return to text}, and is neither before nor after, but co-existent with the Father,
and is the very form \footnote{[\emph{eidos}], ibid. p. 154, note e.

Return to text} of the Godhead,
and is the Creator, and is not created: (for since He is in substance
like \footnote{ibid. p. 210, note e.

Return to text} the Father, He cannot be a
creature, but must be the Creator, as Himself hath said, \emph{My Father
worketh hitherto, and I work} [John v. 17.]:) so being made man,
and bearing our flesh, He is necessarily said to be created and
made, and to have all the attributes of the flesh; howsoever these
men, like Jewish vintners, who mix their wine with water \footnote{Orat. iii. §. 35. also vol. 8. p. 17.

Return to text}, debase the Word, and subject His Godhead to their notions of
created things.

11. Wherefore the Fathers were with reason and
justice indignant, and anathematized this most impious heresy which
these persons are now cautious of and keep back, as being easy to be
disproved and unsound \footnote{[\emph{sathran}].

Return to text} in every part
of it. These that I have set down are but a few of the arguments which
go to condemn their doctrines; but if any one desires to enter more at
large into the proof against them, he will find that this heresy is
not far removed from the Gentile superstitions, and that it is the
lowest and the very dregs of all the other heresies. These last are in
error either concerning the body or the incarnation of the Lord,
falsifying the truth, some in one way and some in another, or else
they deny that the Lord has come at all, as the Jews erroneously
suppose. But this alone more madly than the rest has dared to assail
the very Godhead, and to assert that the Word is not at all \footnote{vol. 8, p. 3, note f.

Return to text}, and that the Father was not always a father; so that one might
reasonably say that that Psalm was written against them; \emph{The fool
hath said in his heart, There is no God} \footnote{ibid. p. 184, note k.

Return to text}. \emph{Corrupt are they, and become abominable in their doings}
[Ps. liii. 1.].

§. 18.

12. ``But,'' say they, ``we are strong, and are able
to defend our heresy by our many devices.'' They would have a better
answer to give, if they were able to defend it, not by artifice nor by
Gentile sophisms, but by the simplicity of the faith. If however they
have confidence in it, and know it to be in accordance with the
doctrines of the Church, let them openly express their sentiments; for
no man when he hath lighted a candle putteth it under a bushel \footnote{vol. 8. p. 193, r. 4.

Return to text}, but on a candlestick, and so it gives light to all that come in.
If therefore they are able to defend it, let them record in writing
the opinions above imputed to them, and expose their heresy bare to
the view of all men, as they would a candle, and let them openly
accuse the Bishop Alexander, of blessed memory, as having unjustly
ejected \footnote{infr. p. 151, note B.

Return to text} Arius for professing these
opinions; and let them blame the Council of Nicæa for putting
forth a written confession of the true faith in opposition to their
impiety. But they will not do this, I am sure, for they are not so
ignorant of the evil nature of those notions which they have invented
and are ambitious of spreading abroad; but they know well enough, that
although they may at first lead astray the simple by vain deceit, yet
their imaginations will soon be extinguished, \emph{as the light of the
ungodly} [Job xviii. 5.] \footnote{vol. 8. p. 193.

Return to text}, and
themselves branded every where as enemies of the Truth.

13. Therefore although they do all things
foolishly, and speak as fools, yet in this at least they have acted
wisely, as \emph{children of this world} [Luke xvi. 8.], hiding their
candle under a bushel, that it may be supposed to give light, and
lest, if it appear, it be condemned and extinguished. Thus when Arius
himself, the author of the heresy, and the associate of Eusebius, was
summoned through the interest of the Eusebians to appear before
Constantine Augustus of blessed memory \footnote{vid. Ep. ad Serap. infr.

Return to text}, and was required to present a written declaration of his faith,
the wily man wrote one, but kept out of sight the peculiar expressions
of his impiety, and pretended, as the Devil did, to quote the simple
words of Scripture, just as they are written. And when the blessed
Constantine said to him, ``If thou holdest no other opinions in thy
mind besides these, take the Truth to witness for thee; the Lord will
repay thee if thou swear falsely:'' the wretched man swore that he held
no other, and that he had never either spoken or thought otherwise
than as he had now written. But as soon as he went out he dropped
down, as if paying the penalty of his crime, and \emph{falling headlong
burst asunder in the midst} [Acts i. 18.].

§. 19.

14. Death, it is true, is the common end of all
men, and we ought not to insult the dead, though he be an enemy, for
it is uncertain whether the same event may not happen to ourselves
before evening. But the end of Arius was not after an ordinary manner,
and therefore it deserves to be related. The Eusebians threatening to
bring him into the Church, Alexander the Bishop of Constantinople
resisted them; but Arius trusted to the violence and menaces of
Eusebius. It was the Sabbath, and he expected to join communion \footnote{[\emph{sunagesthai}].

Return to text} on the following day. There was therefore a great struggle
between them; the Eusebians threatening, Alexander praying. But the
Lord, being judge of the case, decided against the unjust party; for
the sun had not set, when the necessities of nature compelled him to
that place, where he fell down, and was forthwith deprived of
communion with the Church and of his life together. The blessed \footnote{[\emph{makarit\={e}s}].

Return to text} Constantine hearing of this soon after, was struck with wonder to
find him thus convicted of perjury. And indeed it was then evident to
all that the threats of the Eusebians had proved of no avail, and the
hope of Arius had become vain. It was shewn too that the Arian
fanaticism was rejected from communion \footnote{[\emph{akoin\={o}n\={e}tos}].

Return to text} by our Saviour both here and in the Church of the first-born in
heaven.

15. Now who will not wonder to see the
unrighteous ambition of these men, whom the Lord has condemned;—to
see them vindicating the heresy which the Lord has pronounced
excommunicate, (since He did not suffer its author to enter into the
Church,) and not fearing that which is written, but attempting
impossible timings? \emph{For the Lord of hosts hath purposed, and who
shall disannul it?} [Is. xiv. 27.] and whom the Lord hath
condemned, who shall justify? Let them however in defence of their own
imaginations write what they please; but do you, brethren, as \emph{bearing
the vessels of the Lord} [Is. lii. 11.], and vindicating the
doctrines of the Church, examine this matter, I beseech you; and if
they write in other terms than those above recorded as the language of
Arius, then condemn them as hypocrites, who hide the poison of their
opinions, and like the serpent flatter with the words of their lips.
For, though they thus write, they have associated with them those who
were formerly rejected \footnote{p. 151, note B.

Return to text} with Arius.
Such as Secundus \footnote{supr. p. 14. vol. 8. p. 88.

Return to text} of Pentapolis, and
the Clergy who were convicted at Alexandria; and they write to them in
Alexandria. But, what is most astonishing, they have caused us and our
friends to be persecuted, although the most religious Emperor
Constantine sent us back in peace to our country and Church, and
shewed his concern for the harmony of the people. But now they have
caused the Churches to be given up to these men, thus proving to all
that for the sake of the Arians the whole conspiracy against us and
the rest has been carried on from the beginning.

§. 20.

16. Now while such is their conduct, how can they
claim credit for what they write? Had the opinions they have put in
writing been orthodox, they would have expunged from their list of
books the Thalia of Arius, and have rejected the scions of the heresy,
viz. those disciples of Arius, and the partners of his impiety and his
punishment. But since they have not renounced these \footnote{vol. 8. p. 84, note b.

Return to text}, it is manifest to all that their sentiments are not orthodox,
though they write them over ten thousand times \footnote{ibid. p. 2, note c. p. 81, note t.

Return to text}. Wherefore it becomes us to watch, lest some deception be conveyed
under the clothing of their phrases, and they lead away certain from
the true faith. And if they venture to advance the opinions of Arius,
when they see themselves proceeding in a prosperous course, nothing
remains for us but to use great boldness of speech, remembering the
predictions of the Apostle, which he wrote to forewarn us of such like
heresies, and which it becomes us to repeat.

17. For we know that, as it is written, \emph{in the
latter times some shall depart from the sound faith}\footnote{ibid. p. 191, r. 3.

Return to text},\emph{giving heed to seducing spirits, and doctrines of devils, that
turn from the truth} [1 Tim. iv. 1.]; and, \emph{as many as will live
godly in Christ shall suffer persecution. But evil men and seducers
shall wax worse and worse, deceiving and being deceived} [Tit. i.
14. 2 Tim. iii. 12.]. But none of these things shall prevail over us,
nor \emph{separate us from the love of Christ} [Rom. viii. 35.],
though the heretics threaten us with death. For we are Christians, not
Arians \footnote{ibid. p. 180, note f. p. 194 fin.

Return to text}; would that they too, who
have written these things, had not embraced the doctrines of Arius!
Yea, brethren, there is need now of such boldness of speech; for we
have not received \emph{the spirit of bondage again to fear}, but God
hath called us \emph{to liberty} [\emph{v}. 15. Gal. v. 13.]. And it
were indeed disgraceful to us, most disgraceful, were we, on account
of Arius or of those who embrace and advocate his sentiments, to lose
the faith which we have received from our Saviour through His
Apostles. Already very many in these parts, perceiving the craftiness
of these writers, are ready even unto blood to oppose their wiles,
especially since they have heard of your firmness. And seeing that the
refutation of the heresy hath gone forth from you \footnote{vid. infr. p. 151, note B.

Return to text}, and it has been drawn forth from its concealment, like a serpent
from his hole, the Child that Herod sought to destroy is
preserved among you, and the Truth lives in you, and the Faith thrives
among you.

§. 21.

18. Wherefore I exhort you, having always in your
hands the confession which was framed by the Fathers at Nicæa, and
defending it with great zeal and confidence in the Lord, be ensamples
to the brethren every where, and shew them that a struggle is now
before us in support of the Truth against heresy, and that the wiles
of the enemy are various. For a martyr's token lies \footnote{vid. Suicer. Thes. in voc. [\emph{mart}]. iii.

Return to text}, not only in refusing to burn incense to idols; but to refuse to
deny the Faith is also an illustrious testimony \footnote{[\emph{martyrion}].

Return to text} of a good conscience. And not only those who turned aside unto
idols were condemned as aliens, but those also who betrayed the Truth.
Thus Judas was degraded from the Apostolical office, not because he
sacrificed to idols, but because he was a traitor; and Hymenæus and
Alexander fell away not by betaking themselves to the service of
idols, but because they \emph{made shipwreck concerning the faith} [1
Tim. i. 19.]. On the other hand, the Patriarch Abraham received the
crown, not because he suffered death, but because he was faithful unto
God; and the other Saints, of whom Paul speaks, Gedeon, Barak, Samson,
Jephtha, David, and Samuel, and the rest, were not made perfect by the
shedding of their blood, but by faith they were justified; and to this
day they are the objects of our admiration, as being ready even to
suffer death for piety towards the Lord.

19. And if one may add an instance from our own
times, ye know how the blessed Alexander contended even unto death
against this heresy, and what great afflictions and labours, old man
as he was, he sustained, until in extreme age he also was gathered to
his fathers. And how many beside have undergone great toil, in their
teachings against this impiety, and now enjoy in Christ the glorious
reward of their confession! Wherefore, let us also, considering that
this struggle is for our all, and that the choice is now before us,
either to deny or to preserve the faith, let us also make it our
earnest care and aim to guard what we have received, taking as our
instruction the Confession framed at Nicæa, and let us turn away from
novelties, and teach our people not to give heed to \emph{seducing
spirits} [1 Tim. iv. 1.] \footnote{supr. p. 149.

Return to text}, but
altogether to withdraw from the impiety of the Arian fanatics, and
from the coalition which the Meletians have made with them.

§. 22.

20. For you perceive how, though they were
formerly at variance with one another, they have now, like Herod and
Pontius, agreed together in order to blaspheme our Lord Jesus Christ.
And for this, they truly deserve the hatred of every man, because they
were at enmity with one another on private grounds, but have now
become friends and join hands, in their hostility to the Truth and
their impiety towards God. Nay, they are content to do or suffer any
thing, however contrary to their principles, for the satisfaction of
securing their several objects; the Meletians for the sake of
preeminence and the mad \footnote{[\emph{manian}].

Return to text} love of
money, and the Arian fanatics for their own impiety. And thus by this
coalition they are able to assist one another in their malicious
designs, while the Meletians pretend to the impiety of the Arians, and
the Arians from their own wickedness concur in their baseness, so that
by thus mingling together their respective crimes, like the cup of
Babylon, they may carry on their plots against the orthodox
worshippers of our Lord Jesus Christ. The wickedness and falsehood of
the Meletians were indeed even before this evident unto all men; so
too the impiety and godless heresy of the Arians have long been known
every where and to all; for the period of their existence has not been
a short one. The former became schismatics five and fifty years ago,
and it is thirty-six years since the latter were pronounced heretics \footnote{A.
This [\emph{apodeixis}] or declaration is ascribed to S. Alexander,
(as Montfaucon would explain it, supr. p. 125.) supr. p. 43. p. 146,
r. 5. p. 148, r. 3. p. 149, r. 5. vid. also p. 150. It should be
observed that an additional reason for assigning this Letter to the
year 356, is its resemblance in parts to the Orations which were
written not long after.

Return to text

}, and they were rejected from the Church by the judgment of the
whole Ecumenic Council. But by their present proceedings they have
proved at length, even to those who seem openly to favour them, that
they have carried on their designs against me and the rest of the
orthodox Bishops from the very first solely for the sake of advancing
their own impious heresy.

For observe, that which was long ago the great
object of the Eusebians is now brought about. They have caused
the Churches to be snatched out of our hands, they have banished, as
they pleased, the Bishops and Presbyters who refused to communicate
with them; and the laity who withdrew from them they have excluded
from the Churches, which they have given up into the hands of the
Arians who were condemned so long ago, so that with the assistance of
the hypocrisy of the Meletians they can without fear pour forth in
them their impious language, and make ready, as they think, the way of
deceit for Antichrist \footnote{vol. 8. p. 79, note q.

Return to text}, who sowed
among them the seeds \footnote{ibid. p. 5, note k.

Return to text} of this heresy.

§. 23.

Let them however dream and imagine vain things.
We know that when our gracious Emperor shall hear of it, he will put a
stop to their wickedness, and they will not continue long, but
according to the words of Scripture, \emph{the hearts of the impious
shall quickly fail them} [Prov. x. 20. Sept.]. But let us, as it is
written, \emph{put on the words of holy Scripture} [2 Kings xvii. 9.],
and resist them as apostates who would set up fanaticism \footnote{[\emph{manian}].

Return to text} in the house of the Lord. And let us not fear the death of the
body, nor let us emulate their ways; but let the word of Truth be
preferred before all things. I also, as you all know, was formerly
required \footnote{supr. p. 89.

Return to text} by the Eusebians either to
make pretence of their impiety, or to expect their hostility; but I
would not engage myself with them, but chose rather to be persecuted
by them, than to imitate the conduct of Judas. And assuredly they have
done what they threatened; for after the manner of Jezebel, they
engaged the treacherous Meletians to assist them, knowing how the
latter resisted the blessed \footnote{[\emph{makariou}].

Return to text} martyr
Peter, and after him the great Achillas, and then Alexander, of
blessed memory \footnote{[\emph{makaritou}], infr. p. 161, r. 4.

Return to text}, in order that, as
being practised in such matters, the Meletians might pretend against
me also whatever might be suggested to them, while the Eusebians gave
them an opening for persecuting and for seeking to kill me. For this
is what they thirst after; and they continue to this day to desire to
shed my blood.

21. But of these things I have no care; for I
know and am persuaded that they who endure shall receive a reward from
our Saviour; and that ye also, if ye endure as the Fathers did, and
shew yourselves examples to the people, and overthrow these
strange and alien devices of impious men, shall be able to glory, and
say, ``We have \emph{kept the Faith} [2 Tim. iv. 7.];'' and ye shall
receive the \emph{crown of life}, which God \emph{hath promised to them
that love Him} [James i. 12.]. And God grant that I also together
with you may inherit the promises, which were given, not to Paul only,
but also to all them that have loved the appearing of our Lord, and
Saviour, and God, and universal King, Jesus Christ; through whom to
the Father be glory and dominion in the Holy Spirit, both now and for
ever, world without end. Amen.

\xchapter{IV. Apology of our holy Father Athanasius, Archbishop of Alexandria, Addressed to the Emperor Constantius}

———————

[This Apology, which was written with a view to delivery in the Emperor's
presence, (vid. ``stretching out my hand,'' §. 3. ``I have obtained a
hearing,'' §.6. also §. 8 init. ``I see you smile,'' §. 16. also §.
22 fin. §. 27 init.) is the most finished work of its Author. It
professes to answer the new charges with which Athanasius was assailed
after his return from exile upon the Council of Sardica, i.e. between
349, when he was recalled, and 356, which is the date of its
composition. These charges were, 1. that he had influenced the Emperor
Constans to act against his brother Constantius; 2. that he had been a
zealous supporter of Magnentius, who had killed the former; 3. that he
had used a new Church for worship without the Emperor's leave; and 4.
that he had refused to leave Alexandria, which he had been forced to
do since, and to present himself at Court, which he was meditating
when he wrote this Apology. Towards the end of it, he hears news of
his own proscription, which changes his intention, and also his
feelings towards Constantius, though he preserves his respectful tone
in speaking of him to the conclusion.]

———————

§. 1.

1. K\textsc{nowing} that you have been a Christian for many
years \footnote{A. Constantius, though here called a Christian, was not baptized till
his last illness, A.D. 361, and then by the Arian Bishop of Antioch,
Euzoius. At this time he was 39 years of age. Theodoret represents him
making a speech to his whole army on one occasion, exhorting them to
baptism previously to going to war; and recommending all to go thence
who could not make up their mind to the Sacrament. Hist. iii. 1.
Constantius, his grandfather, had rejected idolatry and acknowledged
the One God, according to Eusebius, V. Const. i. 14. though it does
not appear that he had embraced Christianity.

Return to text}, most religious
Augustus, and that you are godly by descent, I cheerfully undertake to
answer for myself at this time;—for I will use the language of
the blessed Paul, and make him my advocate before you, considering
that he was a preacher of the truth, and that you are an attentive
hearer of his words.

2. With respect to those ecclesiastical matters, which have been made
the ground of a conspiracy against me, it is sufficient to refer your
Piety to the testimony of the many Bishops who have written in my
behalf \footnote{Margin Notes

1. supr. p. 14.

Return to text}; enough too is
the recantation of Ursacius and Valens \footnote{2. pp. 14, 86.

Return to text}, to prove to all men, that none of the charges which they set
up against me had any truth in them. For what evidence can others
produce so strong, as what they declared in writing? ``We lied, we
invented these things; all the accusations against Athanasius are full
of falsehood.'' \footnote{3. not supr. In Counc. Milan, 349? Montf.

Return to text} To this
clear proof may be added, if you will vouchsafe to hear it, this
circumstance, that the accusers brought no evidence against Macarius
the presbyter while we were present; but in our absence \footnote{4. pp. 31, 47. \&c.

Return to text}, when they were by themselves, they managed the matter as they
pleased. Now, the Divine Law first of all, and next our own Laws \footnote{5. Const. Apol. ii. 51. Montf.

Return to text}, have expressly declared, that such proceedings are of no force
whatsoever. From these things the piety of your Majesty, as a lover of
God and of the truth, will, I am sure \footnote{6. [\emph{oidas}] qu. [\emph{oida}].

Return to text}, perceive that we are free from all suspicion, and will
pronounce our opponents to be false accusers.

§. 2.

3. But as to the slanderous charge which has been preferred against me
before your Grace, respecting correspondence with the most pious
Augustus, your brother Constans \footnote{B. Constans had so zealously taken the part of S. Athanasius, as to
threaten his brother Constantius with war, if he did not restore him
to his see. vid. Lucifer. Op. p. 91. (ed. Ven. 1778.) This led to the
Council of Sardica. Constantius complains of Athan. in his conference
with Liberius, as ``not ceasing to exasperate Constans to quarrel with
me, had not I with superior meekness sustained the attack both of the
exasperator and the exasperated.'' Theod. Hist. ii. 13. And he says,
infra, Hist. Arian. §. 50. that he only permitted Athan.'s return for
the sake of peace.

Return to text}, of blessed and everlasting memory, (for my enemies report this
of me, and have ventured to assert it in writing,) the result of their
former \footnote{7. vid. Apol. contr. Arian. \emph{pass}.

Return to text} accusation is
sufficient to prove this also to be untrue. Had it been alleged by
another set of persons, the matter would indeed have been a fit
subject of enquiry, but it would have required strong evidence, and
open proof in presence of both parties: but when the same persons who
invented the former charge, are the authors also of this, is it
not reasonable to conclude from the issue of the one, the falsehood of
the other? For this cause they again conferred together in private,
thinking to be able to deceive your Piety before I was aware. But in
this they failed: you would not listen to them as they desired, but
patiently gave me an opportunity to make my defence. And, in that you
were not immediately moved to demand vengeance, you acted only as was
righteous in a Prince, whose duty it is to wait for the defence of the
injured party. Which if you will vouchsafe to hear, I am confident
that in this matter also, you will condemn those reckless men, who
have no fear of that God, who has commanded us not to speak falsely
before the king \footnote{8. vid. Ecclus. vii. 5.

Return to text}.

§. 3.

4. But in truth I am ashamed even to have to defend myself against
charges such as these, which I do not suppose that even the accuser
himself would venture to make mention of in my presence. For he knows
full well that he speaks untruly, and that I was never so mad, so reft
of my senses, as even to be open to suspicion of having conceived any
such thing. So that had I been questioned by any other on this
subject, I would not have answered, lest, while I was making my
defence, my hearers should for a time have suspended their judgment
concerning me. But to your Piety I answer with a loud and clear voice,
and stretching forth my hand, as I have learned from the Apostle, \emph{I
call God for a record upon my soul} [2 Cor. i. 23.], and as it is written in the
book of Kings, (let me be allowed to say the same,) \emph{The Lord is witness, and His
Anointed is witness} [1 Sam. xii. 5.], I have never spoken evil of your
Piety before your brother Constans, the most religious Augustus of
blessed memory. I have never exasperated him against you, as these
falsely accuse me. But whenever in my interviews with him he has
mentioned your Grace, (and he did mention you at the time that
Thalassus \footnote{9. Hist. Arian. 22. vid. supr. p. 79, 80.

Return to text} came to
Pitybion, and I was staying at Aquileia,) the Lord is witness, how I
spoke of your Piety in terms which I would that God would reveal unto
your soul, that you might condemn the falsehood of these my
calumniators.

5. Bear with me, most gracious Augustus, and freely grant me your
indulgence while I speak of this matter. Your most Christian brother
was not a man of so light a temper, nor was I a person of such a
character, that we should communicate together on a subject like this,
or that I should slander a brother to a brother, or speak evil of a
king before a king. I am not so mad, Sire, nor have I forgotten that
divine sentence which says, \emph{Curse
not the king, no, not thy thought; and curse not the rich in thy
bedchamber: for a bird of the air shall carry the voice, and that
which hath wings shall tell the matter} [Eccles. x. 20.]. If then those things, which are spoken in secret
against you that are kings, are not hidden, is it not incredible that
I should have spoken against you in the presence of a king, and of so
many bystanders? For I never saw your brother by myself, nor did he
ever converse with me in private, but I was always introduced in
company with the Bishop of the city, where I happened to be, and with
others that chanced to be there. We entered the presence together, and
together we retired. Fortunatian \footnote{C. All these names of Bishops occur among the subscriptions at Sardica.
supr. pp. 76-78. Fortunatian was raised to the see of Aquileia about
344, signed the condemnation of Athanasius at the Council of Milan in
355, the year before this Apology was written, and in 357 was the
Eusebian tempter in the fall of Liberius. Lucillus, Maximinus, and
Protasius, are in the list of Saints. Maximinus will be mentioned just
below, note G. Vincent, who had been the Pope's legate at Nicæa,
lapsed at Arles so far as to give up S. Athanasius, but recovered
himself by refusing to acknowledge the proceedings at Ariminum. Leis
is Lauda, or Laus Pompeia, \emph{hodie} Lodi; Ughelli, Ital. Sacr. t. 4. p. 656.

Return to text}, Bishop of Aquileia, can testify this, the father Hosius is
able to say the same, as also are Crispinus Bishop of Padua, Lucillus
of Verona, Dionysius of Lëis, and Vincentius of Campania. And
although Maximinus of Treves, and Protasius of Milan, are dead, yet
Eugenius who was Master of the Palace \footnote{D. Or, master of the offices; one of the seven Ministers of the Court
under the Empire; ``He inspected the discipline of the civil and
military schools, and received appeals from all parts of the Empire
… The correspondence between the Prince and his subjects was managed
by the four \emph{scrinia},
or offices of this minister of state … The whole business was
despatched by 148 secretaries, chosen for the most part from the
profession of the law ... But the department of foreign affairs, which
constitutes so essential a part of modern policy, seldom diverted the
attention of the master of the offices; his mind was more seriously
engaged by the general direction of the posts and arsenals of the
Empire.'' Gibbon, ch. 17.

Return to text} can bear witness for me; for he stood before the veil \footnote{E. [\emph{pro
tou b\={e}lou}].
The Veil, which in the first instance was an appendage to the images
of pagan deities, formed at this time a part of the ceremonial of the
imperial Court. It hung over the entrance of the Emperor's bedchamber,
where he gave his audiences. It also hung before the secretarium of
the Judges. vid. Hofman in voc. Gothofred in Cod. Theod i. tit. vii.
1.

Return to text}, and heard what we requested of the Emperor, and what he
vouchsafed to reply to us.

§. 4.

6. This certainly is sufficient for proof, yet suffer me nevertheless to
lay before you an account of my travels, which will further lead you
to condemn the unfounded calumnies of my opponents. When I left
Alexandria, I did not go to your brother's Court \footnote{10. [\emph{stratopedon}], vid. p. 100, note Z.

Return to text}, or to any other persons, but only to Rome \footnote{11. p. 49, §. 29.

Return to text}; and having laid my case before the Church, (for this was my
only concern,) I spent my time in the public worship \footnote{12. [\emph{sunaxesi}].

Return to text}. I did not write to your brother, except when the Eusebians
had written to him to accuse me, and I was compelled while yet at
Alexandria to defend myself; and again when I sent to him volumes \footnote{F. [\emph{puktia}], a bound book, vid. Montf. Coll. Nov. \emph{infr}.
S. Jerome speaks of Hilarion's transcribing a Gospel. Vit. Hilar. 35.
and himself the Psalter, (interpretationem Psalmorum,) ad Florent. Ep.
v. 2. and St. Eusebius of Vercellæ made a copy of the Gospels, which
was extant, as it appears, in the last century. vid. Lami Erud. Apost.
p. 678. Mabillon, Itin. Ital. t. i. p.9. Montfauc. Diar. Ital. xxviii.
p. 445. Tillemont, (t. 8. p. 86.) considers that Athan. alludes in
this passage to the Synopsis Scr. Sacr. which is among his works; but
Montfaucon, Collect. Nov. t. 2. p. xxviii. contends that a copy of the
Gospels is spoken of.

Return to text} containing the holy Scriptures, which he had ordered me to
prepare for him. It behoves me, while I defend my conduct, to tell the
truth to your Piety. When however three years had passed away, he
wrote to me in the fourth year \footnote{13. A.D.
345.

Return to text}, commanding me to meet him, (he was then at Milan;) and upon
enquiring the cause, (for I was ignorant of it, the Lord is my
witness,) I learnt that certain Bishops \footnote{G. Tillemont supposes that Constans was present at the Council of Milan,
at which Eudoxius, Martyrius, and Macedonius, sent to the West with
the Eusebian Creed, (vid. Libr. F. vol. 8. p. 111.) made their
appearance to no purpose. The Bishops principally concerned in
persuading Constans seem to have been Pope Julius, Hosius, and
Maximinus of Treves. Hil. Fragm. 2. p. 16.

Return to text} had gone up and requested him to write to your Piety, desiring
that a Council might be called. Believe me, Sire, this is the truth of
the matter; I lie not. Accordingly I went to Milan, and met with great
kindness from him; for he condescended to see me, and to say that he
had despatched letters to you, requesting that a Council might be
called. And while I remained in that city, he sent for me again into
Gaul; (for the father Hosius was going thither,) that we might travel
from thence to Sardica. And after the Council, he wrote to me while I
continued at Naissus \footnote{H. Naissus was situated in Upper Dacia, and according to some was the
birthplace of Constantine. The Bishop of the place, Gaudentius, whose
name occurs among the subscriptions at Sardica, had protected S. Paul
of Constantinople and incurred the anathemas of the Eusebians at
Philippopolis. Hil. Fragm. iii. 27.

Return to text}, and
I went up, and abode afterwards at Aquileia; where the letters
of your Piety found me. And again, being summoned thence by your
departed brother, I returned into Gaul, and so came at length to your
Piety.

§. 5.

7.Now what place and time does my accuser specify, at which I made use
of these expressions according to his slanderous imputation? In whose
presence was I so mad as to give utterance to the words which he has
falsely charged me with speaking? Who is there ready to support the
charge, and to testify to the fact? What his own eyes have seen that
ought he to speak, as holy Scripture enjoins. But no; he will find no
witnesses of that which never took place. But I take your Piety to
witness, together with the Truth, that I lie not. I request you, for I
know you to be a person of excellent memory, to call to mind the
conversation I had with you, when you condescended to see me, first at
Viminacium \footnote{14. in Mæsia.

Return to text}, a second
time at Cæsarea in Cappadocia, and a third time at Antioch. Did I
speak evil before you even of the Eusebians who have persecuted me?
Did I cast imputations upon any of those that have done me wrong? If
then I imputed nothing to any of those against whom I had a right to
speak; how could I be so possessed with madness as to slander a King
before a King, and to set a brother at variance with a brother? I
beseech you, either cause me to appear before you that the thing may
be proved, or else condemn these calumnies, and follow the example of
David, who says, \emph{Whoso
privily slandereth his neighbour, him will I destroy}
[Ps. ci. 5.]. As much as in them lies, they have slain me; for \emph{the mouth that belieth, slayeth
the soul} [Wisd. i. 11.]. But your long-suffering has prevailed against them, and
given me confidence to defend myself, that they may suffer
condemnation, as contentious and slanderous persons. Concerning your
most religious brother, of blessed memory \footnote{15. [\emph{t\={e}s
makarias mn\={e}m\={e}s}].

Return to text}, this may suffice: for you will be able, according to the
wisdom which God has given you, to gather much from the little I have
said, and to perceive that this accusation is a mere invention.

§. 6.

8. With regard to the second calumny, that I have written letters to the
usurper \footnote{I. Magnentius, a barbarian by origin, securing the troops who were about
the person of Constans, had taken possession of Autun in Gaul, where
the Emperor was, and, on the flight of the latter, had sent a party of
horse after him, by whom he was despatched. Magnentius, after some
successes, was defeated in the great battle of Mursa, and ultimately
destroyed himself at Lyons.

Return to text}, (his name I am
unwilling to pronounce;) I beseech you investigate and try the
matter, in whatever way you please, and by whomsoever you may approve
of. The extravagance of the charge so confounds me, that I am in utter
uncertainty how to act. Believe me, most religious Prince, many times
did I weigh the matter in my mind, but was unable to believe that any
one could be so mad as to utter such a falsehood. But when this charge
was published abroad by the Arians, as well as the former, and they
boasted that they had transmitted to you a copy of the letter, I was
the more amazed, and I have passed sleepless nights contending against
the charge, as if in the presence of my accusers; and suddenly
breaking forth into a loud cry, I have immediately fallen to my
prayers, desiring with groans and tears that I might obtain a
favourable hearing from you. And now that by the grace of the Lord, I
have obtained such a hearing, I am again at a loss how I shall begin
my defence; for as often as I make an attempt to speak, I am prevented
by my horror at the deed.

9. In the case of your departed brother, the slanderers had indeed a
plausible pretence for what they alleged; because I had been admitted
to see him, and he had condescended to write to your brotherly
affection concerning me; and he had often sent for me to come to him,
and had honoured me when I came. But for the traitor \footnote{16. [\emph{diabolon}].

Return to text} Magnentius, \emph{the Lord is witness, and His
Anointed is witness} [1 Sam. xii. 5.], I know him not: I never did know
him. What correspondence then could there be between persons so
entirely unacquainted with each other? What reason was there to induce
me to write to such a man? How could I have commenced my letter, had I
written to him? Could I have said, 'You have done well to murder the
man who honoured me, whose kindnesses I shall never forget?' Or, 'I
approve of your conduct in destroying our Christian friends, and most
faithful brethren?' or, 'I approve of your proceedings in butchering
those who so kindly entertained me at Rome; for instance, your
departed \footnote{17. [\emph{makarias}].

Return to text} Aunt
Eutropia \footnote{K. Nepotian, the son of Eutropia, Constantine's sister, had taken up
arms against Magnentius, got possession of Rome, and enjoyed the title
of Augustus for about a month. Magnentius put him to death, and his
mother, and a number of his adherents, some of whom are here
mentioned.

Return to text}, whose
disposition answered to her name, that worthy man, Abuterius,
the most faithful Spirantius, and many other excellent persons?' §.
7. Is it not mere madness in my accuser even to suspect me of such a
thing? What, I ask again, could induce me to place confidence in this
man? What trait did I perceive in his character on which I could rely?
He had murdered his own master; he had proved faithless to his
friends; he had violated his oath; he had blasphemed God, by
consulting poisoners and sorcerers \footnote{18. Bingh. Antiqu. xvi. 5. §. 5. \&c.

Return to text} contrary to his Law. And with what conscience could I send
greeting to such a man, whose madness and cruelty had afflicted not me
only, but all the world around me? To be sure, I was very greatly
indebted to him for his conduct, that when your departed brother had
filled our churches with sacred offerings, he murdered him. For the
wretch was not moved by the sight of these his gifts, nor did he stand
in awe of the divine grace which had been given to him in baptism: but
like a deadly and devilish spirit, he raged against him, till your
blessed \footnote{19. [\emph{makarit\={e}i}].

Return to text} brother
suffered martyrdom at his hands; while he, henceforth a criminal like
Cain, was driven from place to place, a fugitive and a vagabond, to
the end that he might follow the example of Judas in his death, by
becoming his own executioner, and so bring upon himself a double
weight of punishment in the judgment to come.

§. 8.

10. With such a man the slanderer thought that I had been on terms of
friendship, or rather he did not think so, but like an enemy invented
an incredible fiction: for he knows full well that he has lied. I
would that, whoever he is, he were present here, that I might put the
question to him on the word of Truth itself, (for whatever we speak as
in the presence of God, we Christians consider as an oath \footnote{20. vid. Chrys. in Eph. tr. p. 119, note g.

Return to text};) I say, that I might ask him this question, which of us
rejoiced most in the well-being of the departed \footnote{21. [\emph{makaritou}].

Return to text} Constans? who prayed for him most earnestly? The facts of the
foregoing charge proved this; indeed it is plain how the case stands.
But although he himself knows full well, that no one who was so
disposed towards the departed \footnote{21. [\emph{makaritou}].

Return to text}
Constans, and who truly loved him, could be a friend to his enemy, I
fear that being possessed with other feelings towards him than I
was, he has falsely attributed to me those sentiments of hatred which
were entertained by himself.

§. 9.

11. For myself, I am so surprised at the enormity of the thing, that I
am quite uncertain what I ought to say in my defence. I can only
declare, that I condemn myself to die a thousand deaths, if even the
least suspicion attaches to me in this matter. And to you, Sire, as a
lover of the truth, I confidently make my appeal. I beseech you, as I
said before, to investigate this affair, and especially to call for
the testimony of those who were once sent by him as ambassadors to
you. These are the Bishops Servatius \footnote{L. Sarbatius or Servatius, and Maximus occur in the lists of Gallic
subscriptions at Sardica. The former is supposed to be St. Servatius
or Servatio of Tungri, concerning whom at Ariminum, vid. Sulp. Hist.
ii. 59. vid. also Greg. Turon. Hist. Franc. ii. 5. where however the
Bened. Ed. prefers to read Aravatius, a Bishop, as he considers, of
the fifth century.

Return to text} and Maximus and the rest, with Clementius and Valens. Enquire
of them, I beseech you, whether they brought letters to me. If they
did, this would give me occasion to write to him. But if he did not
write to me, if he did not even know me, how could I write to one with
whom I had no acquaintance? Ask them whether, when I saw Clementius,
and spoke of your brother of blessed memory \footnote{22. [\emph{t\={e}s
makarias mn\={e}m\={e}s}],
supr. p. 159, r. 2.

Return to text}, I did not, in the language of Scripture, wet my garments with
tears, when I remembered his kindness of disposition and his Christian
spirit? Learn of them how anxious I was, on hearing of the cruelty of
that savage beast, and finding that Valens and his company had come by
way of Libya, lest he should attempt a passage also, and like a robber
murder those who held in love and memory the departed \footnote{23. [\emph{makariou}].

Return to text} Prince, among whom I account myself second to none.

§. 10.

12. How with this apprehension of such a design, was there not an
additional probability of my praying for your Grace? Should I feel
affection for his murderer, and entertain dislike towards you his
brother who avenged his death? Should I remember his crime, and forget
that kindness of yours which you vouchsafed to assure me by letter
should remain the same towards me after your brother's death of happy
memory \footnote{24. [\emph{makaritou}].

Return to text}, as it had
been during his lifetime? How could I have borne to look upon the
murderer? Must I not have thought that the blessed Prince beheld
me, when I prayed for your safety? For brothers are by nature the
mirrors of each other. Wherefore as seeing you in him, I never should
have slandered you before him; and as seeing him in you, never should
I have written to his enemy, instead of praying for your safety. Of
this, my witnesses are, first of all, the Lord who has heard and has
given to you entire the kingdom of your forefathers and next those
persons who were present at the time, Felicissimus, who was Duke of
Egypt, Rufinus, and Stephanus, the former of whom was Receiver-general
\footnote{25. supr. pp. 32, 118.

Return to text}, the latter, Master
there; Count Asterius, and Palladius Master of the palace, Antiochus
and Evagrius Official Agents \footnote{M. 1. The Rationales or Receivers, in Greek writers Catholici, ([\emph{logothetai}]
being understood, Vales, ad Euseb. vii. 10.) were the same as the
Procurators, (Gibbon, Hist. ch. xvii. note 148.) who succeeded the
Provincial Quæstors in the early times of the Empire. They were in
the department of the Comes Sacrarum Largitionum, or High Treasurer of
the Revenue, (Gothofr. Cod. Theod. t. 6. p. 327.) Both Gothofr.
however and Pancirolus, p. 134. Ed. 1623. place Rationales also under
the Comes Rerum Privatarum. Pancirolus, p. 120. mentions the Comes
Rationalis Summarum Ægypti as distinct from other functionaries.
Gibbon, ch. xvii. seems to say that there were in all 29, of whom 18
were counts. 2. Stephanus, [\emph{magistros
echei}],
Tillemont translates, ``Master of the camp of Egypt.'' vol. 8. p. 137.
3. The Master of the offices or of the palace has been noticed above,
p. 157, note D. [\emph{agentis\={e}ribous}],
agentes in rebus. These were functionaries under the Master of the
offices, whose business it was to announce the names of the consuls
and the edicts or victories of the Empire. They at length became spies
of the Court, vid. Gibbon, ch. xvii Gothofr. Cod. Th. vi. 27.

Return to text}.
I had only to say, ``Let us pray for the safety of the most religious
Emperor, Constantius Augustus,'' and all the people immediately cried
out with one voice, ``O Christ, send thy help to Constantius;'' and they
continued praying thus for some time \footnote{N. Presbyterum Erachum mihi successorem vo?o. A populo acclamatum est,
Deo gratias, Christo laudes; dictum est vicies terties. Exaudi Christe,
Augustino vita; dictum est sexies decies. Te patrem, te episcopum;
dictum est octies.'' August. Ep. 213.

Return to text}.

§. 11.

13. Now I have already called upon God, and His Word, the Only-begotten
Son our Lord Jesus Christ, to witness for me, that I have never
written to that man, nor received letters from him. And as to my
accuser, give me leave to ask him a few short questions concerning
this charge also. How did he come to the knowledge of this matter?
Will he say that he has got copies of the letter? for this is what the
Arians have declared till they were weary. Now in the first place,
even if he can shew writing resembling mine, the thing is by no means
certain; for there are forgers, who have often imitated the hand \footnote{26. [\emph{cheiros}], supr. p. 107.

Return to text} even of you who are Kings. And the resemblance will not
prove the genuineness of the letter, unless my customary amanuensis
shall testify in its favour. I would then again ask my accusers, Who
provided you with these copies? and whence were they obtained? I had
my writers \footnote{O. vid. Rom. xvi. 22. Lucian is spoken of as the amanuensis of the
Confessors, who wrote to St. Cyprian, Ep. 16. Ed. Ben. Jader perhaps
of Ep. 80. St. Jerome was either secretary or amanuensis to Pope
Damasus, vid. Ep. ad Ageruch. (123. n. 10. Ed. Vallars.) vid. Lami de
Erud. Ap. p. 258.

Return to text}, and he his
servants, who received his letters from the bearers, and gave them
into his hand. My assistants are forthcoming; vouchsafe to summon the
others, (for they are most probably still living,) and enquire
concerning these letters. Search into the matter, as though Truth were
the partner of our throne. She is the defence of Kings, and especially
of Christian Kings; with her you will reign most securely, for holy
Scripture says, \emph{Mercy
and truth preserve the king, and they will encircle his throne in
righteousness} [Prov. xx. 28.]. And the wise Zorobabel gained
a victory over the others by setting forth the power of Truth, and all
the people cried out, \emph{Great
is truth, and mighty above all things}
[1 Esdr. iv. 41.].

§. 12.

14. Had I been accused before any other, I should have appealed to your
Piety; as once the Apostle appealed unto Cæsar, and put an end to the
designs of his enemies against him. But since they have had the
boldness to lay their charge before you, to whom shall I appeal from
you? to the Father of Him who says, \emph{I am the Truth} [John xiv. 6.], that He may incline your heart unto
clemency:—

O Lord Almighty, and King of eternity, the Father of our Lord Jesus
Christ, who by Thy Word hast given this Kingdom to Thy servant
Constantius; do Thou shine into his heart, that he, knowing the
falsehood that is set against me, may both favourably receive this my
defence; and may make known unto all men, that his ears are firmly set
to hearken unto the Truth, according as it is written, \emph{Righteous lips alone are acceptable unto the King}
[Prov. xvi. 13.]. For Thou hast caused it to be said by Solomon, that
thus the throne of a kingdom shall be established.

15. Wherefore at least enquire into this matter, and let the accusers
understand that your desire is to learn the truth and see, whether
they will not shew their falsehood by their very looks; for the
countenance is a test of the conscience, as it is written, \emph{A merry heart maketh a cheerful
countenance, but by sorrow of the heart the spirit is broken}
[Prov. xv. 13.]. Thus they who had conspired against Joseph were
convicted by their own consciences; and the cruelty of Laban towards
Jacob were shewn in his countenance \footnote{27. vid. Vit. Ant. §. 67.

Return to text}. And thus you see the suspicious alarm of these persons, for
they fly and hide themselves; but on my part frankness \footnote{28. supr. pp. 49, 158.

Return to text} in making my defence. And the question between us is not one
regarding worldly wealth, but concerning the honour of the Church. He
that has been struck by a stone, applies to a physician; but sharper
than a stone are the strokes of calumny; for as Solomon has said, \emph{A
false witness is a maul, and a sword, and a sharp arrow} [Prov. xxv. 18.], and its wounds Truth alone is able to cure; and if
Truth be set at nought, they grow worse and worse.

§. 13.

16. It is this that has thrown the Churches every where into such
confusion; for pretences have been devised, and Bishops of great
authority, and of advanced age \footnote{29. Hist. Arian. 72, \&c.

Return to text}, have been banished for holding communion with me. However, if
matters stop here, our prospect is favourable through your gracious
interposition. And that the evil may not extend itself, let Truth
prevail before you; and leave not the whole Church under suspicion, as
though Christian men, nay even Bishops, could be guilty of plotting
and writing in this manner. Or if you are unwilling to investigate the
matter, it is but right that we who offer our defence, should be
believed, rather than our calumniators. They, like enemies, are
occupied in wickedness; we, as earnestly contending for our cause,
present to you our proofs. And truly I wonder how it comes to pass,
that while we address you with fear and reverence, they are possessed
of such an impudent spirit, that they dare even to lie before the King
\footnote{30. supr. p. 156, r. 1. Hist. Ar. §. 52.

Return to text}. But I pray you, for
the Truth's sake, and as it is written, \emph{search diligently}
[Joel i. 7. Sept.] in my presence, on what grounds they affirm these
things, and whence these letters were obtained. But neither will any
of my servants be proved guilty, nor will any of his people be able to
tell whence they came; for they are forgeries. And perhaps one had
better not enquire further. They do not wish it, lest the writer of
the letters should be certain of detection. For the calumniators
alone, and none besides, know who he is.

§. 14.

17. But forasmuch as they have informed against me in the matter of the
great Church, that a congregation was holden there before it was
completed, I will answer to your Piety on this charge also; for the
parties who bear so hearty an enmity against me, constrain me to do
so. I confess this did so happen; for, as in what I have hitherto
said, I have spoken no lie, I will not now deny this. But the facts
are far otherwise than they have represented them. Permit me to
declare to you, most religious Augustus, that we kept no day of
dedication, (it would certainly have been unlawful to do so), before
receiving orders from you,) nor were we led to act as we did through
premeditation. No Bishop or other Clergyman was invited to join in our
proceedings; for much was yet wanting to complete the building. Nay
the congregation was not held on a previous notice, which might give
them a reason for informing against us. Every one knows how it
happened; hear me, however, with your accustomed equity and patience.
It was the feast of Easter, and an exceeding great multitude of
Christians was assembled together, such as Christian kings would
desire to see in all their cities. Now when the Churches were found to
be too few to contain them, there was no little stir among the people,
who desired that they might be allowed to meet together in the great
Church, where they could all offer up their prayers for your safety \footnote{31. supr. p. 163. vol. 8, p. 159.

Return to text}. And this they did; for although I exhorted them to wait
awhile, and to assemble in the other Churches, with whatever
inconvenience to themselves, they would not listen to me; but were
ready to go out of the city, and meet in desert places in the open
air, thinking it better to endure the fatigue of the journey, than to
keep the fast in such a state of discomfort.

§. 15.

18. Believe me, Sire, and let Truth be my witness in this also, when I
declare that in the congregations held during the season of Lent, in
consequence of the narrows limits of our buildings, and the vast
multitude of people assembled, a great number of children, not a few
the younger and very many of the older women, besides several young
men, suffered so much from the pressure of the crowd, that they were
obliged to be carried home; though by the Providence of God,
none perished. All however murmured, and demanded the use of the great
Church. And if the pressure was so great during the days which
preceded the feast, what would have been the case during the feast
itself? Of course matters would have been far worse. It did not
therefore become me to change the people's joy into grief, their
cheerfulness into sorrow, and to make the festival a season of
lamentation.

19. And that the more, because I had a precedent in the conduct of our
Fathers. For the blessed Alexander, when the other places of worship
were too small, and he was engaged in the erection of what was then
considered a very large one, the Church of Theonas \footnote{P. S. Epiphanius mentions 9 Churches in Alexandria. Hær. 69. 2. Athan.
mentions in addition that of Quirinus. Hist. Arian. §. 10. The Church
mentioned in the text was built at the Emperor's expense; and
apparently upon the Emperor's ground, as on the site was or had been a
Basilica, which bore first the name of Hadrian, then Licinius, Epiph. \emph{ibid}. Hadrian, it should be observed, built in many
cities temples without idols, which were popularly considered as
intended by him for Christian worship, and went after his name.
Lamprid. Vit. Alex. See. 43. The Church in question was built in the Cæsareum.
Hist. Arian. 74. There was a magnificent Temple, dedicated to
Augustus, as [\emph{epibat\={e}rios}],
on the harbour of Alexandria, Philon. Legrit. ad Caium, pp. 1013, 4.
ed. 1691. and called the Cæsareum. It was near the Emperor's palace,
vid. Acad. des Inscript. vol. 9. p. 416. As to the Cæsarean Church,
it was begun by Gregory, finished by George, burnt under Julian,
rebuilt by Athanasius. Tillem. vol. 8. pp. 148, 9.

Return to text}, held his congregations there on account of the number of the
people, while at the same time he proceeded with the building. I have
seen the same thing done at Treves and at Aqiuleia, in both which
places, while the building was proceeding, they assembled there during
the feasts, on account of the number of the people; and they never
found any one to accuse them in this manner. Nay, your brother of
blessed memory was present, when a congregation was held under these
circumstances at Aquileia. I also followed this course. There was no
dedication, but only an assembly for the sake of prayer. You, at
least, I am sure, as a lover of God, will approve of the people's
zeal, and will pardon me for being unwilling to hinder the prayers of
so great a multitude.

§. 16.

20. But here again I would ask my accuser, where was it right that the
people should pray? in the desert, or in a place which was in course
of building for the purpose of prayer? Where was it becoming and pious
that the people should answer, Amen \footnote{Q. Bingham, Antiqu. xv. 3. §. 25. Tertullian, (O.T. vol. i. p. 214,
note n.) Suicer, Thesaur. \emph{in
voc}. [\emph{am\={e}n}], Gavanti, Thesaur. vol. i. p. 89. ed. 1763.

Return to text}? in the desert, or in what was already called the Lord's house?
Where would you, most religious Prince, have wished your people to
stretch forth their hands, and to pray for you? where the Greeks, as
they passed by, might stop and listen, or in a place named after
yourself, which all men have long called the Lord's house, even since
the foundations of it were laid? I am sure that you prefer your own
place; for I see you smile, and that tells me so.

21. ``But,'' says the accuser, ``it ought to have been in the Churches.''
They were all, as I said before, too small and confined to admit the
multitude. Then again, in which way was it most becoming that their
prayers should be made? Should they meet together in parts and
separate companies, with danger from the crowded state of the
congregations? or, when there was now a place that would contain them
all, should they assemble in it, and speak as with one and the same
voice in perfect harmony? This was the better course, for this showed
the unanimity of the multitude: in this way God will readily hear
prayer. For if, according to the promise of our Saviour Himself, where
two shall agree together as touching any thing that they shall ask, it
shall be done for them [Mat. xviii. 19.], how shall it be when so
great an assembly of people with one voice utter their Amen to God?
Who indeed was there that did not marvel at the sight? Who but
pronounced you a happy prince, when they saw so great a multitude met
together in one place? I How did the people themselves rejoice to see
each other, having been accustomed heretofore to assemble in separate
places! The circumstance was a source of pleasure to all; of vexation
to the calumniator alone.

§. 17.

22. Now then, I would also meet the other and only remaining objection
of my accuser. He says, the building was not completed, and prayer
ought not to have been made there. But the Lord said, \emph{But thou, when thou prayest,
enter into thy closet, and shut the doors} [Mat. vi. 6.]. What then will the accuser
answer? or rather what will all prudent and true Christians say? Let
your Majesty ask the opinion of such: for it is written of the other, \emph{The
foolish person will speak foolishness} [Is. xxxii. 6. Sept.]; but of these, Ask
counsel of all that are wise [Tob. iv. 18.]. When the Churches were
too small, and the people so numerous as they were, and desirous to go
forth into the desert, what ought I to have done? The desert has no
doors, and all who choose may pass through it, but the Lord's house is
enclosed with walls and doors, and marks the difference between the
pious and the profane. Will not every wise person then, as well as
your Piety, Sire, give the preference to the latter place? For they
know that here prayer is lawfully offered, while a suspicion of
irregularity attaches to it there. Unless indeed, no place proper for
it existed, and the worshippers dwelt only in the desert, as was the
case with Israel; although after the tabernacle was built, they also
had thenceforth a place set apart for prayer.

23. O Christ, Lord and true King of kings, Only-begotten Son of God,
Word and Wisdom of the Father, I am accused because the people prayed
Thy gracious favour, and through Thee besought Thy Father, who is God
over all, to save Thy servant, the most religious Constantius. But
thanks be to Thy goodness, that it is for this that I am blamed, and
for the keeping of Thy laws. Heavier had been the blame, and more true
had been the charge, had we passed by the place which the Emperor was
building, and gone forth into the desert to pray. How would the
accuser then have vented his folly against me! With what apparent
reason would he have said, ``He despised the place which you are
building; he does not approve of your undertaking; he passed it by in
derision; he pointed to the desert to supply the want of room in the
Churches; he prevented the people they wished to offer up their
prayers.'' This is what he wished to say, and sought an occasion of
saying it; and finding none he is vexed, and so forthwith invents a
charge against me. Had he been able to say this, he would have
confounded me with shame; as now he injures me, copying the accuser's
\footnote{32. [\emph{diabolou}], supr. p. 160. r. 1.

Return to text} ways, and watching
for an occasion against those that pray. Thus has he perverted to a
wicked purpose his knowledge of Daniel's history. But he has been
deceived; for he ignorantly imagined, that Babylonian practices were
in fashion with you, and knew not that you are a friend of the blessed
Daniel, and worship the same God, and do not forbid, but wish
all men to pray, knowing that the prayer of all is, that you may
continue to reign in perpetual peace and safety \footnote{33. p. 166. r. 1.

Return to text}.

§. 18.

24. This is what I have to complain of on the part of my accuser. But
may you, most religious Augustus, live through the course of many
years to come, and celebrate the dedication of the Church. Surely the
prayers which have been offered for your safety by all men, are no
hindrance to this celebrity. Let these unlearned persons cease such
misrepresentations, but let them learn from the example of the
Fathers; and let them read the Scriptures. Or rather let them learn of
you, who are so well instructed in such histories, how that Jesus the
son of Josedek the priest, and his brethren, and Zorobabel the wise,
the son of Salathiel, and Ezra the priest and scribe of the law, when
the temple was in course of building after the captivity, the feast of
tabernacles being at hand, (which was a great feast and time of
assembly and prayer in Israel,) gathered the people together with one
accord in the great court within the first gate, which is toward the
East, and prepared the altar to God, and there offered their gifts,
and kept the feast. And so afterwards they brought hither their
sacrifices, on the sabbaths and the new moons, and the people offered
up their prayers. And yet the Scripture says expressly, that when
these things were done, the temple of God was not yet built; but
rather while they thus prayed, the building of the house was set
forward. So that neither were their prayers deferred in expectation of
the dedication, nor was the dedication prevented by the assemblies
held for the sake of prayer. But the people thus continued to pray;
and when the house was entirely finished, they celebrated the
dedication, and brought their gifts for that purpose, and all kept the
feast for the completion of the work.

25. And thus also did the blessed Alexander, and the other Fathers. They
continued to assemble their people, and when they had completed this
work they gave thanks unto the Lord, and celebrated the dedication.
This also it befits you to do, O Prince, most careful in your
inquiries. The place is ready, having been already sanctified by the
prayers which have been offered in it, and requires only the presence
of your Piety. This only is wanting to its perfect beauty. Do
you then supply this deficiency, and there make your prayers unto the
Lord, for whom you have built this house. That you may do so is the
trust of all men.

§. 19.

26. And now, if it please you, let us consider the remaining accusation,
and permit me to answer it likewise. They have dared to charge me with
resisting your commands, and refusing to leave my Church. Truly I
wonder they are not weary of uttering their calumnies, I however am
not yet weary of answering them, I rather rejoice to do so; for the
more abundant my defence is, the more entirely must they be condemned.
I did not resist the commands of your Piety, God forbid; I am not a
man that would resist even the Quæstor \footnote{R. [\emph{logist\={e}i}],
auditor of accounts? vid. Demosth. de Coronâ, p. 290. ed. 1823. Arist.
Polit. vi. 8.

Return to text} of the city, much less so great a Prince. On this matter, I
need not many words, for the whole city will bear witness for me.
Nevertheless, permit me again to relate the circumstances from the
beginning; for when you hear them, I am sure you will be astonished at
the presumption of my enemies.

27. Montanus the officer of the Palace \footnote{34. vid. Cod. Theod. vi. 30.

Return to text}, came and brought me a letter, which purported to be an answer
to one from me, requesting that I might go into Italy, for the purpose
of obtaining a supply of the deficiencies which I thought existed in
the condition of our Churches. Now I desire to thank your Piety, which
condescended to assent to my request, on the supposition that I had
written to you, and made provision \footnote{35. supr. p. 100, note Y.

Return to text} for me to undertake the journey, and to accomplish it without
trouble. But here again I am astonished at those who have spoken
falsehood in your ears, that they were not afraid, seeing that lying
belongs to the Devil, and that liars are alien from Him who says, \emph{I am the Truth} [John xiv. 6.]. For I never wrote to you, nor will
my accuser be able to find any such letter; and though I ought to have
written every day, if I might thereby behold your gracious
countenance, yet it would neither have been pious to desert the
Churches, nor right to be troublesome to your Piety, especially since
you are willing to grant our requests in behalf of the Church,
although we are not present to make them. Now may it please you
to order me to read what Montanus commanded me to do. This is as
follows \footnote{36. lost, or never introduced.

Return to text}. *
* *

§. 20.

28. Now I ask again, whence have my accusers obtained this letter also?
I would learn of them who it was that put it into their hands? Do you
cause them to answer. By this you may perceive that they have forged
this, as they did also the former letter, which they published against
me, with reference to the wretched Magnentius. And being convicted in
this instance also, on what pretence next will they bring me to make
my defence? Their only concern is, to throw every thing into disorder
and confusion; and for this end I perceive they exercise their zeal.
Perhaps they think that by frequent repetition of their charges, they
will at last exasperate you against me. But you ought to turn away
from such persons, and to hate them; for such as themselves are, such
also they imagine those to be who listen to them; and they think that
their calumnies will prevail even before you. The accusation of Doeg
prevailed of old against the priests of God: but it was the
unrighteous Saul, who hearkened unto him [1 Sam. xxii. 9.]. And
Jezebel was able to injure the most religious Naboth by her false
accusations; but then it was the wicked and apostate Ahab who
hearkened unto her [1 Kings xxi.]. But the most holy David, whose
example it becomes you to follow, as all pray that you may, favours
not such men, but was wont to turn away from them and avoid them, as
raging dogs. He says, \emph{Whoso
privily slandereth his neighbour, him have I destroyed}
[Ps. ci. 5.]. For he kept the commandment which says, \emph{Thou
shalt not receive a false report}
[Ex. xxiii. 1. Sept.]. And false are the reports of these men in your
sight. You, like Solomon, have required of the Lord, (and believe
yourself to have obtained your desire,) that it would seem good unto
him to remove far from you vain and lying words.

§. 21.

29. Forasmuch then as the letter was forged by my calumniators, and
contained no order that I should come to you, I concluded that it was
not the wish of your Piety that I should come. For in that you gave me
no absolute command, but merely wrote as in answer to a letter from
me, requesting that I might be permitted to set in order the things
which seemed to be wanting, it was manifest to me (although no one
told me this) that the letter which I had received did not express the
sentiments of your Clemency. All knew, and I also stated in
writing, as Montanus is aware, that I did not refuse to come, but only
that I thought it unbecoming to take advantage of the supposition that
I had written to you to request this favour, fearing also lest my
accusers should find in this a pretence for saying that I made myself
troublesome to your Piety. Nevertheless, I made preparations, as
Montanus also knows, in order that, should you condescend to write to
me, I might immediately leave home, and readily answer your commands;
for I was not so mad as to resist such an order from you. When then in
fact your Piety did not write to me, how could I resist a command
which I never received? or how can they say, that I refused to obey,
when no orders were given me? Is not this again the mere fabrication
of enemies, pretending that which never took place? I fear that even
now, while I am engaged in this defence of myself, they may allege
against me that I am doing that which I have never obtained your
permission to do. So easily is my conduct made matter of accusation by
them, and so ready are they to vent their calumnies in despite of that
Scripture, which says, Love not to slander another, lest thou be cut
off [Prov. xx. 13. Sept.].

§. 22.

30. After a period of six and twenty months, when Montanus had gone
away, there came Diogenes the Notary \footnote{S. Notaries were the immediate attendants on magistrates, whose
judgments, \&c. they recorded and promulgated. Their office was
analogous in the Imperial Court. vid. Gothofred in Cod. Theod. vi. 10.
Ammian. Marcell. tom. 3. p. 464. ed. Erfurdt, 1808. Pancirol. Notit,
p. 143. Hofman \emph{in
voc}.
Scharf enumerates with references the civil officers, \&c. to whom
they were attached in Dissert. 1. de Notariis Ecclesiæ, p. 49.

Return to text}; but he brought me no letter, nor did we see each other, nor
did he charge me with any commands as from you. Moreover when the
General Syrianus entered Alexandria, seeing that certain reports were
spread abroad by the Arians, who declared that matters would now be as
they wished, I enquired whether he had brought any letters on the
subject of these statements of theirs. I confess that I asked for
letters containing your commands. And when he said that he had brought
none, I requested that Syrianus himself; or Maximus the Prefect of
Egypt, would write to me concerning this matter. Which request I made,
because your Grace had written to me, desiring that I would not
suffer myself to be alarmed by any one, nor attend to those who wished
to frighten me, but that I would continue to preside over the Churches
without fear. It was Palladius, the Master of the Palace, and Asterius
Duke of Armenia, who brought me this letter. Permit me to read a copy
of it. It is as follows:

§. 23.

31. A copy \footnote{37. vid. another translation of the Latin, Hist. Arian. §. 24.

Return to text} of the
letter as follows:

Constantius Victor Augustus to Athanasius.

It is not unknown to your Prudence, how constantly I prayed that success
might attend my late brother Constans in all his undertakings, and
your wisdom will easily judge how greatly I was afflicted, when I
learnt that he had been cut off by the treachery of ruffians. Now
forasmuch as certain people are endeavouring at this time to alarm
you, by setting before your eyes that lamentable tragedy, I have
thought good to address to your Reverence this present letter, to
exhort you, that, as becomes a Bishop, you would teach the people to
conform to the established \footnote{38. [\emph{kechre\={o}st\={e}men\={e}n}],
vid. [\emph{kratous\={e}i
pistei}]. infr. §. 31.

Return to text} religion, and, according to your custom, give yourself up to
prayer together with them. For this is agreeable to our wishes; and
our desire is, that you should in every season be a Bishop in our own
place.

And in another hand:—May divine Providence preserve you, beloved
Father, many years.

§. 24.

32. On the subject of this letter, my opponents conferred with the
magistrates. And was it not reasonable that I, having received it,
should demand their letters, and refuse to give heed to mere
pretences? And were they not acting in direct contradiction to the
tenor of your instructions to me, while they failed to shew me the
commands of our Piety? I therefore, seeing they produced no letters
from you, considered it improbable that a mere verbal communication
should be made to them, especially as the letter of your Grace had
charged me not to give ear to such persons. I acted rightly then, most
religious Augustus, that, as I had returned to my country under the
authority of your letters, so I should only leave it by your command;
and might not render myself liable hereafter to a charge of having
deserted the Church, but as receiving your order might have a reason
for my secession. This was demanded for me by all my people, who
went to Syrianus together with the Presbyters, and the greatest part,
to say the least, of the city with them. Maximus the Prefect of Egypt
was also there: and their request was that either he would send me a
declaration of your wishes in writing, or would forbear to disturb the
Churches, while the people themselves were sending a deputation to you
respecting the matter. When they persisted in their demand, Syrianus
at last perceived the reasonableness of it, and consented, protesting
by your life (Hilary was present and witnessed this) that he would put
an end to the disturbance, and refer the case to your Piety. The
guards of the Duke, as well as those of the Prefect of Egypt, know
that this is true; the Prytanis \footnote{39. The Mayor, Tillem. vol. 8. p. 152.

Return to text} of the city also remembers the words; so that you will
perceive that neither I, nor any one else, resisted your commands.

§. 25.

33. All demanded that the letters of your Piety should be exhibited. For
although the bare word of a King is of equal weight and authority with
his written command, especially if he who reports it, boldly affirms
in writing that it has been given him; yet when they neither openly
declared that they had received any command, nor, as they were
requested to do, gave me assurance of it in writing, but acted
altogether as by their own authority; I confess, I say it boldly, I
was suspicious of them. For there were many Arians about them, who
were their companions at table, and their advisers; and while they
attempted nothing openly, they were preparing to assail me, by
stratagem and treachery. Nor did they act at all as under the
authority of a royal command, but, as their conduct betrayed, at the
solicitation of my enemies. This made me demand more urgently that
they should produce letters from you, seeing that all their
undertakings and designs were of a suspicious nature; and because it
was unseemly that after I had entered the Church, under the authority
of so many letters from you, I should retire from it without such a
sanction.

34. When however Syrianus gave his promise, all the people assembled
together in the Churches with feelings of joyfulness and security. But
three and twenty days after, he burst into the Church with his
soldiers, while we were engaged in our usual services, as those
who entered in there witnessed; for it was a vigil, preparatory to a
communion \footnote{40. [\emph{sunaxe\={o}s}].

Return to text} on the
morrow. And such things were done that night as the Arians desired and
had beforehand denounced against us. For the General brought them with
him; and they were the instigators and advisers of the attack. This is
no incredible story of mine, most religious Augustus; for it was not
done in secret, but was noised abroad every where. When therefore I
saw the assault begun, I first exhorted the people to retire, and then
withdrew myself after them, God hiding and guiding me, as those who
were with me at the time witnessed. Since then, I have remained by
myself, though I have all confidence to answer for my conduct, in the
first place before God, and also before your Piety, for that I did not
flee and desert my people, but can point to the attack of the General
upon us, as a proof of persecution. His proceedings have caused the
greatest astonishment among all men; for either he ought not to have
made a promise, or not to have broken it after he had made it.

§. 26.

35. Now why did they form this plot against me, and treacherously lay an
ambush to take me, when it was in their power to enforce the order by
a written declaration? The command of a King is wont to give great
boldness to those entrusted with it; but their desire to act secretly,
made the suspicion stronger that they had received no command. And did
I require any thing so very absurd? Let your Majesty's candour decide
\footnote{41. [\emph{basileu
philal\={e}thes}].

Return to text}. Will not every one
say, that such a demand was reasonable for a Bishop to make? You know,
for you hare read the Scriptures, how great an offence it is for a
Bishop to desert his Church, and to neglect the flock of God. For the
absence of the Shepherd gives the wolves an opportunity to attack the
sheep. And this was what the Arians and all the other heretics
desired, that during my absence they might find an opportunity to
entrap the people into impiety. If then I had fled, what defence could
I have made before true Bishops? or rather before Him who has
committed to me His flock? He it is who judges the whole earth, the
true King of all, our Lord Jesus Christ, the Son of God.

36. Would not every one have rightly charged me with neglect of my
people? Would not your Piety have blamed me, and have justly asked, ``After
you had returned under the authority of our letters, why did you
withdraw without such authority, and desert your people?'' Would not
the people themselves at the day of judgment have reasonably imputed
to me this neglect of them, and have said, ``He that had the oversight
of us fled, and we were neglected, there being no one to put us in
mind of our duty?'' When they said this, what could I have answered?
Such a complaint was made by Ezekiel against the Pastors of old [Ez.
xxxiv. 2. \&c.]; and the blessed Apostle Paul, knowing this, has
charged every one of us, in the person of his disciple, saying, \emph{Neglect
not the gift that is in thee, which was given thee, with the laying on
of the hands of the presbytery} [1 Tim. iv. 14.]. Fearing this, I wished not to
flee, but to receive your commands, if indeed such was the will of
your Piety. But I never obtained what I so reasonably requested, and
now I am falsely accused before you; for I resisted no commands of
your Piety; nor will I now attempt to return to Alexandria, until your
Grace shall desire it. This I say beforehand, lest the slanderers
should again make this a pretence for accusing me.

§. 27.

37. Considering these things, I did not give sentence against myself \footnote{42. [vol. 8. p. 6, note o.

Return to text}, but hastened to come to your Piety, with this my defence,
knowing your goodness, and remembering your faithful promises, and
being confident that, as it is written in the Proverbs of Scripture, \emph{Just
speeches are acceptable to a gracious king} [Prov. xvi. 13.] \footnote{43. quoted otherwise, supr. p. 164.

Return to text}. But when I had already entered upon my journey, and had past
through the desert, a report suddenly reached me \footnote{T. In this chapter he breaks off his Oratorical form, and ends his
Apology much more in the form of a letter. vid. however [\emph{t\={o}n
log\={o}n kairon}], infr. §. 34, 35 init. [\emph{prosth\={o}n\={e}s\={o}}],
§. 35. The events which he here records changed his feelings towards
Constantius, whom henceforth he accounted as a persecutor, worse than
heathen, because an apostate. vid. Lib. F. vol. 8. p. 90, note p.

Return to text}, which at first I thought to be incredible, but which
afterwards proved to be true. It was rumoured every where that
Liberius Bishop of Rome, the great Hosius of Spain, Paulinus of Gaul,
Dionysius and Eusebius of Italy, Lucifer of Sardinia \footnote{44. vid. infr. p. 191.

Return to text}, and certain other Bishops, with theirs Presbyters and
Deacons, had been banished because they refused to subscribe to my
condemnation. These had been banished; and Vincentius \footnote{45. supr. p. 157, note C.

Return to text} of Capua, Fortunatian \footnote{45. supr. p. 157, note C.

Return to text} of
Aquileia, Heremius of Thessalonica, and all the Bishops of the West,
were treated with no ordinary vigour, nay were suffering extreme
violence and grievous injuries, until they could be induced to promise
that they would not communicate with me.

38. While I was astonished and perplexed at these tidings, behold
another report \footnote{46. vid. Hist. Ar. §§. 31, 32, 54, 70, \&c.

Return to text}
overtook me, respecting them of Egypt and Libya, that nearly ninety
Bishops had been under persecution, and that their Churches were given
up to the professors of Arianism; that sixteen had been banished, and
of the rest, some had fled, and others were constrained to dissemble.
For the persecution was said to be so violent in those parts, that at
Alexandria, while the brethren were praying during Easter and on the
Lord's day in a desert place near the cemetery, the General came upon
them with a force of soldiery, more than three thousand in number,
with arms, drawn swords, and spears; whereupon outrages, such as might
be expected to follow so unprovoked an attack, were committed against
women and children, who were doing nothing more than praying to God.
It would perhaps be unseasonable to give an account of them now, lest
the mere mention of such enormities should move us all to tears. But
such was their cruelty, that virgins were stripped, and even the
bodies of those who died from the blows they received were not
immediately given up for burial, but were cast out to the dogs, until
their relatives, with great risk to themselves, came secretly and
stole them away, and much effort was necessary, that no one might know
it.

§. 28.

39. The rest of their proceedings will perhaps be thought incredible,
and will fill all men with astonishment, by reason of their extreme
wickedness. It is necessary however to speak of them, in order that
your Christian zeal and piety may perceive that their slanders and
calumnies against us are framed for no other end, than that they may
drive us out of the Churches, and introduce their own impiety in our
place. For when the lawful Bishops, men of advanced age, had some of
them been banished, and others forced to fly, heathens and
catechumens, those who hold the first places in the senate, and
men who are notorious for their wealth, were straightway commissioned
by the Arians to preach the holy faith instead of Christians \footnote{47. Hist. Ar. §. 73. supr. p. 135, r. 1.

Return to text}. And enquiry was no longer made, as the Apostle enjoined, \emph{if any be blameless}
[1 Tim. iii. 2.]: but according to the practice of the impious
Jeroboam, he who could give most money, was named Bishop; and it made
no difference to them, even if the man happened to be a heathen, so
long as he furnished them with money. Those who had been Bishops from
the time of Alexander, monks and ascetics, were banished: and men
practised only in calumny corrupted, as far as in them lay, the
Apostolic rule, and polluted the Churches. Truly their false
accusations against us have gained them much, that they should be able
to commit iniquity, and to do such things as these in your time \footnote{48. [\emph{kairois}], de Syn. fin. (tr. p. 159.)

Return to text}; so that the words of Scripture may be applied to them, Woe
unto those through whom My name is blasphemed among the Gentiles [vid.
2 Sam. xii. 14. \&c.].

§. 29.

40. These were the rumours that were noised abroad; and although every
thing was thus turned upside down, I still did not relinquish my
earnest desire of coming to your Piety, but was again setting forward
on my journey. And I did so the more eagerly, being confident that
these proceedings were contrary to your wishes, and that if your Grace
should be informed of what was done, you would prevent it for the time
to come. For I could not think that a righteous king could wish
Bishops to be banished, and virgins to be stripped, or the Churches to
be in any way disturbed. While I thus reasoned and hastened on my
journey, behold a third report reached me, to the effect that letters
had been written to the Princes of Auxumis, desiring that Frumentius \footnote{U. Athan. had consecrated Frumentius for the Ethiopian mission, who had
been already successful in introducing Christianity into the heathen
court of Auxumis, where he had held the place first of Minister, then
of Regent. The two Princes to whom Constantius writes in the letter
which is presently to follow were the King's sons, whom Frumentius had
first served.

Return to text}, Bishop of Auxumis, should be brought from thence, and that
search should be made for me even as far as the country of the
Barbarians, that I might be handed over to the Commentaries \footnote{X. That is, the prison. ``The official books;'' Montfaucon (apparently) in
Onomast. vid. Gothofr. Cod. Theod. ix. 3. l. 5. However, in xi. 30. p.
243. he says, Malim pro ipsâ custodiâ accipere. And so Du Cange in
voc. and this meaning is here followed. vid. supr. p. 25. where
commentarius is translated ``jailor.'' Apol. contr. Arian. §. 8.

Return to text} (as they are called) of the Prefects, and that all the
laity and clergy should be compelled to communicate with the Arian
heresy, and that such as would not comply with this order should be
put to death. To shew that these were not merely idle rumours, but
that they were confirmed by facts, since your Grace has given me
leave, I produce the letter. My enemies, who threatened every one with
death, frequently caused it to be read.

§. 30.

41. A copy of the letter.

Victor \footnote{49. pp. 79, 96, 119, \&c.

Return to text} Constantius
Maximus Augustus to the Alexandrians.

Your city, preserving its native spirit and temper, and remembering the
virtue of its founders, has habitually shewn itself obedient unto us,
as it does at this day; and we on our part should consider ourselves
greatly wanting in our duty, did not our good will eclipse even that
of Alexander himself. For as it belongs to a temperate mind, to behave
itself orderly in all respects, so it is the part of royalty, on
account of virtue, permit me to say, such as yours, to embrace you
above all others; you, who rose up as the first teachers of wisdom;
who were the first to acknowledge the God, who is \footnote{50. [\emph{ton
onta}], infr. p. 182, note Y.

Return to text}; who moreover have chosen for yourselves the most consummate
masters; and have cordially acquiesced in our opinion, justly
abominating that impostor and cheat, and dutifully uniting yourselves
to those venerable men who are beyond all admiration. And yet, who is
ignorant, even among those who live in the end of the earth, what
violent party spirit was displayed in the late proceedings? with which
we know not any thing that has ever happened, worthy to be compared.
The majority of the citizens had their eyes blinded, and a man who had
come forth from the lowest dens of infamy obtained authority among
them, entrapping into falsehood, as under cover of darkness, those who
were desirous to know the truth;—one who never provided for them any
fruitful and edifying discourse, but corrupted their minds with
unprofitable subtleties. His flatterers shouted and applauded him;
they were astonished at his powers, and they still probably murmur
secretly; while the majority of the more simple sort took their cue
from them. And thus all went with the stream, as if a flood had
broken in, while every thing was entirely neglected. One of the
multitude was in power;—how can I describe him more truly, than by
saying, that he was superior in nothing to the meanest of the people,
and that the only kindness which he shewed to the city was, that he
did not thrust her citizens down into the pit. This noble-minded and
illustrious person did not wait for judgment to proceed against him,
but sentenced himself to banishment as he deserved. So that now it is
for the interest of the Barbarians to remove him out of the way, lest
he lead some of them into impiety, for he will make his complaint,
like distressed characters in a play, to those who shall first fall in
with him.

42. To him however we will now bid a long farewell. For yourselves there
are few with whom I can compare you: I am bound rather to honour you
separately above all others, for the great virtue and wisdom which
your actions, that are celebrated almost through the whole world,
proclaim you to possess. Go on in this sober course. I would gladly
have repeated to me a description of your conduct in such terms of
praise as it deserves; O ye who have eclipsed your predecessors in the
race of glory, and will be a noble example both to those who are now
alive, and to all who shall come after, and alone have chosen for
yourselves the most excellent guide you could have for your conduct,
both in word and deed, and hesitated not a moment, but manfully
transferred your affections, and gave yourselves up to the other side,
leaving those grovelling \footnote{51. [\emph{t\={o}n
chamai}],
vid. contr. Euseb. Hist. vii. 27.

Return to text}
and earthly teachers, and stretching forth towards heavenly things,
under the guidance of the most venerable George \footnote{52. of Capadocia, vol. 8. p. 134, note f.

Return to text}, than whom no man is more perfectly instructed therein. Under
him you will continue to have a good hope respecting the future life,
and will pass your time in this present world, in rest and quietness.
Would that all the citizens together would lay hold on his words, as a
sacred anchor, so that we might need neither knife nor cautery, for
those whose souls are diseased!

43. Such persons we most earnestly advise to renounce their zeal in
favour of Athanasius, and not even to remember the foolish things
which he spoke so plentifully among them. Otherwise they will bring
themselves before they are aware into extreme peril, from which
we know not any one who will be skilful enough to deliver such
factious persons. For while that pestilent \footnote{53. [\emph{olethron}].

Return to text} fellow Athanasius is driven from place to place, being
convicted of the basest crimes, for which he would only suffer the
punishment he deserves, if one were to kill him ten times over; it
would be inconsistent in us to suffer those flatterers and juggling
ministers of his to exult against us; men of such a character as it is
a shame even to speak of, respecting whom orders have long ago been
given to the magistrates, that they should be put to death. But even
now perhaps they shall not die, if they desist from their former
offences, and repent at last. For that villain Athanasius led them on,
and corrupted the whole state, and laid his impious and polluted hands
upon the most holy things.

§. 31.

44. The following is the letter which was written to the Princes of
Auxumis respecting Frumentius Bishop of that place.

45. Victor Constantius Maximus Augustus, to Æzanes and Sazanes.

It is altogether a matter of the greatest care and concern to us, to
extend the knowledge of the supreme God \footnote{Y. [\emph{h\={e}
tou kreittonos gn\={o}sis}],
vid. [\emph{tou
kreittona}], infr. And so in Arius's Thalia, the Eternal Father, in contrast to
the Son, is called [\emph{ho
kreitt\={o}n}],
de Synod. §. 15. So again, [\emph{theon
ton onta sunientas}],
supr. §. 30. and [\emph{sunet\={o}n theou}] in the Thalia, Orat. i. 5. Again, [\emph{sophias
ex\={e}g\={e}tas}], supr. §. 30. and [\emph{t\={o}n sophias metachont\={o}n}],
in the Thalia, ibid. And [\emph{t\={o}n exeget\={o}n}], supr. §. 30. and [\emph{tout\={o}n
kat}\emph{'}\emph{ichnos \={e}lthon}], in the Thalia.

Return to text}; and I think that the whole race of mankind claims from us
equal regard in this respect, in order that they may pass their lives
in hope, being brought to a proper knowledge of God, and having no
differences with each other in their enquiries concerning justice and
truth. Wherefore considering that you are deserving of the same
provident care as the Romans, and desiring to shew equal regard for
your welfare, we command that the same doctrine be professed in your
Churches as in theirs. Send therefore speedily into Egypt the Bishop
Frumentius, to the most venerab1e Bishop George, and the rest who are
there, who have especial authority to appoint to these offices, and to
decide questions concerning them. For of course you know and remember
(unless you alone pretend to be ignorant of that which all men
are well aware of) that this Frumentius was advanced to his present
rank by Athanasius, a man who is guilty of ten thousand crimes; for he
has not been able fairly to clear himself of any of the charges
brought against him, but was at once deprived of his see, and now
wanders about destitute of any fixed abode, and passes from one
country to another, as if by this means he could escape his own
wickedness.

46. Now if Frumentius shall readily obey our commands, and shall submit
to an enquiry into all the circumstances of his appointment, he will
shew plainly to all men, that he is in no respect opposed to the laws
of the Church and the established \footnote{54. [\emph{kpatous\={e}i}],
supr. p. 174. r. 1.

Return to text} faith. And being brought to trial, when he shall have given
proof of his general good conduct, and submitted an account of his
life to those who are to judge of these things, he shall receive his
appointment from them, if it shall indeed appear that he has any right
to be a Bishop. But if he shall delay and avoid the trial, it will
surely be very evident, that he has been induced by the persuasions of
the wicked Athanasius, thus impiously to act against divine authority,
choosing to follow the course of him whose wickedness has been made
manifest. And our fear is lest he should pass over into Auxumis and
corrupt your people, by setting before them accursed and impious
doctrines, and not only unsettle and disturb the Churches, and
blaspheme the supreme \footnote{55. [\emph{kreittona}].

Return to text}
God, but also thereby cause utter overthrow and destruction to the
several nations whom he visits. But I am sure that Frumentius will
return home, perfectly acquainted with all matters that concern the
Church, having derived much instruction, which will be of great and
general utility, from the conversation of the most venerable George,
and such other of the Bishops, as are excellently qualified to
communicate such knowledge. May God continually preserve you, most
honoured brethren.

§. 32.

47. Hearing, nay almost seeing, these things, through the mournful
representations of the messengers, I confess I turned back again into
the desert, justly concluding, as your Piety will perceive, that if I
was sought after, that I might be sent as soon as I was discovered to
the Prefects \footnote{56. p. 179, 180 init.

Return to text}, I
should be prevented from ever coming to your Grace; and that if those who would not subscribe against me, suffered so severely as they
did, and the laity who refused to communicate with the Arians were
ordered for death, there was no doubt at all but that ten thousand new
modes of destruction would be devised by the calumniators against me;
and that after my death, they would employ against whomsoever they
wished to injure, whatever means they chose, venting their lies
against us the more boldly, for that then there would no longer be any
one left who could expose them. I fled, not because I feared your
Piety, (for I know your long-suffering and goodness,) but because from
what had taken place, I perceived the spirit of my enemies, and
considered that they would make use of all possible means to
accomplish my destruction, from fear that they would be brought to
answer for what they had done contrary to the intentions of your
Excellency. For observe, your Grace commanded that the Bishops should
be expelled only out of the cities and the province. But these worthy
persons presumed to exceed your commands, and banished aged men and
Bishops venerable for their years into desert and unfrequented and
frightful places, beyond the boundaries of three provinces \footnote{Z. Egypt was divided into three Provinces till Hadrian's time, Egypt,
Libya, and Pentapolis; Hadrian made them four; Epiphanius speaks of
them as seven. Hær. 68. 1. By the time of Arcadius they had become
eight. vid. Orlendini Orbis Sacer et Prof. vol. i. p. 118. The
Province specially spoken of seems to be Egypt, which Augustus kept in
his own hands. vid. supr. p. 5, note D. p. 116, r. 1.

Return to text

}. Some of them were sent off from Libya to the great Oasis;
others from the Thebais to Ammoniaca in Libya.

48. Neither was it from fear of death that I fled; let none of them
condemn me as guilty of cowardice; but because it is the injunction of
our Saviour \footnote{57. vid. Apol. de Fug. init. p. 188.

Return to text} that we
should flee when we are persecuted, and hide ourselves when we are
sought after, and not expose ourselves to certain dangers, nor by
appearing before our persecutors inflame still more their rage against
us. For to give one's self up to one's enemies to be murdered, is the
same thing as to murder one's self; but to flee, as our Saviour has
enjoined, is to know our time, and to manifest a real concern for our
persecutors, lest if they proceed to the shedding of blood, they
become guilty of the transgression of the law, \emph{Thou
shalt not kill}
[Exod. xx. 13.]. And yet these men by their calumnies against
me, earnestly wish that I should suffer death.

49. What they have again lately done proves that this is their desire
and murderous intention. You will be astonished, I am sure, most
religious Augustus, when you hear it; it is indeed an outrage worthy
of amazement. What it is, I pray you briefly to hear. §. 33. The Son
of God, our Lord and Saviour Jesus Christ, having become man for our
sakes, and having destroyed death, and delivered our race from the
bondage of corruption, in addition to all His other benefits bestowed
this also upon us, that we should possess upon earth, in the state of
virginity, a picture of the holiness of Angels. Accordingly such as
have attained this virtue, the Catholic Church has been accustomed to
call the brides \footnote{58. [\emph{n}\emph{y}\emph{mphas}].

Return to text} of
Christ. And the heathen who see them express their admiration of them
as the temples of the Word. For indeed this holy and heavenly
profession is no where established, but only among us Christians, and
it is a very strong argument that with us is to be found the genuine
and true religion. Your most religious father Constantine Augustus, of
blessed memory \footnote{59. [\emph{t\={e}s
makarias mn\={e}m\={e}s}].
supr. pp. 159, 162.

Return to text},
honoured the Virgins above all other orders, and your Piety in several
letters has given them the titles of the honourable and holy women.
But now these worthy Arians who have slandered me, and by whom
conspiracies have been formed against most of the Bishops, having
obtained the consent and cooperation of the magistrates, first
stripped them, and then caused them to be suspended upon what are
called the Hermetaries \footnote{60. a rack, or horse, Tillemont. v. Athan. p. 169.

Return to text},
and scourged them on the ribs so severely three several times, that
not even real malefactors have ever suffered the like. Pilate, to
gratify the Jews of old, pierced one of our Saviour's sides with a
spear. These men have exceeded the madness of Pilate, for they have
scourged not one but both His sides; for the limbs of the Virgins are
in an especial manner the Saviour's own.

50. All men shudder at hearing the bare recital of deeds like these.
These men alone, not only did not fear to strip and to scourge those
undefiled limbs, which the Virgins had dedicated solely to our Saviour
Christ; but, what is worse than all, when they were reproached by
every one for such extreme cruelty, instead of manifesting any shame,
they pretended that it was commanded by your Piety. So utterly
presumptuous are they and full of wicked thoughts and purposes. Such a
deed as this was never heard of in past persecutions \footnote{61. vid. Hist. Ar. §. 40. §. 64.

Return to text}: or supposing that it ever occurred before, yet surely it was
not befitting either that Virginity should suffer such outrage and
dishonour, in the time of your Majesty a Christian Prince, or that
these men should impute to your Piety their own cruelty. Such
wickedness belongs only to heretics, to blaspheme the Son of God, and
to do violence to His holy Virgins.

§. 34.

51. Now when such enormities as these were again perpetrated by the
Arians, I surely was not wrong in complying with the direction of Holy
Scripture, which says, \emph{Hide
thyself for a little moment, until the indignation be over-past} [Is. xxvi. 29.]. This was another reason for my
withdrawing myself, most religious Augustus; and I refused not, either
to depart into the desert, or, if need were, to be let down from a
wall in a basket. I endured every thing, I even dwelt among wild
beasts, that your favour might return to me, waiting for an
opportunity to offer to you this my defence, confident as I am that
they will be condemned, and your goodness manifested unto me.

52. O, Augustus, blessed and beloved of God, what would you have had me
to do? to come to you while my calumniators were inflamed with rage
against me, and were seeking to kill me; or, as it is written, to hide
myself a little, that in the mean time they might be condemned as
heretics, and your goodness might be manifested unto me? or would you
have had me, Sire, to appear before your magistrates, in order that
though you had written merely in the way of threatening, they not
understanding your intention, but being exasperated against me by the
Arians, might kill me on the authority of your letters, and on that
ground ascribed the murder to you? It would neither have been becoming
in me to surrender, and give myself up that my blood might be shed,
nor in you, as a Christian King, to have the murder of Christians, and
those too Bishops, imputed unto you.

§. 35.

53. It was therefore better for me to hide myself, and to wait for this
opportunity. Yes, I am sure that from your knowledge of the sacred
Scriptures you will assent and approve of my conduct in this
respect. For you will perceive that, now those who exasperated you
against us have been silenced, your righteous clemency is apparent,
and it is proved to all men that you never persecuted the Christians
at all, but that it was they who made the Churches desolate, that they
might sow the seeds \footnote{62. vol. 8. p. 5. note k.

Return to text}
of their own impiety every where; on account of which I also, had I
not fled, should long ago have suffered from their treachery. For it
is very evident that they who scrupled not to utter such calumnies
against me, before the great Augustus, and who so violently assailed
Bishops and Virgins, sought also to compass my death. But thanks be to
the Lord who has given you this kingdom. All men are confirmed in
their opinion of your goodness, and of their wickedness, from which I
fled at the first, that I might now make this appeal unto you, and
that you might find some one towards whom you may shew kindness. I
beseech you therefore, forasmuch as it is written, \emph{A
soft answer turneth away wrath},
and, \emph{righteous
thoughts are acceptable unto the King} [Prov. xv. 1; xvi. 13. vid. p. 177]; receive
this my defence, and restore all the Bishops and the rest of the
Clergy to their countries and their Churches; so that the wickedness
of my accusers may be made manifest, and that you, both now and in the
day of judgment, may have boldness to say to our Lord and Saviour
Jesus Christ, the King of all, ``\emph{None
of Thine have I lost} [John xviii. 9.], but these are they who designed
the ruin of all, while I was grieved for those who perished, and for
the Virgins who were scourged, and for all other things that were
committed against the Christians; and I brought back them that were
banished, and restored them to their own Churches.''

\xchapter{V. Apology of our Holy Father Athanasius, Archbishop of Alexandria, in Vindication of His Flight, when He was Persecuted by Duke Syrianus}

———————

[This Apology seems to have been written A.D.
357 or 358. The circumstances which led to it are mentioned in the
opening sentences. From what he says to Constantine in the foregoing
work, p. 177, it might almost be said that, in addition to the
considerations insisted on in the following argument, he considered
that the command of the Emperor would in itself have been a sufficient
reason for his leaving his Church; and it was because he had not
received it, that he had not left it before. Now the violence of
Syrianus, acknowledged as it was by Constantius, was of the nature of
a command. The real reason however was, that, if he had been cut off,
there was no one to take his place. vid. supr. p. 184.]

———————

§. 1.

1. I HEAR that Leontius \footnote{Margin Notes

1. infr. §. 26.

Return to text}, now in the see of Antioch, and Narcissus of the city of
Nero, and George \footnote{vol. 8. p. 99, supr. p. 74, \&c.

Return to text}, now
of Laodicea, and the Arians who are leagued with them, are spreading
abroad many slanderous reports concerning me, charging me with
cowardice, because forsooth, when myself was sought by them, I did not
surrender myself into their hands. Now as to their imputations and
calumnies, although there are many things that I could write, which
they are unable to deny, and which all who have heard of their
proceedings know to be true, yet I shall not be prevailed upon to make
any reply to them, except only to remind them of the words of our
Lord, and of the declaration of the Apostle, that \emph{a lie is of the
Devil}, and that, \emph{revilers shall not inherit the kingdom of God}
[vid. 1 John ii. 21. 1 Cor. vi. 10.]. For it is sufficient thereby to
prove, that neither their thoughts nor their words are according to
the Gospel, but that after their own pleasure, whatsoever themselves
desire, that they think to be good.

§. 2.

2. But forasmuch as they pretend to charge me
with cowardice, it is necessary that I should write somewhat
concerning this, whereby it shall be proved that they are men of
wicked minds, who have not read the sacred Scriptures: or if they have
read them, that they do not believe the divine inspiration of the
oracles they contain. For had they believed this, they would never
have dared to act contrary to them, nor have imitated the malice of
the Jews who slew the Lord. For God having given them a commandment, \emph{Honour
thy father and thy mother}, and, \emph{He that curseth father or
mother, let him die the death} [Mat xv. 4.]; that people
established a contrary law, changing the honour into dishonour, and
alienating to other uses the money which was due from the children to
their parents. And though they had read what David did, they acted in
contradiction to his example, and accused the guiltless for plucking
the ears of corn, and rubbing them in their hands on the Sabbath day.
Not that they cared either for the laws, or for the Sabbath, for they
were guilty of greater transgressions of the law on that day: but
being wicked-minded, they grudged the disciples the way of salvation,
and desired that their own private notions should have the sole
pre-eminence. They however have received the reward of their iniquity,
having ceased to be an holy nation, and being counted henceforth as
the rulers of Sodom, and as the people of Gomorrah.

3. And these men likewise, not less than they,
seem to me to have received their punishment already in the ignorance
with which their own folly possesses them. For they understand not
what they say, but think that they know things of which they are
ignorant; while the only knowledge that is in them is to do evil, and
to frame devices more and more wicked day by day. Thus they reproach
me with my present flight, not for the sake of my character, as
wishing me to shew my manliness by coming forward; (how is it possible
that such a wish can be entertained by enemies in behalf of those who
run not with them in the same career of madness?) but being full of
malice, they pretend this, and whisper \footnote{[\emph{peribombein}], vol. 8. p. 22, note y. Greg. Naz. Orat. 27. n.
2.

Return to text} up and down that such is the case, thinking, foolish as indeed
they are, that through fear of their revilings, I shall yet be induced
to give myself up to them. For this is what they desire: to accomplish
this they have recourse to all kinds of schemes: they pretend
themselves to be friends, while they search after me as enemies,
to the end that they may glut themselves with my blood, and put me
also out of the way, because I have always opposed and do still oppose
their impiety, and confute and brand their heresy.

§. 3.

4. For whom have they ever persecuted and taken,
that they have not insulted and injured as they pleased? Whom have
they ever sought after and found, that they have not handled in such a
manner, that either he has died a miserable death, or has been
illtreated in every member \footnote{[\emph{pantachothen}].

Return to text}?
Whatever the magistrates appear to do, it is their work; and the other
are merely the tools of their will and wickedness. In consequence,
where is there a place that has not some memorial of their wickedness?
Who has ever opposed them, without their conspiring against him,
inventing pretexts for his ruin after the manner of Jezebel? Where is
there a Church that is not at this moment lamenting the success of
their plots against her Bishops? Antioch is mourning for the orthodox
Confessor Eustathius \footnote{vid. Hist. Arian. §. 4. also Theodoret Hist.
i. 20. Eustathius was one of the original opponents of Arianism. S.
Alexander wrote to him (then Bishop of Berrhœa) against Arius, as
well as to Philogonius of Antioch and Alexander of Constantinople. He
was deposed by the Arians A.D.
331, on the pretence of Sabellianism and perhaps of incontinency.
Montfaucon, however, doubts whether the latter was ever made a charge,
though Theodoret mentions it. V. Athan. p. 14. Another reason is given
Hist. Arian. loc. cit. The orthodox succession was continued, though
dispossessed, and gave occasion to the schism, after the overthrow of
Arianism, which was not terminated till the time of S. Chrysostom. The
name of Euphration occurs de Syn. 17. (tr. vol.8. p. 99.) as the
Bishop to whom Eusebius of Cæsarea wrote an heretical letter. Balaneæ
is on the Syrian coast. Paltus also and Antaradus are in Syria, and
these persecutions took place about A.D. 340; that of Eutropius, and of Lucius his successor, about
332, shortly after the proceedings against Eustathius. Cyrus too was
banished under pretence of Sabellianism about 340. Asclepas has been
mentioned supr. p. 69. note E. For Theodulus and Olympius vid. Hist.
Arian. §. 19. and supr. p. 71. note G.

Return to text};
Balaneæ for the most admirable Euphration \footnote{Hist. Arian. 5.

Return to text}; Paltus and Antaradus for Cymatius \footnote{Hist. Arian. 5.

Return to text}
and Carterius; Adrianople for that lover of Christ, Eutropius, and
after him for Lucius, who was often loaded with chains by their means,
and so perished; Ancyra mourns for Marcellus, Berrhœa \footnote{Berœa, Hist. Ar. 5.

Return to text} for Cyrus \footnote{Hist. Arian. 5.

Return to text}, Gaza for Asclepas.

5. Of all these, after inflicting many outrages,
they by their intrigues procured the banishment; but for Theodulus and
Olympius, Bishops of Thrace, and for me and my Presbyters, they caused
diligent search to be made, to the intent that if we were discovered
we should suffer capital punishment: and probably we should have
so perished, had we not fled at that very time contrary to their
intentions. For letters to that effect were delivered to the Proconsul
Donatus against Olympius and his friends, and to Philagrius respecting
me. And having raised a persecution against Paul, Bishop of
Constantinople, as soon as they found him, they caused him to be
openly strangled \footnote{infr. Hist. Arian. §. 4.

Return to text} at a
place called Cucusus in Cappadocia, employing as their executioner for
the purpose Philip, who was Prefect. He was a patron of their heresy,
and the tool of their wicked designs.

§. 4.

6. Are they then satisfied with all this, and
content to be quiet for the future? By no means; they have not given
over yet, but like the horseleach \footnote{Hist. Arian. §. 65.

Return to text} in the Proverbs, they revel more and more in their wickedness,
and fix themselves upon the larger dioceses. Who can adequately
describe the enormities they have already perpetrated? who is able to
recount all the deeds that they have done? Even very lately, while the
Churches were at peace, and the people worshipping in their
congregations, Liberius Bishop of Rome, Paulinus \footnote{of Treves.

Return to text} Metropolitan of Gaul, Dionysius \footnote{of Milan.

Return to text} Metropolitan of Italy, Lucifer \footnote{of Cagliari.

Return to text} Metropolitan of the Sardinian islands, and Eusebius \footnote{of Vercellæ.

Return to text} of Italy, all of them excellent Bishops and preachers of the
truth, were seized and banished, on no pretence whatever, except that
they would not unite themselves to the Arian heresy, nor subscribe to
the accusations and calumnies which they had invented against me.

§. 5.

7. Of the great Hosius, who answers to his name,
that confessor of an happy old age \footnote{[\emph{eug\={e}rotatou}], vid. supr. p. 70.

Return to text}, it is superfluous for me to speak, for I suppose it is known
unto all men that they caused him also to be banished; for he is not
an obscure person, but of all men the most illustrious, and more than
this. When was there a Council held, in which he did not take the
lead, and convince every one by his orthodoxy? Where is there a Church
that does not possess some glorious monuments of his patronage? Who
has ever come to him in sorrow, and has not gone away rejoicing? What
needy person ever asked his aid, and did not obtain what he desired?
And yet even on this man they made their assault, because knowing the
calumnies which they invent in behalf of their iniquity, he would not
subscribe to their designs against me. And if afterwards, upon
the repeated blows that were inflicted upon him above measure, and the
conspiracies that were formed against his kinsfolk, he yielded to them
for a time, as being old and infirm in body, yet at least their
wickedness was shewn even in this circumstance; so zealously did they
endeavour by all means to prove that they were not truly Christians \footnote{infr. §. 27 init.

Return to text}.

§. 6.

8. After this they again fastened themselves upon
Alexandria, seeking anew to put me to death: and their proceedings
were now worse than before. For on a sudden the Church was surrounded
by soldiers, and deeds of war took the place of prayers. Then George \footnote{vol. 8. p. 134, note f. supr. p. 181.

Return to text} of Cappadocia who was sent by them, having arrived during the
season of Lent \footnote{supr. p. 7, note H.

Return to text},
brought an increase of evils which they had taught him. For after
Easter week, Virgins were thrown into prison; Bishops were led away in
chains by soldiers; the houses of orphans and widows were plundered,
and their bread taken away; attacks were made upon houses, and
Christians thrust forth in the night, and their dwellings sealed up:
the brothers of clergymen were in danger of their lives on account of
their relations.

9. These outrages were sufficiently dreadful, but
more dreadful than these followed. For on the week that succeeded the
Holy Pentecost, when the people after their fast had gone out to the
cemetery to pray, because that all refused communion with George, that
abandoned person, understanding this to be the case, stirred up
against them the commander Sebastian, a Manichee; who straightway with
a multitude of soldiers with arms, drawn swords, bows, and spears,
proceeded to attack the people, though it was the Lord's day: and
finding a few praying, (for the greater part had already retired on
account of the lateness of the hour,) he committed such outrages as
became a disciple of these men. Having lighted a pile, he placed
certain virgins near the fire, and endeavoured to force them to say
that they were of the Arian faith: and where he saw that they were
getting the mastery, and cared not for the fire, he immediately
stripped them naked, and wounded their faces in such a manner, that
for some time they could hardly be recognised.

§. 7.

10. And having seized upon forty men, he beat
them after a new fashion. Cutting some fresh twigs of the palm
tree with the thorns upon them \footnote{Hist. Arian. §. 72.

Return to text}, he scourged them on the back so severely, that some of them
were for a long time under medical treatment on account of the thorns
which had entered their flesh, and others unable to bear up under
their sufferings died. All those whom they had taken, both the men and
the virgins, they sent away together into banishment to the great
Oasis. And the bodies of those who had perished they would not at
first suffer to be given up to their friends, but concealed them in
any way they pleased, and cast them out without burial \footnote{ibid. §. 72 fin. supr. p. 178.

Return to text}, in order that they might not appear to have any knowledge of
these cruel proceedings. But herein their deluded minds greatly misled
them. For the relatives of the dead, both rejoicing at the confession,
and grieving for the bodies of their friends, published abroad so much
the more this proof of their impiety and cruelty. Moreover they
immediately banished out of Egypt and Libya the following Bishops \footnote{16, ibid. vid. Hist. Ar. §. 72.

Return to text}, Ammonius, Muïus, Gaïus, Philo \footnote{Hieron. V. Hilar. §. 30. perhaps a Luciferian, A.D.
362.

Return to text}, Hermes, Plenius, Psenosiris, Nilammon, Agathus, Anagamphus,
Marcus, Ammonius, another Marcus, Dracontius \footnote{ad Dracont.

Return to text}, Adelphius \footnote{ad Adelph.

Return to text},
Athenodorus, and the Presbyters, Hierax \footnote{ad Dracont. 10.

Return to text}, and Dioscorus; whom they drove forth under such cruel
treatment, that some of them died on the way, and others in the place
of their banishment. They caused also more than thirty Bishops to take
to flight; for their desire was, after the example of Ahab, if it were
possible, utterly to root out the truth. Such are the enormities of
which these impious men have been guilty.

§. 8.

11. But although they have done all this, yet
they are not ashamed of the evils they have already contrived against
me, but proceed now to accuse me, because I have been able to escape
their murderous hands. Nay, they bitterly bewail themselves, that they
have not effectually put me out of the way; and so they pretend to
reproach me with cowardice, not perceiving that by thus murmuring
against me, they rather turn the blame upon themselves. For if it be a
bad thing to flee, it is much worse to persecute; for the one party
hides himself to escape death, the other persecutes with a desire to
kill; and it is written in the Scriptures that we ought to flee, but
he that seeks to destroy transgresses the law, nay, and is himself the
occasion of the other's flight. If then they reproach me with
my flight, let them be more ashamed of their own persecution \footnote{vid. supr. p. 18. §. 4.

Return to text}. Let them cease to compass my destruction, and I shall without
delay cease to flee.

12. But they, instead of giving over their
wickedness, are employing every means to obtain possession of my
person, not perceiving that the flight of those who are persecuted is
a strong argument against them that persecute. For no man flees from
the gentle and the humane, but from the cruel and the evil-minded. \emph{Every
one that was in distress, and every one that was in debt} [1 Sam.
xxii. 2.], fled from Saul, and took refuge with David. But this is the
reason why these men desire to cut off those who are in concealment,
that there may be no evidence forthcoming of their wickedness. But
herein their minds seem to be blinded with their usual error. For the
more the flight of their enemies becomes known, so much the more
notorious will be the destruction or the banishment which their
treachery has brought upon them \footnote{Hist. Arian. §. 34. 35.

Return to text}; so that whether they kill them outright, their death will be
the more loudly noised abroad against them, or whether they drive them
into banishment, they will but be sending forth every where monuments
of their own iniquity.

§. 9.

13. Now if they had been of sound mind, they
would have seen that they were in this strait, and that they were
defeated by their own arguments. But since they have lost all
judgment, they are still led on to persecute, and seek to destroy, and
yet perceive not their own impiety. It may be they even venture to
accuse Providence itself; (for nothing is beyond the reach of their
presumption,) that it does not deliver up to them those whom they
desire, certain as it is, according to the saving of our Saviour, that
not even a sparrow can fall into the net \footnote{p. 199, r. 1.

Return to text} with out our Father which is in heaven. But when these bad
spirits obtain possession of any one, they immediately forget not only
all other, but even themselves; and raising their brow in very
haughtiness, they neither acknowledge times and seasons, nor respect
human nature in those whom they injure. Like the tyrant of Babylon \footnote{pp. 9, 195.

Return to text}, they attack more furiously; they shew pity to none, but
mercilessly \emph{upon the ancient}, as it is written, \emph{they very
heavily lay the yoke}, and \emph{they add to the grief of them that
are wounded} [Is. xlvii. 6. Ps. lxix. 26.].

14. Had they not acted in this manner; had they
not driven into banishment those who spoke in my defence against their
calumnies, their representations might have appeared to some persons
sufficiently plausible. But since they have conspired against so many
other Bishops of high character, and have spared neither the great
confessor Hosius, nor the Bishop of Rome, nor so many others from
Spain and Gaul, and Egypt, and Libya, and the other countries, but
have committed such cruel outrages against all who have in any way
opposed them in my behalf; is it not plain that their designs have
been directed rather against me than against any other, and that their
desire is miserably to destroy me as they have done others? To
accomplish this they vigilantly watch for an opportunity, and think
themselves injured, when they see those safe, whom they wish not to
live. §. 10. Who then does not perceive their profligacy? Is it not
very evident to every one that they do not reproach me with cowardice
from regard to my character, but that being athirst for blood, they
employ these their base devices as a snare, thinking thereby to catch
those whom they seek to destroy? That such is their character is shewn
by their actions, which have convicted them of possessing dispositions
more savage than wild beasts, and more cruel than the Babylonians \footnote{p. 194.

Return to text}. But although the proof against them is sufficiently clear
from all this, yet since they still dissemble with soft words after
the manner of their father the devil, and pretend to charge me with
cowardice, while they are themselves more cowardly than hares; let us
consider what is written in the sacred Scriptures respecting such
cases as this. For thus they will be shewn to fight against the
Scriptures no less than against me, while they detract from the
virtues of the Saints.

15. For if they reproach men for hiding
themselves from those who seek to destroy them, and accuse those who
flee from their persecutors, what will they do when they see Jacob
fleeing from his brother Esau, and Moses withdrawing into Midian for
fear of Pharaoh? What excuse will they make for David, after all this
idle talk, for fleeing from his house on account of Saul, where he
sent to kill him, and for hiding himself in the cave, and for changing
his appearance, until he withdrew from Abimelech \footnote{Achish. 1 Sam. xxi. 13.

Return to text}, and escaped his designs against him? What will they
say, they who are ready to say any thing, when they see the great
Elias, after calling upon God and raising the dead, hiding himself for
fear of Ahab, and fleeing from the threats of Jezebel? At which time
also the sons of the prophets, when they were sought after, hid
themselves with the assistance of Abdias, and lay concealed in caves \footnote{Hist. Ar. §. 53.

Return to text}.

§. 11.

16. Perhaps they have not read these histories;
as being out of date; yet have they no recollection of what is written
in the Gospel? For the disciples also withdrew and hid themselves for
fear of the Jews; and Paul, when he was sought after by the governor
at Damascus, was let down from the wall in a basket, and so escaped
his hands. As the Scripture then relates these things of the Saints,
what excuse will they be able to invent for their wickedness? To
reproach them with cowardice would be an act of madness, and to accuse
them of acting contrary to the will of God, would be to shew
themselves entirely ignorant of the Scriptures. For there was a
command under the Law that cities of refuge should be appointed, in
order that they who were sought after to be put to death, might at
least have some means of saving themselves. And when he who spake unto
Moses, the Word of the Father, appeared in the end of the world, He
also gave this commandment, saying, \emph{But when they persecute you in
this city, flee ye into another} [Mat. x. 23.]: and shortly after
the says, \emph{When ye therefore shall see the abomination of
desolation, spoken of by Daniel the prophet, stand in the holy place,
(whoso readeth, let him understand;) then let them which be in Judea
flee into the mountains: let him which is on the housetop not come
down to take any thing out of his house: neither let him which is in
the field return back to take his clothes}[Mat. xxiv. 15.].
Knowing these things, the Saints regulated their conduct accordingly.
For what our Lord has now commanded, the same also He spoke by His
Saints before His coming in the flesh \footnote{[\emph{t\={e}s ensarkou parousias}], supr. p. 129, r. 3.

Return to text}: and this is the rule which is given unto men to lead them to
perfection,—what God commands, that to do.

§. 12.

17. Wherefore also the Word Himself, being made
man for our sakes, condescended to hide Himself when He was sought
after, as we do: and also when He was persecuted, to flee and avoid
the designs of His enemies. For it became Him, as by hunger and
thirst and suffering, so also by hiding Himself and fleeing, to shew
that He had taken our flesh, and was made man. Thus at the very first,
as soon as He became man, when He was a little child, He Himself by
His Angel commanded Joseph, \emph{Arise, and take the young Child and His
Mother, and flee into Egypt; for Herod will seek the young Child's
life} [Mat. ii. 13.]. And when Herod was dead, we find Him
withdrawing to Nazareth for fear of Archelaus his son. And when
afterwards He was shewing Himself to be God, and made whole the
withered hand, the Pharisees went out, and held a council against Him,
how they might destroy Him; but when Jesus knew it, He withdrew
Himself from thence. So also when He raised Lazarus from the dead, \emph{from
that day forth}, says the Scripture, \emph{they took counsel for to
put Him to death. Jesus therefore walked no more openly among the
Jews; but went thence into a country near to the wilderness} [John
xi. 53, 54.]. Again, when our Saviour said, \emph{Before Abraham was, I
am, the Jews took up stones to cast at Him; but Jesus hid Himself, and
went out of the temple}. And \emph{going through the midst of them, He
went His way, and so passed by} [John viii. 58, 59.]. §. 13. When
they see these things, or rather when they hear of them, for see they
do not, will they not desire, as it is written, to become \emph{fuel of
fire} [Is. ix. 5.], because their counsels and their words are
contrary to what the Lord both did and taught? Also when John was
martyred, and his disciples buried his body, \emph{when Jesus heard of
it, He departed thence by ship into a desert place apart} [Mat.
xiv. 3.].

18. Thus the Lord acted, and thus He taught.
Would that these men were even now ashamed of their conduct, and
confined their rashness to man, nor proceeded to such extreme madness
as even to charge our Saviour with cowardice! for it is against Him
that they now utter their blasphemies. But no one will endure such
madness; nay it will be seen that they do not understand the Gospels.
The cause must be a reasonable and just one, which the Evangelists
represent as weighing with our Saviour to withdraw and to flee; and we
ought therefore to assign the same for the conduct of all the Saints.
(For whatever is written concerning our Saviour in His human nature,
ought to be considered as applying to the whole race of mankind \footnote{vol. 8. p. 241.

Return to text}; because He took our body, and exhibited in Himself
human infirmity.) Now of this cause John has written thus, \emph{They
sought to take Him: but no man laid hands on Him, because His hour was
not yet come} [John vii. 30.]. And before it came, He Himself said
to His Mother, \emph{Mine hour is not yet come} [John ii. 4.]: and to
them who are called His brethren, \emph{My time is not yet come} [John
vii. 6.]. And again, when His time was come, He said to the disciples,
\emph{Sleep on now, and take your rest: for behold, the hour is at hand,
and the Son of man is betrayed into the hands of sinners} [Mat.
xxvi. 45.].

§. 14.

19. Now in so far as He was God and the Word of
the Father, He had no time; for He is Himself the Creator of times \footnote{vol. 8. p. 30, note n.

Return to text}. But being made man, He shews by speaking in this manner that
there is a time allotted to every man; and that not by chance, as some
of the Gentiles imagine in their fables, but a time which He, the
Creator, has appointed to every one according to the will of the
Father. This is written in the Scriptures, and is manifest to all men.
For although it be hidden and unknown to all, what period of time is
allotted to each, and how it is allotted; yet every one knows this,
that as there is a time for spring and for summer, and for autumn and
for winter, so, as it is written, there is a time to die, and a time
to live. And so the time of the generation which lived in the days of
Noah was cut short, and their years were contracted, because the time
of all things was at hand. But to Hezekiah were added fifteen years.
And as God promises to them that serve him truly, \emph{I will fulfil the
number of thy days} [Gen. xxv. 8.]; Abraham dies \emph{full of days},
and David besought God, saying, \emph{Take me not away in the midst of my
days} [Ps. cii. 24.]. And Eliphaz, one of the friends of Job, being
assured of this truth, said, Thou shalt come to thy grave in a full
age, like as a shock of corn cometh in in his season [Job v. 26.]. And
Solomon confirming his words, says, \emph{The souls of the unrighteous
are taken away untimely} [vid. Prov. x. 27.]. And therefore he
exhorts in the book of Ecclesiastes, saying, \emph{Be not overmuch
wicked, neither be thou foolish: why shouldest thou die before thy
time?} [Eccles. vii. 17.]

§. 15.

20. Now as these things are written in the
Scriptures, the case is clear, that the saints \footnote{or sacred writers, supr. p. 128, r. 2.

Return to text} knew that a certain time was allotted to every man, but that no
one knows the end of that time, is plainly intimated by the words of
David, \emph{Declare unto me the shortness of my days} [Ps. cii. 23.
Sept.]. What he did not know, that he desired to be informed of.
Accordingly the rich man also, while he thought that he had yet a long
time to live, heard the words, \emph{Thou fool, this night thy soul shall
be required of thee: then whose shall those things be which thou hast
provided?} [Luke xii. 20.] And the Preacher speaks confidently in
the Holy Spirit, and says, \emph{Man also knoweth not his time}
[Eccles. ix. 12.]. Wherefore the Patriarch Isaac said to his son Esau,
\emph{Behold, I am old, and I know not the day of my death} [Gen.
xxvii. 2.].

21. Our Lord therefore, although as God, and the
Word of the Father, He both knew the period which He had allotted to
all, and was conscious of the time for suffering, which He Himself had
appointed also to His own body; yet since He was made man for our
sakes, He hid Himself when He was sought after before that time came,
as we do; when He was persecuted, He fled; and avoiding the designs of
His enemies He passed by, and so \emph{went through the midst of them}
[Luke iv. 30.]. But when He had brought on that time which He Himself
had appointed, at which He desired to suffer in the body for all men,
He announces it to the Father, saying, \emph{Father, the hour is come;
glorify Thy Son} [John xvii. 1.]. And then He no longer hid Himself
from those who sought Him, but stood willing to be taken by them; for
the Scripture says, He said to them that came unto Him, \emph{Whom seek
ye?} and when they answered, \emph{Jesus of Nazareth}, He saith
unto them, \emph{I am He whom ye seek} [John xviii. 4, 5.]. And this
He did even more than once; and so they straightway led Him away to
Pilate. He neither suffered Himself to be taken before the time came,
nor did He hide Himself when it was come; but gave himself up to them
that conspired against Him, that He might shew to all men that the
life and death of man depends upon the divine sentence; and that
without our Father which is in heaven, neither a hair of man's head
can become white or black, nor a sparrow fall into the net \footnote{p. 194, r. 3.

Return to text}.

§. 16.

22. Our Lord therefore, as I said before, thus
offered Himself for all; and the Saints having received this example
from their Saviour, (for all of them before His coming, nay always,
were under His teaching \footnote{vol. 8. p. 236, note c.

Return to text},) in their
conflicts with their persecutors acted lawfully in flying, and hiding
themselves when they were sought after. And being ignorant, as men, of
the end of the time which Providence had appointed unto them,
they were unwilling at once to deliver themselves up into the power of
those who conspired against them. But knowing on the other hand what
is written, that \emph{the times} of man \emph{are in God's hand}
[Ps. xxxi. 15.], and that \emph{the Lord killeth}, and the Lord \emph{maketh
alive} [1 Sam. iii. 6.], they the rather endured unto the end, \emph{wandering
about}, as the Apostle has spoken, \emph{in sheepskins, and goatskins,
being destitute, tormented, wandering in deserts}, and hiding
themselves \emph{in dens and caves of the earth} [Heb. xi. 37, 38.];
until either the appointed time of death arrived, or God who had
appointed their time spoke unto them, and stayed the designs of their
enemies, or else delivered up the persecuted to their persecutors,
according as it seemed to Him to be good. This we may well learn
respecting all men from David: for when Joab instigated him to slay
Saul, he said, \emph{As the Lord liveth, the Lord shall smite him; or his
day shall come to die; or he shall descend into battle, and perish;
the Lord forbid that I should stretch forth my hand against the
Lord's anointed} [ 1 Sam. xxvi. 10, 11.].

§. 17.

23. And if ever in their flight they voluntarily
came unto those that sought alter them, they did not do so without
reason: but when the Spirit spoke unto them, then as righteous men
they went and met their enemies; by which they also shewed their
obedience and zeal towards God. Such was the conduct of Elias, when,
being commanded by the Spirit, he shewed himself unto Ahab; and of
Micaiah the prophet when he came to the same Ahab; and of the prophet
who cried against the altar in Samaria, and rebuked Jeroboam; and of
Paul when he appealed unto Cæsar. It was not certainly through
cowardice that they fled: God forbid. The flight to which they
submitted was rather a conflict and war against death. For with wise
caution they guarded against these two things; either that they should
offer themselves up without reason, (for this would have been to kill
themselves, and to become guilty of death, and to transgress the
saying of the Lord, \emph{What God hath joined, let not men put asunder}
[Mat. xix. 6.];) or that they should willingly subject themselves to
the reproach of negligence, as if they were unmoved by the
tribulations which they met with in their flight, and which brought
with them sufferings greater and more terrible than death. For he that
dies, ceases to suffer \footnote{vid. supr. p. 19.

Return to text}; but he
that flies, while he expects daily the assaults of his enemies,
esteems death a lighter evil. They therefore whose course was
consummated in their flight did not perish dishonourably, but attained
as well as others the glory of martyrdom. Therefore it is that Job is
accounted a man of mighty fortitude, because he endured to live under
so many and such severe sufferings, of which he would have had no
perception, had he come to his end.

24. Wherefore the blessed Fathers thus regulated
their conduct also; they shewed no cowardice in fleeing from the
persecutor, but rather manifested their fortitude of soul in shutting
themselves up in close and dark places, and living a hard life. Yet
did they not desire to avoid the time of death when it arrived; for
their concern was neither to shrink from it when it came, nor to
forestall the sentence determined by Providence, nor to resist His
dispensation, for which they knew themselves to be preserved; lest by
acting hastily, they should become to themselves the cause of terror:
for thus it is written, \emph{He that is hasty with his lips, shall bring
terror upon himself}. [Prov. xiii. 3. Sept.]

§. 18.

25. Of a truth no one can possibly doubt that
they were well furnished with the virtue of fortitude. For the
Patriarch Jacob who had before fled from Esau, feared not death when
it came, but at that very time blessed the Patriarchs, each according
to his deserts. And the great Moses who previously had hid himself
from Pharaoh, and had withdrawn into Midian for fear of him, when he
received the commandment, \emph{Go into Egypt} [vid. Ex. iii. 10.],
feared not to do so. And again when he was bidden to go up into the
mountain Abarim and die, he delayed not through cowardice, but even
joyfully proceeded thither. And David who had before fled from Saul,
feared not to risk his life in war in defence of his people; but
having the choice of death or of flight set before him, when he might
have fled and lived, he wisely preferred death. And the great Elias
who had at a former time hid himself from Jezebel, shewed no cowardice
when he was commanded by the spirit to meet Ahab, and to reprove
Ochozias. And Peter who had hid himself for fear of the Jews, and the
Apostle Paul into was let down in a basket, and fled, when they were
told, \emph{Ye must bear witness at Rome} [vid. Acts xxiii. 11.],
deferred not the journey; yea, rather, they departed rejoicing \footnote{vid. Euseb. Hist. ii. 25.

Return to text}; the one as hastening to meet his friends, received his death with
exultation; and the other shrunk not from the time when it came, but
gloried in it, saying, \emph{For I am now ready to be offered, and the
time of my departure is at hand} [2 Tim. iv. 6.].

§. 19.

26. These things both prove that their previous
flight was not the effect of cowardice; and testify that their after
conduct also was of no ordinary character: and they loudly proclaim
that they possessed in a high degree the virtue of fortitude. For
neither did they withdraw themselves to gratify a slothful timidity,
seeing they were at such times under the practice of a severer
discipline \footnote{[\emph{ton tonon t\={e}s ask\={e}se\={o}s}].

Return to text} than at others; nor were
they condemned for their flight, or charged with cowardice, by such as
are now so fond of criminating others. Nay they were blessed through
that declaration of our Lord, \emph{Blessed are they which are persecuted
for righteousness' sake} [Mat. v. 10.]. Nor yet were these their
sufferings without profit to themselves; for having tried them as \emph{gold
in the furnace}, as the Book of Wisdom has said, God found them
worthy for Himself. And then they shone the more \emph{like sparks},
being saved from them that persecuted them, and delivered from the
designs of their enemies, and preserved to the end that they might
teach the people, that their flight and escape from the rage of them
that sought after them, was according to the dispensation of the Lord.
And so they became dear in the sight of God, and obtained the most
glorious testimony to their fortitude.

§. 20.

27. Thus for example, the Patriarch Jacob was
favoured in his flight with many, even divine visions, and remaining
quiet himself; he had the Lord on his side, rebuking Laban, and
hindering the designs of Esau; and afterwards he became the father of
Judah, of whom sprang the Lord according to the flesh; and he
dispensed the blessings to the Patriarchs. And Moses the beloved of
God, when he was in exile, then it was that he saw that great sight,
and being preserved from his persecutors, was sent as a prophet into
Egypt, and being made the minister of those mighty wonders and of the
Law, he led that great people in the wilderness. And David when he was
persecuted wrote the Psalm, \emph{My heart is inditing a good matter}
[Ps. xlv. 1.]; and, \emph{Our God shall come, and shall not keep}
silence [Ps. l. 3.]. And again he speaks more confidently,
saying, \emph{Mine eye hath seen his desire upon mine enemies} [Ps.
liv. 7.]; and again, \emph{In God have I put my trust; I will not be
afraid what man can do unto me} [Ps. lvi. 11.]. And when he fled
and escaped from the face of Saul to the cave, he said, \emph{He hath
sent from heaven, and hath saved me. He hath given them to reproach
that would tread me under their feet. God hath sent His mercy and
truth, and hath delivered my soul from the midst of lions} [Ps.
lvii. 3.]. Thus he too was saved according to the dispensation of God,
and afterwards became king, and received the promise, that from his
seed our Lord should spring.

28. And the great Elias, when he withdrew himself
to mount Carmel, called upon God, and destroyed at once more than four
hundred prophets of Baal; and when there were sent to take him two
captains of fifty with their hundred men, he said, \emph{Let fire come
down from heaven} [2 Kings i. 10.], and thus rebuked them. And he
too was preserved, so that he anointed Elisha in his own stead, and
became a pattern of virtue for the sons of the prophets. And the
blessed Paul, after writing these words, \emph{what persecutions I
endured; but out of them all the Lord delivered me, and will deliver}
[2 Tim. iii. 11.]; could speak more confidently and say, \emph{But in all
these things we are more than conquerors, for nothing shall separate
us from the love of Christ} [Rom. viii. 35, 37.] \footnote{pp. 149, 220.

Return to text}. For then it was that he was caught up to the third heaven, and
admitted into paradise, where he heard \emph{unspeakable words, which it
is not lawful for a man to utter} [2 Cor. xii. 4.]. And for this
end was he then preserved, that \emph{from Jerusalem even unto Illyricum}
he might \emph{fully preach the Gospel} [Rom. xv. 19.].

§. 21.

29. The flight of the saints therefore was
neither blameable nor unprofitable. If they had not avoided their
persecutors, how would it have come to pass that the Lord should
spring from the seed of David? Or who would have preached the glad
tidings of the word of truth? It was for this that the persecutors
sought after the saints, that there might be no one to teach, as the
Jews charged the Apostles; but for this cause they endured all things,
that the Gospel might be preached \footnote{p. 184.

Return to text}.
Behold, therefore, in that they were thus engaged in conflict with
their enemies, they passed not the time of their flight unprofitably,
nor while they were persecuted, did they forget the welfare of others:
but as being ministers of the good word, they grudged not to
communicate it to all men; so that even while they fled, they preached
the Gospel, and gave warning of the wickedness of those who conspired
against them, and confirmed the faithful by their exhortations.

30. Thus the blessed Paul, having found it so by
experience, declared beforehand, \emph{As many as will live godly in
Christ, shall suffer persecution} [2 Tim. iii. 12.]. And so he
straightway prepared them that fled for the trial, saying, \emph{Let us
run with patience the race that is set before us} [Heb. xii. 1.];
for although there be continual tribulations, \emph{yet tribulation
worketh patience, and patience experience, and experience hope, and
hope maketh not ashamed} [Rom. v. 4.]. And the Prophet Esaias when
such-like affliction was expected, exhorted and cried aloud, \emph{Come,
my people, enter thou into thy chambers, and shut thy doors: hide
thyself as it were for a little moment, until the indignation be
overpast} [Is. xxvi. 29.] \footnote{p. 186.

Return to text}, And
the Preacher who knew the designs of the wicked against the righteous,
and said, \emph{If thou seest the oppression of the poor, and violent
perverting of judgment and justice in a province, marvel not at the
matter: for He that is higher than the highest regardeth, and there be
higher than they: moreover there is the profit of the earth}
[Eccles. v. 8, 9]. He had his own father David for an example, who had
himself experienced the sufferings of persecution, and who supports
them that suffer by the words, \emph{Be of good courage, and He shall
strengthen your heart, all ye that put your trust in the Lord} [Ps.
xxxi. 24.]; for them that so endure, not man, but the Lord Himself,
(he says,) \emph{shall help them, and deliver them, because they put
their trust in Him}[Ps. xxxviii. 40.]: for I also \emph{waited
patiently for the Lord, and He inclined unto me, and heard my calling;
He brought me up also out of the lowest pit, and out of the mire and
clay} [Ps. xl. 1.]. Thus is shewn how profitable to the people and
productive of good is the flight of the Saints, howsoever the Arians
may think otherwise.

§. 22.

31. Thus the Saints, as I said before, were
abundantly preserved in their flight by the Providence of God, as
physicians for the sake of them that had need. And to all men
generally, even to us, is this law given, that we should flee when we
are persecuted, and hide ourselves when we are sought after, and not
rashly tempt the Lord, but should wait, as I said above, until the
appointed time of death arrive, or the Judge determine something
concerning us, according as it shall seem to Him to be good: that we
should be ready, that, when the time calls for us, or when we are
taken, we may contend for the truth even unto death. This rule the
blessed Martyrs observed in their several persecutions. When
persecuted they fled, while concealing themselves they shewed
fortitude, and when discovered they submitted to martyrdom. And if
some of them came and presented themselves to their persecutors \footnote{A.
Vid. instances and passages collected in Pearson's Vind. Ignat. part
ii. c. 9. also Gibbon, ch. xvi. p. 438. Mosheim de Reb. Ante Const. p.
941.

Return to text}, they did not do so without reason; for immediately in that case
they were martyred, and thus made it evident to all that their zeal,
and this offering up of themselves to their enemies, were from the
Spirit.

§. 23.

32. Seeing therefore that such are the commands
of our Saviour, and that such is the conduct of the Saints, let these
persons, to whom one cannot give a name suitable to their
character,—let them, I say, tell us, from whom they learnt to
persecute? They cannot say, from the Saints \footnote{Hist. Arian. §§. 33, 67.

Return to text}. No, but from the Devil; (that is the only answer which is left
them;)—from him who says, \emph{I will pursue, I will overtake} [Ex.
xv. 9.]. Our Lord commanded to flee, and the saints fled: but
persecution is a device of the Devil, and one which he desires to
exercise against all. Let them say then, to which we ought to submit
ourselves; to the words of the Lord, or to their fabrications? Whose
conduct ought we to imitate, that of the Saints, or that of those
whose example they have adopted? But since it is likely they cannot
determine this question, (for, as Esaias said, their minds and their
consciences are blinded, and they think \emph{bitter to be sweet}, and
\emph{light darkness} [Is. v. 20.] \footnote{p. 220. vol. 8. p. 9.

Return to text},)
let some one come forth from among us Christians, and put them to
rebuke, and cry with a loud voice, “It is better to trust in the
Lord, than to attend to the foolish sayings of these men; for the \emph{words}
of the Lord have \emph{eternal life} [John vi. 68.], but the things
which these utter are full of iniquity and blood.”

§. 24.

33. This were sufficient to put a stop to the
madness of these impious men, and to prove that their desire is for
nothing else, but only through a love of contention to utter revilings
and blasphemies. But forasmuch as having once dared to fight
against Christ, they have now become officious, let them enquire and
learn into the manner of my withdrawal from their own friends. For the
Arians were mixed with the soldiers in order to exasperate them
against me, and, as they were unacquainted with my person, to point me
out to them. And although they are destitute of all feelings of
compassion, yet when they hear the circumstances they will surely be
quiet for very shame.

34. It was now night \footnote{p. 176.

Return to text}, and some of the people were keeping a vigil preparatory to a
communion on the morrow, when the General Syrianus suddenly came upon
us with more than five thousand soldiers, having arms and drawn
swords, bows, spears, and clubs, as I have related above. With these
he surrounded the Church, stationing his soldiers near at hand, in
order that no one might be able to leave the Church and pass by them.
Now I considered that it would be unfair in me to desert the people
during such a disturbance, and not to endanger myself in their behalf;
therefore I sat down upon my throne, and desired the Deacon to read
the Psalm, and the people to answer, \emph{For His mercy endureth for
ever} [Ps. cxxxvi. 1.], and then all to withdraw and depart home.
But the General having now made a forcible entry, and the soldiers
having surrounded the Chancel for the purpose of apprehending me, the
Clergy and those of the laity, who were still there, cried out, and
demanded that I should withdraw. But I refused, declaring that I would
not do so, until they had retired one and all. Accordingly I stood up,
and having bidden prayer, I then made my request of them, that all
should depart before me, saying that it was better that my safety
should be endangered, than that any of them should receive hurt. So
when the greater part had gone forth, and the rest were following, the
monks who were there with me and certain of the Clergy came up and
dragged me away. And thus, (Truth is my witness,) while some of the
soldiers stood about the Chancel, and others were going round the
Church, I passed through, under the Lord's guidance, and with His
protection withdrew without observation, greatly glorifying God, that
I had not betrayed the people, but had first sent them away, and
then had been able to save myself, and to escape the hands of them
which sought after me.

§. 25.

35. Now when Providence had delivered me in such
an extraordinary manner, who can justly lay any blame upon me, because
I did not give myself up into the hands of them that sought after me,
nor return and present myself before them? This would have been
plainly to shew ingratitude to the Lord, and to act against His
commandment, and in contradiction to the practice of the Saints. He
who censures me in this matter must presume also to blame the great
Apostle Peter, because though he was shut up and guarded by soldiers,
he followed the angel that summoned him, and when he had gone forth
from the prison and escaped in safety, he did not return and surrender
himself, although he heard what Herod had done. Let the Arian in his
madness censure the Apostle Paul, because when he was let down from
the wall and had escaped in safety, he did not change his mind, and
return and give himself up; or Moses, because he returned not out of
Midian into Egypt, that he might be taken of them that sought after
him; or David, because when he was concealed in the cave, he did not
discover himself to Saul. As also the sons of the prophets remained in
their caves, and did not surrender themselves to Ahab. This would have
been to act contrary to the commandment, since the Scripture says, \emph{Thou
shall not tempt the Lord thy God} [Deut. vi. 16. Mat. iv. 7.]. §.
26. Being careful to avoid such an offence, and instructed by these
examples, I so ordered my conduct; and I do not undervalue the favour
and the help which have been shewn me of the Lord, howsoever these
madmen may gnash their teeth \footnote{Sent. Dion. 16. Hist. Ar. §. 68. 72.

Return to text} against
me. For since the manner of my retreat was such as I have described, I
do not think that any blame whatever can attach to it in the minds of
those who are possessed of a sound judgment: seeing that according to
holy Scripture, this pattern has been left us by the Saints for our
instruction. But there is no atrocity, it would seem, which these men
neglect to practise, nor will they leave any thing undone, which may
shew their own wickedness and cruelty.

36. And indeed their lives are only in accordance
with their spirit and the follies of their doctrine; for there are no
sins that one could charge them with, how heinous soever, that
they do not commit without shame. Leontius \footnote{Hist. Arian. §. 28.

Return to text}, for instance, being censured for his intimacy with a certain
young woman, named Eustolium, and prohibited from living with her,
mutilated himself for her sake, in order that he might be able to
associate with her freely. He did not however clear himself from
suspicion, but rather on this account he was degraded from his rank as
Presbyter, although the heretic Constantius by violence caused him to
be named a Bishop. Narcissus \footnote{p. 60. vol. 8. p. 99.

Return to text},
besides being charged with many other transgressions, was degraded
three times by different Councils; and now he is the most wicked among
them. And George \footnote{p. 25.

Return to text} who was a
Presbyter, was degraded on account of his vices, and although he had
nominated himself a Bishop, he was nevertheless a second time degraded
in the great Council of Sardica. And besides all this, his dissolute
life is notorious, for he is condemned even by his own friends, as
making the end of existence and happiness to consist in the commission
of the most disgraceful crimes.

§. 27.

37. Thus each surpasses the other in his own
peculiar vices. But there is a common blot that attaches to them all,
in that through their heresy they are enemies of Christ, and are no
longer called Christians \footnote{B.
Vid. supr. p. 149, r. 4. infr. Hist. Arian. §§. 17. 34 fin. 41 init.
59 fin. 64 init. vol. 8. p. 27, note h. pp. reff. 85, 1. 179, 4. 182.
188, 4. 194, 2.

Return to text

}, but Arians.
They ought indeed to accuse each other of the sins they are guilty of,
for they are contrary to the faith of Christ; but they rather conceal
them for their own sakes. And it is no wonder, that being possessed of
such a spirit, and implicated in such vices, they persecute and seek
after those who follow not the same impious heresy as themselves; that
they delight to destroy them, and are grieved if they fail of
obtaining their desires, and think themselves injured, as I said
before, when they see those alive, whom they wish to perish. May they
continue to be injured in such sort, that they may lose the power of
inflicting injuries, and that those whom they persecute may give
thanks unto the Lord, and say in the words of the twenty-sixth Psalm, \emph{The
Lord is my light and my salvation; whom then shall I fear? The Lord is
the strength of my life; of whom then shall I be afraid? When the
wicked, even mine enemies and my foes,}\emph{come upon me to
eat up my flesh, they stumbled and fell} [Ps. xxvii. 1.]: and again
in the thirtieth Psalm, \emph{Thou hast known my soul in adversities;
Thou hast not shut me up into the hands of my enemies; Thou hast set
my feet in a large room} [Ps. xxxi. 7, 8.]; in Christ Jesus our
Lord, through whom to the Father in the Holy Spirit be glory and power
for ever and ever. Amen.

\xchapter{VI. An Epistle of our Holy Father Athanasius, Archbishop of Alexandria, to his brother Serapion, concerning the death of Arius}

———————

[S. Serapion, Bishop of Thmuis, was a friend of
St. Anthony’s; to him the Saint on his death, which took place
shortly before the following Letter from Athanasius, left one of his
sheepskins, leaving the other to S. Athanasius himself. His fellowship
with Athanasius in persecution, has gained him the title of Confessor,
and his accomplishments and talents that of Scholasticus. Jerom. de
Vir. Illustr. 99. At his suggestion Athanasius about the same date
wrote his work upon the divinity of the Holy Spirit, addressing it to
him. He seems also to have booms a correspondent of Apollinaris. His
name is found in the Roman Martyrology under March 21. It appears from
the commencement of the following Letter, written A.D. 358-360, that Serapion had asked
Athanasius, first for a history of his times, next for a refutation of
Arianism, and thirdly for an account of the death of Arius. The death
of Arius is the subject of this Letter itself; for the history of his
times he refers him to his history of Arianism addressed to the Monks,
which he sent him at the same time; and the refutation of Arianism,
which was also addressed to the Monks, has sometimes been supposed to
be the four celebrated Orations which are his principle dogmatic work.
Though in strict order of time the Epistles both to Serapion and to
the Monks are later than the history, and the latter Epistle, as
containing scarcely an allusion to the History, might easily be
detached from it, yet it seems best in a matter of this kind to follow
the arrangement adopted in the Benedictine Edition.]

———————

1. Athanasius to Serapion a brother and
fellow-minister sends health in the Lord.

§. 1.

1. I HAVE
read the letters of your Piety, in which you have requested roe to
make known to you the events of my times relating to myself, and
to give an account of that most impious heresy of the Arians, in
consequence of which I have endured these sufferings, and also of the
manner of the death of Arius. With two out of your three demands I
have readily undertaken to comply, and have sent to your Godliness the
letter which I wrote to the Monks; from which you will be able to
learn my own history as well as that of the heresy. But with respect
to the other matter, I mean the Death, I debated with myself for a
long time, fearing lest any one should suppose that I was exulting in
the death of that man. But yet, since a disputation which has taken
place amongst you concerning the heresy, has issued in this question,
whether Arius died in communion with the Church; I therefore was
necessarily desirous of giving an account of his death, as thinking
that the question would thus be set at rest, considering also that by
making this known I should at the same time silence those who are fond
of contention. For I conceive that when the wonderful \footnote{Margin Notes

1. [\emph{thaumatos}]. vid. p. 217, r. 5.

Return to text} circumstances connected with his death become known, even those
who before questioned it will no longer venture to doubt that the
Arian heresy is hateful in the sight of God \footnote{[\emph{theostug\={e}s}], p. 217, r. 6.

Return to text}.

§. 2.

2. I was not at Constantinople when he died, but
Macarius the Presbyter was, and I heard the account of it from him.
Arius had been summoned by the Emperor Constantine, through the
interest of the Eusebians; and when he entered the presence the
Emperor enquired of him, whether he held the Faith of the Catholic
Church? And he declared upon oath that he held the right \footnote{[\emph{orth\={o}s}].

Return to text} Faith, and gave in an account of his Faith in writing,
suppressing the points for which he had been cast out of the Church by
the Bishop Alexander, and speciously \footnote{[\emph{hypokrinomenos}].

Return to text} alleging expressions out of the Scriptures. When therefore he
swore that he did not profess the opinions for which Alexander had
excommunicated him, the Emperor dismissed him, saying, “If thy Faith
be right, thou hast done well to swear; but if thy Faith be impious,
and thou hast sworn, God judge thee according to thy oath.” When he
thus came forth from the presence of the Emperor, the Eusebians with
their accustomed violence desired to bring him into the Church \footnote{p. 147, 8.

Return to text}. But Alexander the Bishop of Constantinople of blessed
memory \footnote{p. 162, r. 3.

Return to text} resisted them,
saying that the inventor of the heresy ought not to be admitted to
communion; whereupon the Eusebians threatened, declaring, “As we
have caused him to be summoned \footnote{[\emph{kl\={e}th\={e}nai}], vid. p. 49, r. 4. p. 70, r. 1.

Return to text} by the Emperor, in opposition to your wishes, so tomorrow,
though it be contrary to your desire, Arius shall have communion with
us in this Church.” It was the Sabbath when they said this.

§. 3.

3. When the Bishop Alexander heard this, he was
greatly distressed, and entering into the Church, he stretched forth
his hands unto God, and bewailed himself; and casting himself upon his
face in the Chancel \footnote{[\emph{heiratei\={o}i}], p. 20.

Return to text}, he
prayed, lying upon the pavement. Macarius also was present, and prayed
with him, and heard his words. And he besought these two things,
saying, “If Arius is brought to communion tomorrow, let me Thy
servant depart, and destroy not the pious with the impious; but if
Thou wilt spare Thy Church, (and I know that Thou wilt spare,) look
upon the words of the Eusebians, and give not Thine inheritance to
destruction and reproach, and take off Arius \footnote{[\emph{aron areion}].

Return to text}, lest if he enter into the Church, the heresy also may seem to
enter with him, and henceforward impiety \footnote{p. 29, r. 1.

Return to text} be accounted for piety.” When the Bishop had thus prayed, he
retired in great anxiety; and a wonderful and extraordinary
circumstance took place. While the Eusebians threatened, the Bishop
prayed; but Arius, who had great confidence in the Eusebians, and
talked very wildly, urged by the necessities of nature withdrew \footnote{[\emph{eis thakas}].

Return to text}, and suddenly, in the language of Scripture, \emph{falling
headlong he burst asunder in the midst} [Acts i. 18.], and
immediately expired as he lay, and was deprived both of communion and
of his life together.

§. 4.

4. Such was the end of Arius and the Eusebians,
overwhelmed with shame, buried \footnote{vid. Acts v. 6. 10.

Return to text} their accomplice, while the blessed Alexander, amidst the
rejoicings of the Church, celebrated the Communion with piety and
orthodoxy, praying with all the brethren, and greatly glorifying God,
not as exulting in his death, (God forbid!) for it is appointed unto
all men once to die, but because this thing had been shewn forth in a
manner surpassing the expectations of all men. For the Lord Himself
judging between the threats of the Eusebians and the prayer of
Alexander, condemned the Arian heresy, shewing it to be unworthy
of communion with the Church, and making manifest to all, that
although it receive the support of the Emperor and of all mankind, yet
it has been condemned by the Church herself.

5. Thus this antichristian workshop \footnote{[\emph{ergast\={e}rion}].

Return to text} of the Arian fanatics has been shewn to be unpleasing to God
and impious; and many of those who before were deceived by it have
changed their opinions. For none other than the Lord Himself who was
blasphemed by them has condemned the heresy which rose up against Him,
and has again shewn, that howsoever the Emperor Constantius may now
use violence to the Bishops in behalf of it, yet it is excluded from
the communion of the Church, and alien from the kingdom of heaven \footnote{[\emph{allotria t\={o}n ouran\={o}n}].

Return to text}. Wherefore also let the question which has arisen among you be
henceforth set at rest; (for this is the agreement that was made among
you,) and let no one join himself to the heresy, but let even those
who have been deceived repent. For who shall receive a heresy which
the Lord has condemned? And will not he who takes up the support of
that which He has made excommunicate, be guilty of great impiety, and
manifestly an enemy of Christ?

§. 5.

6. Now this is sufficient to confound the
contentious; read it therefore to those who before raised this
question, as well as what I have briefly \footnote{p. 216, r. 2.

Return to text} addressed to the Monks against the heresy, in order that they
may be led thereby more strongly to condemn the impiety and wickedness
of the Arian fanatics. Do not however consent to give a copy of these
letters to any one, neither transcribe them for yourself, (I have
signified the same to the Monks also \footnote{p. 217, r. 7.

Return to text};) but as a sincere friend, if any thing is wanting in what I
have written, add it, and immediately send them back to me. For you
will be able to learn from the letter which I have written to the
Brethren, what pains it has cost me to write it \footnote{p. 215, r. 2.

Return to text}, and also to perceive that it is not safe for the writings of
an individual \footnote{[\emph{idi\={o}tou}], p. 218, r. 1.

Return to text} to be
published, (especially if they relate to the highest and chief \footnote{[\emph{koruphaiotat\={o}n}].

Return to text} doctrines,) for this reason;—lest what is imperfectly
expressed through infirmity or the obscurity of language, do hurt to
the reader. For the majority of men do not consider the faith or the
aim of the writer \footnote{p. 130, r. 2. p. 134, r. 4.

Return to text},
but either through envy or a spirit of contention, receive what is
written as themselves choose, according to an opinion which they
have previously formed, and misinterpret it to suit their pleasure.
But the Lord grant that the Truth and a sound \footnote{[\emph{hygiainousan}], vid. Alex. Ep. Encycl. §. 5 fin.

Return to text} faith in our Lord Jesus Christ may prevail among all, and
especially among those to whom you read this. Amen.

\xchapter{VII. An Epistle of our Holy Father Athanasius, Archbishop of Alexandria, to the Monks}

———————

[The beautiful and striking Letter which follows
formed the introduction to a work, which the Author, as he says in the
course of it, thought unworthy of being preserved for posterity. Some
critics have supposed it to be the Orations against the Arians, which
form his greatest work; but this opinion can hardly be maintained,
though the discussion of it does not belong to this place. The Epistle
to the Monks was written in 358, or later, but before the foregoing
Epistle to Serapion.]

———————

§. 1.

1. To those in every place who are living a
monastic life, who are established in the faith of God, and sanctified
in Christ, and who say, \emph{Behold, we have forsaken all, and I
followed Thee} [Mat. xix. 27.], brethren dearly beloved and longed
for, a full greeting in the Lord.

1. IN
compliance with your affectionate request, which you have frequently
urged upon me, I have written a short account of the sufferings which
ourselves and the Church have undergone, refuting, according to my
ability, the accursed \footnote{Margin Notes

1. [\emph{musaran}].

Return to text}
heresy of the Arian fanatics, and proving how entirely it is alien
from the Truth. And I thought it needful to represent to your Piety
what pains \footnote{p. 213, r. 4.

Return to text} the writing
of theses things has cost me, in order that you may understand thereby
how truly the blessed Apostle has said, O the depth of the \emph{riches
both of the wisdom and knowledge of God} [Rom xi. 33.]; and may
kindly bear with a weak man such as I am by nature. For the more I
desired to write, and endeavoured to force myself to understand the
Divinity of the Word, so much the more did the knowledge thereof
withdraw itself from me; and in proportion as I thought that I
apprehended it, in so much I perceived myself to fail of doing so.
Moreover also I was unable to express in writing even what I seemed to
myself to understand; and that which I wrote was unequal to the
imperfect shadow of the truth which existed in my conceptions.

§. 2.

2. Considering therefore how it is written in the
Book of Ecclesiastes, \emph{I said, I will be wise, but it was far from
me; That which is far off, and exceeding deep, who shall find it out?}
[Eccles. vii. 23, 24.] and what is said in the Psalms, \emph{The
knowledge of Thee is too wonderful for me; it is high, I cannot attain
unto it} [Ps. cxxxix. 6.]; and that Solomon says, \emph{It is the
glory of God to conceal a thing} [Prov. xxv. 2.]; I frequently
designed to stop and to cease writing; believe me \footnote{pp. 240, 158.

Return to text}, I did. But lest I should be found to disappoint you, or by my
silence to lead into impiety those who have made enquiry of you, and
are given to disputation, I constrained myself to write briefly \footnote{p. 213, r. 2.

Return to text}, what I have now sent to your Piety. For although a perfect
apprehension of the truth is at present far removed from us by reason
of the infirmity of the flesh; yet it is possible, as the Preacher
himself has said, to perceive the madness of the impious, and having
found it, to say that it is \emph{more bitter than death} [Eccles.
vii. 26.]. Wherefore for this reason, as perceiving this and able to
find it out, I have written, knowing that to the faithful the
detection of impiety is a sufficient information wherein piety
consists. For although it be impossible to comprehend what God is, yet
it is possible to say, what He is not \footnote{This negative character of our knowledge,
whether of the Father or of the Son, is insisted on by other writers.
“When we speak of the substance of any being, we have to say what it
is, not what it is not; however, as relates to God, it is impossible
to say what He is as to His substance. All we can know about the
Divine Nature is, that it is not to be known; and whatever positive
statements we make concerning God, relate not to His Nature, but to
the accompaniments of His Nature.” Damasc. F. O. i. 4. S. Basil ad
Eunom. i. 10 speaks similarly of the negative attributes, (so to
speak,) of the Divine Nature, adding, however, the positive. And St.
Austin says, “Totum ab animo rejicite; \emph{quidquid occurrerit,
negate} ... dicite \emph{non est illud}.” August. Enarrat. 2. in
Psalm xxvi. 8. “How,” says St. Cyril, “the Father begat the Son,
we profess not to tell; \emph{only} we insist upon its \emph{not}
being in \emph{this} manner or \emph{that}.” Catech. xi. 11.
“Patrem \emph{non} esse Filium, sed habere Filium qui Pater \emph{non}
sit; Filium \emph{non} esse Patrem, sed Filium Dei esse natum; sanctum
quoque Paracletum esse, qui \emph{nec} Pater sit ipse, \emph{nec}
Filius, sed a Patre Filioque procedat. Anonym. in Append. Aug. Oper.
t. 5. p. 383.

Return to text}. And we know that He is not as man; and that it is not
lawful to conceive of any created \footnote{[\emph{t\={o}n genn\={e}t\={o}n}], vol. 8. p. 261, note e.

Return to text} nature as existing in Him. So also respecting the Son of God,
although we are by nature very far from being able to apprehend Him;
yet it is possible and easy to condemn the assertions of the heretics
concerning Him, and to say, that the Son of God is not such; nor is it
lawful even to conceive in our minds such things as they speak,
concerning His divinity; much less to utter them with the lips.

§. 3.

3. Accordingly I have written as well as I was
able; and you, dearly beloved, receive these communications not as
containing a perfect exposition of the doctrine of the divinity of the
Word, but as being merely a refutation of the impiety of the enemies
of Christ, and as containing and affording to those who desire it,
suggestions \footnote{[\emph{aphorm\={e}n}].

Return to text} for
arriving at a pious and sound \footnote{p. 214, r. 1.

Return to text} faith in Christ. And if in any thing they are defective, (and I
think they are defective in all respects,) pardon it with a pure
conscience, and only receive favourably the boldness \footnote{[\emph{to tolm\={e}ron}].

Return to text} of my good intentions in support of godliness. For an utter
condemnation of the heresy of the Arians, it is sufficient for you to
know the judgment which has been given by the Lord in the death of
Arius, of which you have already been informed by others. \emph{For the
Lord of Hosts hath purposed, and who shall disannul it?} [Is. xiv.
27.] and whom the Lord hath condemned who shall justify \footnote{so quoted p. 148.

Return to text}? After such a sign \footnote{[\emph{s\={e}meion}], vid. p. 211, r. 1.

Return to text}
has been given, who does not now acknowledge, that the heresy is hated
of God \footnote{[\emph{theomis\={e}tos}], vid. p. 211, r. 2.

Return to text}, however it
may have men for its patrons?

4. Now when you have read this account, pray for
me, and exhort one another so to do. And immediately send it back to
me, and suffer no one whatever to take a copy of it, nor transcribe it
for yourselves \footnote{p. 213, r. 3.

Return to text}. But
like good money-changers \footnote{On this celebrated text, as it may be called, which is cited so
frequently by the Fathers, vid. Coteler. in Const. Apol. ii. 36. in
Clement. Hom. ii. 51. Potter in Clem. Strom. i. p. 425. Vales. in
Euseb. Hist. vii. 7. vid. also S. Cyril, Catech. tr. p. 78, note o.

Return to text

}
be satisfied with the reading; but read it repeatedly if you desire to
do so. For it is not safe that the writings of us babblers and
private persons \footnote{[\emph{idi\={o}t\={o}n}],
p. 213, r. 5. Apol. contr. Ar. §. 9. supr. p. 27. §. 12. p. 30.

Return to text}
should fall into the hands of them that shall come after. Salute one
another in love, and also all that come unto you in piety and faith.
For \emph{if any man}, as the Apostle has said, \emph{love not the Lord,
let him be anathema. The grace of our Lord Jesus Christ be with you.
Amen} [1 Cor. xvi. 22.].

\xchapter{VIII. The History of the Arians, [Down to the Year 357, the beginning being lost.]}

———————

[The earlier portion of this History, which seems
to have commenced with the Author's elevation to his see, has not been
preserved, because, as Montfaucon conjectures, it was considered but a
repetition of the second part of the Apology against the Arians, §.
59–84. pp. 88–116. supr. He notices a correspondence even in the
words employed in the two works, at the place in the Apology where the
line of narrative may be considered to be taken up by the opening but
broken sentence of the following History. In the beginning of §. 84.
of the Apology, supr. p. 116, towards the end of its second part,
Athanasius says, ``As such is the nature of their machinations, so they
\emph{very soon} shewed plainly the reasons of their conduct. For,
when they went away, they took the Arians with them to Jerusalem, and
there \emph{admitted them to communion};'' and in the beginning, as
extant, of the History. ``And \emph{not long after}, they proceeded to
put in execution the designs for the sake of which they had had
recourse to these artifices; for they \emph{no sooner} had formed
their plans, but they \emph{immediately admitted the Arians to communion}.''
vid. also infr. p. 220, r. 2. Papebroke, whom Tillemont in the main
follows, considers that the whole Apology formed a sort of third part
of the Work addressed to the Monks, (the dogmatic treatise being the
first of the three.) And in maintenance of this opinion he proposes an
ingenious though untenable emendation of some words in the text of
Athanasius, or rather in the notes added to the text by his copyists.
(in Maii 2. p. 187.) A question has been raised about the genuineness
of the work before us, under the idea that it probably was the writing
of a companion of Athanasius, not of the Saint himself. It cannot be
denied that in parts it is written in a livelier and terser, not to
say freer, style than his other works, and he speaks of himself in the
third person. And there is a passage, where, if the text be not
corrupt, the writer distinguishes himself from Athanasius, §. 52. But
on the other hand, there is a passage in which he speaks in the first
person where none but Athanasius can be meant. vid. §. 21. p. 236.
And he speaks of himself in other works in the third person, e.g. Orat.
i. §. 3. Moreover, it is plain that the very circumstance that he was
not writing in his own person would make a considerable alteration in
his mode of writing, not to dwell on the difference between an apology
and what is a history and invective. Some instances of agreement in
words, phrases, texts, \&c. are pointed out in the margin and
notes.]

———————

\xsection{Chapter 1. Arian Persecution under Constantine}

§. 1.

1. … AND
not long after they proceeded to put in execution the designs for the
sake of which they had had recourse to these artifices; for they
no sooner had formed their plans, but they immediately admitted the
Arians to communion. They set aside the repeated condemnations which
had been passed upon them, and again pretended the imperial authority
\footnote{Margin Notes

1. p. 246.

Return to text} in their behalf. And
they were not ashamed to say in their letters, ``since Athanasius has
suffered, all opposition \footnote{[\emph{ph?ones}]. vid. twice p. 116 fin.

Return to text}
has ceased, and let us henceforward receive the Arians;'' adding, in
order to frighten their hearers, 'because the Emperor has commanded
it.' Moreover they were not ashamed to add, ``for these men profess
orthodox opinions;'' not fearing that which is written, \emph{Woe unto
them that call bitter sweet, that put darkness for light} [Is. v.
20.] \footnote{vol. 8. p. 9. supr. p. 205, r. 2.

Return to text}; for they are
ready to undertake any thing in support of their heresy. Now is it not
hereby plainly proved to all men, that we both suffered heretofore,
and that you now persecute us, not under the authority of an
Ecclesiastical sentence \footnote{infr. §. 76.

Return to text},
but on the ground of the Emperor's threats, and on account of our
Piety towards Christ? As also they conspired in like manner against
the Bishops, fabricating charges against them also; some of whom are
fallen asleep in the place of their exile, having attained the glory
of Christian confession; and others are at this day banished from
their country, and contend still more and more manfully against their
heresy, saying, \emph{Nothing shall separate us from the love of Christ}
[Rom. viii. 35.] \footnote{pp. 149, 203.

Return to text}.

§. 2.

2. And hence also you may discern its character,
and be able to condemn it more confidently. The man who is their
friend and their associate in impiety, although he is open to ten
thousand charges for other enormities which he has committed; although
the evidence and proof against him are most clear; he is approved of
by them, and straightway becomes time friend of the Emperor, obtaining
favour by his impiety; and making large gains, he acquires confidence
before the magistrates to do whatever he desires. But he who exposes
their impiety, and honestly advocates the cause of Christ, though he
is pure in all things, though he is conscious of no delinquencies,
though he meets with no accuser; yet on the false pretences which they
have framed against him, is immediately seized and sent into
banishment under a sentence of the Emperor, as if he were guilty
of the crimes which they wish to charge upon him, or as if, like
Naboth, he had blasphemed the king. While he who advocates the cause
of their heresy, is sought for and immediately sent to take possession
of the other's Church; and henceforth confiscations and insults, and
all kinds of cruelty are exercised against those who do not receive
him. And what is the strangest thing of all \footnote{[\emph{to paradoxotaton}], vid. p. 32. §. 14 fin.

Return to text}, the man whom the people desire, and know to be blameless [1
Tim. iii. 2.] \footnote{p. 179, infr. p. 222.

Return to text}, the
Emperor takes away and banishes; but him whom they neither desire, nor
know, he sends to them from a distant place \footnote{vid. p. 133, r. 10.

Return to text} with soldiers and letters \footnote{p. 8, r. 3.

Return to text} from himself. And henceforward a strong necessity is laid upon
them, either to hate him whom they love; who has been their teacher,
and their father in godliness; and to love him whom they do not
desire, and to trust their children to one of whose life and
conversation and character they are ignorant; or else certainly to
suffer punishment, if they disobey the Emperor.

§. 3.

3. In this manner the impious are now proceeding,
as heretofore, against the orthodox; giving proof of their malice and
impiety amongst all men every where. For granting \footnote{[\emph{est\={o}}], vid. Apol. contr. Ar. 35. supr. p. 56.

Return to text} that they have justly accused Athanasius; yet what have the
other Bishops done? On what grounds can they charge them? Has there
been found in their case too the dead body of an Arsenius? Is there a
Presbyter Macarius, or has a chalice been broken amongst them? Is
there a Meletian to play the hypocrite? No: but as their proceedings
against the other Bishops shew the charges which they have brought
against Athanasius, in all probability, to be false; so their attacks
upon Athanasius make it plain, that their accusations of the other
Bishops are unfounded likewise. This heresy has come forth upon the
earth like some wild monster, which not only injures the innocent with
its words, as with teeth \footnote{vid. Dan. vii. 5. 7.

Return to text};
but it has also hired external power to assist it in its designs.

4. And strange it is that, as I said before, no
accusation is brought against any of them; or if any be accused, he is
not brought to trial; or if a shew of enquiry be made, he is acquitted
against evidence, while the convicting party is plotted against,
rather than the criminal put to shame. Thus the whole party of them is
full of vileness \footnote{[\emph{rhupou}].

Return to text} and
their spies \footnote{[\emph{kataskopoi}].

Return to text}, for
Bishops \footnote{[\emph{episkopoi}].

Return to text} they
are not, are the vilest of them all. And if any one among them desires
to become a Bishop, he is not told, \emph{a Bishop must be blameless}
[1 Tim. iii. 2.] \footnote{p. 221.

Return to text};
but only, ``Take up opinions contrary to Christ, and care not for
manners. This will be sufficient to obtain favour for you, and
friendship with the Emperor.'' Such is the character of those who
support the tenets of Arius. And they who are zealous for the truth,
however holy and pure they shew themselves, are yet, as I said before,
made criminals, whenever these men choose, and on whatever pretences
it may seem good to them to invent. The truth of this, as I before
remarked, you may clearly gather from their proceedings.

§. 4.

5. There was one Eustathius \footnote{p. 190, note A.

Return to text}, Bishop of Antioch, a Confessor, and sound in the Faith. This
man, because he was very zealous for the truth, and hated the Arian
heresy, and would not receive those who adopted its tenets, is falsely
accused before the Emperor Constantine, and a charge invented against
him, that he had insulted his mother \footnote{If the common slander of the day concerning
St. Helena was imputed to St. Eustathius, Constantine was likely to
feel it keenly. ``Stabulariam,'' says St. Ambrose, ``hanc primò fuisse
asserunt, sic cognitam Constantio.'' de Ob. Theod. 42. Stabularia, i.e.
an innkeeper; so Rahab is sometimes considered to be ``cauponaria sive
tabernaria et meretrix,''Cornel. à Lap. in Jos. ii. 1. [\emph{ex homilias
gunaikos ou semn\={e}s oude kata nomon sunelthou s\={e}s}].
Zosim. Hist. ii. p. 78.Constantinus ex concubinâ Helenâ procreatus.
Hieron. in Chron. Euseb. p. 773. (ed. Vallars.) Tillemont however
maintains, (Empereurs, t. 4. p. 613.) and Gibbon fully admits (Hist.
ch. 14. p. 190.) the legitimacy of Constantine. The latter adds, ``Entropius
(x. 2.) expresses in a few words the real truth, and the occasion of
the error, ex \emph{obscuriori matrimonio} ejus filius.''

Return to text}. And immediately he is driven into banishment, and a great
number of Presbyters and Deacons with him. And immediately after the
banishment of the Bishop, those whom he would not admit into the
clerical order on account of their impiety were not only received into
the Church by them, but were even appointed the greater part of them
to be Bishops, in order that they might have accomplices in their
impiety. Among these was Leontius the eunuch \footnote{[\emph{ho apokopos}]. pp. 208, 241, note A.

Return to text}, now of Antioch, and before him Stephanus, George of Laodicea,
and Theodosius who was of Tripolis, Eudoxius of Germanicia, and
Eustathius \footnote{p. 133.

Return to text} now of
Sebastia.

§. 5.

6. Did they then stop here? No. For Eutropius \footnote{p. 190, note A.

Return to text} who was Bishop of Adrianople, a good man, and excellent in all
respects, because he had often convicted Eusebius, and had advised them who came that way, not to comply with his impious
dictates, suffered the same treatment as Eustathius,—and was cast
out of his city and his Church. Basilina \footnote{Julian's mother.

Return to text} was the most active in the proceedings against him. And
Euphration of Balanea, Cymatius of Paltus, another Cymatius of Taradus,
Asclepas of Gaza, Cyrus of Berea \footnote{qu. Berrhœa?

Return to text} in Syria, Diodorus of Asia, Domnion of Sirmium, and Ellanicus
of Tripolis, were merely known to hate \footnote{21.
p. 217, r. 7.

Return to text} the heresy; and some of them on one pretence or another, some
without any, they removed under the authority of royal letters \footnote{22.
p. 221, r. 3.

Return to text}, drove them out of their cities, and appointed others whom
they knew to be impious men, to occupy the Churches in their stead.

§. 6.

7. Of Marcellus \footnote{23.
p. 52, note L.

Return to text} the Bishop of Galatia it is perhaps superfluous for me to
speak; for all men have heard how the Eusebians, who had been first
accused by him of impiety, brought a counter-accusation against him,
and caused him to be banished in his old age. He went up \footnote{24.
[\emph{anelth\={o}n}], vid. Acts xxi. 15. infr. pp. 239, r. 3. 242,
r. 4.

Return to text} to Rome, and there made his defence, and being required by
them, he offered a written declaration of his faith, of which the
Council of Sardica approved. But the Eusebians made no defence, nor,
when they were convicted of impiety out of their writings, were they
put to shame, but rather assumed greater boldness against all. For
they had interest with the Emperor through the women \footnote{25.
i.e. Constantia, Const.'s sister.

Return to text}, and were formidable to all men.

§. 7.

8. And I suppose no one is ignorant of the case
of Paul \footnote{26.
p. 191, r. 1.

Return to text}, Bishop of
Constantinople; for the more illustrious any city is, so much the more
that which takes place in it is not concealed. A charge was fabricated
against him also. For Macedonius his accuser, who has now become
Bishop in his stead, (I was present myself at the accusation,)
afterwards held communion with him, and was a Presbyter under Paul
himself. And yet when Eusebius with an evil eye \footnote{27.
[\emph{ep' ophthalmiai}], supr. p. 23.

Return to text} wished to seize upon the Bishopric of that city, (he had been
translated in the same manner from Berytus to Nicomedia,) the charge
was revived against Paul; and they did not give up their plot, but
persisted in the calumny. And he was banished first into Pontus by
Constantine, and a second time by Constantius he was sent bound with
iron chains to Singara in Mesopotamia, and from thence transferred to
Emesa, and a fourth time he was banished to Cucusus in
Cappadocia, near the deserts of mount Taurus; where, as those who were
with him have declared, he died by strangulation \footnote{28.
p. 191, r. 1.

Return to text} at their hands. And yet these men who never speak the truth,
though guilty of this, were not ashamed after his death to invent
another story, representing that he had died from disease; although
all who live in that place know the circumstances. And even Philagrius
\footnote{It is remarkable that this Philagrius, who has been so often mentioned
with dishonour in these Tracts of St. Athanasius, as an apostate and a
persecutor, vid. supr. pp. 5, 3l, \&c. is represented by St. Greg.
Naz. as very popular in Alexandria, and as on that account appointed
to the prefecture there a second time. He compares his entry into the
city on this occasion to that of St. Athan.'s after banishment. vid.
Greg. Orat. 21. 28. St. Athan. however wrote on the spot and at the
time, and there is nothing inconsistent in his being a popular
magistrate and an enemy of the Church.

Return to text} who was then
Deputy-Governor \footnote{Vicarius, i.e. ``vicarius Præfecti, agens vicem Præfecti;'' Gothofred
in Cod. Theod. i. tit. 6. vid. their office, \&c. drawn out at
length, ibid. t. 6. p. 334.

Return to text

} of those
parts, and represented all their proceedings in such manner as they
desired, was yet astonished at this; and being grieved perhaps that
another, and not himself, had done the evil deed, he informed Serapion
the Bishop as well as many other of our friends, that Paul was shut up
by them in a very confined and dark place, and left to perish of
hunger; and when after six days they went in and found him still
alive, they immediately set upon him, and strangled him.

9. This was the end of his life; and they said
that Philip who was Prefect was their agent in the perpetration of
this murder. Divine Justice however did not overlook this; for not a
year had past, when Philip was deprived of his office in great
disgrace, so that being reduced to a private station, he became the
mockery of those whom he least desired to be the witnesses of his
fall. For in extreme distress of mind, \emph{a fugitive and a vagabond}
[Gen. iv. 12.], like Cain \footnote{29.
supr. p. 161.

Return to text},
and expecting every day that some one would destroy him, far from his
country and his friends, he died, like one astounded at his
misfortunes, in a manner that he least desired. Moreover these men
spare not even after death those against whom they have invented
charges whilst living. They are so eager to shew themselves formidable
to all, that they banish the living, and shew no mercy on the dead;
but alone of all the world they manifest their hatred to them
that are departed, and conspire against their friends, truly inhuman
as they are,—and haters of that which is good, savage in temper
beyond mere enemies, in behalf of their impiety, who eagerly plot the
ruin of me and of all the rest, with no regard to truth, but by false
charges.

§. 8.

10. Perceiving this to be the case, the three
brothers, Constantine, Constantius, and Constans, caused us all after
the death of their father to return to our own country and Church; and
while they wrote letters concerning the rest to their respective
Churches, concerning Athanasius they wrote the following; which
likewise shews the violence of the whole proceedings, and proves the
murderous disposition, of the Eusebian party.

11. A copy of the Letter of Constantine Cæsar to
the people of the Catholic Church in the city of the Alexandrians

I suppose that it has not escaped the knowledge
of your pious minds \footnote{vid. Apol.
contr. Arian. §.
87. supr. p. 121.

Return to text},
\&c.

12. This is his letter; and what more credible
witness of their conspiracy could there be than this, who knowing
these circumstances has thus written of them?

\xsection{Chapter 2. First Arian persecution under Constantius}

§. 9.

1. THE
Eusebians however, seeing the declension of their heresy, wrote to
Rome, as well as to the Emperors Constantine and Constans, to accuse
Athanasius: but when the persons who were sent by Athanasius disproved
the statements which they had written, they were repulsed with
disgrace by the Emperors; and Julius, Bishop of Rome, wrote to say
that a Council ought to be held, wherever we should desire, in order
that they might exhibit the charges which they had to make, and might
also freely defend themselves concerning those things of which they
themselves were accused. The Presbyters also who were sent by them,
when they saw themselves making an exposure, requested that this might
be done. Whereupon these men, whose conduct is suspicious in all that
they do, when they see that they are not likely to get the better in
an Ecclesiastical trial, betake themselves to Constantius alone, and
there bewail themselves, as to the patron \footnote{Margin Notes

1. [\emph{prostat\={e}n}],
de Syn. §. 31. tr. p. 127.

Return to text} of their heresy. ``Spare,'' they say, ``the heresy; you see that
all men have withdrawn from us; and very few of us are now left. Begin
to persecute, for we are being deserted even of those few, and are
left destitute. Those persons whom we forced over to our side, when
these men were banished, they now by their return have persuaded again
to take part against us. Write letters therefore against them all, and
send out Philagrius a second time \footnote{p. 224, note B.

Return to text} as Prefect of Egypt, for he is able to carry on a persecution
favourably for us, as he has already shewn upon trial, and the more
so, as he is an apostate \footnote{[\emph{parabat\={e}s}], p. 5.

Return to text}.
Send also Gregory as Bishop to Alexandria, for he too is able to
strengthen our heresy.''

§. 10.

2. Accordingly Constantius at once writes
letters, and commences a persecution against all, and sends Philagrius
as Prefect with one Arsacius an eunuch; he sends also Gregory with a
military force. And the same consequences followed as before \footnote{upon the Commission, p. 33.

Return to text}. For gathering together \footnote{pp. 5, \&c. 53, \&c.

Return to text} a multitude of herdsmen and shepherds, and other dissolute
youths belonging to the town, armed with swords and clubs, they
attacked in a body the Church which is called the Church of Cyrinus \footnote{Quirinus.

Return to text}; and some they slew, some they trampled under foot, others they
beat with stripes and cast into prison or banished. They haled away
many women also, and dragged them openly into the court, and insulted
them, dragging them by the hair. Some they proscribed; from some they
took away their bread \footnote{vid. infr. §. 63.

Return to text}
for no other reason, but that they might be induced to join the
Arians, and receive Gregory who had been sent by the Emperor.

§. 11.

3. Athanasius however, before these things
happened, at the first report of their proceedings, sailed to Rome,
knowing the rage of the heretics, and for the purpose of having the
Council held as had been determined. And Julius wrote letters to them,
and sent the Presbyters Elpidius and Philoxenus \footnote{p. 39.

Return to text}, appointing a day \footnote{[\emph{prothesmian}], supr. p. 45,
r. 4.

Return to text},
and saying, that either they must come, or consider themselves as
altogether suspected persons. But as soon as the Eusebians heard that
the trial was to be an Ecclesiastical one, at which no Count would be
present \footnote{pp. 25, 249, r. 8.

Return to text}, nor
soldiers stationed before the doors, and that the proceedings would
not be regulated by royal order, (for they have always depended upon
these things to support them against the Bishops, and without them
they have no boldness even to speak;) they were so alarmed that they
detained the Presbyters till after the appointed time, and pretended
this indecent excuse, that they were not able to come now on account
of the war which was begun by the Persians \footnote{p. 46, r. 1.

Return to text}. But this was not the true cause of their delay, but the fears
of their own consciences. For what have Bishops to do with war? Or if
they were unable on account of the Persians to come to Rome, although
it is at a distance and beyond sea, why did they like lions \footnote{1 Pet. v. 8.

Return to text} traverse \footnote{[\emph{peri\={e}rchonto}], p.
126. infr. p. 235, r. 2. Orat. i. §. 22. iii. fin. supr.

Return to text} the
parts of the East and those which are near the Persians, seeking who
was opposed to them, that they might falsely accuse and banish them?

§. 12.

4. However when they had dismissed the Presbyters
with this improbable excuse, they said to one another, ``Since we are
unable to get the advantage in an Ecclesiastical trial, let us
exhibit our usual audacity.'' Accordingly they write to Philagrius, and
cause him after a while to go out with Gregory into Egypt. Whereupon
the Bishops are severely scourged and cast into chains \footnote{pp. 51, 53.

Return to text}; Sarapammon, for instance, Bishop and Confessor, they drive
into banishment; Potammon, Bishop and Confessor, who had also lost an
eye in the persecution, they beat with stripes on the back so cruelly,
that he appeared to be dead before they came to an end. In which
condition he was cast aside, and hardly after some hours, being
carefully attended and fanned, he revived, God granting him his life;
but a short time after he died of the sufferings caused by the
stripes, and attained in Christ to the glory of a second martyrdom.
And besides these, how many monks were scourged, while Gregory sat by
with Balacius the Duke! how many Bishops were wounded! how many
virgins were beaten!

§. 13.

5. After this the wretched Gregory called upon
all men to have communion with him \footnote{p. 8.

Return to text}; but if thou didst demand of them communion, they were not
worthy of stripes: and if thou did scourge them as if evil persons,
why didst thou ask it of them as if holy? But he had no other end in
view, except to fulfil the designs of them that sent him, and to
establish the heresy. Wherefore he became in his folly a murderer and
an executioner \footnote{[\emph{d\={e}mios}], pp. 133
fin. 247, r. 2.

Return to text},
injurious, crafty, and profane; in one word, an enemy of Christ. He so
cruelly persecuted the Bishop's aunt, that even when she died he would
not suffer her to be buried \footnote{p. 178, §. 27 fin.

Return to text}. And this would have been her lot; she would have been cast
away without burial, had not they who attended on the corpse carried
her out as one of their own kindred. Thus even in such things he
shewed his profane temper. And again when the widows and other
mendicants \footnote{[\emph{anexod\={o}n}],
vid. infr. §. 60. Tillemont translates it, prisoners. Montfaucon has
been here followed; vid. Collect. Nov. t. 2. p. xliii.

Return to text} had received
alms, he commanded what had been given them to be seized, and the
vessels in which they carried their oil and wine to be broken, that he
might not only show impiety by robbery, but in his deeds dishonour the
Lord; from whom very shortly \footnote{[\emph{hoson oudep\={o}}], vid.
p. 245, r. 4. George was pulled to pieces by the populace, A.D.
362. This was written A.D. 358,
or later.

Return to text}
he will hear those words, \emph{Inasmuch
as thou hast dishonoured these, thou hast dishonoured Me} [vid.
Mat. xxv. 45.].

6. And many other things he did, which exceed the
power of language to describe, and which whoever should hear would
think to be incredible. And the reason why he acted thus was, because
he had not received his ordination according to ecclesiastical rule,
nor had been called to be a Bishop by apostolical tradition \footnote{He had neither apostolical calling, nor canonical ordination, for he
was a layman, nominated to his see by the Emperor, and that, when
there was a lawful occupant, and consecrated by heretics. ``Tradition''
and ``Canon'' seem used nearly as synonymous. p. 249, r. 6.

Return to text}; but had been sent out from court with military power and pomp,
as one entrusted with a secular government. Wherefore he boasted
rather to be the friend of Governors, than of Bishops and Monks.
Whenever therefore Father Antony wrote to him from the mountains, as
godliness is an abomination to a sinner, so he abhorred the letters of
the Saint. But whenever the Emperor, or a General, or other
magistrate, sent him a letter, he was as much overjoyed as those in
the Proverbs, of whom the Word has said indignantly, \emph{Woe
unto them who leave the paths of uprightness; who rejoice to do evil,
and delight in the frowardness of the wicked} [Prov. ii. 13, 14.
Sept.]. And so he honoured with presents the bearers of these letters;
but once when Antony wrote to him he caused Duke Balacius to spit upon
the letter, and to cast it from him. But Divine Justice did not
overlook this; for no long time after, when the Duke was on horseback,
and on his way to the first halt \footnote{[\emph{mon\={e}n}]. vid. supr.
p. 50, note H. This halt or station which lay up the Nile was called
Cereu, (Athan. V. Ant. §. 86.) or Chæreus, or the land or property
of Chæreus, vid. Naz. Orat. 21, 29. who says it was the place where
the people met Athanasius on his return from exile on Constantius's
death.

Return to text

}, the horse turned his head, and biting him on the thigh, threw
him off; and after three days he died.

\xsection{Chapter 3. Restoration of the Catholics on the Council of Sardica}

§. 15.

1. WHILE
they were proceeding in like measures towards all, at Rome about fifty
Bishops assembled \footnote{Margin Notes

1. p. 14, note B.

Return to text}, and
denounced the Eusebians, as persons suspected, afraid to come, and
also condemned as unworthy of credit the written statement they had
sent; but us they received, and gladly embraced \footnote{and p. 232, r. 3. [\emph{\={e}gap\={e}san}]. vid. [\emph{agap\={e}n}],
p. 39, r. 5. p. 60, r. 2.

Return to text} our communion. While these things were taking place, a report
of the Council held at Rome, and of the proceedings against the
Churches at Alexandria, and through all the East, came to the hearing
of the Emperor Constans \footnote{p. 158, note G.

Return to text}.
He writes to his brother Constantius, and immediately they both
determine \footnote{infr. §. 50.

Return to text} that a
Council shall be called, and matters be brought to a settlement, so
that those who had been injured may be released from further
suffering, and the injurious be no longer able to perpetrate such
outrages. Accordingly there assemble at the city of Sardica both from
the East and West to the number of one hundred and seventy Bishops \footnote{vid. supr. p. 78, note O. Montfaucon argues in
the Addenda in his Præf. p. xxxiv. from the subscriptions in St.
Hilary, p. 1292. that the Bishops whose signatures occur without
provinces annexed, (supr. p. 76.) were the Bishops present—whereas
those who were absent signed with a mention of their provinces.
Accordingly he considers the number of orthodox to be 86; to which if
we add the 70 or 80 Eusebians, we approximate to the number170. If the
parties were so nearly matched, it is very remarkable that the
Eusebians should withdraw. But they had the Pope, Athanasius, and
Hosius against them.

Return to text}, more or less; those who came from the West were Bishops only,
having Hosius for their father \footnote{vid. p. 158, and [\emph{ho pat\={e}r t\={o}n episkop\={o}n}],
infr. pp. 255, 256.

Return to text}, but those from the East brought with them instructors of youth
and advocates, Count Musonianus \footnote{p. 59, notes A, B.

Return to text}, and Hesychius \footnote{p. 59, notes A, B.

Return to text} the Castrensian;
on whose account they came with great alacrity, thinking that every
thing would be again managed by their authority. For thus by means of
these persons they have always shown themselves formidable to any whom
they wished to intimidate, and have prosecuted their designs against
whomsoever they chose. But when they arrived and saw that the cause was to be conducted as simply an ecclesiastical one, without the
interference of the Count or of soldiers; when they saw the accusers
who came from every church and city, and the evidence which was
brought against them, when they saw the venerable Bishops Arius and
Asterius \footnote{p. 233, note C.

Return to text}, who came up
in their company, withdrawing from them and siding with us, and giving
an account of their profligate conduct; as their whole proceedings had
been suspicious, so now they fear the consequences of a trial, lest
they should be convicted by us of being false informers, and it should
be discovered by those whom they produced in the character of
accusers, that they had themselves suggested all they were to say, and
were the contrivers of the plot.

2. Perceiving this to be the case, although they
had come with great zeal, as thinking that we should be afraid to meet
them, yet now when they saw our alacrity, they shut themselves up in
the Palace \footnote{The word Palatium sometimes stands for the space or limits set apart
in cities for the Emperor, Cod. Theod. xv. 1. 47. sometimes for the
buildings upon it, ibid. vii. 10, 2. which wore one of the four public
works mentioned in the Laws. ibid. xv. 1. 35. and 36. None but great
officers of state were admitted into it. xv. 1. 47. Even the judges
might not lodge in it, except there was no Prætorium, vii. 10. 2.
Gothofr. in vii. 10, 1. enumerates (with references) the Palatia in
Antioch, Daphne, Constantinople, Heraclea, Milan, Treves, \&c. It
was a great mark then of imperial favour that the Eusebians were
accommodated in the Palatium at Sardica.

Return to text}, (for they had
their abode there,) and proceeded to confer with one another in the
following manner, ``We came hither for one result; and we see in
prospect another; we arrived in company with Counts, and the trial is
proceeding without them. We are certainly condemned. You all know the
orders that have been given. The Athanasians have the reports of the
proceedings in the Mareotis \footnote{pp. 47, r. 2. 116, §. 83.

Return to text},
by which he is cleared, and we are covered with disgrace. Why then do
we delay? why are we so slow? Let us invent some excuse and be gone,
or we shall be condemned if we remain. It is better to suffer the
shame of fleeing, than the disgrace of being convicted as false
accusers. If we flee, we shall find some means of defending our
heresy; and even if they condemn us for our flight, still we have the
Emperor as our patron \footnote{p. 226, r. 1.

Return to text},
who will not suffer the people to expel us from the Churches.''

§. 16.

3. They reasoned with themselves in this manner:
and Hosius and all the other Bishops repeatedly signified to
them the alacrity of the Athanasians, saying, 'They are ready with
their defence, and pledge themselves to prove you false accusers.'
They said also, 'If you fear the trial, why did you come to meet us?
either you ought not to have come, or now that you have come, not to
flee.' When they heard this, being still more alarmed, they had
recourse to an excuse even more indecent than that they pretended at
Antioch \footnote{p. 227, r. 8.

Return to text}, viz. that
they betook themselves to flight because the Emperor had written to
them the news of his victory over the Persians. And this excuse they
were not ashamed to send by Eustathius a Presbyter of the Sardican
Church. Nevertheless their flight did not succeed according to their
wishes; for immediately the holy Council, of which the great Hosius
was president, wrote to them plainly, saying, ``Either come forward and
answer the charges which are brought against you, for the false
accusations which, you have made against others, or know that the
Council will condemn you as guilty, and declare Athanasius and his
friends free and clear from all blame.'' Whereupon they were rather
impelled to flight by the alarms of conscience, than to compliance
with the proposals of the letter: for when they saw those who had been
injured by them, they did not even turn their faces to listen to their
words, but fled with greater speed.

§.17.

4. Under these disgraceful and unseemly
circumstances their flight took place. And the holy Council, which had
been assembled out of more than five and thirty provinces \footnote{pp. 14, 60, r. 9.

Return to text}, perceiving the malice of the Arians, admitted the Athanasian
party to answer to the charges which they had brought against them,
and to declare the sufferings which they had undergone. And when they
had thus made their defence, as we said before, they approved and so
highly admired their conduct, that they gladly embraced \footnote{p. 230, r. 2.

Return to text} our communion, and wrote letters to all quarters, to the
diocese of each, and especially to Alexandria, and Egypt, and the
Libyas, declaring Athanasius and his friends to be innocent, and free
from all blame, and their opponents to be calumniators, evil-doers,
and every thing rather than Christians \footnote{p. 208, note B.

Return to text}. Accordingly they dismissed them in peace; but deposed \footnote{p. 75.

Return to text} Stephanus and Menophantus, Acacius and George of Laodicea,
Ursacius and Valens, Theodorus and Narcissus. And against
Gregory who had been sent to Alexandria by the Emperor, they put forth
a proclamation to the effect that he had never been made a Bishop, and
that he ought not to be called a Christian \footnote{p. 68.

Return to text}. They therefore declared the ordinations which he professed to
have conferred to be void, and commanded that they should not be even
named in the Church, on account of their novel and illegal nature.
Thus Athanasius and his friends were dismissed in peace (the letters
concerning them are inserted at the end on account of their length \footnote{not found there, but in Apol. contr. Ar. supr. p. 60-76.

Return to text}); and the Council was dissolved.

§. 18.

5. But the deposed persons, who ought now to have
remained quiet, with those who had separated after so disgraceful a
flight, were guilty of such conduct, that their former proceedings
appear trifling in comparison of these. For when the people of
Adrianople would not have communion with them, as men who had fled
from the Council, and had been declared guilty, they carried their
complaints to the Emperor Constantius, and succeeded in causing ten of
the laity to be beheaded, belonging to the Manufactory of arms \footnote{de Fabricis, vid. Gothofr. in Cod. Theod. x. 21.

Return to text} there, Philagrius, who was there again as Count, assisting
their designs in this matter also. The tombs of these persons, which
we have seen in passing by, are in front of the city.

6. Then as if they had been quite successful,
because they had fled lest they should be convicted of false
accusation, they prevailed with the Emperor to command whatsoever they
wished to be done. Thus they caused two Presbyters and three Deacons
to be banished from Alexandria into Armenia. As to Arius and Asterius,
the former Bishop of Petræ \footnote{This seems to be the famous Petra, the capital of Edom, which has
lately been discovered by travellers Edom being formed into the
Province Tertia Palestina, or at least called Palestine, about or soon
after the time of St. Athanasius. But a difficulty arises from a
passage in the Tomus ad Antioch. §. 10. where Asterius, the
fellow-sufferer with Arius, (or Macarius, as he is called supr. p.
74.) is said to be Bishop of Petræ of Arabia, as if the Petræ of
Palestine were distinct from it. Reland observes, Palestine, p. 928,
(in answer to Cellarius, who considered in consequence that there were
two Petræ, Le Quien Orien. Christ. t. 3. pp. 665. 666.) that as there
is already one error of text in this passage (as it appears), of Arius
for Macarius, so the word Petra may have fallen into the wrong place,
instead of ``the one of Palestine, the other of Petra of Arabia;'' or
that Petra may be a marginal note, which has been incorporated with
the text in the wrong place, as is confirmed by the run of the passage
supr. p. 74. and by passages to which he refers in St. Hilary. He
observes moreover, on the improbability of the silence of Eusebius,
St. Jerome, the acts of Councils, and ancient Notitiæ, supposing
there were two Petræ. Dr. Robinson, who in his recent elaborate and
useful work on Palestine, referring to Beland, observes, that ``the
passage [in the text] is usually referred to as contained in the 'Epist.
ad solitariam Vitam agentes,' though in the Benedictine Edition at
least, it is found, not in that Epistle, but in the Historia Arianor.
§. 18.'' Palest. vol. 2. p. 655. But these were but two titles of the
same work, till Montfaucon more correctly confined the former to the
short introduction to a lost dogmatic work which is prefixed to the
History, (vid. supr. pp. 210. 215, and note of Amanuensis in Calc.
Hist.) yet even Montf. calls the History, ``Ep. ad Sol.'' Notes, tom. 1.
p.150, 151. correcting himself in Præf. xxxiii. And it is called ``Epistle
to the Solitaries'' by Ceillier even since Montfaucon.

Return to text

}
in Palestine, the latter Bishop in Arabia, who had withdrawn
from their party, they not only banished into upper Libya, but also
caused them to be treated with insult; §. 19. and as to Lucius \footnote{p. 71, note F. p. 190.

Return to text}, Bishop of Adrianople, when they saw that he used great
boldness of speech against them, and exposed their impiety, they
again, as they had done before, caused him to be bound with iron
chains on the neck and hands, and so drove him into banishment, where
he died, as they know. And Diodorus the Bishop \footnote{of Tenedos, vid. pp. 76, 223.

Return to text} they transported from his place; but against Olympius of Æni,
and Theodulus of Trajanople \footnote{p. 71, note G.

Return to text}, both Bishops of Thrace, good and orthodox men, when they
perceived their hatred of the heresy, they brought false charges.

7. This the party of Eusebius had done first of
all, and the Emperor Constantius wrote letters on the subject; and
next these men \footnote{Acacians, \&c. p. 241, r. 1. vol. 8. p. 7, note p.

Return to text}
revived the accusation. The purport of the letters was, that they
should not only be expelled from their cities and churches, but should
suffer capital punishment, wherever they were discovered. However
surprising this conduct may be, it is only in accordance with their
principles; for as being instructed by the Eusebians in such
proceedings, and as heirs of their impiety and evil principles, they
wished to shew themselves formidable at Alexandria, as their fathers
had done in Thrace. They caused an order to be written, that the ports
and gates of the cities should be watched, lest availing themselves of
the permission granted by the Council, the banished persons should
return to their churches. They also cause orders to be sent to the
magistrates at Alexandria, respecting Athanasius and certain
Presbyters, named therein, that if either the Bishop \footnote{this accounts for Ath.'s caution, pp. 79, 80, and 236, r. 5.

Return to text}, or any of the others, should be found coming to the city or
its borders, the magistrate should have power to behead those who were
so discovered. Thus this new Jewish heresy \footnote{vol. 8. p. 4.

Return to text} does not only deny the Lord, but has also learnt to commit
murder.

§. 20.

8. Yet even after this they did not rest; but as
the father of their heresy goeth about \footnote{[\emph{perierchetai}], vid. supr. p. 227.

Return to text} like a lion, seeking whom he may devour, so these obtaining
the use of the public posts \footnote{p. 100, note Y.

Return to text} went about, and whenever they found any that reproached them
with their flight, and that hated the Arian heresy, they scourged
them, cast them into chains, and caused them to be banished from their
country; and they rendered themselves so formidable, as to induce many
to dissemble, many to fly into the deserts, rather than willingly even
to have any dealings with them. Such were the enormities which their
madness prompted them to commit after their flight.

9. Moreover they perpetrate another outrageous
act, which is indeed in accordance with the character of their heresy,
but is such as was never heard of before, nor is likely soon to take
place again, even among the more dissolute of the Gentiles \footnote{pp. 9, 195, 235. infr. §. 64.

Return to text}, much less among Christians. The holy Council had sent as
Legates the Bishops Vincentius \footnote{p. 157, note C.

Return to text} of Capua, (this is the Metropolis of Campania,) and Euphrates
of Agrippina \footnote{Cologne.

Return to text}, (this
is the Metropolis of Upper Gaul,) that they might obtain the Emperor's
consent to the decision of the Council, that the Bishops should return
to their Churches, inasmuch as he was the author of their expulsion.
The most religious Constans had also written to his brother \footnote{infr. §. 50.

Return to text}, and supported the cause of the Bishops. But these admirable
men, who are equal to any act of audacity, when they saw the two
Legates at Antioch, consulted together and formed a plot, which
Stephanus \footnote{Bishop of Antioch, p. 60, r. 6.

Return to text} undertook
by himself to execute, as being a suitable instrument for such
purposes. Accordingly they hire a common harlot, even at the season of
the most holy Easter, and stripping her introduce her by night into
the apartment of the Bishop Euphrates. The harlot who thought that it
was a young man who had sent to invite her, at first willingly
accompanied them; but when they thrust her in, and she saw the man
asleep and unconscious of what was going on, and when presently she
distinguished his features, and beheld the face of an old man, and the
figure of a Bishop, she immediately cried aloud, and declared that
violence was used towards her. They desired her to be silent,
and to lay a false charge against the Bishop; and so when it was day,
the matter was noised abroad, and all the city ran together; and those
who came from the Palace \footnote{p. 231, note B.

Return to text}
were in great commotion, wondering at the report which had been spread
abroad, and demanding that it should not be passed by in silence. An
enquiry therefore was made, and her master \footnote{[\emph{heta rotrophos}].

Return to text} gave information concerning those who came to fetch the
harlot, and these informed against Stephanus; for they were his
Clergy. Stephanus therefore is deposed, and Leontius the eunuch \footnote{[\emph{ho akokopos}], p. 241, note A.

Return to text} appointed in his place, only that the Arian heresy may not
want a supporter.

§. 21.

10. And now the Emperor Constantius, feeling some
compunctions, returned to a right \footnote{[\emph{eis heauton}].

Return to text} mind; and concluding from their conduct towards Euphrates,
that their attacks upon the others were of the same kind, he gives
orders that the Presbyters and Deacons who had been banished from
Alexandria into Armenia should immediately be released. He also writes
publicly to Alexandria commanding that the clergy and laity who were
friends of Athanasius should suffer no further persecution. And when
Gregory died about ten months after, he sends for Athanasius with
every mark of honour, writing to him no less than three times a very
friendly letter \footnote{pp. 79, 80.

Return to text}, in
which he exhorted him to take courage and come. He sends also a
Presbyter and a Deacon, that he may be still further encouraged to
return; for he thought that, through alarm at what had taken place
before, I \footnote{vid. p. 219.

Return to text} did not
care to return. Moreover he writes to his brother Constans, that he
also would exhort me to return. And he affirmed that he had been
expecting Athanasius a whole year, and that he would not permit any
change to be made, or any ordination to take place, as he was
preserving the Churches for Athanasius their Bishop.

§. 22.

11. When therefore he wrote in this strain, and
encouraged him by means of many, (for he caused Polemius, Datianus,
Bardion, Thalassus \footnote{p. 156, r. 2.

Return to text},
Taurus \footnote{At Ariminum.

Return to text}, and
Florentius, his Counts, in whom Athanasius could best confide, to
write also;) Athanasius committing the whole matter to God, who had
stirred the conscience of Constantius to do this, came with his
friends to him; and he gave him a favourable audience, and sent him
away to go to his country and his Churches, writing at the same
time to the magistrates in the several places, that whereas he had
before commanded the ways to be guarded, they should now grant him a
free passage. Then when the Bishop complained of the sufferings he had
undergone, and of the letters which the Emperor had written against
him, and besought him that the false accusations against him might not
be revived by his enemies after his departure, saying, ``If you please,
summon these persons; for as far as we are concerned they are at
liberty to stand forth, and we will expose their conduct;'' he would
not do this, but commanded that whatever had been before slanderously
written against him should all be destroyed and obliterated, affirming
that he would never again listen to any such accusations, and that his
purpose was fixed and unalterable. This he did not simply say, but
sealed his words with an oath, calling upon God to be witness of them.
And so encouraging him with many other words, and desiring him to be
of good courage, he sends the following letters to the Bishops and
Magistrates.

§. 23.

12. Constantius Augustus, the Great, the
Conqueror, to the Bishops and Clergy of the Catholic Church.

The most Reverend Athanasius has not been
deserted by the grace of God \footnote{vid. Apol. contr. Arian. §. 54. supr. p. 82.

Return to text}, \&c.

Another
Letter.

From Constantius to the people
of Alexandria.

Desiring as we do your welfare in all respects \footnote{vid. Apol. contr. Arian. §. 55. supr. p. 83.

Return to text}, \&c.

Another Letter.

Constantius Augustus, the
Conqueror, to Nestorius,

Prefect of Egypt.

It is well known that an order was heretofore
given by us, and that certain documents are to be found prejudicial to
the character of the most reverend Bishop Athanasius; and that these
exist among the Orders  of your worship. Now we desire your
Prudence, of which we have good proof, to transmit to our Court, in
compliance with this our order, all the letters respecting the
fore-mentioned person, which are found in your Order-Book \footnote{or Acta Publica, vid. supr. p. 84.

Return to text}.

§. 24.

13. The following is the letter which he wrote
after the death of the blessed Constans. It was written in Latin, and
is here translated into Greek \footnote{another translation, p. 174.

Return to text}.

Constantius Augustus, the
Conqueror, to Athanasius.

It is not unknown to your Prudence, that it was
my constant prayer, that prosperity might attend my late brother
Constans in all his undertakings; and your wisdom may therefore
imagine how greatly I was afflicted when I learnt that he had been
taken off by most unhallowed hands. Now whereas there are certain
persons who at the present time endeavour to alarm you by that so
melancholy event, I have therefore thought it right to address this
letter to your Constancy, to exhort you that, as becomes a Bishop, you
would teach the people those things which pertain to the service of
God, and that, as you are accustomed to do, you would employ your time
in prayers together with them, and not give credit to vain rumours,
whatever they may be. For our fixed determination is, that you should
continue, agreeably to our desire, to perform the office of a Bishop
in your own place. May Divine Providence preserve you, most beloved
Father \footnote{[\emph{goneu}].

Return to text}, many years.

§. 25.

14. Under these circumstances, when they had at
length taken their leave, and commenced their journey, those who were
friendly to them rejoiced to see their friend; but of the other party,
some were confounded at the sight of him; others not having the
confidence to appear, hid themselves; and others repented of what they
had written against the Bishop. Thus all the Bishops of Palestine,
except some two or three, and those men of suspected character, so
willingly received Athanasius, and embraced communion with him \footnote{p. 85.

Return to text}, that they wrote to excuse themselves, on the ground that in
what they had formerly written, they had acted, not according to their
own wishes \footnote{[\emph{kata proairesis}].

Return to text}, but by
compulsion. Of the Bishops of Egypt and the Libyan provinces, of the
laity both of those countries and of Alexandria, it is superfluous for
me to speak. They all ran together, and were possessed with
unspeakable delight, that they had not only received their friends
alive contrary to their hopes; but that they were also delivered
from the heretics who were as tyrants and as raging dogs towards them.
Accordingly great was their joy, the people in the congregations
encouraging one another in virtue. How many unmarried women, who were
before ready to enter upon marriage, now remained virgins to Christ!
How many young men, seeing the examples of others, embraced the
monastic life! How many fathers persuaded their children, and how many
were urged by their children, to submit themselves to Christian
discipline \footnote{[\emph{ask\={e}se\={o}s}]. vid. p. 202, r. 2.

Return to text}! how many
wives persuaded their husbands, and how many were persuaded by their
husbands, to give themselves to prayer, as the Apostle has spoken! How
many widows and how many orphans, who were before hungry and naked,
now through the great zeal of the people, were no longer hungry, and
went forth clothed! In a word, so great was their emulation in virtue,
that you would have thought every family and every house a Church, by
reason of the goodness of its inmates, and the prayers which were
offered to God. And in the Churches there was a profound and wonderful
peace, while the Bishops wrote from all quarters, and received from
Athanasius the customary letters of peace.

§. 26.

15. Moreover Ursacius and Valens, as if suffering
the scourge of conscience, came to another mind, and wrote to the
Bishop himself a friendly and peaceable letter \footnote{p. 86, note Q.

Return to text}, although they had received no communication from him. And
going up \footnote{[\emph{anelthontes}], p. 26, r. 2. 39. and p. 242, r. 4.

Return to text} to Rome
they repented, and confessed that all their proceedings and assertions
against him were founded in falsehood and mere calumny. And they not
only voluntarily did this, but also anathematized the Arian heresy,
and presented a written declaration of their repentance, addressing to
the Bishop Julius the following letter in Latin, which has been
translated into Greek. The Latin copy was sent to us by Paul \footnote{Paulinus, supr. p. 86. Paulinus? p. 78.

Return to text} Bishop of Tibur.

Translation from the Latin.

Ursacius and Valens to my \footnote{[\emph{kuri\={o}i mou}], supr. p. 113.

Return to text} Lord the most blessed Pope Julius.

Whereas it is well known that we \footnote{vid. Apol. contr. Ar. §. 58. supr. p. 86.

Return to text}, \&c.

Translation from the Latin.

The Bishops Ursacius and Valens to
my \footnote{[\emph{kuri\={o}i mou}], supr. p. 113.

Return to text} Lord and Brother,

the Bishop Athanasius.

Having an opportunity of sending \footnote{vid. Apol. contr. Ar. §. 58. supr. p. 87.

Return to text}, \&c.

After writing these, they also subscribed the
letters of peace which were presented to them by Peter and Irenæus,
Presbyters of Athanasius, and by Ammonius a layman, who were passing
that way, although Athanasius had sent no communication to them by
these persons.

§. 27.

16. Now who was not filled with admiration at
witnessing these things, and the great peace that prevailed in the
Churches? who did not rejoice to see the concord of so many Bishops?
who did not glorify the Lord, beholding the delight of the people in
their assemblies? How many enemies repented! How many excused
themselves who had formerly accused him falsely! How many who formerly
hated him, now shewed affection for him! How many of those who had
written against him, recanted \footnote{[\emph{palin\={o}dian \={e}san}].

Return to text} their assertions! Many also who had sided with the Arians, not
through choice but by necessity, came by night and excused themselves.
They anathematized the heresy, and besought him to pardon them,
because, although through the plots and calumnies of these men they
appeared bodily on their side, yet in their hearts they held communion
with Athanasius, and were always with him. Believe me, this is true \footnote{pp. 158, 216.

Return to text}.

\xsection{Chapter 4. Second Arian persecution under Constantius}

§. 28.

1. BUT
the inheritors \footnote{Margin Notes

1. p. 234, r. 4.

Return to text} of the
opinions and impiety of the Eusebians, the eunuch Leontius \footnote{Various writers have treated on the subject of
that great scandal of the early centuries, the [\emph{suneisaktai}].
The most charitable account of it is to be found in the unprotected
state of women dedicated to a single life when or where Convents did
not exist. ``She says that she has no protector, husband, guardian,
father, nay, nor brother,'' \&c. Chrysost. ap. Basn. Dissert. vii.
19. ad Ann. Eccles. t. 2. And the example of the Holy and Blessed
Virgin was adduced, whom our Lord consigned to the care of St. John,
Epiph. Hær. 78. 11. which the Nicene Council so far acknowledged that
it dispensed with its prohibition in favour of mother, sister, aunt,
or other person, to whom no suspicion could attach. Nay, even in the
case of the atrocious extravagance, which St. Cyprian reprobates, Ep.
62. (ed. Ben.) and which in a still more perverted and shocking form
is spoken of in the text, it must be recollected that it was not
unknown to primitive times for husband and wife to vow continency and
yet to cohabit. Theodoret gives an instance in which a youth persuades
his bride, [\emph{en aut\={e}i t\={e}i pastadi, t\={e}i pr\={o}t\={e}i
t\={o}n gam\={o}n h\={e}merai}]. Hist. iv. 12. Another is
the instance so beautifully related by St. Gregory of Tours, in which
the bride persuades her husband; ``puella, graviter contristata, aversa
ad parietem, amarissime flebat,'' till ``tunc ille, armatus crucis
vexillo, ait, Faciam quæ hortaris, et datis inter se dextris,
quieverunt.'' He adds, ``Multos postea in uno strato recumbentes annos,
vixerunt cum castitate laudabili.'' Hist. Franc. i. 42. What was found
possible in the married, others had the indecency and wildness to
attempt in the single state. On the [\emph{suneisaktai}],vid. Mosheim
de Rebus Ante Const. p. 599. Routh, Reliqu. Sacr. t. 2. p. 506. t. 3.
p. 445. Basnag. Diss. vii. 19. in Ann. Eccles. t. 2. Muratori Anecdot.
Græc. p. 218. Dodwell, Dissert. Cyprian. iii. Bevereg. in Can. Nic.
3. Suicer. Thesaur. in voc. \&c. \&c. it is conjectured by
Beveridge, Dodwell, Van Espen, \&c. that Leontius gave occasion to
the first Canon of the Nicene Council, [\emph{peri t\={o}n solm\={o}nt\={o}n
heautous ektemnein}].

Return to text}, who ought not to remain in communion even as a layman \footnote{Can. Ap. 17. but vid. Morin. de Pæn. p. 185.

Return to text}, because he mutilated himself that he might henceforward be at
liberty to sleep with one Eustolium \footnote{p. 208.

Return to text}, who is a wife as far as he is concerned, but is called a
virgin; and George and Acacius, and Theodorus, and Narcissus, who were
deposed by the Council; when they heard and saw these things, were
greatly ashamed. And when they perceived the unanimity and peace that
existed between Athanasius and the Bishops; (they were more than four
hundred \footnote{after Sardica, vid. p. 78, note O.

Return to text}, from great
Rome, and all Italy, from Calabria, Apulia, Campania, Bruttia, Sicily,
Sardinia, Corsica, and time whole of Africa; and those from
Gaul, Britain, and Spain, with the great Confessor Hosius; and also
those from Pannonia, Noricim, Siscia, Dalmatia, Dardana, Dacia, Mysia,
Macedonia, Thessaly, and all Achaia, and from Crete, Cyprus, and Lycia,
with most of those from Palestine, Isauria, Egypt, the Thebais, the
whole of Libya, and Pentapolis;) when I say they perceived these
things, they were possessed with envy and fear; with envy, on account
of the communion of so many together; and with fear, lest those who
had been entrapped by them should be brought over by the unanimity of
so great a number, and henceforth their heresy should be triumphantly
exposed, and every where proscribed.

§. 29.

2. First of all they persuade Ursacius and Valens
to change sides again, and like dogs to return to their own vomit
[vid. 2 Pet. ii. 22.], and like swine to wallow again in the former
mire of their impiety; and they make this excuse for their
retractation, that they did it through fear of the most religious
Constans. And yet even had there been cause for fear, yet if they had
confidence in what they had done, they ought not to have become
traitors to their friends. But when there was no cause for fear, and
yet they were guilty of a lie, are they not deserving of utter
condemnation? For no soldier was present, no Palatine \footnote{p. 171, r. 1.

Return to text} or Notary \footnote{p. 173, note S.

Return to text} had
been sent, as they now send them, nor yet was the Emperor there, nor
had they been summoned \footnote{[\emph{kl\={e}thentes}], p. 212, r. 2.

Return to text}
by any one, when they wrote their recantation. But they voluntarily
went up \footnote{p. 223, r. 6.

Return to text} to Rome, and of
their own accord recanted and wrote it down in the Church, where there
was no fear from without, where the only fear is the fear of God, and
where every one has liberty of conscience \footnote{infr. p. 245, note B.

Return to text}. And yet although they have a second time become Arians, and
then have devised this indecent excuse for their conduct, they are
still without shame.

§. 30.

3. In the next place they went in a body to the
Emperor Constantius, and besought him, saying, ``When we first made our
request to you, we were not believed; for we told you, when you sent
for Athanasius, that by inviting him to come forward, you were
expelling our heresy. For he has been opposed to it from the very
first, and never ceases to anathematize it. He has already written
letters against us into all of the world, and the majority of
men have embraced communion with him; and even of those who seemed to
be on our side, some have been gained over by him, and others are
likely to be. And we are left alone, so that the fear is, lest the
character of our heresy become known, and henceforth both we and you
gain the name of heretics. And if this come to pass, you must take
care that we be not classed with the Manichæns. Therefore begin again
to persecute, and support the heresy, for it accounts you its king.''
Such was the language of their iniquity. And the Emperor, when in his
passage through the country on his hasty march against Magnentius \footnote{p. 159, note I.

Return to text}, he saw the communion of the Bishops with Athanasius, like one
set on fire, suddenly changed his mind, and no longer remembered his
oaths, but was alike forgetful of what he had written, and regardless
of the duty he owed his brother. For in his letters to him, as well as
in his interview with Athanasius, he took an oath that he would not
act otherwise than as the people should wish, and as should be
agreeable to the Bishop. But his zeal for impiety caused him at once
to forget all these things. And yet one ought not to wonder that after
so many letters and so many oaths Constantius had altered his mind,
when we remember that Pharaoh \footnote{p. 246, r. 5.

Return to text} of old, the tyrant of Egypt, after frequently promising and by
that means obtaining a remission of his punishments, likewise changed,
until he at last perished together with his associates in wickedness.

§. 31.

4. He compelled then the people in every city to
change their party; and on arriving at Arles and Milan, he proceeded
to act entirely in accordance with the designs and suggestions of the
heretics; or rather they acted themselves, and receiving authority
from him, furiously attacked every one. Letters and orders were
immediately sent hither to the Prefect, that for the future the corn
should be taken from Athanasius and given to those who favoured the
Arian doctrines, and that whoever pleased might freely insult them
that held communion with him; and a threat was held out to the
magistrates, if they did not hold communion with the Arians. These
things were but the prelude to what afterwards took place under the
direction of the Duke Syrianus.

5. Orders were sent also to the more distant
parts, and Notaries despatched to every city, and Palatines,
with threats to the Bishops and Magistrates, directing the Magistrates
to urge on the Bishops, and informing the Bishops that either they
must subscribe against Athanasius, and hold communion with the Arians,
or themselves undergo the punishment of exile, while the people who
took part with them were to understand that chains, and insults, and
scourgings, and the loss of their possessions, would be their portion.
These orders were not neglected, for the commissioners had in their
company the Clergy of Ursacius and Valens, to inspire them with zeal,
and to inform the Emperor if the Magistrates neglected their duty. The
other heresies, as younger sisters of their own \footnote{vol. 8. p. 89, note M. p. 189, note B.

Return to text}, they permitted to blaspheme the Lord, and only conspired
against the Christians, not enduring to hear orthodox language
concerning Christ. How many Bishops in consequence, according to the
words of Scripture, were brought before rulers and kings, and received
this sentence from magistrates, ``Subscribe, or withdraw from your
churches, for the Emperor has commanded you to be deposed!'' How many
in every city were made to waver, lest they should accuse them as
friends of the Bishops! Moreover letters were sent to the city
authorities, and a threat of a fine was held out to them, if they did
not compel the Bishops of their respective cities to subscribe. In
short, every place and every city was full of fear and confusion,
while the Bishops were dragged along to trial, and the magistrates
witnessed the lamentations and groans of the people.

§. 32.

6. Such were the proceedings of the Palatine
commissioners; on the other hand, those admirable persons, confident
in the patronage which they had obtained, display great zeal, and
cause some of the Bishops to be summoned before the Emperor, while
they persecute others by letters, inventing charges against them; to
the intent that the one might be overawed by the presence of
Constantius, and the other, through fear of the commissioners and the
threats held out to them in these pretended accusations, might be
brought to renounce their orthodox and pious opinions \footnote{[\emph{mn\={e}m\={e}s. l. gn\={o}m\={e}s}]. Montf. after
Nann.

Return to text}. In this manner it was that the Emperor forced so great a
multitude of Bishops, partly by threats, and partly by promises, to
declare, ``We will no longer hold communion with Athanasius.'' For
those who came for an interview, were not admitted to his presence,
nor allowed any relaxation, not so much as to go out of their
dwellings, until they had either subscribed, or refused and incurred
banishment thereupon. And this he did because he saw that the heresy
was hateful \footnote{p. 217, r. 7. p. 223, r. 3. p. 248, r. 3. p. 259.

Return to text} to all
men. For this reason especially he compelled so many to add their
names to the small number \footnote{p. 132, r. 5.

Return to text}
of the Arians, his earnest desire being to collect together a crowd of
names, both from envy of the Bishop, and for the sake of making a shew
in favour of the Arian impiety, of which he is the patron; supposing
that he will be able to alter the truth, as easily as he can influence
the minds of men. He knows not, nor has ever read, how that the
Sadducees and the Herodians, taking unto them the Pharisees, were not
able to obscure the truth; rather it shines out thereby more brightly
every day, while they crying out, We have no king but Cæsar [John xix.
25.] \footnote{vol. 8. p. 190.

Return to text}, and obtaining
the judgment of Pilate in their favour, are nevertheless left
destitute, and wait in utter shame, expecting shortly \footnote{[\emph{hoson oudep\={o}}]. p. 228, note B. Const. died in 362, aged
45.

Return to text} to become bereft, like the partridge, when they shall see
their patron near his death [vid. Jer. xvii. 11. Sept.].

§. 33.

7. Now if it was altogether unbecoming in any of
the Bishops to change their opinions merely from fear of these things,
yet it was much more so \footnote{p. 193 fin.

Return to text},
and not the part of men who have confidence in what they believe, to
force and compel the unwilling. In this manner it is that the Devil,
when he has no truth on his side \footnote{The fault consists in substituting persecution \emph{for} the power of
truth. Vid. p. 279, note C.

Return to text

}, attacks and breaks down the doors of them that admit him with
axes and hammers [vid. Ps. lxxiv. 6.]. But our Saviour is so gentle
that He teaches thus, \emph{If any man wills to come after Me} [Mat.
xvi. 24.], and, \emph{Whoever wills to be My disciple}; and coming to
each He does not force them, but knocks at the door and says, \emph{Open
unto Me, My sister, My spouse} [Cant. v. 2.]; and if they open to
Him, He enters in, but if they delay and will not, He departs from
them. For the truth is not preached with swords or with darts, nor by
means of soldiers; but by persuasion and counsel. But what persuasion
is there where fear of the Emperor prevails? or what counsel is there,
when he who withstands them receives at last banishment and
death? Even David, although he was a king, and had his enemy in his
power, prevented not the soldiers by an exercise of authority when
they wished to kill his enemy, but, as the Scripture says, David
persuaded his men by arguments, and suffered them not to rise up and
put Saul to death [vid. 1 Sam. xxvi. 9.]. But he, being without
arguments of reason, forces all men by his power, that it may be shewn
to all, that their wisdom is not according to God, but merely human,
and that they who favour the Arian doctrines have indeed no king but Cæsar;
for by his means it is that these enemies of Christ accomplish
whatsoever they wish to do.

8. But while they thought that they were carrying
on their designs against many by his means, they knew not that they
were making many to be confessors, of whom are those who have lately
made so glorious a confession, religious men, and excellent Bishops,
Paulinus \footnote{p. 239, r. 4.

Return to text} Bishop of
Treves the Metropolis of Gaul, Lucifer \footnote{p. 191, r. 3-6.

Return to text} Bishop of the Metropolis of Sardinia, Eusebius of Vercelli in
Italy, and Dionysius of Milan, which is the Metropolis of Italy. These
the Emperor summond before him, and commanded them to subscribe
against Athanasius, and to hold communion with the heretics; and when
they were astonished at this novel procedure, and said that there was
no Ecclesiastical Canon \footnote{p. 3.

Return to text}
to this effect, he immediately said, ``Whatever I will, be that
esteemed a Canon; the Bishops of Syria let me thus speak. Either then
obey, or go into banishment.

§. 34.

9. When the Bishops heard this they were utterly
amazed, and stretching forth their hands to God, they used great
boldness of speech against him, teaching him that the kingdom was not
his, but God’s who had given it to him, whom also they bid him fear,
lest he should suddenly take it away from him. And they threatened him
with the day of judgment, and warned him against infringing
Ecclesiastical order, and mingling Roman sovereignty with the
constitution \footnote{[\emph{diatag\={e}i}], p. 249, r. 10.

Return to text} of the
Church, nor to introduce the Arian heresy into the Church of God. But
he would not listen to them, nor permit them to speak further, but
threatened them so much the more, and drew his sword against them, and
gave orders for some to be led to punishment; although afterwards,
like Pharaoh \footnote{p. 243, r. 2.

Return to text}, he
repented. The holy men therefore shaking off the dust, and
looking up to God, neither feared the threats of the Emperor, nor
betrayed their cause before his drawn sword; but received their
banishment, as a service pertaining to their ministry. And as they
passed along, they preached the Gospel in every place and city \footnote{infr. p. 253, r. 2. vid. Acts viii. 4. Phil. i. 12.

Return to text}, although they were in bonds, proclaiming the orthodox faith,
anathematizing the Arian heresy, and stigmatizing the recantation of
Ursacius and Valens. But this was contrary to the intention of their
enemies; for the greater was the distance of their place of
banishment, so much the more was the hatred against them increased,
while the wanderings of these men were but the heralding of their
impiety. For who that saw them as they passed along, did not greatly
admire them as Confessors, and renounce and abominate the others,
calling them not only impious men, but executioners \footnote{[\emph{d\={e}mious}], vid. p. 133, r. 12.

Return to text} and murderers, and every thing rather than Christians \footnote{supr. p. 208, note B.

Return to text}.

\xsection{Chapter 5. Persecution and lapse of Liberius}

§. 35.

1. N\textsc{ow} it had been better
if from the first Constantius had never become connected with this
heresy at all; or being connected with it, if he had not yielded so
much to those impious men; or having yielded to them, if he had stood
by them only thus far, so that judgment might come upon them all for
these atrocities alone. But as it would seem, like madmen, having
entangled themselves in the bonds of impiety, they are drawing down
upon their own heads a more severe judgment. Thus from the first \footnote{Margin Notes

1. in contrast to date of his fall, p. 255, r. 6.

Return to text} they spared not even Liberius Bishop of Rome, but extended \footnote{[\emph{t\={e}n manian exeteinan}]; vid. [\emph{ekteinai
t\={e}n manian}]. infr. p. 254. r. 1. And so in the letter of
the Council of Chalcedon to Pope Leo; which says that Dioscorus, [\emph{kat'
autou t\={e}s ampelou t\={e}n phulak\={e}n para tou s\={o}t\={e}ros
epitetrammenou t\={e}n manian exe teine. legomen d\={e}. t\={e}s
s\={e}s hosiot\={e}tos}]. Hard. Conc. t. 2. p. 656. As to the
words [\emph{hoti apostolikos esti thronos}], the phrase ``Apostolical
throne or see,'' is given also, though not as an appellative, to the
sees of Antioch, Ephesus, \&c. vid. Tertull. de Præscript. 36.
August. Ep. 43. 7. Even were it to be here construed ``because it is
the Apostolical see,'' yet perhaps Athanasius uses it from his
familiarity with Latin ideas during his frequent exiles in the West,
just as he also adopts some of their theological terms. The Eusebians
had in the first instance resisted the authority of Rome, though with
expressions of respect. supr. p. 40, note C.

Return to text} their fury \footnote{[\emph{manian}].

Return to text} even to those parts;
they respected not his bishopric, because it was an Apostolical
throne; they felt no reverence for Rome, because she is the Metropolis
of Romania \footnote{By Romania is meant the Roman Empire, according to Montfaucon after
Nannius. vid. Præfat. xxxiv. xxxv. And so Epiph. Hær. lxvi. 1 fin.
p. 618. and lxviii. 2 init. p. 728. Nil. Ep. i. 75. vid. Ducange
Gloss. Græc. in voc.

Return to text}; they remembered not that
formerly in their letters they had spoken of her Bishops as
Apostolical men. But confounding all things together, they at once
forgot every thing, and cared only to shew their zeal in behalf of
impiety. When they perceived that he was an orthodox man, and hated \footnote{pp. 245 r. 1. 254, r. 2.

Return to text} the Arian heresy, and earnestly endeavoured to persuade all persons
to renounce and withdraw from it, these impious men reasoned thus with
themselves: ``If we can persuade Liberius, we shall soon prevail over
all.''

2. Accordingly they accuse him falsely before the
Emperor; and he, expecting easily to draw over all men to his side by
means of Liberius, writes to him, and sends a certain eunuch called
Eusebius with letters and offerings, to cajole him with the presents,
and to threaten him with the letters. The eunuch accordingly went to
Rome, and first proposed to Liberius to subscribe against Athanasius,
and to hold communion with the Arians, saying, ``The Emperor wishes it,
and commands you to do so.'' And then showing him the offerings, he
took him by the hand, and again besought him, saying, ``Be persuaded to
comply with the Emperor's request, and receive these.'' §. 36. But the
Bishop endeavoured to convince him, reasoning with him thus: ``How is
it possible for me to do this against Athanasius? how can we condemn a
man, whom not one \footnote{at Alexandria.

Return to text} Council only, but a
second \footnote{at Sardica.

Return to text} assembled from all parts of the
world \footnote{[\emph{pantachoden}].

Return to text}, has fairly acquitted, and whom
the Church of Rome dismissed in peace? who will approve of our
conduct, if we reject in his absence one, whose presence \footnote{vid. p. 49 fin.

Return to text} amongst us we gladly welcomed \footnote{p. 230, r. 2.

Return to text}, and
admitted him to our communion? There is no Ecclesiastical Canon \footnote{pp. 41, 49, 55.

Return to text} which can authorize such a proceeding; nor have we had transmitted
to us any such tradition \footnote{[\emph{paradosis}], vid. p. 229, note B.

Return to text} from the
Fathers, which they might have received from the great and blessed
Apostle Peter \footnote{p. 57, note U.

Return to text}.

3. ``But if the Emperor is really concerned for
the peace of the Church, if he requires our decrees respecting
Athanasius to be reversed, let their proceedings both against him and
against all the others be reversed also; and then let an
Ecclesiastical Council be called at a distance from the Court \footnote{or Palace, supr. pp. 25, 227, 246.

Return to text}, at which the Emperor shall not be present, nor any Count be
admitted, nor magistrate to threaten us, but where only the fear of
God, and the Apostolical rule \footnote{[\emph{t\={o}n apostol\={o}n diataxis}], supr. pp. 57, 246.

Return to text} shall
prevail; that so in the first place, the faith of the Church may be
secured, as the Fathers defined it in the Council of Nicæa, and the
supporters of the Arian doctrines may be cast out, and their heresy
anathematized. And then after that, an enquiry being made into the
charges brought against Athanasius, and any other beside, as well as
into those things of which the other party is accused, let the guilty
be cast out, and the innocent receive encouragement and support.
For it is impossible that they who maintain an impious creed can be
admitted as members of a Council; nor is it fit that an enquiry into
matters of conduct should precede the enquiry concerning the faith \footnote{vid. Pallavicin. Conc. Trid. vi. 7. Sarpi. Hist. ii.

Return to text}; but all diversity of opinion on points of faith ought first to be
eradicated, and then the enquiry made into matters of conduct. Our
Lord Jesus Christ did not heal them that were afflicted, until they
shewed and declared what faith they had in Him. These things we have
received from the Fathers; these report to the Emperor; for they are
both profitable for him and edifying to the Church. But let not
Ursacius and Valens be listened to, for they have retracted their
former assertions, and in what they now say they are not to be
trusted.''

§. 37.

4. These were the words of the Bishop Liberius.
And the eunuch \footnote{[\emph{eunouchos}].

Return to text}, who was vexed, not
so much because he would not subscribe, as because he found him an
enemy to the heresy, forgetting that he was in the presence of a
Bishop \footnote{[\emph{pros episkopon \={e}n}].

Return to text}, after threatening him
severely, went away with the offerings; and proceeded to perpetrate an
offence, which is foreign from a Christian, and too audacious for a
eunuch \footnote{[\emph{spadont\={o}n}].

Return to text}. In imitation of the
transgression of Saul, he went to the Martyry \footnote{``Under this canopy,'' [the Baldacchino in the present St. Peter's
Church,] ``is the high altar, which is only used on the most solemn
ceremonies, and beneath it repose the bodies of St. Peter and St.
Paul. That of St. Peter lies in the place where it was first buried.
It is said that Pope Anacletus, while he was only a priest,
constructed a chapel here in 106, which was called the Confessional of
St. Peter, and inclosed the body of the Apostle in a marble urn.
Constantine is reported to have covered the urn with metal, so that it
can never be seen.'' Burton's Rome, p. 425.

Return to text} of the Apostle Peter, and then presented the offerings. But
Liberius having notice of it, was very angry with the person who kept
the place, that he had not prevented him, and cast out the offerings
as an unlawful sacrifice, which increased the anger of the mutilated \footnote{[\emph{ton thladian}].

Return to text} creature against him. Consequently he exasperates the Emperor
against him, saying, ``The matter that concerns us is no longer the
obtaining the subscription of Liberius, but the fact that he is so
resolutely opposed to the heresy, that he anathematizes the Arians by
name.'' He also stirs up the other eunuchs to say the same; for many of
those who are about Constantius, or rather the whole number of them,
are eunuchs \footnote{vid. Gibbon, Hist. ch. 19 init.

Return to text}, who engross all
the influence with him, and it is impossible to do any thing there
without them. The Emperor accordingly writes to Rome, and again
Palatines, and Notaries, and Counts are sent off with letters to the
Prefect, in order that either they may inveigle Liberius by stratagem
away from Rome and send him to the Court to him, or else persecute him
by violence.

§. 38.

5. Such being the tenor of the letters, there
also fear and treachery forthwith prevailed throughout the whole city.
How many were the families against which threats were held out! How
many received great promises on condition of their acting against
Liberius! How many Bishops hid themselves when they saw these things!
How many noble women retired to their estates in consequence of the
calumnies of the enemies of Christ! how many ascetics were made the
objects of their plots! How many who were sojourning there, and had
made that place their home, did they cause to be persecuted! How often
and how strictly did they guard the harbour \footnote{Ostia, vid. Gibbon, Hist. ch. 31. p. 303.

Return to text} and the approaches to the gates, lest any orthodox person should
enter and visit Liberius! Rome also had trial of the enemies of
Christ, and now experienced what before she would not believe, when
she heard how the other Churches in every city were ravaged by them.

6. It was the eunuchs who instigated these
proceedings against all. And the most remarkable circumstance in the
matter is this; that the Arian heresy which denies the Son of God,
receives its support from eunuchs, who, as both their bodies are
fruitless, and their souls barren of the seeds of virtue, cannot bear
even to hear the name of son. The Eunuch of Ethiopia indeed, though he
understood not what he read, believed the words of Philip, when he
taught him concerning our Saviour [Acts viii. 27.]; but the eunuchs of
Constantius cannot endure the confession of Peter \footnote{Mat. xvi. 16. allusion to Liberius? vid. p. 57, note U. Hard. Conc. t.
2. p. 305, E.

Return to text}, nay, they turn away when the Father manifests the Son, and madly
rage against those who say, that the Son of God is His genuine Son,
thus claiming as a heresy of eunuchs, that there is no genuine and
true offspring of the Father. On these grounds it is that the law
forbids such persons to be admitted into any ecclesiastical Council \footnote{Can. Nic. 1.

Return to text}; notwithstanding which these have now regarded them as competent
judges of ecclesiastical causes, and whatever seems good to
them, that Constantius decrees, while men with the name of Bishops
dissemble with them. Oh! who shall be their historian? who shall
transmit the record of these things to future generations? who indeed
would believe it, were he to hear it, that eunuchs who are scarcely
entrusted with household services (for theirs is a pleasure-loving \footnote{[\emph{phil\={e}donon}], this the key to his severity towards them.

Return to text} race, that has no serious concern but that of hindering in others
what nature has taken from them); that these, I say, now exercise
authority in ecclesiastical matters, and that Constantius in
submission to their will treacherously conspired against all, and
banished Liberius!

§. 39.

7. For after the Emperor had frequently written
to Rome, had threatened, sent commissioners, devised schemes, on the
persecution subsequently breaking out at Alexandria, Liberius is
dragged before him, who uses great boldness of speech towards him. ``Cease,''
he said, ``to persecute the Christians; attempt not by my means to
introduce impiety into the Church. We are ready to suffer any thing
rather than to be called Arian fanatics. We are Christians; compel us
not to become enemies of Christ. We also give you this counsel: fight
not against Him who gave you this empire, nor show impiety towards Him
instead of thankfulness \footnote{p. 246. §. 34.

Return to text}; persecute
not them that believe in Him, lest you also hear the words, \emph{It is
hard for thee to kick against the pricks} [Acts ix. 5.]. Nay, I
would that you might hear them, that you might obey, as the holy Paul
did. Behold, here we are; we are come, before they fabricate charges.
For this cause we hastened hither, knowing that banishment awaits us
at your hands, that we might suffer before a charge encounters us, and
that all may clearly see that all the others too have suffered as we
shall suffer, and that the charges brought against them were
fabrications of their enemies, and all their proceedings are mere
calumny and falsehood.''

§. 40.

8. These were the words of Liberius at that time,
and he was admired by all men for them. But the Emperor instead of
answering, only gave orders for their banishment, separating each of
them from the rest, as he had done in the former cases. For he had
himself devised this plan in the banishments which he inflicted, that
so the severity of his punishments might be greater than that of
former tyrants and persecutors \footnote{p. 235, r. 4. infr. §. 60. §. 64.

Return to text}. In
the former persecution Maximian who was then Emperor commanded a
number of Confessors to be banished together, and thus lightened their
punishment by the consolation which he gave them in each other's
society. But this man was more savage than he; he separated those who
had spoken boldly and confessed together, he put asunder them who were
united by the bond of faith, that when they came to die they might not
see one another; thinking that bodily separation can disunite also the
affections of the mind, and that being severed from each other, they
would forget the concord and unanimity which existed among them. He
knew not that however each one may remain apart from the rest, he has
nevertheless with him that Lord, whom they confessed in one body
together, who will also provide, (as he did in the case of the prophet
Elisha,) that more shall be with each of them, than there are soldiers
with Constantius. Of a truth iniquity is blind; for in that they
thought to afflict the Confessors, by separating them from one
another, they rather brought thereby a great injury upon themselves.
For had they continued in each other's company, and abode together,
the pollutions of those impious men would have been proclaimed from
one place only; but now by putting them asunder, they have made their
impious heresy and wickedness to spread abroad and become known in
every place \footnote{p. 247, r. 1.

Return to text}.

§. 41.

9. Who that shall hear what they did in the
course of these proceedings will not think them to be any thing rather
than Christians \footnote{pp. 247, r. 3. 208, note B.

Return to text}? When Liberius sent
Eutropius a Presbyter and Hilarius a Deacon with letters to the
Emperor, at the time that Lucifer and his friends made their
confession, they banished the Presbyter on the spot, and after
stripping Hilarius \footnote{This Hilary afterwards followed Lucifer of Cagliari in his schism. He
is supposed to be the author of the Comments on St. Paul's Epistles
attributed to St. Ambrose, who goes under the name of Ambrosiaster.

Return to text

} the Deacon and
scourging him on the back, they banished him too, exclaiming, ``Why
didst thou not resist Liberius instead of being the bearer of letters
from him.'' Ursacius and Valens with the eunuchs who sided with them
were the authors of this outrage. The Deacon, while he was being
scourged, praised the Lord, remembering his words, \emph{I gave My back
to the smiters} [Is. l. 6.]; but they while they scourged him
laughed and mocked him, feeling no shame that they were insulting a
Levite. Indeed they acted but consistently in laughing while he
continued to praise God; for it is the part of Christians to endure
stripes, but to scourge Christians is the outrage of a Pilate or a
Caiaphas \footnote{p. 194, r. 1.

Return to text}.

10. Thus they endeavoured at the first to corrupt
the Church of the Romans, wishing to introduce impiety into it as well
as others. But Liberius after he had been in banishment two years gave
way, and from fear of threatened death was induced to subscribe. Yet
even this only shews their violent conduct, and the hatred \footnote{p. 217, r. 7.

Return to text} of Liberius against the heresy, and his support of Athanasius, so
long as he was suffered to exercise a free choice. For that which men
are forced by torture to do contrary to their first judgment, ought
not to be considered the willing deed of those who are in fear, but
rather of their tormentors \footnote{p. 245, note B.

Return to text}. They
however attempted every thing in support of their heresy, while the
people in every Church, preserving the faith which they had learnt,
waited for the return of their teachers, and cast from them, and all
avoided, as they would a serpent, the Antichristian heresy.

\xsection{Chapter 6. Persecution and lapse of Hosius}

§. 42.

1. B\textsc{ut} although they had
done all this, yet these impious men thought they had accomplished
nothing, so long as the great Hosius escaped their wicked
machinations. And now they undertook to extend their fury \footnote{Margin Notes

1. [\emph{ekteinai t\={e}n manian}], p. 248,
note A.

Return to text} to that venerable old man. They felt no shame at the thought
that he is the father of the Bishops \footnote{pp. 158, 230, 256.

Return to text}; they regarded not that he had been a Confessor \footnote{under Maximian.

Return to text}; they reverenced not the length of his Episcopate, in which he
had continued more than sixty years; but they set aside every thing,
and looked only to the interests of their heresy, as being of a truth
such as neither fear God, nor regard man [vid. Luke xviii. 2.].
Accordingly they went to Constantius, and again employed such
arguments as the following, ``We have done every thing; we have
banished the Bishop of the Romans; and before him a very great number
of other Bishops, and have filled every place with alarm. But these
strong measures of yours are as nothing to us, nor is our success at
all more secure, so long as Hosius remains. While he is in his own
place, the rest also continue in their Churches, for he is able by his
arguments and his faith to persuade all men against us. He is the
president of Councils \footnote{of Nicæa and Sardica.

Return to text},
and his letters are every where attended to. He it was who put forth
the Nicene Confession, and proclaimed every where that the Arians were
heretics. If therefore he is suffered to remain, the banishment of the
rest is of no avail, for our heresy will be destroyed. Begin then to
persecute him also, and spare him not, ancient \footnote{[\emph{archaios}], vid. p. 284.

Return to text} as he is. Our heresy knows not to honour the hoary hairs of the
aged.''

§. 43.

2. Upon hearing this, the Emperor no longer
delayed, but knowing the man, and the weight of his years, wrote to
summon him. This was when he first \footnote{supr. p. 248, r. 1. i.e. two years before his fall.

Return to text} began his attempt upon Liberius. Upon his arrival he
desired him, and urged him with the usual arguments, with which he
thought also to deceive the others, that he would subscribe against
us, and hold communion with the Arians. But the old man, scarcely
bearing to hear the words, and grieved that he had even ventured to
utter such a proposal, severely rebuked him, and after endeavouring to
convince him of his error, withdrew to his own country and Church. But
the heretics still complaining, and instigating him to proceed, (he
had the eunuchs also to remind him and to urge him further,) the
Emperor again wrote in threatening terms; but still Hosius, while he
endured their insults, was unmoved by any fear of their designs
against him, and remaining firm to his purpose, as one who had built
the house of his faith upon the rock, he spoke boldly against the
heresy, regarding the threats held out to him in the letters but as
drops of rain and blasts of wind. And although Constantius wrote
frequently, sometimes flattering him with the title of Father \footnote{p. 255, r. 2.

Return to text}, and sometimes threatening and recounting the names of those
who had been banished, and saying, ``Will you continue the only person
to oppose the heresy? Be persuaded and subscribe against Athanasius;
for whoever subscribes against him thereby embraces with us the Arian
cause;'' still Hosius remained fearless, and while suffering these
insults, wrote an answer in such terms as these. We have read the
letter, which is placed at the end \footnote{transferred by copyists hither.

Return to text}.

§. 44.

3. Hosius to Constantius the Emperor sends health
in the Lord.

I was a Confessor at the first, when a
persecution arose in the time of your grandfather Maximian; and if you
shall persecute me, I am ready now too to endure any thing rather than
to shed innocent blood and to betray the truth. But I cannot approve
of your conduct in writing after this threatening manner. Cease to
write thus; adopt not the cause of Arius, nor listen to those in the
East, nor give credit to Ursacius and Valens. For whatever they
assert, it is not on account of Athanasius, but for the sake of their
own heresy. Believe my statement, O Constantius, who am of an age to
be your grandfather. I was present at the Council of Sardica,
when you and your brother Constans of blessed memory assembled us all
together; and on my own account I challenged the enemies of
Athanasius, when they came to the Church where I abode \footnote{Corduba.

Return to text}, that if they had any thing against him they might declare it;
desiring them to have confidence, and not to expect otherwise than
that a right judgment would be passed in all things. This I did once
and again, requesting them, if they were unwilling to appear before
the whole Council, yet to appear before me alone; promising them also,
that if he should be proved guilty, he should certainly be rejected by
us; but if he should be found to be blameless, and should prove them
to be calumniators, that if they should then refuse to hold communion
with him, I would persuade him to go with me into Spain. Athanasius
was willing to comply with these conditions, and made no objection to
my proposal; but they, altogether distrusting their cause, would not
consent. And on another occasion Athanasius came to your Court \footnote{[\emph{stratopedon}], p. 100, note Z.

Return to text}, when you wrote for him, and his enemies being at the time in
Antioch, he requested that they might be summoned either altogether or
separately, in order that they might either convict him, or be
convicted, and might either in his presence prove him to be what they
represented, or cease to accuse him when absent. To this proposal also
you would not listen, and they equally rejected it.

4. Why then do you still give ear to them that
speak evil of him? How can you endure Ursacius and Valens, although
they have retracted, and made a written confession of their calumnies?
For it is not true, as they pretend, that they were forced to confess;
there were no soldiers at hand to influence them; your brother was not
cognizant of the matter \footnote{p. 15, note F. p. 242.

Return to text}.
No, such things were not done under his government, as are done now;
God forbid. But they voluntarily went up \footnote{p. 223, r. 6.

Return to text} to Rome, and in the presence of the Bishop and Presbyters
wrote their recantation, having previously addressed to Athanasius a
friendly and peaceable letter. And if they pretend that force was
employed towards them, and acknowledge that this is an evil thing,
which you also disapprove of; then do you cease to use force \footnote{pp. 19, 205, 221, n. 3 fin. 242, r. 5. 245, note B. 267, r. 2. 279,
note C.

Return to text}; write no letters, send no Counts; but release those that have
been banished, lest while you are complaining of violence, they
do but exercise greater violence. When was any such thing done by
Constans? What Bishop suffered banishment at his hands? When did he
appear in presence at an Ecclesiastical trial? When did any Palatine
of his compel men to subscribe against any one, that Valens and his
fellows should be able to affirm this?

5. Cease these proceedings, I beseech you, and
remember that you are a mortal man. Be afraid of the day of judgment,
and keep yourself pure thereunto. Intrude not yourself into
Ecclesiastical matters, neither give commands unto us concerning them;
but learn them from us. God hath put into your hands the kingdom; to
us He hath entrusted the affairs of his Church; and as he who should
steal the empire from you would resist the ordinance of God, so
likewise fear on your part lest by taking upon yourself the government
of the Church, you become guilty of a great offence. It is written, \emph{Render
unto Cæsar the things that are Cæsar's, and unto God the things that
are God's} [Mat. xxii. 21.]. Neither therefore is it permitted unto
us to exercise an earthly rule, nor have you, Sire, any authority to
burn incense \footnote{Incense is mentioned in the Apostolical Canon
iii. but. apparently no where else till this date. Hippol. de Consumm.
Mund. adduced by Beveridge on the Canon is not genuine. At the same
time it must be recollected, that Hosius was at this time 100 years
old, and a rite which he singles out (if he does not speak figurately)
to desciibe the Eucharistic Sacrifice, could not be a recent one. From
Tertull. Apol. 42. and Arnobius, contr. Gent. vii. 27. it appears to
have been unknown to the African Churches in their day. vid. Bon. Rer.
Lit. i. 25. n. 9. Bellarm. de Miss. ii. 15. Bevereg. Cod. Can. Vind.
ii. 2. r. 5. Dall. de Pseudepig. Apost. iii. 14. §. 4. Dodwell, Use
of Incense.

Return to text}. These
things I write unto you out of a concern for your salvation. With
regard to the subject of your letters, this is my determination: I
will not unite myself to the Arians; I anathematize their heresy.
Neither will I subscribe against Athanasius, whom both we and the
Church of the Romans, and the whole Council pronounced to be
guiltless. And yourself also, when you understood this, sent for the
man, and gave him permission to return with honour to his country and
his Church. What reason then can there be for so great a change in
your conduct? The same persons who were his enemies before, are so now
also; and the things they now whisper to his prejudice, (for they
do not declare them openly in his presence,) the same they spoke
against him, before you sent for him; the same they spread abroad
concerning him when they came to the Council. And when I required them
to come forward, as I have before said, they were unable to produce
their proofs; had they possessed any, they would not have fled so
disgracefully. Who then has persuaded you so long after to forget your
own letters and declarations? Forbear, and be not influenced by evil
men, lest while you act for the mutual advantage of yourself and them,
you bring guilt upon yourself. For here you comply with their desires,
hereafter in the judgment you will have to answer for doing so alone.
These men desire by your means to injure their enemy, and wish to make
you the minister of their wickedness, in order that through your help
they may sow the seeds \footnote{vid. vol. 8. p. 5. note k. It is remarkable, this letter having so
much its own character, and being so unlike Athanasius's writings in
style, that a phrase characteristic of him should here occur in it.
Did Athan. translate it from Latin?

Return to text} of
their accursed heresy in the Church. Now it is not a prudent thing to
cast one's self into manifest danger for the pleasure of others. Cease
then, I beseech you, O Constantius, and be persuaded by me. These
things it becomes me to write, and you not to despise.

§. 45.

6. Such were the sentiments, and such the letter,
of the Abraham-like old man, Hosius \footnote{i.e. sacred, saintly.

Return to text}, truly so called \footnote{[\emph{ho al\={e}th\={o}s Hosios, kataskopoi, ou gar episkopoi}],
supr. §. 3. infr. §§. 48, 75 fin. and so [\emph{al\={e}th\={o}s
Eusebie}], Theod. Hist. i. 4. [\emph{On\={e}simon, ton pote soi
achr\={e}ston, nuni de euchr\={e}ston}], Ep. ad Phil. 10.
vid. vol. 8. p. 114, note b.

Return to text}.
But the Emperor desisted not from his designs, nor ceased to seek an
occasion against him; but continued to threaten him severely, with a
view either to bring him over by force, or to banish him if he refused
to comply. And as the Officers and Satraps of Babylon \footnote{p. 195, r. 1.

Return to text} seeking an occasion against Daniel, found none except in the
law of his God; so likewise these present Satraps of impiety were
unable to invent any charge against the old man, (for this true
Hosius, and his blameless life were known to all,) except the charge
of hatred \footnote{p. 260, r. 1.

Return to text} to their
heresy. They therefore proceeded to accuse him; though not under the
same circumstances as those others accused Daniel to Darius, for
Darius was grieved to hear the charge, but as Jezebel accused
Naboth, and as the Jews applied themselves to Herod. And they said, ``He
not only will not subscribe against Athanasius, but also on his
account condemns us; and his hatred \footnote{p. 245, r. 1.

Return to text} to the heresy is so great, that he also writes to others, that
they should rather suffer death, than become traitors to the truth.
For, he says, our beloved Athanasius also is persecuted for the Truth's
sake, and Liberius Bishop of Rome, and all the rest, are treacherously
assailed.''

7. When this patron of impiety, and Emperor of
heresy \footnote{vid. pp. 226, r. 1. 243, 267, r. 3.

Return to text}, Constantius,
heard this, and especially that there were others also in Spain of the
same mind as Hosius, after he had tempted them also to subscribe, and
was unable to compel them to do so, he sent for Hosius, and instead of
banishing him, detained him a whole year in Sirmium. Godless, unholy,
without natural affection, he feared not God, he regarded not his
father's love for Hosius, he reverenced not his great age, for he was
now a hundred years old \footnote{[\emph{oute ton Theon phob\={e}theis ho atheos oute tou patros t\={e}n
diathesin aidestheis ho anosios, oute to g\={e}ras aischuntheis ho
astorgos}].

Return to text

};
but all these things this modern Ahab, this second Belshazzar of our
times, disregarded for the sake of impiety. He used such violence
towards the old man, and confined him so straitly, that at last,
broken by suffering, he was brought, though hardly, to hold communion
with Valens and Ursacius, though he would not subscribe against
Athanasius. Yet even thus he forgot not his duty, for at the approach
of death, as it were by his last testament, he bore witness to the
force which had been used towards him, and anathematized the Arian
heresy, and gave strict charge that no one should receive it.

§. 46.

8. Who that witnessed these things, or that has
merely heard of them, will not be greatly amazed, and cry aloud unto
the Lord, saying, \emph{Wilt Thou make a full end of the remnant of
Israel?} [Ez. xi. 13.] Who that is acquainted with these
proceedings, will not with good reason cry out and say, \emph{A wonderful
and horrible thing is committed in the land}; and, \emph{The heavens
are astonished at this, and the earth is even more horribly afraid}
[Jer. v. 30; ii. 12.]. The fathers of the people and the teachers of
the faith are taken away, and the impious are brought into the
Churches? Who that saw when Liberius Bishop of Rome was banished, and
when the great Hosius the father \footnote{p. 230, r. 5.

Return to text} of the Bishops suffered these things, or who that saw so many
Bishops banished out of Spain and the other parts, could fail to
perceive, however little sense he might possess, that the charges \footnote{vid. in Apol. contr. Ar. and ad Const.

Return to text} against Athanasius also and the rest were false, and
altogether mere calumny? For this reason those others also endured all
suffering, because they saw plainly that the conspiracies laid against
these were founded in falsehood. For what charge was there against
Liberius? or what accusation against the aged Hosius? who bore even a
false witness against Paulinus, and Lucifer, and Dionysius, and
Eusebius? or what sin could be laid to the account of the rest of the
banished Bishops, and Presbyters, and Deacons? None whatever; God
forbid. There were no charges against them on which a plot for their
ruin might be formed; nor was it on the ground of any accusation that
they were severally banished. It was a breaking out of impiety against
godliness \footnote{[\emph{asebeias, eusebeias}], vol. 8. p. 1, note a.

Return to text}; it was
zeal for the Arian heresy, and a prelude to the coming of Antichrist,
for whom Constantius is thus preparing the way.

\xsection{Chapter 7. Persecution at Alexandria}

§. 47.

1. A\textsc{fter} he had accomplished all that he desired
against the Churches in Italy, and the other parts; after he had
banished some, and violently oppressed others, and filled every place
with fear, he at last turned his fury, as it had been some
pestilential disorder, against Alexandria. This was artfully contrived
by the enemies of Christ; for in order that they might have a show of
the signatures of many Bishops, and that Athanasius might not have a
single Bishop in his persecution to whom he could even complain, they
therefore anticipated his proceedings, and filled every place with
terror, which they kept up to second \footnote{Margin Notes

1. [\emph{ephedron}].

Return to text} them in the prosecution
of their designs. But herein they perceived not through their folly
that they were not exhibiting the free sentiments \footnote{p. 141, note A, p. 257, r. 5.

Return to text} of the
Bishops, but rather the violence which themselves had employed; and
that, although his brethren should desert him, and his friends and
acquaintance stand afar off, and no one be found to sympathise with
him and console him, yet far above all these, a refuge with his God
was sufficient for him. For Elias also was alone in his persecution,
and God was all in all to the holy man. And our Saviour has given us
an example herein, who also was left alone, and exposed to the designs
of His enemies, to teach us, that when we are persecuted and deserted
by men, we must not faint, but place our hope in Him, and not betray
the Truth. For although at first it may seem to be afflicted, yet even
they who persecute shall afterwards acknowledge it.

§. 48.

2. Accordingly they urge on the Emperor, who
first writes a menacing letter, which he sends to the Duke and the
soldiers. The Notaries Diogenius \footnote{vid. pp. 173, 175, 294.

Return to text} and Hilarius \footnote{vid. pp. 173, 175, 294.

Return to text}, and
certain Palatines with them were the bearers of it; upon whose arrival
those terrible and cruel outrages were committed against the Church,
which I have briefly related a little above \footnote{p. 243, \&c.

Return to text}, and which are
known to all men from the protests put forth by the people, which are
inserted at the end of this history \footnote{vid. p. 293, note A.

Return to text}, so that any one may read
them. Then after these proceedings on the part of Syrianus, after
these enormities had been perpetrated, and violence offered to the
Virgins, as approving of such conduct and the infliction of these
evils upon us, he writes again to the senate and people of Alexandria,
instigating the younger men, and requiring them to assemble together,
and either to persecute Athanasius, or consider themselves as his
enemies. He however had withdrawn before these instructions reached
them, and from the time when Syrianus broke into the Church; for he
remembered that which is written, \emph{Hide thyself as it were for a
little moment, until the indignation be overpast} [Is. xxvi. 20.] \footnote{pp. 186, 204.

Return to text}.

§. 50. [there is no §. 49 in Montf.]

3. One Heraclius, by rank a Count, was the bearer
of this letter, and the precursor of a certain George that was
dispatched by the Emperor as a spy, for one that was sent from him
cannot be a Bishop \footnote{[\emph{kataskopou, ouk episkopos}], vid. p.
259, note C.

Return to text}; God forbid. And so indeed his conduct and
the circumstances which preceded his entrance sufficiently prove. Heraclius then published the letter, which reflected great
disgrace upon the writer. For whereas, when the great Hosius wrote to
Constantius, he had been unable to make out any plausible pretext for
his change of conduct, he now invented an excuse much more
discreditable to himself and to his advisers. He said, ``From regard
to the affection I entertained towards my brother of divine and pious
memory, I endured for a time the coming of Athanasius among you.''
This proves that he has both broken his promise, and behaved
ungratefully to his brother after his death. He then declares him to
be, as indeed he is, ``deserving of sacred and pious remembrance;''
yet as regards a command of his, or to use his own language, the ``affection'' he bore him, even though he complied merely
``for the
sake'' of the blessed Constans, he ought to deal fairly by his
brother, and make himself heir to his sentiments as well as to
the Empire. But, although, when seeking to obtain his just rights, he
deposed Vetranio, with the question, ``To whom does the inheritance
belong after a brother's death?'' \footnote{``It was an easy task to deceive the
frankness and simplicity of Vetranio, who, fluctuating some time
between the opposite views of power and interest, displayed to the
world the insincerity of his temper, and was insensibly engaged in the
snares of an artful negociation. Constantius acknowledged him as a
legitimate and equal colleague in the Empire, on condition that he
would renounce his disgraceful alliance with Magnentius, and appoint a
place of interview on the frontiers of their respective provinces ...
The united armies were commanded to assemble in a large plain near the
city [Sardica]. In the centre, according to the rules of ancient discipline,
a military tribunal, or rather scaffold, was erected, from whence the
Emperors were accustomed, on solemn and important occasions, to
harangue the troops ... The first part of his [C.'s] Oration seemed
to be pointed only against the tyrant of Gaul [Magnentius], but while
he tragically lamented the cruel murder of Constans, he insinuated,
that \emph{none, except a brother, could claim a right to the succession
of his brother}. He displayed, with some complacency, the glories
of his Imperial race, \&c. … The contagion of loyalty and
repentance was communicated from rank to rank; till the plain of
Sardica resounded with the universal acclamation of 'Away with these
upstart usurpers!' \&c. Gibbon, Hist. ch. xviii.

Return to text} yet for the sake of the accursed
heresy of the enemies of Christ, he disregards the claims of justice,
and behaves undutifully towards his brethren.

4. Nay, for the sake of this heresy, he would not
consent to observe his father's wishes without infringement; but, in
what he may gratify those impious men, he pretends to adopt his
intention, while in order to distress the others, he cares not to shew
the reverence which is due unto a father. For in consequence of the
calumnies of the Eusebians, his father sent the Bishop for a time into
Gaul to avoid the cruelty of his persecutors, (this was shewn by the
blessed Constantine, the brother of the former, after their father's
death, as appears by his letters \footnote{p. 121.

Return to text},) but he would not be
persuaded by the Eusebians to send the person whom they desired for a
Bishop, but prevented the accomplishment of their wishes, and put a
stop to their attempts with severe threats.

§. 51.

5. If therefore, as he declares in his letters,
he desired to observe his father's practice, why did he first send
out Gregory, and now this George, who eats his own stores \footnote{George had been pork-contractor to the army,
and had been detected in peculation. vid. vol. 8. p. 89, r. 1. p. 134,
note f. and infr. p. 286.

Return to text}?
Why does he endeavour so earnestly to introduce into the Church these
Arians, whom his father named Porphyrians \footnote{Constantine called the Arians by this title
after the philosopher Porphyry, the great enemy of Christianity.
Socrates has preserved the Edict. Hist. i. 9.

Return to text}, and banish others
while he patronises them? Although his father admitted Arius to
his presence, yet when Arius perjured himself and burst asunder \footnote{pp. 147, 212.

Return to text}, he lost the compassion of his father; who, on learning the truth,
condemned him as a heretic.

6. Why moreover, while pretending to respect the
Canons of the Church, has he ordered the whole course of his conduct
in opposition to them? For where is there a Canon that a Bishop should
be appointed from Court? Where is there a Canon \footnote{p. 249, r. 6. p. 268, r. 1.

Return to text} that permits
soldiers to invade Churches? What tradition \footnote{p. 249, r. 7.

Return to text} is there
allowing counts and ignorant \footnote{[\emph{alogistos}], vid. vol. 8. p. 2, note e

Return to text} eunuchs to exercise authority
in Ecclesiastical matters, and to make known by their edicts the
decisions of those who bear the name of Bishops? He is guilty of all
manner of falsehood for the sake of this unholy heresy. At a former
time he sent out Philagrius as Prefect a second time \footnote{p. 224, note B.

Return to text}, in
opposition to the opinion of his father, and we see what has taken
place now.

7. Nor ``for his brother's sake'' does he
speak the truth. For after his death he wrote as often as three times
to the Bishop, and repeatedly promised him that he would not change
his behaviour towards him, but exhorted him to be of good courage, and
not suffer any one to alarm him, but to continue to abide in his
Church in perfect security \footnote{pp. 174, 238.

Return to text}. He also sent his commands by
Count Asterius, and Palladius the Notary, to Felicissimus who was then
Duke, and to the Prefect Nestorius, that if either Philip the Prefect,
or any other should venture to form any plot against Athanasius, they
should prevent it. §. 52. Wherefore when Diogenes came, and Syrianus laid in
wait for us, both he \footnote{p. 219.

Return to text} and we and the people demanded to see
the Emperor's letters, supposing that, as it is written, \emph{Let not
a falsehood be spoken before the king} \footnote{vid. p. 165, r. 4.

Return to text}; so when a king
has made a promise, he will not lie, nor change. If then ``for his
brother's sake he complied,'' why did he also write those letters
upon his death? And if he wrote them for ``his memory's sake,''
why did he afterwards behave so very unkindly towards him, and
persecute the man, and write what he did, alleging a judgment of
Bishops, while in truth he acted only to please himself \footnote{p. 267, r. 4.

Return to text}?

8. Nevertheless his craft has not escaped
detection, but we have the proof of it ready at hand. For if a
judgment had been passed by Bishops, what concern had the
Emperor with it? Or if it was only a threat of the Emperor, what need
in that case was there of the so-named Bishops? When was such a thing
heard of before from the beginning of the world? When did a judgment
of the Church receive its validity \footnote{[\emph{to kuros}].

Return to text} from the Emperor? or
rather when was his decree ever recognised by the Church? There have
been many Councils held heretofore; and many judgments passed by the
Church; but the Fathers never sought the consent of the Emperor
thereto, nor did the Emperor busy himself with the affairs of the
Church. The Apostle Paul had friends among them of Cæsar's
household, and in his Epistle to the Philippians he sent salutations
from them; but he never took them as his associates in Ecclesiastical
judgments \footnote{p. 249, r. 9.

Return to text}. Now however we have witnessed a novel sight,
which is a discovery of the Arian heresy. Heretics have assembled
together with the Emperor Constantius, in order that he, alleging the
authority of the Bishops, may exercise his power against whomsoever he
pleases, and while he persecutes may avoid the name of persecutor \footnote{p. 279, note C.

Return to text}; and that they, supported by the Emperor's government, may
conspire the ruin of whomsoever they will \footnote{[\emph{hois an ethel\={o}si}], and just
before [\emph{h\={o}n an etheloi}]. [And more strikuuigly just
below, §. 53 fin. [\emph{ha thelousi prattei, epei kai autos haper \={e}thelen
\={e}kousi par aut\={o}n}].] This is a very familiar phrase
with Athan. i.e. [\emph{h\={o}s ethel\={e}sen, haper ethel\={e}san,
hotan thel\={o}sin, hous ethel\={e}san}], \&c. \&c.
Some instances have been given supr. p. 15, note E. and vol. 8. p. 92,
note r. Among the many passages that might be noticed, are the
following, de Decr. §. 3 A. de Syn. §. 13 A. Apol. contr. Arian.
§§. 2 C. 14 D. 35 D. 36 D. 73 A. B. 74 F. 77 D. Ep. Æg. §§. 5 B.
19 A. 22 B. Ap. ad Const. §. 1 C. de Fug. §§. 3 C. 7 E. ad Serap.
fin. And so in this History, besides the above passage, the phrase is
found in §§. 2 D. 3 fin. 7 C. ib. D twice. 47 C. 54 init. 59 A. 60
fin. In like manner, [\emph{h\={o}s \={e}boulonto, ha boulontai}],
\&c. Ep. Encycl. §. 7 D. Apol. contr. Arian. §§. 36 D. 73 A. 74
A. 77 B. twice. ibid. D. 82 init. 83 F. ibid. B. Ep. Æg. §. 6 B. C.
Apol. ad Const. §. 32 D. de Fug. §. 1 fin. And so in this History,
§§. 2 D. 15 D. 18 C.

Return to text}; and these are all
such as are not as impious as themselves. One might look upon their
proceedings as a comedy which they are performing on the stage, in
which the pretended Bishops are actors \footnote{p. 34, r. 6.

Return to text}, and Constantius the
performer of their behests, who makes promises to them, as Herod did
to the daughter of Herodias, and they dancing before him, accomplish,
through false accusations \footnote{[\emph{orchoumenou; tas diabolas epi}], vid.
Herod. Hist. vi. 129 fin.

Return to text}, the banishment and death of the
true believers in the Lord.

§. 53.

9. Who indeed has not been injured by their
calumnies? Whom have not these enemies of Christ conspired to destroy?
Whom has Constantius failed to banish upon charges which they
have brought against them? When did he refuse to hear them willingly?
And what is most strange \footnote{p. 221, r. 1. de Decr. §. 3. F.

Return to text}, when did he permit any one to
speak against them, and did not more readily receive their testimony,
of whatever kind it might be? Where is there a Church which now enjoys
the privilege of worshipping Christ freely \footnote{p. 262, r. 2.

Return to text}? If a Church be a
maintainer of true piety, it is in danger; if it dissemble, it abides
in fear. Every place is full of hypocrisy and impiety, so far as he is
concerned; and wherever there is a pious person and a lover of Christ,
(and there are many such every where, as were the prophets and the
great Elias,) they hide themselves, if so be that they can find a
faithful friend like Abdias, and either they withdraw into caves and
dens of the earth, or pass their lives in wandering about in the
deserts. These men in their madness prefer such calumnies against
them, as Jezebel invented against Naboth, and the Jews against our
Saviour; while the Emperor, who is the patron \footnote{p. 260, r. 2.

Return to text} of the heresy,
and wishes to pervert the truth, as Ahab wished to change the vineyard
into a garden of herbs, does whatever they desire him to do, for the
suggestions he receives from them are agreeable to his own wishes \footnote{p. 265, r. 9.

Return to text}.

§. 54.

10. Accordingly he banished, as I said before,
the genuine Bishops, because they would not profess impious doctrines,
to suit his own pleasure; and now he has sent Count Heraclius to
proceed against Athanasius, who has publicly made known his decrees,
and announced the commands of the Emperor to be, that unless they
complied with the instructions contained in his letters, their bread \footnote{p. 243, §. 31. p. 276, note A.

Return to text} should be taken away, their idols overthrown, and the
persons of many of the city-magistrates and people delivered over to
certain slavery. After threatening them in this manner, he was not
ashamed to declare publicly with a loud voice, ``The Emperor
disclaims Athanasius, and has commanded that the Churches be given up
to the Arians.'' And when all wondered to hear this, and made signs
to one another, exclaiming, ``What! has Constantius become a
heretic?'' instead of blushing as he ought, this man the more
strictly obliged the senators and heathen magistrates and wardens of
the idol temples to subscribe to these conditions, and to agree
to receive as their Bishop whomsoever the Emperor should send them. Of
course Constantius was strictly upholding the Canons \footnote{p. 249, r. 6.

Return to text} of the
Church, when he caused this to be done; when, instead of requiring
letters \footnote{p. 221, r. 4.

Return to text} from the Church, he demanded them of the
market-place \footnote{infr. note F.

Return to text}, and instead of the people he asked them of the
wardens of the temples. He was conscious that he was not sending a
Bishop to preside over Christians, but a certain pragmatical person
for those who subscribed to his terms.

§. 55.

11. The Gentiles accordingly, as purchasing by
their compliance the safety of their idols, and certain of the trades
\footnote{[\emph{t\={o}n ergasi\={o}n}],—trades,
or workmen. vid. supr. p. 33, r. 2. Montfaucon has a note upon the
word in the Collect. Nov. t. 2. p. xxvi. where he corrects his Latin \emph{in
loc}. of the former passage very nearly in conformity to the
rendering given of it above, p. 33. ``In Onomastico monuimus, hic [\emph{ergasias}]
officinarum operas commodius exprimere.'' And he quotes an
inscription discovered by Spon, [\emph{touto to \={e}r\={o}on
stephanoi h\={e} ergasia t\={o}n baphe\={o}n}].

Return to text}, subscribed, though unwillingly, from fear of the threats
which he had held out to them; just as if the matter had been the
appointment of a general, or other magistrate. Indeed what, as
heathen, were they likely to do, except whatever was pleasing to the
Emperor? But the people having assembled in the great Church \footnote{vid. p. 167, note P.

Return to text}, (for it was the fourth day of the week,) Count Heraclius on the
following day takes with him Cataphronius the Prefect of Egypt, and
Faustinus the Receiver-General \footnote{Catholic, p. 163, note M.

Return to text}, and Bithynus a heretic; and
together they stir up the younger men of the common multitude \footnote{[\emph{t\={o}n agorai\={o}n}], vid. Acts
xvii. 5. [\emph{agora}] has been used just above. vid. Suicer. Thesaur.
\emph{in voc}.

Return to text}
who worshipped idols, to attack the Church, and stone the people,
saying that such was the Emperor's command. As the time of
separation \footnote{[\emph{apolusi\={o}s}], vid. Suicer. \emph{in
voc}.

Return to text} however had arrived, the greater part had already
left the Church, but there being a few women still remaining, they did
as these men had charged them, whereupon a piteous spectacle ensued.
The few women had just risen from prayer and had sat down, when the
youths having stripped themselves suddenly came upon them with stones
and clubs. Some of them the godless \footnote{[\emph{hoi atheoi}], vid. p. [sic] vol. 8. p. 3,
note f.

Return to text} wretches stoned to
death; they lacerated with stripes the holy persons of the Virgins,
tore off their veils \footnote{p. 7, r. 3.

Return to text} and exposed their heads, and when they
resisted the insult, the cowards kicked them with their feet. This was
dreadful, exceedingly dreadful; but what ensued was worse, and more intolerable than any outrage. Knowing the holy character of the
virgins, and that their ears were unaccustomed to pollution, and that
they were better able to bear stones and swords than expressions of
obscenity, they assailed them with such language. This the Arians
suggested to the young men, and laughed at all they said and did;
while the holy Virgins and other godly women fled from such words as
they would from the bite of asps, but the enemies \footnote{p. 270, note L.

Return to text} of Christ assisted them in the work, nay even, it may be, gave utterance to the
same; for they were well-pleased with the obscenities which the youths
vented upon them.

§. 56.

12. After this, that they might fully execute the
orders they had received, (for this was what they earnestly desired,
and what the Count and the Receiver-General instructed them to do,)
they seized upon the seats, the throne, and the table which was of
wood \footnote{vid. Fleury's Church History, xxii. 7. p.
129, note k. [Oxf. tr. 1843.] By specifying the material, Athan.
implies that altars were sometimes not of wood.

Return to text}, and the curtains \footnote{Curtains were at the entrance, and before the
chancel. vid. Bingh. Antiqu. viii. 6. §. 8. Hofman. Lex. in voc. \emph{velum}.
also Chrysost. Hom. iii. in Eph. [tr. p. 133, note o.]

Return to text} of the Church, and whatever
else they were able, and carrying them out burnt them before the doors
in the great street, and cast frankincense upon the flame. Alas! who
will not weep to hear of these things, and, it may be, close his ears
\footnote{p. 140 fin. vol. 8. p. 188 init.

Return to text}, that he may not have to endure the recital, esteeming it
hurtful merely to listen to the accounts of such enormities? Moreover
they sang the praises of their idols, and said, ``Constantius hath
become a heathen, and the Arians have acknowledged our customs;'' for
indeed they scruple not even to pretend heathenism, if only their
heresy may be established. They even were ready to sacrifice a heifer
which drew the water for the gardens at the Cæsareum \footnote{The royal quarter in Alexandria, vid. supr. p.
167, note P. In other Palatia an aqueduct was necessary, e.g. vid.
Cod. Theod. xv. 2. even at Daphne, though it abounded in springs,
ibid. l. 2.

Return to text}; and
would have sacrificed it, had it not been a female \footnote{vid. Herodot. ii. 41. who says that cows and
heifers were sacred to Isis. vid. Jablonski Pantheon Æg. i. 1. §.
15. who says that Isis was worshipped in the shape of a cow, and
therefore the cows received divine honours. Yet bulls were sacrificed
to Apis, ibid. iv. 2. §. 9. vid. also Schweighæuser \emph{in loc}.
Herod.

Return to text}; for they
said that it was unlawful for such to be offered among them.

§. 57.

13. Thus acted the impious \footnote{[\emph{dussebeis}], misbelieving, as passim.

Return to text} Arians in
conjunction with the heathens, thinking that these things tended to
our dishonour. But Divine justice reproved their iniquity, and wrought
a great and remarkable miracle \footnote{[\emph{s\={e}meion}], p. 217, r. 6.

Return to text}, thereby plainly shewing to
all men, that as in their acts of impiety \footnote{vid. vol. 8. p. 1, note 1. This is a
remarkable instance of the special and technical sense of the words, [\emph{eusebeia,
asebountes}], \&c. being here contrasted with pagan blasphemy,
\&c. vid. also p. 269, r. 1.

Return to text} they had dared to
attack none other but the Lord, so in these proceedings also, they
were again attempting to do dishonour unto Him. This was more
manifestly proved by the marvellous \footnote{[\emph{thaumatos}], p. 211, r. 1.

Return to text} event which now came to
pass. One of these licentious youths ran into the Church, and ventured
to sit down upon the throne; and as he sat there the wretched man
uttered with a nasal sound some lascivious song. Then rising up he
attempted to pull away the throne, and to drag it towards him; he knew
not that he was drawing down vengeance upon himself. For as of old the
inhabitants of Azotus, when they ventured to touch the Ark, which it
was not lawful for them even to look upon, were immediately destroyed
by it, being first grievously tormented by emerods; so this unhappy
person who presumed to drag the throne, drew it upon himself, and, as
if Divine justice had sent the wood to punish him, he struck it into
his own bowels; and instead of carrying out the throne, he brought out
by the blow his own entrails, so that the throne took away his life,
instead of his taking it away. For, as it is written of Judas, his
bowels gushed out, and he fell down and was carried away, and the day
after died. Another also entered the Church with boughs of trees, and,
as in the Gentile manner he waved them in his hands and mocked, he was
immediately struck with blindness, so as straightway to lose his
sight, and to know no longer where he was; but as he was about to
fall, he was taken by the hand and supported by his companions out of
the place, and when on the following day he was with difficulty
brought to his senses, he knew not either what he had done or suffered
in consequence of his audacity.

§. 58.

14. The Gentiles, when they beheld these things,
were seized with fear, and ventured on no further outrage; but the
Arians were not yet touched with shame, but, like the Jews when
they saw the miracles, were faithless and would not believe, nay, like
Pharaoh, they were hardened; they too having placed their hopes below,
on the Emperor and his eunuchs. They permitted the Gentiles, or rather
the more abandoned of the Gentiles, to act in the manner before
described; for they found that Faustinus, who is the Receiver-General
by style, but is a vulgar \footnote{[\emph{agoraion}], p. 268, note F.

Return to text} person in habits, and profligate
in heart, was ready to play his part with them in these proceedings,
and to stir up the heathen. Nay, they undertook to do the like
themselves, that as they had struck off their heresy from all other
heresies together \footnote{p. 244, r. 1.

Return to text}, so they might divide their wickedness
with the more depraved part of mankind. What they did through the
instrumentality of others I have described above; the enormities they
committed themselves, surpass the bounds of all wickedness; and they
exceed the vileness of any hangman \footnote{[\emph{d\={e}miou}], p. 247, r. 3.

Return to text}. Where is there a house
which they did not ravage? where is there a family they did not
plunder on pretence of searching for their opponents? where is there a
garden they did not trample under foot? what tomb \footnote{vid. Socr. Hist. iv. 13.

Return to text} they did
not open, pretending they were seeking for Athanasius, though their
sole object was to plunder and spoil all that came in their way? How
many men's houses were sealed up! From how many did they accept
hospitality to give it to the soldiers who assisted them! Who had not
experience of their wickedness? Who that met them in the market-place
but was obliged to hide himself? Did not many an one leave his house
from fear of them, and pass the night in the desert? Did not many an
one, while anxious to preserve his property from them, lose the
greater part of it? And who, however inexperienced, did not choose
rather to commit himself to the sea, and to risk all its dangers, than
to witness their threatenings? Many also changed their residences, and
removed from street to street, and from the city to the suburbs. And
many submitted to severe fines, and when they were unable to pay,
borrowed of others, merely that they might escape their machinations.

§. 59.

15. For they made themselves formidable to all
men, and treated all with great arrogance, using the name of the
Emperor, and threatening them with his displeasure. They had to assist them in their wickedness the Duke Sebastianus, a Manichee, and
a profligate young man; the Prefect, the Count, and the
Receiver-General to play his part. Many Virgins who condemned their
impiety, and professed the truth, they threw down from the houses;
others they insulted as they walked along the streets, and caused
their heads to be uncovered \footnote{p. 268, r. 8.

Return to text} by their young men. They also
gave permission to the females of their party to insult whom they
chose; and although the holy and faithful women withdrew on one side,
and gave them the way, yet they gathered round them like Bacchanals
and Furies \footnote{vid. vol. 8. p. 91, note q. also Greg. Naz.
Orat. 35. 3. Epiph. Hær. 69. 3. Theod. Hist. i. 3. (p. 730. ed.
Schulze.)

Return to text}, and esteemed it a misfortune if they found no
means to injure them, and spent that day sorrowfully on which they
were unable to do them some mischief. In a word, so cruel and bitter
were they against all, that all men called them hangmen \footnote{p. 271, r. 3.

Return to text},
murderers, lawless, intruders, evil-doers, and by any other name
rather than that of Christians \footnote{p. 208, note B.

Return to text}.

§. 60.

16. Moreover, imitating the savage practices of
Scythians \footnote{p. 253, r. 1. vid. Hofman. in voc. \emph{fin}.

Return to text}, they seized upon Eutychius the Sub-deacon, a man
who had served the Church honourably, and causing him to be scourged
on the back with a heathen whip, till he was at the point of death,
they demanded that he should be sent away to the mines; and not simply
to any mine, but to that of Phæno \footnote{The mines of Phæno lie almost in a direct
line between Petræ and Zoar, which is at the southern extremity of
the Dead Sea. They formed the place of punishment of Confessors in the
Maximinian Persecution, Euseb. de Mart. Pal. 7. and in the Arian
Persecution at Alexandria after Athan. Theod. Hist. iv. 19. p. 996. Phænon
was once the seat of a Bishopric, which sent a Bishop to the Councils
at Ephesus, the Ecumenical, and the Latrocinium. vid. Reland,
Palestine, pp. 951, 952. Montfaucon \emph{in loc}. Athan. Le Quien.
Or. Christ. t. 3. p. 745.

Return to text}, where even a condemned
murderer is hardly able to live a few days. And what was most
unreasonable in their conduct, they would not permit him even a few
hours to have his wounds dressed, but caused him to be sent off
immediately, saying, ``If this is done, all men will be afraid, and
henceforward will be on our side.'' After a short interval however,
being unable to accomplish his journey to the mine on account of the
pain of his wounds, he died on the way. He perished rejoicing, having
obtained the glory of martyrdom.

17. But the miscreants \footnote{[\emph{asebeis}].

Return to text} were not even
yet ashamed, but in the words of Scripture, \emph{having bowels without
mercy} [Prov x. 12.], they acted accordingly, and now again
perpetrated a devilish \footnote{[\emph{satanikon}], vol. 8. p. 9, note s.

Return to text} deed. When the people prayed them to
spare Eutychius and besought them for him, they caused four honourable
and free citizens to be seized, one of whom was Hermias who washed the
beggars' feet \footnote{[\emph{Hermeian, louonta tous anexodous}],
``Inauspicato verterat Hermantius, qui angiportos non pervios lavabat.''
Montfaucon, Coll. Nov. t. 2. p. xliii. who translates as above, yet
not satisfactorily, especially as there is no article before [\emph{louonta}].
Tillemont says, ``qui avait \emph{quelle charge} dans la police de la
ville,'' understanding by [\emph{anexodoi}] ``inclusi sive
incarcerati homines;'' whereas they are ``ii qui [\emph{ana tas exodous}]
in exitibus viarum, stipem cogunt.'' Montf. ibid. For the custom of
washing the feet, vid. Bingh. Antiqu. xii. 4. §. 10. Justinian in 1
Ep. ad Trin. v. 10.

Return to text}; and after scourging them very severely, the
Duke cast them into the prison. But the Arians, who are more cruel
even than Scythians \footnote{pp. 272, r. 4. 275, r. 4.

Return to text}, when they saw that they did not die
from the stripes they had received, complained of the Duke and
threatened, saying, ``We will write and tell the eunuchs, that he
does not flog as we wish.'' Hearing this he was afraid, and was
obliged to beat the men a second time; and they being beaten, and
knowing for what cause they suffered and by whom they had been
accused, said only, ``We are beaten for the sake of the Truth, but we
will not hold communion with the heretics; beat us now as thou wilt;
God will judge thee for this.'' The impious heretics \footnote{miscreants.

Return to text} wished
to expose them to danger in the prison, that they might die there; but
the people of God observing their time, besought him for them, and
after seven days or more they were set at liberty.

§. 61.

18. But the Arians, as being grieved at this,
again devised another yet more cruel and unholy deed; cruel in the
eyes of all men, but well suited to their antichristian heresy. Our
Lord commanded that we should remember the poor; He said, \emph{Sell that
ye have, and give alms} [Luke xii. 33.]; and again, \emph{I was an
hungred, and ye gave Me meat; I was thirsty, and ye gave Me drink; for
inasmuch as ye have done it unto one these little ones, ye have done
it unto Me} [Mat. xxv. 35, 40.]. But these men, as being in truth
opposed to Christ, have presumed to act contrary to His will in this
respect also. For when the Duke gave up the Churches to the Arians,
and the destitute persons and widows were unable to continue any
longer in them, the widows sat down in places which the Clergy
entrusted with the care of them appointed. And when the Arians saw
that the brethren readily ministered unto them and supported them,
they persecuted them also, beating them on the feet \footnote{p. 284, r. 10.

Return to text}, and
accused those who gave to them before the Duke. This was done by means
of a certain soldier named Dynamius. And it was well-pleasing to
Sebastian, for there is no mercy in the Manichæans; nay, it is
considered a hateful thing among them to shew mercy to a poor man.
Here then was a novel subject of complaint; and a new kind of court
now first invented by the Arians. Persons were brought to trial for
acts of kindness which they had performed; he who shewed mercy was
accused, and he who had received a benefit was beaten; and they wished
rather that a poor man should suffer hunger, than that he who was
willing to shew mercy should give to him. Such sentiments these modern
Jews, for such they are, have learned from the Jews of old, who when
they saw him who had been blind from his birth recover his sight, and
him who had been a long time sick of the palsy made whole, accused our
Lord who had bestowed these benefits upon them, and judged them to be
transgressors who had experienced his goodness \footnote{vid. de Decr. §. 1. tr. p. 3.

Return to text}.

§. 62.

19. Who was not struck with astonishment at these
proceedings? Who did not execrate both the heresy, and its defenders?
Who failed to perceive that the Arians are indeed more cruel than wild
beasts? For they had no prospect of gain \footnote{vid.
vol. 8. p. 191, note c.

Return to text} from their
iniquity, for the sake of which they might have acted in this manner;
but they rather increased the hatred \footnote{p. 248, r. 3.

Return to text} of all men against
themselves. They thought by treachery and terror to force certain
persons into their heresy, so that they might be brought to
communicate with them; but the event turned out quite the contrary.
The sufferers endured as martyrdom whatever they inflicted upon them,
and neither betrayed nor denied the true faith in Christ. And those
that were without and witnessed their conduct, and at last even the
heathen when they saw these things, execrated them as antichristian \footnote{[\emph{antichristous}].

Return to text}, as cruel executioners \footnote{[\emph{d\={e}mious}], p. 247, r. 2.

Return to text}; for human nature is prone
to pity and sympathise with the poor. But these men have lost even the
common sentiments of humanity; and that kindness which they
would have desired to meet with at the hands of others, had themselves
been sufferers, they would not permit others to receive, but employed
against them the severity and authority of the magistrates, and
especially of the Duke.

§. 63.

20. What they did to the Presbyters and Deacons;
how they drove them into banishment under sentence passed upon them by
the Duke and the Magistrates, causing the soldiers to throw down their
kinsfolk from the houses \footnote{p. 272 init.

Return to text}, and Gorgonius the commander of the
police \footnote{[\emph{strategou}], infr. p. 295, note B.

Return to text} to beat them with stripes; and how (most cruel act of
all) with much insolence they plundered the bread \footnote{A. [\emph{tous artous}], the word occurs above,
pp. 7, 192, 267. in this sense; but Nannius, Hermant, and Tillemont,
with some plausibility understand it as a Latin term naturalized, and
translate ``most cruel of all, with much insolence they tore the \emph{limbs}
of the dead,'' alleging that merely to take away \emph{loaves} was
not so ``cruel'' as to take away \emph{lives}, which the Arians had
done.

Return to text

} of these
and of those who were now dead; these things it is impossible for
words to describe, for their cruelty surpasses all the powers of
language. What terms could one employ which might seem equal to the
subject? What circumstances could one mention first, so that those
next recorded would not be found more dreadful, and the next more
dreadful still? All their attempts and iniquities \footnote{[\emph{aseb\={e}mata}].

Return to text} were full
of murder and impiety; and so unscrupulous and artful are they, that
they endeavour to deceive by promises of protection, and by bribing
with money \footnote{pp. 135, r. 1. 285, r. 2.

Return to text}, that so, since they cannot recommend themselves
by fair means, they may thereby appear to the simple to make some
show.

\xsection{Chapter 8. Persecution in Egypt}

§. 64.

1. W\textsc{ho} would call them even
by the name of Gentiles; much less by that of Christians \footnote{Margin Notes

1. p. 208, note B.

Return to text}?
Would any one regard their habits and feelings as human, and not
rather those of wild beasts, seeing their cruel and savage conduct?
They are more malignant than public hangmen \footnote{p. 274, r. 6.

Return to text}; more audacious
than all other heretics. To the Gentiles they are much inferior, and
stand far apart and separate from them \footnote{pp. 235, r. 4. 253, r. 1.

Return to text}. I have heard from our
fathers, and I believe their report to be a faithful one, that long
ago, when a persecution arose in the time of Maximian, the grandfather
of Constantius, the Gentiles concealed our brethren the Christians,
who were sought after, and frequently suffered the loss of their own
substance, and had trial of imprisonment, solely that they might not
betray the fugitives. They protected those who fled to them for
refuge, as they would have done their own persons, and were determined
to run all risks on their behalf. But now these admirable persons, the
inventors of a new heresy, act altogether the contrary part \footnote{p. 275 init.

Return to text},
and are distinguished for nothing, but their treachery. They have
appointed themselves as executioners \footnote{p. 274, r. 6.

Return to text}, and seek to betray all
alike, and make those who conceal others the objects of their plots,
esteeming equally as their enemy both him that conceals and him that
is concealed. So murderous are they; so emulous in their
evil-doings of the wickedness of Judas.

§. 65.

2. The crimes these men have committed cannot
worthily be described. I would only say, that as I write and wish to
enumerate all their deeds of iniquity, the thought enters my mind,
whether this heresy be not the fourth daughter of the horse-leach \footnote{p. 191, r. 2.

Return to text} in the Proverbs, since after so many acts of injustice, so
many murders, it hath not yet said, 'It is enough.' No; it still
rages, and goes about \footnote{[\emph{perierchetai}], p. 235, r. 2. ad Adelph.
§. 2 fin.

Return to text} seeking after those whom it has not yet
discovered, while those whom it has already injured, it is eager to
injure anew. After the midnight attack, after the evils committed in
consequence of it, after the persecution brought about by Heraclius,
they cease not yet to accuse us falsely before the Emperor, (and they
are confident that as impious persons they will obtain a hearing,)
desiring that something more than banishment may be inflicted upon us,
and that hereafter those who do not consent to their impieties may be
destroyed. Accordingly, being now emboldened in an extreme degree,
that most abandoned Secundus \footnote{p. 133, r. 2.

Return to text} of Pentapolis, and Stephanus
\footnote{p. 235.

Return to text} his accomplice, conscious that their heresy was a defence of
any injustice they might commit, on discovering a Presbyter at Barea
who would not comply with their desires, (he was called Secundus,
being of the same name, but not of the same faith with the heretic,)
they kicked till he died \footnote{In like manner the party of Dioscorus at the
Latrocinium, or the Eutychian Council of Ephesus, A.D. 449. kicked to
death Flavian, Patriarch of Constantinople.

Return to text}. While he was thus suffering he
imitated the Saint and said, ``Let no one avenge my cause before human
judges; I have the Lord for my avenger, for whose sake I suffer these
things at their hands.'' They however were not moved with pity at these
words, nor did they feel any awe of the sacred season; for it was
during the time of Lent \footnote{p. 7, note U.

Return to text} that they thus kicked the man to
death.

§. 66.

3. O new heresy, that hast put on the whole devil
in impiety and wicked deeds! For in truth it is but a lately invented
evil; and although certain heretofore appear to have adopted its
doctrines, yet they concealed them and were not known to hold them.
But Eusebius and Arius, like serpents coming out of their holes,
have vomited \footnote{de Syn. tr. p. 96. Orat. i. tr. p. 232.

Return to text} forth the poison
\footnote{Orat. i. tr. pp. 177, 189, 218.

Return to text} of this impiety;
Arius daring to blaspheme openly, and Eusebius defending his
blasphemy. He was not however able to support the heresy, until, as I
said before, he found a patron \footnote{p. 260, r. 2.

Return to text} for it in the Emperor. Our
fathers called an Ecumenical Council, when three hundred of them, more
or less, met together and condemned the Arian heresy, and all declared
that it was alien and strange to the faith of the Church \footnote{[\emph{ekkl\={e}siastik\={e}s}], Orat.
i. tr. p. 242, r. 4. Ep. Æg. §. 18 init. Vales. in Eus. Hist. ii.
25.

Return to text}.
Upon this its supporters, perceiving that they were dishonoured and
had now no good ground of argument to insist upon, devised a different
method, and attempted to vindicate it by means of external power \footnote{p. 279, r. 1.

Return to text}.

4. And herein one may especially admire the
novelty as well as wickedness of their device, and how they go beyond
all other heresies. For these support their fond \footnote{[\emph{t\={e}n manian}].

Return to text} inventions
by persuasive arguments calculated to deceive the simple; the Greeks,
as the Apostle has said, make their attack with sublime and enticing
words, and with plausible fallacies; the Jews, leaving the divine
Scriptures, now, as the Apostle again has said, contend about \emph{fables
and endless genealogies} [1 Tim. i. 4.]; and the Manichees and
Valentinians with them, and others, corrupting the divine Scriptures,
put forth fables in terms of their own invention. But the Arians are
bolder than them all, and have shewn that the other heresies are but
their younger sisters \footnote{p. 244.

Return to text}, whom, as I have said, they surpass in
impiety, emulating them all, and especially the Jews, in their
iniquity. For as the Jews, when they were unable to prove the charges
which they pretended to allege against Paul, straightway led him to
the chief captain and the governor; so likewise these men, who surpass
the Jews in their devices, make use only of the power of the judges;
and if any one so much as speaks against them, he is dragged before
the Governor or the General. §. 67. The other heresies also, when the
very Truth has refuted them on the clearest evidence, are wont to be
silent, being simply confounded by their conviction. But this modern
and accursed heresy, when it is overthrown by argument, when it is
cast down and covered with shame by the very Truth, forthwith
endeavours to reduce by violence and stripes and imprisonment those whom it has been unable to persuade by argument, thereby
acknowledging itself to be any thing rather than godly. For it is the
part of true godliness not to compel \footnote{The early theory about persecution seems to
have been this,—that that was a bad cause which \emph{depended} upon
it, but that, when a \emph{cause} was good, there was nothing wrong in
using force in due \emph{subordination} to argument; that there was as
little impropriety in the civil magistrate's inducing \emph{individuals}
by force, when they were incapable of higher motives, as by those
secular blessings which follow on Christianity. Our Lord's kingdom was
not of this world, that is, it did not depend on this world; but, as
subduing, engrossing, and swaying this world, it at times condescended
to make use of this world's weapons against itself. The simple
question was \emph{whether a cause depended enforce for its existence}.
St. Athanasius declared, and the event proved, that Arianism was so
dependent. When Emperors ceased to persecute, Arianism ceased to be;
it had no life in itself. Again, all cruel persecution, or long
continued, or on a large scale, was wrong, as arguing \emph{an absence}
of moral and rational grounds in the \emph{cause} so maintained.
Again, there was an evident \emph{impropriety} in ecclesiastical
functionaries using secular weapons, as there would be in their
engaging in a secular pursuit, or forming secular connections; whereas
the soldier might as suitably, and should as dutifully, defend
religion with the sword, as the scholar with his pen. And further
there was an abhorrence of cruelty natural to us, which it was a duty
to cherish and maintain. All this being considered, there is no
inconsistency in St. Athanasius denouncing persecution, and in
Theodosius decreeing that ``the heretical teachers, who usurped the
sacred titles of Bishops or Presbyters,'' should be ``exposed to the
heavy penalties of exile and confiscation.'' Gibbon, Hist. ch. 27. For
a list of passages from the Fathers on the subject, vid. Limborch on
the Inquisition, vol. 1. Bellarmin. de Laicis, c. 21. 22. and of
authors in favour of persecution, vid. Gerhard de Magistr. Polit. p.
741, \&c.

Return to text}, but to persuade, as I
said before \footnote{p. 257, r. 5.

Return to text}. Thus our Lord Himself, not as employing force,
but as offering to their free choice, has said to all, \emph{If any man
will follow after Me} [Mat. xvi. 24.]; and to His disciples, \emph{Will
ye also go away?} [John vi. 67.]

5. This heresy however is altogether alien from
godliness; and therefore how otherwise should it act, than contrary to
our Saviour, seeing also that it has enlisted that enemy of Christ
Constantius, as it were Antichrist himself \footnote{vid. vol. 8. p. 79, note q.

Return to text}, to be its leader
in impiety? He for its sake has earnestly endeavoured to emulate Saul
in savage cruelty. For when the priests gave victuals to David, Saul
commanded, and they were all destroyed, in number three hundred and
five \footnote{85 priests, rec. text.

Return to text}; and this man, now that all avoid the heresy, and
confess a sound faith in the Lord, overthrows a Council of full three
hundred Bishops, banishes the Bishops themselves, and hinders the
people from the practice of piety, and from their prayers to God,
preventing their public assemblies. And as Saul overthrew Nob, the
city of the priests, so this man, advancing even further in
wickedness, has given up the Churches to the impious. And as he
honoured Doeg the accuser before the true priests, and
persecuted David, giving ear to the Ziphites; so this man prefers
heretics to the godly, and even persecutes them that flee from him,
giving ear to his own eunuchs, who falsely accuse the orthodox. He
does not perceive that whatever he does or writes in behalf of the
heresy of the Arians, amounts to an attack upon his Saviour.

§. 68.

6. Ahab himself did not act so cruelly towards
the priests of God, as this man has acted towards the Bishops. For he
was at least pricked in his conscience when Naboth had been murdered,
and was afraid at the sight of Elias; but this man neither reverenced
the great Hosius, nor was wearied or pricked in conscience, after
banishing so many Bishops; but like another Pharaoh, the more he is
afflicted, the more he is hardened, and imagines greater wickedness
day by day. And the most extraordinary instance of his iniquity was
the following. It happened that when the Bishops were condemned to
banishment, certain other persons also received their sentence on
charges of murder or sedition or theft, each according to the quality
of his offence. These men after a few months he released, on being
requested to do so, as Pilate did Barabbas; but the servants of Christ
he not only refused to set at liberty, but even sentenced them to more
unmerciful punishment in the place of their exile, proving himself a
perpetual torment to them. To the others through congeniality of
disposition he became a friend; but to the orthodox he was an enemy on
account of their true faith in Christ. Is it not clear to all men from
hence, that the Jews of old when they demanded Barabbas, and crucified
the Lord, acted but the part which these present enemies of Christ are
acting together with Constantius? nay, that he is even more bitter
than Pilate. For Pilate when he perceived the injustice of the deed,
washed his hands; but this man, while he banishes the saints, gnashes
\footnote{infr. p. 284, r. 7. p. 207, r. 1.

Return to text} his teeth against them more and more.

§. 69.

7. But what wonder is it if, after he has been
led into impious errors, he is so cruel towards the Bishops, since the
common feelings of humanity could not induce him to spare even his own
kindred? his uncles \footnote{The brothers of Constantine were Julius
Constantius, and Dalmatius; of these Julius Constantius was father of
Gallus and Julian, and Dalmatius of Dalmatius and Hannibalianus. (vid.
supr. p.94, note S. p.108, note C.) Constantine had put his two
last-mentioned nephews almost on an equality with his three sons;
Dalmatius being a Cæsar, and Hannibalianus ``King,'' the only prince
with that title in any age of the Empire. On the Emperor's death some
of his great officers as well as the soldiers and people came to a
resolution that none but his sons should be their masters. Constantius
promised his kinsmen his protection under an oath; but Eusebius of
Nicomedia produced a last will of Constantine's, in which he declared
his suspicions that he had been poisoned by his brothers, and called
on his sons to avenge him. Vid. Gibbon, ch. 18. who continues, ``The
spirit, and even the forms of legal proceedings were repeatedly
violated in a promiscuous massacre; which involved the two uncles of
Constantius, seven of his cousins, of whom Dalmatius and Hannibalianus
were the most illustrious, the Patrician Optatus, who had married a
sister of the late Emperor, and the Prefect Ablavius, whose power and
riches had inspired him with some hope of obtaining the purple.'' p.
132. Constantius had married the daughter of his uncle Julius
Constantius, and had given his sister in marriage to his cousin
Hannibalianus. ``Of so numerous a family,'' continues Gibbon, ``Gallus
and Julian alone, the two youngest children of Constantius, were saved
from the hands of the assassins.'' Constantius married Gallus to his
sister, and made him Cæsar. Gallus abused his power, was recalled
from the seat of his government, and beheaded in prison. Olympias was
the daughter of Ablavius, who was betrothed to the Emperor Constans;
about the time of Ath.'s writing, Constantius married her to Arsaces,
king of Armenia. Amm. Marcell. xx. 11 init. We may suppose Athan. in
the text expresses the feeling of the day at this alliance, or of
Constantius's enemies. Arsaces was a Christian. St. Olympias was niece
to this Olympias. Tillem. Empereurs, t. 4. p 219.

Return to text} he slew; his cousins he put out of
the way; he commiserated not the sufferings of his father-in-law,
though he had married his daughter, or of his kinsmen; but he has ever
been a transgressor of his oath towards all. So likewise he treated
his brother in an unholy manner; and now he pretends to build his
sepulchre, although he delivered up to the barbarians his betrothed
wife Olympias, whom he had protected till his death, and had brought
up as his intended consort. Moreover he attempted to set aside his
wishes, although he boasts to be his heir \footnote{p. 264, note A.

Return to text}; for so he writes,
in terms which any one possessed but of a small measure of sense would
be ashamed of. But when I compare his letters, I find that he does not
possess common understanding, but that his mind is solely regulated by
the suggestions of others, and is by no means in his own power. Now
Solomon says, \emph{If a ruler hearken to lies, all his servants are
wicked} [Prov. xxix. 12.]. This man proves by his actions that he
is such an unjust one, and that those about him are wicked.

§. 70.

8. How then, being such an one, and taking
pleasure in such associates, can he ever design any thing just or
reasonable, entangled as he is in the iniquity of his followers, men
given to sorcery, who have trampled his brains under the soles
of their feet? Wherefore he now writes letters, and then repents that
he has written them, and after repenting is again stirred up to anger,
and then again laments his fate, and being undetermined what to do, he
shews a soul destitute of understanding. Being then of such a
character, one would rather pity him, because that under the semblance
and name of freedom he is the slave of those who drag him on to
gratify their own impious pleasure. In a word, while through his folly
and inconstancy, as the Scripture saith, he is willing to comply with
the desires of others, he has given himself up to condemnation, to be
consumed by fire in the future judgment; at once consenting to do
whatever they wish, and gratifying them in their designs against the
Bishops, and in their exertion of authority over the Churches.

9. For behold, he has now again thrown into
disorder all the Churches of Alexandria and of Egypt and Libya, and
has publicly given orders, that the Bishops of the Catholic Church and
faith be cast out of them, and that they be given up to the professors
of the Arian doctrines. The General began to carry this order into
execution; and straightway Bishops were sent off in chains, and
Presbyters and Monks bound with iron, after being almost beaten to
death with stripes. Disorder prevails in every place; all Egypt and
Libya are in danger, the people being indignant at this unjust
command, and seeing in it the preparation for the coming of
Antichrist, and beholding their property plundered by others, and
given up into the hands of the heretics.

§. 71.

10. When was ever such iniquity heard of? when
was such an evil deed ever perpetrated, even in times of persecution?
They were heathens who persecuted formerly; but they did not bring
their idols into the Churches. Zenobia was a Jewess, and a supporter
of Paul of Samosata; but she did not give up the Churches to the Jews
for Synagogues. This is a new piece of iniquity \footnote{[\emph{musos}].

Return to text}. It is not
simply persecution, but more than persecution, it is a prelude and
preparation \footnote{vol. 8. p. 79, note q.

Return to text} for the coming of Antichrist. Even if it be
admitted \footnote{[\emph{est\={o}}], p. 221, r. 5.

Return to text} that they invented false charges against Athanasius
and the rest of the Bishops whom they banished, yet what is this to
their later practices? What charges have they to allege against
the whole of Egypt and Libya and Pentapolis \footnote{p. 221, §. 3.

Return to text}? For they have
begun no longer to lay their plots against individuals, in which case
they might be able to frame a lie against them; but they have set upon
all in a body, so that, however they may wish to invent accusations
against them, they must be condemned. Thus their wickedness has
blinded their understanding; and they have required, without any
reason assigned, that the whole body of the Bishops shall be expelled,
and thereby they shew that the charges they framed against Athanasius
and the rest of the Bishops whom they banished were false, and
invented for no other purpose than to support the accursed heresy of
the Arian enemies of Christ.

11. This is now no longer concealed, but has
become most manifest to all men. He commanded Athanasius to be
expelled out of the city, and gave up the Churches to them. And the
Presbyters and Deacons that were with him, who had been appointed by
Peter and Alexander, were also expelled and driven into banishment;
and the real Arians, who not through any suspicions arising from
circumstances \footnote{[\emph{ex\={o}then}].

Return to text}, but on account of the heresy had been
expelled at first together with Arius himself by the Bishop Alexander,
Secundus in Libya, in Alexandria Euzöius \footnote{infr. Dep.
Ar.

Return to text} the Chananean,
Julius, Ammon, Marcus, Irenæus, Zozimus, and Serapion surnamed
Pelycon, and in Libya Sisinnius, and the younger men with him,
associates in his impiety; these obtained possession of the Churches.
§. 72. And the General Sebastian wrote to the governors and military
authorities in every place; and the true Bishops were persecuted, and
those who professed impious doctrines were brought in in their stead.
They banished Bishops who had grown old in orders \footnote{[\emph{kl\={e}r\={o}i}].

Return to text}, and had
been many years in the Episcopate, having been ordained by the Bishop
Alexander; Ammonius \footnote{p. 193.

Return to text}, Hermes, Anagamphus, and Marcus, they
sent to the Upper Oasis; Muis, Psenosiris, Nilammon, Plenes, Marcus,
and Athenodorus to Ammoniaca, with no other intention than that they
should perish in their passage through the deserts. They had no pity
on them though they were suffering from disease, and indeed proceeded
on their journey with so much difficulty on account of their weakness,
that they were obliged to be carried in litters, and their
sickness was so dangerous that the materials for their burial
accompanied them. One of them indeed died, but they would not even
permit the body to be given up to his friends for interment \footnote{pp. 228, 193, r. 2.

Return to text}.
With the same purpose they banished also the Bishop Dracontius \footnote{p. 193.

Return to text} to the desert places about Clysma, Philo to Babylon, Adelphius to
Psinabla in the Thebais, and the Presbyters Hierax and Dioscorus to
Syene. They likewise drove into exile Ammonius, Agathus, Agathodæmon,
Apollonius, Eulogius, Apollo, Paphnutius, Gaius, and Flavius, ancient
\footnote{[\emph{archaiois}], p. 255, r. 5.

Return to text} Bishops, as also the Bishops Dioscorus, Ammonius, Heraclides,
and Psais; some of whom they gave up to work in the stone-quarries,
others they persecuted with an intention to destroy, and many others
they plundered.

12. They banished also forty of the laity, with
certain virgins whom they had before exposed to the fire \footnote{p. 192.

Return to text};
beating them so severely with rods taken from the palm-tree, that
after lingering five days some of them died, and others had recourse
to medical treatment on account of the thorns left in their limbs,
from which they suffered torments worse than death \footnote{p. 193. of the 40 men.

Return to text}. But what
is most dreadful to the mind of any man of sound understanding, though
characteristic of these miscreants \footnote{misbelievers.

Return to text}, is this: When the
Virgins during the scourging called upon the Name of Christ, they
gnashed their teeth against them with increased fury \footnote{p. 280, r. 1.

Return to text}. Nay
more, they would not give up the bodies of the dead to their friends
for burial, but concealed them that they might appear to be ignorant
of the murder. They did not however escape detection; the whole city
perceived it, and all men withdrew from them as executioners \footnote{p. 275, r. 4.

Return to text}, as malefactors and robbers. Moreover they overthrew monasteries
\footnote{[\emph{monast\={e}ria}].

Return to text}, and endeavoured to cast the Monks into the fire; they
plundered houses, and breaking into the house of certain free citizens
where the Bishop had deposited a treasure, they plundered and took it
away. They scourged the widows on the soles of their feet \footnote{p. 274, r. 1.

Return to text},
and hindered them from receiving their alms.

§. 73.

13. Such were the iniquities practised by the
Arians; and as to their further deeds of impiety, who could hear the
account of them without shuddering? They had caused these venerable
old men and aged Bishops to be sent into banishment; they now
appointed in their stead profligate heathen youths, whom they thought
to raise at once to the highest dignity, though they were not even
Catechumens \footnote{vid. Hallier de Ordin. part 2. i, l. art. 2.

Return to text}. And others who were accused of bigamy
\footnote{[\emph{digunaiois}], [\emph{digamois}]. on the
latter, vid. Suicer, Thes. in voc. [\emph{digamia}]. Tertull. Works,
tr. vol. i. p. 419, note N.

Return to text},
and even of worse crimes, they nominated Bishops on account of the
wealth and civil power which they possessed, and sent them out as it
were from a market, upon their giving them gold \footnote{p. 5, r. 1. p. 135, r. 1.

Return to text}. And now
more dreadful calamities befel the people. For when they rejected
these mercenary dependents of the Arians, so alien from themselves,
they were scourged, they were proscribed, they were shut up in prison
by the General, (who did all this readily, being a Manichee,) in order
that they might no longer seek after their own Bishops, but be forced
to accept those whom they abominated, men who were now guilty of the
same mockeries as they had before practised among their idols.

§. 74.

14. Will not every just person break forth into
lamentations at the sight or hearing of these things, at perceiving
the arrogance and extreme injustice of these impious men? \emph{The
righteous lament in the place of the impious} [Prov. xxviii. 28.
Sept.]. After all these things, and now that the impiety has reached
such a pitch of audacity, who will any longer venture to call this
Costyllius \footnote{An irregularly formed diminutive, or a quasi
diminutive from Constantius, as Agathyllus from Agathocles, Heryllus
from Heracles, \&c. vid. Matth. Gr. Gramm. §. 102. ed. 1820.

Return to text} Christian, and not rather the image of Antichrist?
For what mark of Antichrist is yet wanting to him? How can he in any
way fail to be regarded as he? or how can the latter fail to be
supposed such a one as he is? Did not the Arians and the Gentiles
offer those sacrifices in the great Church in the Cæsareum \footnote{p. 269, note I.

Return to text},
and utter their blasphemies against Christ as by His command? And does
not the vision of Daniel thus describe Antichrist; that he shall make
war with the saints, and prevail against them, and exceed all that
have been before him in evil deeds, and shall humble three kings, and
speak words against the Most High, and shall think to change times and
laws? Now what other person besides Constantius has ever attempted to
do these things? He is surely such a one as Antichrist would be. He
speaks words against the Most High by supporting this impious
heresy: he makes war against the saints by banishing the Bishops;
although indeed he exercises this power but for a little while \footnote{Short lives are generally considered the
destiny of the Church's persecutors, and length of days the token of
her protectors. What of old was said of pain, applies to
persecution—si gravis, brevis, Antichrist's oppression seems to be
marked out as three years and a half. Constantius died at 45. having
openly apostatized for about six years. Julian died at 32, after a
reign of a year and a half. vid. supr. p. 245, r. 4. vid. also
Bellarmin. de Notis Eccl. 17. and 18.

Return to text} to his own destruction. Moreover he has surpassed those before him
in wickedness, having devised a new mode of persecution; and after he
had overthrown three kings, namely Vetranio, Magnentius, and Gallus,
he straight-way undertook the patronage \footnote{p. 278, r. 3.

Return to text} of impiety; and like
a giant \footnote{vid. de Decr. §. 32. tr. p. 58, note m. Orat.
ii. §. 32. Naz. Orat. 43, 26. Socr. Hist. v. 10. p. 268.

Return to text} he has dared in his pride to set himself up against
the Most High.

15. He has thought to change laws, by
transgressing the ordinance of the Lord given us through His Apostles,
by altering the customs of the Church, and inventing a new kind of
ordinations. For he sends from strange places distant a fifty days'
journey \footnote{p. 133, r. 10.

Return to text}, Bishops attended by soldiers to people unwilling to
receive them; and instead of an introduction to the acquaintance of
their people, they bring with them threatening messages, and letters
to the magistrates. Thus he has sent Gregory from Cappadocia to
Alexandria; he has transferred Germinius \footnote{vol. 8. p. 74, note e.

Return to text} from Cyzicus to
Sirmium; he has removed Cecropius \footnote{p. 133, r. 8.

Return to text} from Laodicea to
Nicomedia. §. 75. Again he transferred from Cappadocia to Milan one
Auxentius \footnote{vol. 8. p. 82, note x.

Return to text}, a man pragmatical rather than Christian, whom he
commanded to stay there after he had banished for his piety towards
Christ, Dionysius the Bishop of the place, a godly man. But this
person was as yet even ignorant of the Latin language, and unskilful
in every thing except impiety. And now one George a Cappadocian, who
was contractor of stores \footnote{[\emph{hypodekt\={e}n}], pp. 192, 264,
note B.

Return to text} at Constantinople, and having
embezzled all monies that he received, was obliged to fly, he
commanded to enter Alexandria with military pomp, and supported by the
authority of the General. And he, finding there one Epictetus \footnote{Epictetus is mentioned above, p. 133, where he
is called [\emph{hypokrit\={e}s}], which after Montfaucon was
translated ``stage-player.'' It is a question, however, especially
considering the correspondence between that passage and the present,
whether more than 'actor' is meant by it, alluding to the mockery of
an ordination in which he seems to have taken part. Though an Asiatic
apparently by birth, he was made Bishop of Civita Vecchia. We hear of
him at the conference between Constantius and Liberius. Theod. Hist.
ii. 13. Then he assists in the ordination of Felix. Afterwards he made
a martyr of S. Ruffinian by making him run before his carriage; and he
ends his historical career by taking a chief part among the Arians at
Ariminum, vid. Tillem. t. 6. p. 380, \&c. Ughell. Ital. t. 10. p.
56.

Return to text}
a novice, a bold young man, made him his friend \footnote{The Greek is [\emph{Epikt\={e}ton tina … ne\={o}teron
… \={e}gap\={e}sen, hor\={o}n k.t.l}.] So in the account
of the [\emph{neaniskos, Ho de I\={e}sous emblepsas aut\={o}i, \={e}gap\={e}sen
auton}]. Mark x. 21.

Return to text},
perceiving that he was ready for any wickedness; and by his means he
carries on his designs against those of the Bishops whom he desires to
ruin. For he is prepared to do every thing that the Emperor wishes;
who accordingly availing himself of his assistance, has committed at
Rome a strange act, but one truly resembling the malice of Antichrist.
Having made preparations in the Palace instead of the Church, and
caused some three of his own eunuchs to attend instead of the people,
he then compelled three \footnote{i.e. to keep up the form of the canonical
number; and so a century earlier, in the case of Novatian, in the same
see, while the capital was still heathen, we read in Eusebius that he
brought from some obscure part of Italy ``three Bishops,'' ``rustic and
ignorant,'' who after a full meal, when they were not themselves,
consecrated him. Hist. vi. 43. On the custom itself, vid. Bingh.
Antiqu. ii. 11. §. 4.

Return to text} ill-conditioned spies
\footnote{pp. 221 fin. 263, r. 5.

Return to text}, (for
one cannot call them Bishops,) to ordain forsooth as Bishop one Felix
\footnote{This Felix has been in after times accounted a
true Pope and Martyr, and has been supposed to have condemned
Constantius. The circumstances will be found in Tillemont, Mem. t. 6.
p. 778. Bolland. Catal. Pontif. Gibbon, ch. 21. p. 390.

Return to text}, a man worthy of them, then in the Palace. For the people
perceiving the iniquitous proceedings of the heretics would not allow
them to enter the Churches, and withdrew themselves entirely from
them.

§. 76.

16. Now what is yet wanting to make him
Antichrist? or what more could Antichrist do at his coming than this
man has done? Will he not find when he comes that the way has been
already prepared for him by this man easily to deceive the people?
Again, he claims to himself the right of deciding causes, which he
refers to the Court instead of the Church, and presides at them in
person. And strange it is to say, when he perceives the accusers at a
loss, he takes up the accusation himself, so that the injured party
may no longer be able to defend himself on account of the violence
which he displays. This he did in the proceedings against Athanasius.
For when he saw the boldness of the Bishops Paulinus, Lucifer,
Eusebius, and Dionysius, and how out of the recantation of
Ursacius and Valens they confuted those who spoke against the Bishop,
and advised that Valens and his associate should no longer be believed
since they had already retracted what they now asserted, he
immediately stood up and said, ``I am now the accuser of Athanasius; on
my account you must believe what these assert.'' And then, when they
said,—``But how can you be an accuser, when the accused person is not
present? and if you are his accuser, yet he is not present, and
therefore cannot be tried. And the cause is not one that concerns
Rome, so that you should be believed as being the Emperor; but it is a
matter that concerns a Bishop; and the trial ought to be conducted on
equal terms both to the accuser and the accused. And besides, how can
you accuse him? for you could not be present to witness the conduct of
one who lived at so great a distance from you; and if you speak but
what you have heard from these, you ought also to give credit to what
he says; but if you will not believe him, while you do believe them,
it is plain that they assert these things for your sake, and accuse
Athanasius only to gratify you \footnote{p. 267, r. 4.

Return to text}?''—when he heard this,
thinking that what they had so truly spoken was an insult to himself,
he sent them into banishment; and being exasperated against
Athanasius, he wrote in a more savage strain, requiring that he should
suffer what has now befallen him, and that the Churches should be
given up to the Arians, and that they should be allowed to do whatever
they pleased.

§. 77.

17. Terrible indeed, and worse than terrible are
such proceedings; and yet is this conduct suitable to him who
represents the character of Antichrist. Who that beheld him bearing
sway over his pretended Bishops, and presiding in Ecclesiastical
causes, would not justly exclaim that this was the abomination of
desolation spoken of by Daniel? For having put on the profession of
Christianity, and entering into the holy places, and standing therein,
he lays waste the Churches, transgressing their Canons, and enforcing
the observance of his own decrees. Will any one now venture to say
that this is a peaceful time with Christians, and not a time of
persecution? A persecution indeed, such as never arose before, and
such as no one perhaps will again stir up, except \emph{the son of
lawlessness} [2 Thess. ii. 8.], do these enemies of Christ exhibit,
who already present a picture of him in their own persons. Wherefore
it especially behoves us to be sober, lest this heresy which has
reached such a height of impudence, and has diffused itself abroad
like the \emph{poison of an adder} [Ps. lviii. 4.], as it is written
in the Proverbs, and which teaches doctrines contrary to the Saviour;
lest, I say, this be that \emph{falling away} [2 Thess. ii. 3.], after
which He shall be revealed, of whom Constantius is surely the
forerunner \footnote{[\emph{prodromos}], vid. vol. 8. p. 79, note
q.

Return to text}. Else wherefore is he so mad against the godly?
wherefore does he contend for it as his own heresy, and call every one
his enemy who will not comply with the madness of Arius, and admit
gladly the allegations of the enemies of Christ, and dishonour so many
venerable Councils? why did he command that the Churches should be
given up to the Arians? was it not that, when that other comes, he may
thus find a way to enter into them, and may take to himself him who
has prepared those places for him?

18. For the ancient Bishops who were ordained by
Alexander, and by his predecessor Achilles, and by Peter before him,
have been cast out; and those introduced whom the companions of
soldiers nominated; and they nominated only such as promised to adopt
their doctrines. §. 78. This was an easy proposition for the
Meletians to comply with; for the greater part, or rather the whole of
them, have never had a religious education, nor are they acquainted
with the \emph{sound faith} \footnote{p. 149, r. 3.

Return to text} in Christ, nor do they know at
all what Christianity is, or what writings we Christians possess. For
having come out, some of them from the worship of idols, and others
from the senate, or from the first civil offices, for the sake of the
miserable exemption \footnote{pp. 84, 85.

Return to text} from duty and for the patronage they
gained, and having bribed \footnote{pp. 89, 151, 291.

Return to text} the Meletians who preceded them,
they have been advanced to this dignity even before they were
Catechumens. And even if they pretended to have been such, yet what
kind of instruction \footnote{catechising.

Return to text} is to be obtained among the Meletians?
But indeed without even pretending to have been instructed, they came
at once, and immediately were called Bishops, just as children receive
a name. Being then persons of this description, they thought the thing
of no great consequence, nor even supposed that piety \footnote{[\emph{eusebeia}].

Return to text}
was different from impiety. Accordingly from being Meletians they
readily and speedily became Arians; and if the Emperor should command
them to adopt any other profession, they are ready to change again to
that also. Their ignorance of true godliness \footnote{[\emph{eusebeia}].

Return to text} quickly brings
them to submit to the prevailing folly, and that which happens to be
first taught them. For it is nothing to them to be carried about by
every wind and tempest, so long as they are only exempt from duty, and
obtain the patronage of men; nor would they care probably to change
again \footnote{pp. 88, 92.

Return to text} to what they were before, even to become such as they
were when they were heathens.

19. Any how, being men of such an easy temper,
and considering the Church as a civil senate \footnote{[\emph{politeian boul\={e}s}].

Return to text}, and like
heathen, being infected with the worship of idols, they have put on
the honourable name of our Saviour, under which they have polluted the
whole of Egypt, were it only that they have caused the name of the
Arian heresy to be known therein. For Egypt has heretofore been the
only country, throughout which the profession of the orthodox faith
was boldly maintained \footnote{p. 81.

Return to text}; and therefore these misbelievers have
striven to introduce jealousy there also, or rather not they, but the
Devil who has stirred them up, in order that when his herald
Antichrist shall come, he may find that the Churches in Egypt also are
his own, and that the Meletians have already been instructed in his
principles, and may recognise himself as already formed \footnote{[\emph{morph\={o}thenta}].

Return to text} in
them. §. 79. Such is the effect of that iniquitous \footnote{[\emph{paranomon}].

Return to text} order
which was issued by Constantius. On the part of the people there was
displayed a ready alacrity to submit to martyrdom, and an increased
hatred of this most impious heresy; and yet lamentations for their
Churches, and groans burst from all, while they cried unto the Lord, ``Spare
Thy people, O Lord, and give not Thine heritage unto Thine enemies to
reproach [Joel ii. 17.]; but make haste to deliver us out of the hand
of the lawless \footnote{[\emph{anom\={o}n}], 2 Thess. ii. 8.

Return to text}. For behold, they have not spared Thy
servants, but are preparing the way for Antichrist.''

20. For the Meletians will never resist him, nor
will they care for the truth, nor will they esteem it an evil thing to
deny Christ. They are men who have not approached the Lord with
sincerity; like the chameleon \footnote{vol. 8. p. 2. note c.

Return to text} they assume every various
appearance; they are hirelings \footnote{p. 289, r. 4.

Return to text} of any who will make use of
them. They make not the truth their aim, but prefer before it their
present pleasures; they say only, \emph{Let us eat and drink, for
tomorrow we die} [1 Cor. xv. 32.]. Such a profession and faithless
temper is more worthy of the Epicritian \footnote{histrionum genus.
Montf.

Return to text} players than of the
Meletians. But the faithful servants of our Saviour, and the true
Bishops who believe with sincerity, and live not for themselves, but
for the Lord; they faithfully believing in our Lord Jesus Christ, and
knowing, as I said before, that the charges which were alleged against
the truth were false, and plainly fabricated for the sake of the Arian
heresy, (for by the recantation \footnote{p. 86.

Return to text} of Ursacius and Valens they
detected the calumnies which were devised against Athanasius, for the
purpose of removing him out of the way, and of introducing into the
Churches the impieties of the enemies of Christ;) they, I say,
perceiving all this, as defenders and preachers of the truth, chose
rather, and endured to be insulted and driven into banishment, than to
subscribe against him, and to hold communion with the Arian fanatics.
They forgot not the lessons they had taught to others; yea, they know
well that great dishonour remains for the traitors, but for them which
confess the truth, the kingdom of heaven \footnote{supr. p. 213, r. 1.

Return to text}; and that to the
careless and such as fear Constantius will happen no good thing; but
for them that endure tribulations here, as sailors reach a quiet haven
after a storm, as wrestlers receive a crown after the combat, so these
shall obtain great and eternal joy and delight in heaven;—such as
Joseph obtained after his tribulations; such as the great Daniel had
after his temptations and the manifold conspiracies of the courtiers
against him; such as Paul now enjoys having received a crown from his
Saviour; such as the people of God every where expect. They, seeing
these things, were not infirm of purpose, but strong in faith, and
increased in their zeal more and more. Being fully persuaded of the
calumnies and impieties of the heretics, they condemn the persecutor,
and in heart and mind run together the same course with them that are
persecuted, that they also may obtain the crown of Confession.

§. 80.

21. One might say much more against this accursed
and antichristian heresy, and might demonstrate by many
arguments that the practices of Constantius are a prelude to the
coming of Antichrist. But seeing that, as the Prophet has said, from
the feet even to the head there is no soundness in it, but it is full
of all filthiness and all impiety, so that the very name \footnote{p. 138.

Return to text} of
it ought to be avoided as a dog's vomit or the poison of serpents; and
seeing that Costyllius openly exhibits the image of the adversary \footnote{[\emph{antikeimenou}], 2 Thess. ii. 4.

Return to text}; in order that our words may not be too many, it will be
well to content \footnote{[\emph{kalon arkesth\={e}nai, touto arkei}],
and so [\emph{\={e}rkei men gar}], Apol. contr. Ar. 2 init. [\emph{hikana
men oun tauta}]. de Decr. 15 init. [\emph{kai \={e}rkei men tauta}],
de Sent. D. 4 init. [\emph{arkei gar autous}], Apol. de Fug. 1 fin. [\emph{hikana
men oun tauta}], ibid. 24 init. [\emph{hikanon men oun touto}], ad
Serap. de M. A. 5 init. [\emph{esti men oun touto hikanon}], Orat. i.
17. [\emph{hikana men oun}], Ep. ad Serap. iii. 2 init. [\emph{arkei
tauta}]. ad Serap. iv. 7 init. [\emph{arkei hoti}], ad Epict. Vid.
also Orat. i. 7. B. Orat. ii init. Orat. iii. 47. Ep. Æg. 9 init. ad
Serap. iv. 1 init. ad Max. 5. \&c.

Return to text} ourselves with the divine Scripture, and that
we all obey the precept \footnote{supr. p. 148.

Return to text} which it has given us both in regard
to other heresies, and especially respecting this. That precept is as
follows; \emph{Depart ye, depart ye, go ye out from thence, touch no
unclean thing; go ye out of the midst of them, and be ye clean, that
bear the vessels of the Lord} [Is. lii. 11.]. This may suffice [Note
N] to instruct us all, so that if any one has been deceived by
them, he may go out from them, as out of Sodom, and not return again
unto them, lest he suffer the fate of Lot's wife; and if any one has
continued from the beginning pure from this impious heresy, he may
glory in Christ and say, \emph{We have not stretched out our hands to a
strange god} [Ps. xliv. 20.]; neither have we worshipped the works
of our own hands, nor served the creature \footnote{supr. p. 141, r. 1.

Return to text}, more than Thee,
the God that hast created all things through Thy Word, the
Only-begotten Son our Lord Jesus Christ, through whom to Thee the
Father together with the same Word in the Holy Spirit be glory and
power for ever and ever. Amen.''

The Second Protest \footnote{A. Of the two Protests referred to supr. p. 263,
the first was omitted by the copyists, as being already contained, as
Montfaucon seems to say, in the Apology against the Arians; yet if it
be the one to which allusion is made in the beginning of the Protest
which follows, it is not found there, nor does it appear what document
of A.D. 356. could properly have a place in a set of papers which
end with A.D.
350.

Return to text}

———————

§. 81.

1. The people of the Catholic Church in
Alexandria, which is under the government of the most Reverend Bishop
Athanasius, make this public protest by those whose names are
underwritten.

We have already protested against the nocturnal
assault which was committed upon ourselves and the Lord's house \footnote{[\emph{kuriakon}].

Return to text}; although in truth there needed no protest in respect to
proceedings with which the whole city has been already made
acquainted. For the bodies of the slain which were discovered were
exposed in public, and the bows and arrows and other arms found in the
Lord's house loudly proclaim the iniquity.

2. But whereas after our Protest already made,
the most illustrious Duke Syrianus endeavours to force all men to
agree with him, as though no tumult had been made, nor any had
perished, (wherein is no small proof that these things were not done
according to the wishes of the most gracious Emperor Augustus
Constantius; for he would not have been so much afraid of the
consequences of this transaction, had he acted therein by command;)
and whereas also, when we went to him, and requested him not to do
violence to any, nor to deny what had taken place, he ordered us,
being Christians, to be beaten with clubs; thereby again giving proof
of the nocturnal assault which has been directed against the
Church:—

We therefore make also this present Protest,
certain of us being now about to travel to the most religious Emperor
Augustus: and we adjure Maximus the Prefect of Egypt, and the
Controllers \footnote{Curiosi, p. 105, note A.

Return to text}, in the name of Almighty God, and for the sake
of the salvation of the most religious Augustus Constantius, to relate
all these things to the piety of Augustus, and to the authority of the
most illustrious Prefects \footnote{i.e. the
Prætorian.

Return to text}. We adjure also all the masters of
vessels, to publish these things every where, and to carry them to the
ears of the most religious Augustus, and to the Prefects and the
Magistrates in every place, in order that it may be known that a war
has been waged against the Church, and that, in the times \footnote{[\emph{kairois}], supr. p. 179, r. 2.

Return to text} of
Augustus Constantius, Syrianus has caused Virgins and many others to
become martyrs.

3. As it dawned upon the fifth before the Ides of
February \footnote{Febr. 9.

Return to text}, that is to say, the fourteenth of the month Mechir,
while we were keeping vigil \footnote{supr. pp. 176 init. 206.

Return to text} in the Lord's house, and engaged in our
prayers (for there was to be a communion on the Preparation \footnote{Friday, vid. p. 7, note I.

Return to text}); suddenly about midnight, the most illustrious Duke Syrianus
attacked us and the Church with many legions of soldiers \footnote{i.e. more than 5000, p. 206.

Return to text}
armed with naked swords and javelins and other warlike instruments,
and wearing helmets on their heads; and even while we were praying,
and while the lessons were being read, they broke down the doors. And
when the doors were burst open by the violence of the multitude, he
gave command, and some of them shot their arrows; others shouted;
their arms rattled, and their swords flashed in the light of the
lamps; and forthwith the Virgins were slain, many men were trampled
down, and fell over one another as the soldiers came upon them, and
several were pierced with arrows and perished. Some of the soldiers
also betook themselves to plundering, and stripped the Virgins naked,
who were more afraid of being even touched by them than they were of
death.

4. The Bishop continued sitting upon his throne,
and exhorted all to pray. The Duke led on the attack, having with him
Hilarius the notary, whose part in the proceedings was shown in the
sequel. The Bishop was seized, and hardly escaped being torn to
pieces; and having fallen into a state of insensibility, and appearing
as one dead, he disappeared from among them, and has gone we know not
whither. They were eager to kill him. And when they saw that many had
perished, they gave orders to the soldiers to remove out of
sight the bodies of the dead. But the most holy Virgins who were left
there were buried in the tombs, having attained the glory of martyrdom
in the times \footnote{p. 294, r. 2.

Return to text} of the most religious Constantius. Deacons also
were beaten with stripes even in the Lord's house, and were shut up
there.

5. Nor did matters stop even here: for after all
this had happened, whosoever pleased broke open any door that he
could, and searched, and plundered what was within. They entered even
into those places, which not even all Christians are allowed to enter.
Gorgonius the commander of the city force \footnote{B. [\emph{strat\={e}gou}]. There were two [strategoi]
or duumvirs at the head of the police force at Alexandria; they are
mentioned in the plural in Euseb. vii. 11. where S. Dionysius speaks
of their seizing him. We read of them at Philippi in Luke xvi. 35.
vid. Vales. \emph{in loc}. Euseb. et in Amm. Marc. xxxi. 6. The word
is translated in the Justinian Code, Prætor. vid. Du Cange, Gloss. Græc.
\emph{in voc}.

Return to text} knows this, for he was
present. And no unimportant evidence of the nature of this hostile
assault is afforded by the circumstance, that the armour and javelins
and swords borne by those who entered were left in the Lord's house.
They have been hung up in the Church until this time, that they might
not be able to deny it: and although they sent several times Dynamius
the soldier \footnote{[\emph{ton t\={e}s taxe\={o}s}], i.e.
supr. p. 274. [\emph{strati\={o}ton}].

Return to text}, as well as the Commander of the city police,
desiring to take them away, we would not allow it, until the
circumstance was known to all.

6. Now if an order has been given that we should
be persecuted, we are all ready to suffer martyrdom. But if it be not
by order of Augustus, we desire Maximus the Prefect of Egypt and all
the city magistrates to request of him that they may not again be
suffered thus to assail us. And we desire also that this our petition
may be presented to him, that they may not attempt to bring in hither
any other Bishop: for we have resisted unto death \footnote{pp. 63, 81.

Return to text}, desiring
to have the most Reverend Athanasius, whom God gave us at the
beginning, according to the succession of our fathers; whom also the
most religious Augustus Constantius himself sent to us with letters
and oaths. And we believe that when his Piety is informed of what has
taken place, he will be greatly displeased, and will do nothing
contrary to his oath, but will again give orders that our Bishop
Athanasius shall remain with us.

To the Consuls to be elected \footnote{C. Since the Consuls came into office on the
first of January, and were proclaimed in each city, vid. p. 153, note
M, it is strange that the Alexandrians here speak in February as if
ignorant of their names. The phrase, however, is found elsewhere. Thus
in this very year the Anonymous Maffeianus, (who is spoken of in the
Preface of this Volume,) dates Jan. 5. as ``post Consulatum Arbitionis
et Loliani.'' And in Socr. Hist. ii. 29. in the instance of the year
351, when there were no Consuls, and in 346, when there was a
difference on the subject between the Emperors who were eventually
themselves Consuls, the first months are dated in like manner from the
Consuls of the foregoing year.

Return to text

} after the Consulship
of the most illustrious Arbæthion and Collianus \footnote{Lollianus.

Return to text}; on the
seventeenth Mechir, which is the day before the Ides of February \footnote{Febr. 12.

Return to text}.

\xsection{Appendix}

S. Alexander's Deposition of Arius
and his companions,

and Encyclical Letter on the subject

———————

[As Montfaucon has introduced the two documents
which follow into his Edition, it has been thought that, though not
Athanasius's, they might occupy a place in a volume, like the present,
which already contains so large a collection of the ecclesiastical
tracts and papers of the day in which it belongs. Should the internal
character of the Encyclical Letter lead to the suspicion that it is
probably Athan.'s own composition, in his situation of Deacon to St.
Alexander, or at least as being in his intimate confidence, there will
be a further reason for introducing it here. The grounds of this
conjecture are such as the following. 1. it is written in a style
altogether unlike S. Alexander's, which, (as we see in his Epistle to
S. Alexander of Constantinople contained in Theod. Hist. i. 3.) is
elaborate and involved and abounding in compound words, with nothing
of the simplicity and vigor of St. Athan.'s; with which, 2. the style
of this document is identical, using the very same words and terms of
expression for which Athan. is so remarkable. 3. The theological
terms, nay the theological view, of St. Alex., is proper to himself,
and could not suitably be ascribed to S. Athan., who, to say no more,
has far fewer technical phrases than his predecessor; and here the
Encyclical Epistle answers to S. Athan.'s writings, not to St. Alex.'s.
4. Certain texts quoted in the course of it, are used as Athan. quotes
and uses them in his acknowledged works. Some of these points of
resemblance and dissimilarity shall be mentioned in the notes. The
date of St. Alexander's document is 321.]

———————

Alexander, being assembled \footnote{Margin Notes

1. [\emph{par\={o}n parousin}].

Return to text} with his beloved brethren, the Presbyters and Deacons of
Alexandria, and the Mareotis, greets them in the Lord.

Although you have already subscribed to the
letter I addressed to the followers of Arius, exhorting them to
renounce his impiety, and to submit themselves to the sound Catholic
Faith, and have shewn your right-mindedness \footnote{[\emph{orth\={e}n}].

Return to text} and agreement in the doctrines of the Catholic Church; yet
forasmuch as I have written also to our fellow-ministers in every
place concerning the Arians, and especially since some of you, as the
Presbyters Chares and Pistus \footnote{pp. 37, 44.

Return to text},
and the Deacons Serapion, Parammon, Zozimus, and Irenæus, have
joined the Arian party, and been content to suffer deposition with
them, I thought it needful to assemble together you, the Clergy of the
city, and to send for you the Clergy of the Mareotis, in order that
you may understand what I have now written, and may testify your
agreement thereto, and give your concurrence in the deposition of the
followers of Arius and Pistus. For it is desirable that you should be
made acquainted with the sentiments I have expressed, and that each of
you should heartily embrace them, as though he had written them
himself.

A Copy

To his dearly beloved and most honoured
fellow-ministers \footnote{[\emph{sulleitouzgois}],
colleagues.

Return to text} of the
Catholic Church in every place, Alexander sends health in the Lord.

§. 1.

1. AS
there is one body \footnote{St. Alexander in Theod. begins his Epistle to
his namesake of Constantinople with some moral reflections, concerning
ambition and avarice. Athan. indeed uses a similar introduction to his
Ep. Æg. but it is not addressed to an individual.

Return to text} of the
Catholic Church, and a command is given us in the sacred Scriptures to
preserve the bond of unity and peace [Eph. iv. 3.], it is agreeable
thereto, that we should write and signify to one another whatever is
done by each of us individually; so that whether one member suffer or
rejoice, we may either suffer or rejoice with one another. Now there
are gone forth in this diocese, at this time, certain lawless \footnote{[\emph{paranomoi}] vid. Hist. Ar. §. 71 init. §. 75 fin. 79. A.

Return to text} men, enemies of Christ, teaching an apostasy, which one may
justly suspect and designate as the forerunner \footnote{[\emph{prodromon Antichristou}]. vid. Orat. i. 7. B. Vit. Ant. 69. A.
vol. 8. p. 79, note q.

Return to text} of Antichrist. I was desirous \footnote{[\emph{kai eboulom\={e}n men si\={o}t\={e} … epeide de …
anagk\={e}n eschon}]. vid. Apol. contr. Ar. §. 1 init. de Decr.
§. 2. F. Orat. i. 23 init. Orat. ii init. Orat. iii. 1. A. ad Serap.
i. 1. C. 16. C. ii. 1 init. iii init. iv. 8 init. Ep. ad Mon. §. 2.
E. ad Epict. 3 fin, ad Max. §. 1. contr. Apollin. i. 1 init.

Return to text} to pass such a matter by without notice, in the hope that
perhaps the evil would spend itself among its supporters, and not
extend to other places to defile \footnote{[\emph{rhup\={o}s\={e}i}] and infr. [\emph{rhupon}]. vid. Hist.
Ar. §. 3. C. §. 80. B. de Decr. §. 2. C. Ep. Æg. 11 fin. Orat. i.
10. C.

Return to text} the ears \footnote{[\emph{akoas}], and infr. [\emph{akoas buei}]. vid. Ep. Æg. §. 13.
A. Orat. i. §. 7. A. Hist. Ar. § 56. B.

Return to text} of the
simple \footnote{[\emph{akerai\={o}n}] Apol. contr. Ar. §. 1. A. Ep. Æg. §. 18.
E. ad Epict. §. 1. fin. ad Adelph. §. 2. fin. Orat. i. 8. E.

Return to text}. But seeing that
Eusebius now of Nicomedia, who thinks that the government of the
Church rests with him, because retribution has not come upon him for
his desertion of Berytus, when he had cast an eye \footnote{[\emph{epophthalmisas}] also used of Eusebius. Apol. contr. Ar. §. 6.
D. Hist. Ar. §. 7. A.

Return to text} of desire on the Church of the Nicomedians, begins to support
these apostates, and has taken upon him to write letters every where
in their behalf, if by any means he may draw in certain ignorant
persons to this most base and antichristian heresy; I am therefore
constrained, knowing what is written in the law, no longer to hold my
peace, but to make it known to you all; that you may understand who
the apostates are, and the unhappy terms \footnote{[\emph{rh\={e}matia}]. vid. de Decr. §. 8. A. 18. E. Orat. i. 10.
D. de Sent. D. §. 23. init. S. Dionysius also uses it. ibid. §. 18.
A.

Return to text} which their heresy has adopted, and that, should Eusebius write
to you, you may pay no attention to him, for he now desires by means
of these men to exhibit anew his old malevolence \footnote{[\emph{kakonoian}]. vid. Hist. Ar. §. 75. E. de Decr. §. 1. D. et
al.

Return to text}, which has so long been concealed, pretending to write in their
favour, while in truth it clearly appears, that he does it to forward
his own interests.

§. 2.

2. Now the apostates are these, Arius, Achilles,
Anthales, Carpones, another Arius, and Sarmates, sometime Presbyters;
Euzoïus, Lucius, Julius, Menas, Helladus, and Gaius, sometime
Deacons; and with them Secundus and Theonas, sometime called Bishops.
And the novelties they have invented and put forth contrary to the
Scriptures are these following:—God was not always a Father \footnote{[\emph{ouk aei pat\={e}r}]. This enumeration of Arius's tenets, and
particularly the mention of the first, corresponds to de Decr. §. 6.
Ep. Æg. t. 12. as being taken from the Thalia. Orat. i. §. 5. and
far less with Alex. ap. Theod. p. 731, 2. vid. also Sent. D. §. 16. [\emph{katachr\={e}stik\={o}s}],
which is found here, occurs de Decr. §. 6. B.

Return to text}, but there was a time when God was not a Father. The Word of
God was not always, but was made of things that were not: for God that
is, made Him that was not, of things that were not; wherefore there
was a time when He was not; for the Son is a creature and a work.
Neither is He like in substance to the Father; neither is He the true
and natural Word of the Father; neither is He His true Wisdom; but He
is one of the things made and created, and is called the Word and
Wisdom by an abuse of terms, since He Himself was made by the proper
Word of God, and by the Wisdom that is in God, by which God made not
only all other things but Him also. Wherefore He is by nature subject
to change and variation, as are all rational creatures. And the Word
is foreign from the substance \footnote{[\emph{ousian; ousia tou logou}] or [\emph{tou huiou}] is a familiar
expression with Athan. e.g. Orat. i. 45. ii. 7. B. 9. B. 11. B. 12. A.
13. B. C. 18 init. 22. E. 47 init. 56 init. \&c. for which Alex.
Theod. uses the word [\emph{hypostasis}], e.g. [\emph{t\={e}n
idiotropon autou hypostasin; t\={e}s hypostase\={o}s autou
aperiepgastou; ne\={o}teran t\={e}s hypostase\={o}s genesin;
h\={e} tou monogenous anakdi\={e}g\={e}tos hypostasis; t\={e}n
tou logou hypostasin}].

Return to text}
of the Father, and is alien and separate therefrom. And the Father
cannot be described by the Son, for the Word does not know the Father
perfectly and accurately, neither can He see Him perfectly. Moreover,
the Son knows not His own substance as it really is; for He was
created for us, that God might create us by Him, as by an instrument;
and He would not have existed, had not God wished to create us.
Accordingly, when some one asked them, whether the Word of God can
possibly change as the devil changed, they were not afraid to say that
He can; for being something made and created, His nature is subject to
change.

§. 3.

3. Now when the Arians made these assertions, and
shamelessly avowed them, we being assembled with the Bishops of Egypt
and Libya, nearly a hundred in number, anathematized both them and
their followers. But the Eusebians admitted them to communion,
being desirous to mingle falsehood with the truth, and impiety with
piety. But they will not be able to do so, for the truth must prevail;
neither is there any \emph{communion of light with darkness} [2 Cor.
vi. 14.], nor any \emph{concord of Christ with Belial} \footnote{[\emph{koin\={o}nia ph\={o}ti}]. This is quoted Alex. ap. Theod.
Hist. i. 3. p. 738; by S. Athan. in the Letter published by Maffei,
ed. Patav. t. 3. p. 87. It seems to have been a received text in the
controversy, as the Sardican Council uses it supr. p. 76. and S. Athan.
seems to put it into the mouth of St. Anthony, Vit. Ant. 69. A.

Return to text}. For who ever heard such assertions before \footnote{[\emph{tis gar \={e}kouse}]. Ep. Æg. §. 7 init. ad Epict. §. 2
init. Orat. i. 8. B. C. Apol. contr. Ar. 85 init. Hist. Ar. §. 46
init. §. 73 init. §. 74 init. ad Serap. iv. 2 init.

Return to text}? or who that hears them now is not astonished and does not stop
his ears lest their filthy language should touch them? Who that has
heard the words of John, \emph{In the beginning was the Word} [John i.
1.], will not denounce the saying of these men, that ``there was a time
when He was not?'' Or who that has heard in the Gospel, \emph{the
Only-begotten Son}, and \emph{by Him were all things made} [Ib.
xiv. and xviii. 3.], will not detest their declaration that He is ``one
of the things that were made.'' For how can He be one of those things
which were made by Himself? or how can He be the Only-begotten, when,
according to them, He is counted as one among the rest, since He is
Himself a creature and a work? And how can He be ``made of things that
were not,'' when the Father saith, \emph{My heart hath brought forth a
good Word}, and, \emph{Out of the womb I have begotten Thee before the
morning star} [Ps. xlv. 1. Ib. cx. 3.]? Or again, how is He ``unlike
in substance to the Father,'' seeing He is the perfect \emph{image} and
\emph{brightness} of the Father [Heb. i. 3.], and that He saith, \emph{He
that hath seen Me hath seen the Father} [John xiv. 9.] \footnote{On the concurrence of these three texts in Athan. (though other
writers use them too, and Alex. ap. Theod. has two of them,) vid. vol.
8. p. 208, note b.

Return to text}? And if the Son is the Word and Wisdom of God, how was there ``a
time when He was not?'' It is the same as if they should say that God
was once without Word and without Wisdom \footnote{[\emph{alogon kai asophon ton theon}]. de Decr. §. 15. Orat. i. §.
19. vid. vol. 8. p. 25, note c. p. 208, note b.

Return to text}. And how is He ``subject to change and variation,'' who says, by
Himself, \emph{I am in the Father, and the Father in Me} \footnote{On the concurrence of these three texts in Athan. (though other
writers use them too, and Alex. ap. Theod. has two of them,) vid. vol.
8. p. 208, note b.

Return to text}, and, \emph{I and the Father are one} [v. 10. Ib. x. 30.] \footnote{On the concurrence of these three texts in Athan. (though other
writers use them too, and Alex. ap. Theod. has two of them,) vid. vol.
8. p. 208, note b.

Return to text}, and by the Prophet, \emph{Behold Me, for I am, and I change not}
[Mal. iii. 6.] \footnote{This text is thus applied by Athan. Orat. i. 36. D. ii. 10. A. In the
first of these passages he uses the same apology, nearly in the same
words, which is contained in the text.

Return to text}? For
although one may refer this expression to the Father, yet it may now
be more aptly spoken of the Word, viz. that though He has been made
man, He has not changed; but as the Apostle has said, \emph{Jesus Christ
is the same yesterday, today, and for ever} [Heb. xiii. 8.]. And
who can have persuaded them to say, that He was made for us, whereas
Paul writes, \emph{for whom are all things, and by whom are all things}
[Ib. ii. 10.]? §. 4. As to their blasphemous position that ``the Son knows not the Father perfectly,'' we ought not to wonder at it; for
having once set themselves to fight against Christ, they contradict
even His express words, since He says, \emph{As the Father knoweth Me,
even so know I the Father} [John x. 15.]. Now if the Father knows
the Son but in part, then it is evident that the Son does not know the
Father perfectly; but if it is not lawful to say this, but the Father
does know the Son perfectly, then it is evident that as the Father
knows His own Word, so also the Word knows His own Father whose Word
He is.

§. 5.

4. By these arguments and appeals to the sacred
Scriptures we frequently overthrew them; but they changed like
chameleons \footnote{[\emph{chamaileontes}]. vid. de Decr. §. 1. D. Hist. Ar. §. 79.

Return to text}, and again
shifted their ground, striving to bring upon themselves that sentence,
\emph{when the impious falleth into the depth of evils, he is filled with
contempt} [Prov. xviii. 3.]. There have been many heresies before
them, which, venturing further than they ought, have fallen into
folly; but these men by endeavouring in all their positions to
overthrow the Divinity of the Word, have justified the other in
comparison of themselves, as approaching nearer to Antichrist.
Wherefore they have been excommunicated and anathematized by the
Church. We grieve for their destruction, and especially because,
having once been instructed in the doctrines of the Church, they have
now fallen away. Yet we are not greatly surprised; for Hymeneus and
Philetus did the same, and before them Judas, who followed our Saviour,
but afterwards became a traitor and an apostate [2 Tim. ii. 17.]. And
concerning these same persons, we have not been left without
instruction; for our Lord has forewarned us; \emph{Take heed lest any man
deceive you: for many shall come in My name, saying, I am Christ, and
the time draweth near, and they shall deceive many; go ye not after
them} [Luke xxi. 8.]. And Paul, who was taught these things by our
Saviour, wrote, that \emph{in the latter times some shall depart from the
sound faith, giving heed to seducing spirits and doctrines of devils,
which reject the truth} [1 Tim. iv. 1.] \footnote{Into this text which Athan. also applies to the Arians, (vid. vol. 8.
p. 191, note e.) Athan. also introduces, like Alexander here, the word
[\emph{hygiainous\={e}}], e.g. Ep. Æg. §. 20. Orat. i. 8 fin. de
Decr. 3. E. Hist. Arian. §. 78 init. \&c. It is quoted without the
word by Origen contr. Cels. v. 64. but with [\emph{hygious}] in Matth.
t. xiv. 16. Epiphan. has [\emph{hygiainous\={e}s didaskalias}], Hær.
78. 2. [\emph{hygious did}]. ibid. 23. p. 1055.

Return to text}.

§. 6.

5. Since then our Lord and Saviour Jesus Christ
hath instructed us by His own mouth, and also hath signified to us
concerning such men by the Apostle, we accordingly being personal
witnesses of their impiety, have anathematized, as we said, all such,
and declared them to be alien from the Catholic Faith and Church. And
we have made this known to your piety, dearly beloved and most
honoured fellow- ministers, in order that should any of them have the
boldness \footnote{[\emph{propeteusainto}]. vid. de Decr. §. 2. B.

Return to text} to come unto you, you may not receive them, nor comply with the desires of
Eusebius, or any other person writing in their behalf. For it becomes
us who are Christians to turn away from all who speak or think any
thing against Christ, as being enemies of God, and destroyers \footnote{[\emph{phthoreas t\={o}n psuch\={o}n}], but S. Alex. in Theod.
uses the compound word [\emph{phthoropoios}]. p. 731. Other compound
or recondite words (to say nothing of the construction of sentences)
found in S. Alexander's Letter in Theod., and unlike the style of the
Circular under review, are such as [\emph{h\={e} philarchos kai
philarguros prothesis; christemporian phrenoblabous idiotropon;
homostoichois syllabais; the\={e}gorous apostolous antidiastol\={e}n
t\={e}s patrik\={e}s maieuse\={o}s melancholik\={e}n
philotheos saph\={e}neia anosiourgias; phl\={e}naph\={o}n
myth\={o}n}]. Instances of theological language in S. Alex. to
which the Letter in the text contains no resemblance are [\emph{ach\={o}rista
pragmata duo; ho huios t\={e}n kata panta homoiot\={e}ta autou
ex physe\={o}s apomaxamenos; di esoptrou ak\={e}lid\={o}tou
kai empsuchou theias eikonos; mesiteuouoa physis monogen\={e}s; tas
t\={e}i hypostasei duo physeis}].

Return to text

} of souls; and not even to bid such God speed, lest we become
partakers of their sins, as the blessed John hath charged us. Salute
the brethren that are with you. They that are with me salute you.

Presbyters of Alexandria

§. 7.

6. I, Colluthus, Presbyter, agree with what is
here written, and give my assent to the deposition of Arius and his
associates in impiety.

Alexander \footnote{vid.
Presbyters p. 105.

Return to text},
Presbyter, likewise
Nemesius, Presbyter

Dioscorus \footnote{vid.
Presbyters p. 105.

Return to text}, Presbyter,
likewise
Longus \footnote{vid.
Presbyters p. 105.

Return to text}, Presbyter

Dionysius \footnote{vid.
Presbyters p. 105.

Return to text}, Presbyter,
likewise
Silvanus, Presbyter

Eusebius, Presbyter, likewise
Perous, Presbyter

Alexander, Presbyter, likewise
Apis, Presbyter

Nilanus \footnote{Nilaras? p.
105

Return to text},
Presbyter, likewise
Proterius, Presbyter

Arpocration, Presbyter, likewise
Paulus, Presbyter

Agathus, Presbyter
Cyrus, Presbyter, likewise

Deacons

Ammonius \footnote{vid.
Presbyters p. 105.

Return to text},
Deacon, likewise
Ambytianus, Deacon

Macarius, Deacon
Gaius \footnote{vid.
Presbyters p. 105.

Return to text}, Deacon, likewise

Pistus \footnote{vid.
Presbyters p. 105.

Return to text}, Deacon, likewise
Alexander, Deacon

Athanasius \footnote{vid.
Presbyters p. 105.

Return to text}, Deacon
Dionysius, Deacon

Eumenes, Deacon
Agathon, Deacon

Apollonius \footnote{vid.
Presbyters p. 105.

Return to text}, Deacon
Polybius, Deacon, likewise

Olympius, Deacon
Theonas, Deacon

Aphthonius \footnote{vid.
Presbyters p. 105.

Return to text}, Deacon
Marcus, Deacon

Athanasius, Deacon
Comodus, Deacon

Macarius, Deacon, likewise
Serapion \footnote{vid.
Presbyters p. 105.

Return to text}, Deacon

Paulus, Deacon
Nilus, Deacon

Petrus, Deacon
Romanus, Deacon, likewise

Presbyters of the Mareotis

I, Apollonius, Presbyter, agree with what is here
written, and give my assent to the deposition of Arius and his
associates in impiety.

Ingenius \footnote{p. 107.

Return to text},
Presbyter, likewise
Serenus, Presbyter

Ammonius, Presbyter
Didymus, Presbyter

Dioscorus \footnote{p. 107.

Return to text}, Presbyter
Heracles \footnote{Heraclius? p.
107.

Return to text},
Presbyter

Sostras, Presbyter
Boccon \footnote{p. 107.

Return to text}, Presbyter

Theon \footnote{p. 107.

Return to text}, Presbyter
Agathus, Presbyter

Tyrannus, Presbyter
Achillas, Presbyter

Copres, Presbyter
Paulus, Presbyter

Ammonas \footnote{p. 107.

Return to text}, Presbyter
Thalelæus, Presbyter

Orion, Presbyter
Dionysius, Presbyter, likewise

Deacons

Serapion \footnote{p. 107.

Return to text},
Deacon, likewise
Didymus, Deacon

Justus, Deacon, likewise
Ptollarion \footnote{p. 107.

Return to text}, Deacon

Didymus, Deacon
Seras, Deacon

Demetrius \footnote{p. 107.

Return to text}, Deacon
Gaius \footnote{p. 107.

Return to text}, Deacon

Maurus \footnote{p. 107.

Return to text},
Deacon
Hierax \footnote{p. 107.

Return to text}, Deacon

Alexander, Deacon
Marcus, Deacon

Marcus \footnote{p. 107.

Return to text}, Deacon
Theonas, Deacon

Comon, Deacon
Sarmaton, Deacon

Tryphon \footnote{p. 107.

Return to text}, Deacon
Carpon, Deacon

Ammonius \footnote{p. 107.

Return to text}, Deacon
Zoilus, Deacon. likewise
